% Options for packages loaded elsewhere
\PassOptionsToPackage{unicode}{hyperref}
\PassOptionsToPackage{hyphens}{url}
%
\documentclass[
  11pt,
]{book}
\usepackage{amsmath,amssymb}
\usepackage{iftex}
\ifPDFTeX
  \usepackage[T1]{fontenc}
  \usepackage[utf8]{inputenc}
  \usepackage{textcomp} % provide euro and other symbols
\else % if luatex or xetex
  \usepackage{unicode-math} % this also loads fontspec
  \defaultfontfeatures{Scale=MatchLowercase}
  \defaultfontfeatures[\rmfamily]{Ligatures=TeX,Scale=1}
\fi
\usepackage{lmodern}
\ifPDFTeX\else
  % xetex/luatex font selection
\fi
% Use upquote if available, for straight quotes in verbatim environments
\IfFileExists{upquote.sty}{\usepackage{upquote}}{}
\IfFileExists{microtype.sty}{% use microtype if available
  \usepackage[]{microtype}
  \UseMicrotypeSet[protrusion]{basicmath} % disable protrusion for tt fonts
}{}
\makeatletter
\@ifundefined{KOMAClassName}{% if non-KOMA class
  \IfFileExists{parskip.sty}{%
    \usepackage{parskip}
  }{% else
    \setlength{\parindent}{0pt}
    \setlength{\parskip}{6pt plus 2pt minus 1pt}}
}{% if KOMA class
  \KOMAoptions{parskip=half}}
\makeatother
\usepackage{xcolor}
\usepackage{color}
\usepackage{fancyvrb}
\newcommand{\VerbBar}{|}
\newcommand{\VERB}{\Verb[commandchars=\\\{\}]}
\DefineVerbatimEnvironment{Highlighting}{Verbatim}{commandchars=\\\{\}}
% Add ',fontsize=\small' for more characters per line
\newenvironment{Shaded}{}{}
\newcommand{\AlertTok}[1]{\textcolor[rgb]{1.00,0.00,0.00}{\textbf{#1}}}
\newcommand{\AnnotationTok}[1]{\textcolor[rgb]{0.38,0.63,0.69}{\textbf{\textit{#1}}}}
\newcommand{\AttributeTok}[1]{\textcolor[rgb]{0.49,0.56,0.16}{#1}}
\newcommand{\BaseNTok}[1]{\textcolor[rgb]{0.25,0.63,0.44}{#1}}
\newcommand{\BuiltInTok}[1]{\textcolor[rgb]{0.00,0.50,0.00}{#1}}
\newcommand{\CharTok}[1]{\textcolor[rgb]{0.25,0.44,0.63}{#1}}
\newcommand{\CommentTok}[1]{\textcolor[rgb]{0.38,0.63,0.69}{\textit{#1}}}
\newcommand{\CommentVarTok}[1]{\textcolor[rgb]{0.38,0.63,0.69}{\textbf{\textit{#1}}}}
\newcommand{\ConstantTok}[1]{\textcolor[rgb]{0.53,0.00,0.00}{#1}}
\newcommand{\ControlFlowTok}[1]{\textcolor[rgb]{0.00,0.44,0.13}{\textbf{#1}}}
\newcommand{\DataTypeTok}[1]{\textcolor[rgb]{0.56,0.13,0.00}{#1}}
\newcommand{\DecValTok}[1]{\textcolor[rgb]{0.25,0.63,0.44}{#1}}
\newcommand{\DocumentationTok}[1]{\textcolor[rgb]{0.73,0.13,0.13}{\textit{#1}}}
\newcommand{\ErrorTok}[1]{\textcolor[rgb]{1.00,0.00,0.00}{\textbf{#1}}}
\newcommand{\ExtensionTok}[1]{#1}
\newcommand{\FloatTok}[1]{\textcolor[rgb]{0.25,0.63,0.44}{#1}}
\newcommand{\FunctionTok}[1]{\textcolor[rgb]{0.02,0.16,0.49}{#1}}
\newcommand{\ImportTok}[1]{\textcolor[rgb]{0.00,0.50,0.00}{\textbf{#1}}}
\newcommand{\InformationTok}[1]{\textcolor[rgb]{0.38,0.63,0.69}{\textbf{\textit{#1}}}}
\newcommand{\KeywordTok}[1]{\textcolor[rgb]{0.00,0.44,0.13}{\textbf{#1}}}
\newcommand{\NormalTok}[1]{#1}
\newcommand{\OperatorTok}[1]{\textcolor[rgb]{0.40,0.40,0.40}{#1}}
\newcommand{\OtherTok}[1]{\textcolor[rgb]{0.00,0.44,0.13}{#1}}
\newcommand{\PreprocessorTok}[1]{\textcolor[rgb]{0.74,0.48,0.00}{#1}}
\newcommand{\RegionMarkerTok}[1]{#1}
\newcommand{\SpecialCharTok}[1]{\textcolor[rgb]{0.25,0.44,0.63}{#1}}
\newcommand{\SpecialStringTok}[1]{\textcolor[rgb]{0.73,0.40,0.53}{#1}}
\newcommand{\StringTok}[1]{\textcolor[rgb]{0.25,0.44,0.63}{#1}}
\newcommand{\VariableTok}[1]{\textcolor[rgb]{0.10,0.09,0.49}{#1}}
\newcommand{\VerbatimStringTok}[1]{\textcolor[rgb]{0.25,0.44,0.63}{#1}}
\newcommand{\WarningTok}[1]{\textcolor[rgb]{0.38,0.63,0.69}{\textbf{\textit{#1}}}}
\usepackage{longtable,booktabs,array}
\usepackage{calc} % for calculating minipage widths
% Correct order of tables after \paragraph or \subparagraph
\usepackage{etoolbox}
\makeatletter
\patchcmd\longtable{\par}{\if@noskipsec\mbox{}\fi\par}{}{}
\makeatother
% Allow footnotes in longtable head/foot
\IfFileExists{footnotehyper.sty}{\usepackage{footnotehyper}}{\usepackage{footnote}}
\makesavenoteenv{longtable}
\usepackage{graphicx}
\makeatletter
\def\maxwidth{\ifdim\Gin@nat@width>\linewidth\linewidth\else\Gin@nat@width\fi}
\def\maxheight{\ifdim\Gin@nat@height>\textheight\textheight\else\Gin@nat@height\fi}
\makeatother
% Scale images if necessary, so that they will not overflow the page
% margins by default, and it is still possible to overwrite the defaults
% using explicit options in \includegraphics[width, height, ...]{}
\setkeys{Gin}{width=\maxwidth,height=\maxheight,keepaspectratio}
% Set default figure placement to htbp
\makeatletter
\def\fps@figure{htbp}
\makeatother
\setlength{\emergencystretch}{3em} % prevent overfull lines
\providecommand{\tightlist}{%
  \setlength{\itemsep}{0pt}\setlength{\parskip}{0pt}}
\setcounter{secnumdepth}{-\maxdimen} % remove section numbering
\ifLuaTeX
\usepackage[bidi=basic]{babel}
\else
\usepackage[bidi=default]{babel}
\fi
\babelprovide[main,import]{american}
% get rid of language-specific shorthands (see #6817):
\let\LanguageShortHands\languageshorthands
\def\languageshorthands#1{}
\usepackage[letterpaper,left=1.20in,right=1.20in,top=1.00in,bottom=1.05in,includehead,includefoot,headsep=14pt,footskip=26pt]{geometry}
\usepackage{microtype}
\usepackage{mathtools}
\usepackage{amssymb}
\usepackage{mathrsfs}
\usepackage{centernot}
\usepackage{float}
\usepackage{caption}
\usepackage{fancyhdr}
\usepackage{titlesec}
\usepackage{enumitem}
\usepackage{titling}
\usepackage[most]{tcolorbox}
\usepackage{etoolbox}
\usepackage{booktabs}
\usepackage{tabularx}
\usepackage{array}
\usepackage{pdflscape}
\usepackage{adjustbox}
\usepackage{needspace}
\usepackage{xurl}   % long URLs break at any character
\graphicspath{{src/}{./}}

\definecolor{tnnavy}{HTML}{173B5E}
\definecolor{tnblue}{HTML}{275DAD}
\definecolor{tnteal}{HTML}{2A9D8F}
\definecolor{tnrust}{HTML}{B84A2B}
\definecolor{tncream}{HTML}{F7F2E8}
\definecolor{tnlightblue}{HTML}{EDF4FA}
\definecolor{tngray}{HTML}{4B5563}


% ---- status vocabulary -------------------------------------------------
\definecolor{tngreen}{HTML}{1F7A4D}
\definecolor{tnamber}{HTML}{8A6D1F}

\newcommand{\tnstatusbox}[2]{%
  {\setlength{\fboxsep}{2.1pt}%
   \colorbox{#1}{\textcolor{white}{\scriptsize\bfseries\hspace{0.5pt}#2\hspace{0.5pt}}}}}
\newcommand{\PROVED}{\tnstatusbox{tngreen}{PROVED}}
\newcommand{\CERT}{\tnstatusbox{tnblue}{CERTIFIED}}
\newcommand{\EQUIV}{\tnstatusbox{tnnavy}{RH-EQUIVALENT}}
\newcommand{\OPENSTAT}{\tnstatusbox{tnrust}{OPEN}}
\newcommand{\EVID}{\tnstatusbox{tnamber}{EVIDENCE}}

\newcommand{\statuslegend}{%
\begin{center}\small
\begin{tabular}{@{}l@{\hspace{7pt}}p{0.78\textwidth}@{}}
\PROVED{} & proved here or cited as proved: an unconditional theorem or an exact identity.\\[1.6pt]
\CERT{} & rigorous but computer-assisted, or dependent on verified zero heights; the finite region it certifies is stated with it.\\[1.6pt]
\EQUIV{} & logically equivalent to the Riemann hypothesis. An equivalence is not a proof in either direction.\\[1.6pt]
\OPENSTAT{} & not proved anywhere in this work, and not implied by any finite result in it.\\[1.6pt]
\EVID{} & numerical or empirical support only. It never upgrades a status.\\
\end{tabular}
\end{center}}

\newtcolorbox{tnclosed}[1]{enhanced,breakable,
  colback=white,colframe=tngreen,boxrule=1.1pt,arc=2pt,
  left=7pt,right=7pt,top=6pt,bottom=6pt,
  fonttitle=\bfseries\scriptsize,coltitle=white,colbacktitle=tngreen,
  title={CLOSED FORM\quad\normalfont\itshape #1}}

\newtcolorbox{tnequiv}[1]{enhanced,breakable,
  colback=tnlightblue,colframe=tnnavy,boxrule=1.1pt,arc=2pt,
  left=7pt,right=7pt,top=6pt,bottom=6pt,
  fonttitle=\bfseries\scriptsize,coltitle=white,colbacktitle=tnnavy,
  title={RH-EQUIVALENT\quad\normalfont\itshape #1}}

\newtcolorbox{tnopen}[1]{enhanced,breakable,
  colback=white,colframe=tnrust,boxrule=1.1pt,arc=2pt,
  left=7pt,right=7pt,top=6pt,bottom=6pt,
  fonttitle=\bfseries\scriptsize,coltitle=white,colbacktitle=tnrust,
  title={OPEN\quad\normalfont\itshape #1}}
% -----------------------------------------------------------------------------


% Page fill: leave slack at the foot rather than stretching every gap.
\raggedbottom
\setlength{\parskip}{3.4pt plus 1.0pt minus 0.8pt}
\AtBeginDocument{%
  \setlength{\abovedisplayskip}{6.0pt plus 1.8pt minus 2.2pt}%
  \setlength{\belowdisplayskip}{6.0pt plus 1.8pt minus 2.2pt}%
  \setlength{\abovedisplayshortskip}{2.5pt plus 1pt}%
  \setlength{\belowdisplayshortskip}{4.5pt plus 1pt}%
}

% Figures were sized for a narrower text block; hold them to their old measure.
\setkeys{Gin}{width=0.86\linewidth,keepaspectratio}
\floatplacement{figure}{H}
\DeclareCaptionFont{tnnavyfont}{\color{tnnavy}}
\DeclareCaptionFont{tngrayfont}{\color{tngray}}
\captionsetup{
  font=small,
  labelfont={bf,tnnavyfont},
  textfont=tngrayfont,
  justification=centering,
  singlelinecheck=false,
  skip=5pt
}

\setlist{itemsep=2pt,topsep=4pt}
\renewcommand{\arraystretch}{1.13}
\AtBeginEnvironment{longtable}{\scriptsize\setlength{\tabcolsep}{3pt}}
\let\cleardoublepage\clearpage

\titleformat{name=\chapter,numberless}[display]
  {\normalfont\huge\bfseries\color{tnnavy}}
  {}
  {0pt}
  {\titlerule[1.1pt]\vspace{0.7ex}\filright}
  [\vspace{0.6ex}\titlerule]

\titleformat{\section}
  {\normalfont\Large\bfseries\color{tnnavy}}
  {}
  {0pt}
  {}

\titleformat{\subsection}
  {\normalfont\large\bfseries\color{tnteal}}
  {}
  {0pt}
  {}

\titlespacing*{\chapter}{0pt}{-42pt}{16pt}
\titlespacing*{\section}{0pt}{17pt}{7pt}
\titlespacing*{\subsection}{0pt}{12.5pt}{5pt}

\setlength{\droptitle}{-3.3em}
\pretitle{\begin{center}\color{tnnavy}\Huge\bfseries}
\posttitle{\par\end{center}\vspace{0.35em}}
\preauthor{\begin{center}\large\lineskip .5em\begin{tabular}[t]{c}}
\postauthor{\end{tabular}\par\end{center}}
\predate{\begin{center}\normalsize}
\postdate{%
  \par\vspace{0.35em}
  {\color{tnrust}\bfseries RH STATUS: OPEN}
  \par\end{center}\thispagestyle{empty}
}

\pagestyle{fancy}
\fancyhf{}
\fancyhead[R]{\small\color{tngray}\thepage}
\fancyhead[L]{\small\color{tngray}\nouppercase{\rightmark}}
\fancyfoot[L]{\scriptsize\color{tngray}$s=\sigma+it$}
\fancyfoot[C]{\scriptsize\color{tngray}Volume I -- Technical Dossier, v5.1 -- RH OPEN}
\renewcommand{\headrulewidth}{0.4pt}
\renewcommand{\footrulewidth}{0pt}
\setlength{\headheight}{15pt}

% A two-column contents list: the same depth, a third of the pages.
\usepackage{multicol}
\makeatletter
\renewcommand{\tableofcontents}{%
  \chapter*{\contentsname}%
  \@mkboth{\contentsname}{\contentsname}%
  \begingroup\small\setlength{\columnsep}{20pt}%
  \begin{multicols}{2}\@starttoc{toc}\end{multicols}%
  \endgroup}
\makeatother

\renewenvironment{quote}
  {\begin{tcolorbox}[
      enhanced,
      breakable,
      colback=tncream,
      colframe=tnrust,
      boxrule=0.8pt,
      arc=1.5pt,
      left=7pt,
      right=7pt,
      top=6pt,
      bottom=6pt]}
  {\end{tcolorbox}}

\AtBeginDocument{
  \hypersetup{
    linkcolor=tnblue,
    urlcolor=tnteal,
    citecolor=tnrust,
    pdfauthor={Jeffery Huckstead; Cerebral Graphix},
    pdftitle={Triangular-n Postmaster: Technical Dossier and Empirical Context},
    pdfsubject={v5.1 technical dossier: triangular coordinates, Euler constant bridge, completed folding, Widder sectors, Li-Loewner geometry, Bernstein-Laguerre energy, exponential prime cutoffs and the open fixed-scale prime estimate; Riemann hypothesis open}
  }
}
\ifLuaTeX
  \usepackage{selnolig}  % disable illegal ligatures
\fi
\usepackage{bookmark}
\IfFileExists{xurl.sty}{\usepackage{xurl}}{} % add URL line breaks if available
\urlstyle{same}
\hypersetup{
  pdftitle={Triangular-n Postmaster: Technical Dossier and Empirical Context},
  pdfauthor={Jeffery Huckstead; ORCID 0009-0007-0234-2177; Cerebral Graphix},
  pdflang={en-US},
  hidelinks,
  pdfcreator={LaTeX via pandoc}}

\title{Triangular-n Postmaster: Technical Dossier and Empirical Context}
\usepackage{etoolbox}
\makeatletter
\providecommand{\subtitle}[1]{% add subtitle to \maketitle
  \apptocmd{\@title}{\par {\large #1 \par}}{}{}
}
\makeatother
\subtitle{Volume I --- the proofs behind the Reading Volume -- v5.1}
\author{Jeffery Huckstead \and ORCID 0009-0007-0234-2177 \and Cerebral
Graphix}
\date{September 2026}

\begin{document}
\frontmatter
\maketitle

{
\setcounter{tocdepth}{1}
\tableofcontents
}
\mainmatter
\newpage

\chapter{Abstract}\label{abstract}

This Dossier supplies the proofs, finite certificates, countermodels,
and separately marked empirical material behind \emph{Triangular-n
Postmaster, Volume I}. The analytic coordinate is
\(q_\rho=\rho(1-\rho)=-2T_{\rho-1}\), indexed by reflection orbits
\([\rho]=\{\rho,1-\rho\}\) with multiplicity \(m_\rho\). An off-line
quartet contributes two conjugate folded nodes. This convention controls
the canonical product, resolvent, Widder and heat identities, and the Li
sum. Conjugation is not silently included in the quotient.

The chapters develop the triangular arithmetic, the completed source,
finite positivity certificates, and exact obstructions to several
proposed shortcuts. The inherited analytic source result covers \(W_1\);
the v5.0 computer-assisted certificates cover \(W_2\)--\(W_6\) on their
stated domains. The independent source proof at \(W_7\), the all-node
matrix inequality, and the specified uniform energy and arithmetic
bounds remain open. This edition clarifies the interface and closes with
a local summary; it adds no RH implication and invokes no companion
entropy paper.

The v5.0 certificate programs and receipts are retained unchanged as
provenance. They have not been re-executed for this editorial revision.
The current receipt distinguishes preservation and interface checks from
those inherited mathematical certificates.

\chapter{How to use this Dossier}\label{how-to-use-this-dossier}

This is the second of two documents. The \textbf{Reading Volume} carries
the continuous argument from the triangular primitive to the research
boundary, in reader sections \textbf{§R1--§R25}. This Dossier carries
the proofs, the finite certificates, the countermodels, and the
empirical context that supports none of them.

\textbf{A section number without an R is here.} References of the form
§34, §104 or §166C name sections of this Dossier; references of the form
§R12 or §R23 name sections of the Reading Volume. Equation tags are
shared: (33.1), (SW.4b) and (166C.25) name the same equations in both
documents, and the identifiers under which the proofs were derived have
not been renumbered.

The signed coefficient grid \(c_{k,j}=(-1)^j\binom{k}{j}\) referred to
throughout is printed in full, for \(0\le k\le17\), as the
coefficient-grid figure of the Reading Volume.

Chapters follow mathematical dependency, not reading order:

\begin{longtable}[]{@{}
  >{\raggedright\arraybackslash}p{(\columnwidth - 4\tabcolsep) * \real{0.3333}}
  >{\raggedright\arraybackslash}p{(\columnwidth - 4\tabcolsep) * \real{0.3333}}
  >{\raggedright\arraybackslash}p{(\columnwidth - 4\tabcolsep) * \real{0.3333}}@{}}
\toprule\noalign{}
\begin{minipage}[b]{\linewidth}\raggedright
chapter
\end{minipage} & \begin{minipage}[b]{\linewidth}\raggedright
subject
\end{minipage} & \begin{minipage}[b]{\linewidth}\raggedright
supports
\end{minipage} \\
\midrule\noalign{}
\endhead
\bottomrule\noalign{}
\endlastfoot
I--II & discrete primitives, factorial families, the fractional grid,
worked triangular arithmetic (§§169--172) & §§R1--R5E \\
III & centered geometry and operator models & §§R6--R8 \\
IV & completion, folded support, the Gamma source & §§R9--R10B \\
V--VI & flux geometry, rational witnesses, Widder source proofs &
§§R11--R14 \\
VII--VIII & one-scale moment criteria, finite Löwner certificates &
§§R14--R16 \\
IX--XI & cyclic arithmetic, prime wave packets, finite Weil models &
§§R17--R17E \\
XII--XIV & dilation, completion forms, completed Bernstein rows, prime
cutoffs & §§R17F--R17H \\
XV & remaining proof obligations and countermodels & §§R22--R25 \\
Part E & empirical context and historical transcriptions & nothing \\
\end{longtable}

Status tags in this Dossier use the \textbf{single master key in the
Reading Volume front matter}. The certificate that carries the source
rungs of §74A is bundled with this edition and re-executed by the v5.0
build; software versions, precisions and run dates are recorded in the
package's QA receipt rather than in the text. The Riemann hypothesis is
open, and no result in this Dossier is offered as a proof of it.

\textbf{Labels and vocabulary.} Labels such as PP85 or PP190-T name the
author's numbered working papers, where a result or certificate was
first produced. They are kept as provenance tags only: each result is
stated here with its own status, and the working papers are not part of
this publication. An \emph{adversarial replay} reruns a certificate with
every material numerical parameter changed at once (precision, partition
and expansion orders), so that agreement is not an artifact of one
configuration. The other working terms (source, owner, rung, carrier,
rail) are defined in the glossary of the Reading Volume.

\textbf{Classical sources.} Classical results are cited where they are
used, and the lineage table of the Reading Volume maps each part of the
argument to its sources. Reference numbers are shared with the Reading
Volume; entries 83--120 were added in the previous edition.

\textbf{What changed in v5.1.} The abstract, the folded-multiplicity
explanation in Section 32, and the technical summary make the object and
the open boundary explicit. No theorem status changed.

\textbf{Retained from v5.0.} Section 155A now prints the diagonal
specializations (155A.4a) of its two product identities, so the square
form of every reflected Li rung, odd and even, appears in one place, and
cross-references Reading Volume §R10C, the new one-orbit coordinate
dictionary. Section 74A is unchanged and still carries the five
certified source rungs \(W_2,\ldots,W_6\) in one certificate --- one
partition of \(1\le s\le128\), one exact tail for \(s\ge128\), and two
programs that share no arithmetic --- re-executed by the v5.0 build on
both backends. No status changed, and the Riemann hypothesis remains
open.

\newpage

\chapter{Technical chapter I. Discrete primitives and factorial
families}\label{technical-chapter-i.-discrete-primitives-and-factorial-families}

Primitive identities generate discrete ladders. Factorial and Pell
classifications retain their exact integer conditions.

\section{5. Discrete ladders: squaring, perfect numbers, and
Pell}\label{discrete-ladders-squaring-perfect-numbers-and-pell}

Reader §R1A proves the primitive identities (TDP.1)--(TDP.5). This
section derives their repeated-square and Pell consequences. All
statements here are arithmetic; none assumes a zero-location theorem.

\subsection{Theorem TDP-A - Square-index triangular
conjugacy}\label{theorem-tdp-a---square-index-triangular-conjugacy}

Fix a positive integer \(n\) and define

\[
m_0=n,\qquad m_{k+1}=m_k^2=n^{2^{k+1}}
\qquad
U_k=T_{m_k-1},\quad V_k=T_{m_k}
\tag{TDP.A1}
\]

Then

\[
m_{k+1}=U_k+V_k,\qquad
V_{k+1}=U_k^2+V_k^2
\]

and therefore

\[
\boxed{
(U_{k+1},V_{k+1})
=\bigl(U_k^2+V_k^2-U_k-V_k,\ U_k^2+V_k^2\bigr)
}
\tag{TDP.A2}
\]

This recurrence is exactly one-dimensional on its invariant curve. Put

\[
D_k=V_k-U_k,\qquad S_k=U_k+V_k
\]

Consecutiveness gives \(D_k=m_k\) and \(S_k=m_k^2=D_k^2\). Under
(TDP.A2),

\[
\boxed{D_{k+1}=S_k=D_k^2,\qquad S_{k+1}=D_k^4}
\tag{TDP.A3}
\]

Thus the two-component polynomial recurrence is conjugate to
\(D\mapsto D^2\) on the parabola \(S=D^2\).

For \(n=2\) the first steps are

\begin{longtable}[]{@{}rrr@{}}
\toprule\noalign{}
\(k\) & \(m_k\) & \((U_k,V_k)\) \\
\midrule\noalign{}
\endhead
\bottomrule\noalign{}
\endlastfoot
0 & 2 & \((1,3)\) \\
1 & 4 & \((6,10)\) \\
2 & 16 & \((120,136)\) \\
3 & 256 & \((32640,32896)\) \\
\end{longtable}

Equivalently, every step gives

\[
T_{n^{2^{k+1}}}
=T_{n^{2^k}-1}^{\,2}+T_{n^{2^k}}^{\,2}
\tag{TDP.A4}
\]

\subsection{Corollary TDP-B - Exact Euclid--Euler
intersection}\label{corollary-tdp-b---exact-euclideuler-intersection}

The Euclid--Euler theorem says that an even perfect number is exactly

\[
2^{p-1}(2^p-1)=T_q
\qquad q=2^p-1\ \text{prime}
\tag{TDP.B1}
\]

On the lower \(n=2\) ladder, \(q_k=m_k-1=2^{2^k}-1\). At \(k=0\),
\(q_0=1\) is not prime; at \(k=1\), \(q_1=3\) gives \(T_3=6\). For every
\(k\ge2\), the exponent \(2^k\) is composite, and \(2^{2^k}-1\) is
correspondingly composite by the factorization of \(X^{ab}-1\). Hence

\[
\boxed{
\{T_{q_k}:q_k\text{ is a Mersenne-prime index}\}=\{T_3\}=\{6\}
}
\tag{TDP.B2}
\]

The next even perfect number \(T_7=28\) lies in the full Euclid--Euler
family, not in this repeated-square subfamily. This corollary concerns
even perfect numbers only.

\begin{figure}
\centering
\includegraphics{figures/figure63_triangular_square_ladder.png}
\caption{The square-index triangular recurrence is scalar squaring on
the invariant parabola, while its lower n=2 branch meets the
Euclid--Euler family only at 6.}
\end{figure}

\subsection{Theorem TDP-C - Pell companion and the unique fourth-power
lock}\label{theorem-tdp-c---pell-companion-and-the-unique-fourth-power-lock}

For nonnegative integers \(a,b\) with \(a+b>0\),

\[
b^2-a^2=T_{a+b}
\quad\Longleftrightarrow\quad
b=3a+1
\tag{TDP.C1}
\]

Indeed, cancellation of \(a+b\) reduces the equation to
\(2(b-a)=a+b+1\). If \(a=T_r\) and \(b=T_s\), (TDP.C1) becomes

\[
T_s=3T_r+1
\quad\Longleftrightarrow\quad
(2s+1)^2-3(2r+1)^2=6
\tag{TDP.C2}
\]

All nonnegative solutions form one Pell orbit. From \((r_0,s_0)=(0,1)\),

\[
\boxed{
r_{k+1}=2r_k+s_k+1
\qquad
s_{k+1}=3r_k+2s_k+2
}
\tag{TDP.C3}
\]

The first pairs are

\[
(0,1),(2,4),(9,16),(35,61),(132,229),(494,856),\ldots
\]

and every pair satisfies

\[
\boxed{T_s^2-T_r^2=T_{T_s+T_r}}
\tag{TDP.C4}
\]

Completeness follows by writing \(x=2s+1\), \(y=2r+1\) and descending
every solution with \(y>1\) by

\[
(x,y)\longmapsto(2x-3y,\ 2y-x)
\]

The first nontrivial pair carries one extra identity, and it is unique.
Put \(t=T_r\), so \(T_s=3t+1\). Equality of the fourth-power and
square-difference constructions factors as

\[
t^4+3t+1=8t^2+6t+1
\quad\Longleftrightarrow\quad
\boxed{t(t-3)(t^2+3t+1)=0}
\tag{TDP.C5}
\]

Thus \(t=0\) is trivial and \(t=3\) is the unique nontrivial nonnegative
solution. Consequently

\[
\boxed{
T_2^4+T_4
=T_4^2-T_2^2
=T_{13}
=91
}
\tag{TDP.C6}
\]

This Pell ladder is independent of the repeated-square dynamical ladder
in Theorem TDP-A. Neither has an RH implication.

\begin{figure}
\centering
\includegraphics{figures/figure68_pell_companion.png}
\caption{The repeated-square ladder and the Pell companion are distinct
exact triangular ladders: repeated squaring on one side, and a complete
Pell orbit with a unique fourth-power lock on the other.}
\end{figure}

\section{6. Simplex hierarchy and factorial
absorption}\label{simplex-hierarchy-and-factorial-absorption}

The \(q\)th simplex polynomial is

\[
S_q(x)=\binom{x+q}{q}=\frac1{q!}\prod_{j=1}^q(x+j)
\]

Every factor is an adjacent triangular difference. Consequently every
simplex polynomial is a normalized product of consecutive triangular
legs. Its reflection-completed centered roots form a symmetric
arithmetic progression, producing central-factorial parity formulas.

For integer \(n\ge1\),

\[
\frac{T_n}{n!}=\frac{n+1}{2(n-1)!}
\]

With

\[
g_n=\gcd(n+1,2(n-1)!)
\]

the reduced numerator and denominator are

\[
A_n=\frac{n+1}{g_n},\qquad
B_n=\frac{2(n-1)!}{g_n}
\]

For \(n\ge2\), the reduced numerator \(U_n\) obeys the exact successor
gate

\[
\boxed{
U_n=
\begin{cases}
n+1,&n+1\text{ prime}\\
1,&n+1\text{ composite}
\end{cases}}
\]

The prime \(2\) is the explicit base exception. Therefore
\(\epsilon_n=(U_n-1)/n\) is the successor-prime indicator and

\[
\zeta(s)=\frac1{1-2^{-s}}
\prod_{n\ge2}\left(1-(n+1)^{-s}\right)^{-\epsilon_n}
\qquad \Re s>1
\]

This is an exact encoding of existing prime support, not an estimate of
prime distribution.

\subsection{Proposition 6.1 - Ideal-lattice and radical
form}\label{proposition-6.1---ideal-lattice-and-radical-form}

The divisibility reversal in the ideal lattice of \(\mathbb Z\) is

\[
\boxed{d\mid n\iff(n)\subseteq(d)}
\]

Suppose the triangular-factorial equation

\[
T_m=n!
\]

holds. Then

\[
m(m+1)=2n!
\]

Since consecutive integers are coprime,

\[
(m)+(m+1)=\mathbb Z
\]

For comaximal ideals, product and intersection agree, so

\[
\boxed{
(m)\cap(m+1)=(m)(m+1)=(2n!)
}
\]

For \(n\ge2\), taking the ideal radical yields

\[
\boxed{
\sqrt{(2n!)}
=\bigcap_{p\le n}(p)
=\left(\prod_{p\le n}p\right)
}
\]

The logarithm of the positive generator is therefore Chebyshev's
function

\[
\log\prod_{p\le n}p=\sum_{p\le n}\log p=\vartheta(n)
\]

Thus the exact structural lane is

\[
\begin{gathered}
\text{triangular successor pair}\longrightarrow
\text{comaximal ideal factorization}\\
\longrightarrow\text{radical prime support}\longrightarrow\vartheta(n)
\end{gathered}
\]

This reformulation does not estimate \(\vartheta(n)\) and hence does not
supply prime cancellation. Its durable lesson is categorical: in
extensions where element factorization can fail, prime ideals rather
than prime elements are the stable transport objects.

\begin{figure}
\centering
\includegraphics{figures/radical_prime_support.png}
\caption{The ideal radical preserves prime support while discarding
factorial multiplicity.}
\end{figure}

\section{156. The factorial triangular switch and its diagonal
defect}\label{the-factorial-triangular-switch-and-its-diagonal-defect}

For integer \(x\ge1\),

\[
\boxed{
\frac{(x+1)!}{2(x-1)!}
=\frac{x(x+1)}2
=T_x}
\tag{156.1}
\]

The same cancellation can be displayed in the factored form noticed in
the exploratory round:

\[
\boxed{
\left(\frac{x!}{2}\right)
\left(\frac{x+1}{(x-1)!}\right)
=\frac12\left(\frac{x!}{(x-1)!}\right)(x+1)
=\frac{x(x+1)}2}
\tag{156.2}
\]

This is elementary, but structurally useful: inserting the factorial
turns a reduced rational successor gate into an exact diagonal
polynomial. Gamma continuation shows its natural domain:

\[
\frac{\Gamma(x+2)}{2\Gamma(x)}=\frac{x(x+1)}2
\tag{156.3}
\]

where the quotient is regular, and everywhere after filling its
removable values.

The mixed expression

\[
F_\triangle(x,y):=\frac{(y+1)!}{2(x-1)!}
=\frac{\Gamma(y+2)}{2\Gamma(x)}
\tag{156.4}
\]

restricts to the triangular polynomial on its diagonal:

\[
\boxed{F_\triangle(x,x)=T_x}
\tag{156.5}
\]

The diagonal is therefore an equality or balance locus. At regular
points with \(T_x\ne0\), the exact signed off-diagonal defect and
relative defect are

\[
\boxed{
\Delta_\triangle(x,y)
=F_\triangle(x,y)-T_x
=T_x\left(\frac{\Gamma(y+2)}{\Gamma(x+2)}-1\right),\qquad
\frac{F_\triangle(x,y)}{T_x}
=\frac{\Gamma(y+2)}{\Gamma(x+2)}}
\tag{156.6}
\]

Near \(y=x\), with \(y=x+\varepsilon\),

\[
\frac{\Delta_\triangle(x,x+\varepsilon)}{T_x}
=\varepsilon\psi(x+2)+O(\varepsilon^2)
\tag{156.7}
\]

Thus the graph's immediate departure from the diagonal is controlled by
the digamma slope. For integer \(y>x\) the ratio in (156.6) is the
product \(\prod_{k=x+2}^{y+1}k\), explaining the rapid visual growth.

The successor-prime information in the reduced rational

\[
\frac{x+1}{2(x-1)!}
\tag{156.8}
\]

does remain visible \textbf{inside that factor before the product is
reduced}: for \(x\ge2\), its reduced numerator is \(x+1\) when \(x+1\)
is prime and \(1\) when \(x+1\) is composite. It is not preserved as a
separate invariant of (156.2), because multiplying by \(x!\) introduces
the cancellation \(x!/(x-1)!=x\) and the whole expression becomes the
integer \(T_x\). The switch is real; the prime recognizer and triangular
polynomial are two different reductions of the same displayed
factorization.

\section{76. Dyadic triangular factorials and the reciprocal simplex
operator}\label{dyadic-triangular-factorials-and-the-reciprocal-simplex-operator}

The unrestricted equation \(T_m=n!\) remains open. Its dyadic branch,
however, closes completely.

\begin{quote}
\textbf{Theorem 76.1 (complete dyadic classification).} The positive
integer solutions of \[
2^{k-1}(2^k-1)=n!
\tag{76.1}
\] are exactly \[
\boxed{(k,n)=(1,1),(2,3),(4,5)}
\tag{76.2}
\]
\end{quote}

\textbf{Proof.} Taking \(2\)-adic valuations and using Legendre's
identity gives

\[
k-1=v_2(n!)=n-s_2(n)
\qquad
k=n-s_2(n)+1\le n
\tag{76.3}
\]

For \(n\ge3\), \(k\) must be even because \(3\mid n!\). The
lifting-the-exponent formula then gives

\[
v_3(2^k-1)=1+v_3(k/2)\le1+\log_3 n
\tag{76.4}
\]

On the other hand \(v_3(n!)\ge\lfloor n/3\rfloor\). Hence a solution
must satisfy

\[
\left\lfloor\frac n3\right\rfloor\le1+\log_3 n
\tag{76.5}
\]

The inequality fails for \(n=12,13,14\). Replacing \(n\) by \(n+3\)
raises its left side by one and its right side by \(\log_3(1+3/n)<1\),
so it fails for every \(n\ge12\). Exact evaluation of (76.3) for
\(1\le n\le11\) leaves precisely (76.2). \(\square\)

The square form is exactly equivalent:

\[
\boxed{(2^{k+1}-1)^2=8n!+1}
\tag{76.6}
\]

because the difference between the left side and \(1\) is eight times
the left side of (76.1). All three solutions, including the unit seed,
therefore give

\[
\boxed{8\cdot1!+1=3^2,\qquad
8\cdot3!+1=7^2,\qquad
8\cdot5!+1=31^2}
\tag{76.6a}
\]

They are \(1!=T_1\), \(3!=T_3\), and \(5!=T_{15}\). The next odd
factorial fails the square test:

\[
\boxed{200^2=40000<8\cdot7!+1=40321<40401=201^2}
\tag{76.6b}
\]

There is also a local arithmetic obstruction. Since
\(7!\equiv2\pmod{11}\), \(8\cdot7!+1\equiv6\pmod{11}\), whereas the
square residues modulo \(11\) are \(0,1,3,4,5,9\). Thus \(7!=5040\) is
not triangular; indeed \(T_{99}=4950<5040<T_{100}=5050\).

The odd-factorial step itself is exact: \((2r+1)!=2T_{2r}(2r-1)!\). At
\(r=3\) it gives \(7!/240=21=T_6\). This multiplication identity does
not preserve triangularity: the square test remains necessary. Theorem
76.1 closes only the subfamily \(m=2^k-1\). Neither that theorem nor the
rejection of \(7!\) proves that \(1!,3!,5!=1,6,120\) exhaust all
triangular factorials.

The reciprocal arithmetic connects the same triangular coordinates:

\[
\boxed{\frac1{7!}=\frac1{56\cdot90}=\frac1{4T_7T_9},\qquad
\frac{7!}{T_{44}}=\frac{5040}{990}=\frac{56}{11}=5+\frac1{11}}
\tag{76.6c}
\]

Here \(56=7\cdot8\), \(90=9\cdot10\), and \(990=11\cdot90=T_{44}\). Thus
the recurring decimal \(5.\overline{09}\) is the exact quotient
\(56/11\); it is not another solution of \(T_m=n!\). The two factors
\(56\) and \(90\) are also the real source coordinates \(x=s(s-1)\) at
\(s=8\) and \(s=10\) used in §147. Their analytic sign certificates have
the scope stated there. The same \(7!\) normalizes the fourth Widder row
in (32.4): \(W_4(0)/7!=\mu_3\). This is the factorial rung
normalization, distinct from asserting that \(7!\) itself is triangular.

The independent near-integer \(\pi^3\approx31.0062766803\) has a
convergent triangular owner, rather than a factorial square-test proof.
From (4.4f) and Basel, \(\pi^2=10-\tfrac18\sum_{n\ge1}T_n^{-3}\). Since
\(T_n\ge n^2/2\) and the tail is strictly positive,

\[
\sum_{n=1}^3T_n^{-3}=\frac{25}{24},\qquad
0<\sum_{n>N}T_n^{-3}<\frac8{5N^5},\qquad
\frac{1895}{192}-\frac1{1215}<\pi^2<\frac{1895}{192}
\tag{76.6d}
\]

The tail estimate follows by comparison with
\(8\int_N^\infty t^{-6}dt\). Exact rational comparison, with all
quantities positive, gives

\[
31^2<\left(\frac{1895}{192}-\frac1{1215}\right)^3
<\pi^6<\left(\frac{1895}{192}\right)^3
<\left(\frac{3101}{100}\right)^2
\tag{76.6e}
\]

Taking positive square roots proves \(31<\pi^3<31.01\). Meanwhile the
factor \(960/961\) in (PR.10d) comes from \(T_{15}=5!\) and
\(8T_{15}+1=31^2\). These are two exact, distinct connections to the
earlier triangular formulas. The denominator \(7!\) reappears in
\(W_4\), and the next Widder normalization grows by \(2T_{2k}\) in
(SW.4d).

The neighbouring value \(91\) has its own proved identity in (TDP.C6):
\(T_{13}=T_2^4+T_4=T_4^2-T_2^2=91=2T_9+1\). That unique nonzero
intersection uses the triangular value \(T_2=3\) and \(T_4=10\);
\(T_3=6\) belongs to a different index. This connects the reciprocal
scale \(90\) to the arithmetic branch already established in Reading
Volume §R1A.

The reciprocal digit mechanism is elementary:
\(1/5040=0.0001\overline{984126}\) has period six after a four-digit
prefix, whereas \(5040/990=5.\overline{09}\) has period two. Their
denominator orders explain these distinct repetitions.

\begin{figure}
\centering
\includegraphics{figures/dyadic_factorial_classification_v1_2.png}
\caption{The dyadic subfamily is completely classified by a valuation
tail and an eleven-case exact core.}
\end{figure}

The same packet contains a useful simplex identity. Put

\[
\Sigma_d(m)=\binom{m+d-1}{d}
\]

Then for \(d\ge2\),

\[
\boxed{
\frac1{\Sigma_d(m)}
=\frac d{d-1}
\left(
\frac1{\Sigma_{d-1}(m)}-
\frac1{\Sigma_{d-1}(m+1)}
\right)}
\tag{76.7}
\]

Consequently

\[
\boxed{
\sum_{m=1}^{N}\frac1{\Sigma_d(m)}
=\frac d{d-1}
\left(1-\frac1{\Sigma_{d-1}(N+1)}\right)
\qquad
\sum_{m\ge1}\frac1{\Sigma_d(m)}=\frac d{d-1}}
\tag{76.8}
\]

For \(d=2\), this is the familiar handshake identity

\[
\frac1{T_m}=2\left(\frac1m-\frac1{m+1}\right)
\]

The coefficient \(d/(d-1)\) in higher dimension is algebraic; it is not
an orientation multiplicity.

The value \(120\) has the exact face-count packet

\[
120=\Sigma_2(15)=\Sigma_3(8)=\Sigma_7(4)=\Sigma_{14}(3)=5!
\tag{76.9}
\]

but this is a cardinal resonance, not a prime selector or a zeta-zero
theorem.

\section{120. Scope of the repeated-square perfect-number
intersection}\label{scope-of-the-repeated-square-perfect-number-intersection}

Theorem TDP-A and Corollary TDP-B are the controlling square-index
statements. They use only the primitive identities

\[
T_n-T_{n-1}=n
\qquad
T_n+T_{n-1}=n^2
\]

so the cubic difference and the square-index identity are derived,
rather than presented as isolated coincidences. Exact integer replay
verifies (TDP.1)--(TDP.A4) for \(1\le n\le1000\) and the recurrence
through every staged ladder step.

The Euclid--Euler corollary is narrower than the full perfect-number
family. It says only that the special lower indices
\(1,3,15,255,\ldots\) meet the Mersenne-prime index set at \(3\). In
particular, \(T_7=28\) remains the next even perfect number and is not
discarded; it simply does not lie on the repeated-square ladder. No part
of PP176 changes an analytic, Löwner, source, zero-density, or RH grade.

\section{128. Scope of the Pell fourth-power
intersection}\label{scope-of-the-pell-fourth-power-intersection}

Theorem TDP-C supplies the exact arithmetic companion to the
square-index ladder. Its square-difference identity is not merely a list
of examples: (TDP.C2) is a complete Pell classification, with forward
recurrence (TDP.C3) and a strict descent proof. The extra fourth-power
lock is not a scan. Its exact factored equation is

\[
t(t-3)(t^2+3t+1)=0
\]

so \((T_2,T_4)=(3,10)\) is the unique nontrivial nonnegative
intersection. PP176A is an elementary side rail and changes no analytic
grade.

\section{162. Two Pell orbits and the dyadic
restriction}\label{two-pell-orbits-and-the-dyadic-restriction}

The finite-dyadic terminal of (139.4) gives
\((s_6,n_6,x_6)=(32,31,992)\). Its arithmetic companion is

\[
2T_{31}=992=T_{44}+2
\qquad T_{31}=496
\qquad T_{44}=990
\tag{162.1}
\]

The value \(496=2^4\cdot31\) belongs to the even-perfect-number family
because \(31=2^5-1\) is prime. This is separate from the repeated-square
subfamily of Theorem TDP-B.

For \(m\ge0,n\ge1\), the complete companion problem is

\[
2T_n=T_m+2
\iff X^2-2Y^2=-17
\qquad X=2m+1,\quad Y=2n+1
\tag{162.2}
\]

\Needspace{10\baselineskip}

\begin{quote}
\textbf{Theorem 162.1 --- two complete Pell orbits.} All solutions of
(162.2) form exactly two forward families, with seeds \((m,n)=(0,1)\)
and \((4,3)\), under \[m'=3m+4n+3,\qquad n'=2m+3n+2\] For strictly
positive \(m\), omit only \((0,1)\).
\end{quote}

The first members are

\begin{longtable}[]{@{}ll@{}}
\toprule\noalign{}
family & first four pairs \((m,n)\) \\
\midrule\noalign{}
\endhead
\bottomrule\noalign{}
\endlastfoot
first & \((0,1),(7,5),(44,31),(259,183)\) \\
second & \((4,3),(27,19),(160,113),(935,661)\) \\
\end{longtable}

\textbf{Proof.} Multiplication of \(X+Y\sqrt2\) by \(3+2\sqrt2\) gives
the displayed recurrence and preserves the norm. Conversely, every
positive solution has odd \(X,Y\) by reduction modulo 2 and 4. For
\(Y\ge9\), the norm equation implies \(4/3<X/Y<\sqrt2\) and \(X>Y\). The
inverse transformation

\[
(X,Y)\mapsto(3X-4Y,\,3Y-2X)
\tag{162.3}
\]

therefore keeps both coordinates positive and decreases \(Y\). Descent
terminates at odd \(Y\le7\). Checking \(Y=1,3,5,7\) leaves exactly
\((X,Y)=(1,3)\) and \((9,7)\), proving completeness. \(\square\)

The dyadic intersection is an additional restriction:

\[
Y=2^N-1
\qquad X^2=2^{2N+1}-2^{N+2}-15
\tag{162.4}
\]

The exact hits \(N=2,3,6\) give \(m=0,4,44\) respectively. The first
uses \(T_0=0\). A finite scan reported in the PP230 directive found no
others through \(N=100\); that scan is not a completeness proof. Modulo
5 excludes \(N\equiv0\pmod4\), and modulo 7 excludes \(N\equiv1\pmod3\).
Thus

\[
N\bmod12\in\{2,3,5,6,9,11\}
\tag{162.5}
\]

is necessary, not sufficient. An additional modulus \(9\) test sharpens
it. With \(u=2^N\), the right side of (162.4) is \(2u^2-4u-15\). For
\(N=0,1,\ldots,5\pmod6\), its residues modulo \(9\) are \(1,3,1,0,1,6\)
respectively. Since squares modulo \(9\) are \(0,1,4,7\), the classes
\(N\equiv1,5\pmod6\) are excluded. Thus

\[
\boxed{N\bmod12\in\{2,3,6,9\}}
\tag{162.5a}
\]

This remains only a necessary condition. Global completeness of the
dyadic intersection is open in the present work.

Two scalar relations have exact family explanations. For any offset
\(c\),

\[
2T_n=T_m+c
\iff (2m+1)^2-2(2n+1)^2=-(8c+1)
\tag{162.6}
\]

Thus the number 17 is generated by offset 2; this does not independently
select the modulus \(1009=992+17\). Also, with \(x=2T_n\),

\[
2j_n^2=8x+2
\qquad 2(j_n+1)^2-2j_n^2=4j_n+2
\tag{162.7}
\]

For the dyadic terminal, \(j_n=2^N-1\) and \(2(j_n+1)^2=2^{2N+1}\). At
\(N=6\) these endpoints are 7938 and 8192. This is a centered-coordinate
boundary identity. It adds no sign theorem to the source obstruction
(139.3).

The \(\sqrt2\) norm \(-17\) family is distinct from the \(\sqrt3\) norm
6 family in Theorem TDP-C. The complete arithmetic classification is
retained as a sidecar; no transport from it to the all-node Gamma/prime
inequality has been proved.

The completed derivative array determines signed Bernstein rows, their
exact Laguerre energy, and a limiting measure under a variation bound.
The proofs below establish the ordinary growth rate and identify the
remaining fixed-scale source estimate. That estimate remains open.

\Needspace{11\baselineskip}

\chapter{Technical chapter II. The fractional grid, reciprocal moments,
and decimal
coordinates}\label{technical-chapter-ii.-the-fractional-grid-reciprocal-moments-and-decimal-coordinates}

Native addresses retain their gcd and denominator. Reciprocal sums give
exact measures and scale laws; spacing experiments are assessed
separately.

\section{7. Native Triangular-Fractional
Grid}\label{native-triangular-fractional-grid}

The native cell and cumulative address are

\[
a(n,k)=n-1+\frac{k}{n},\qquad 1\le k\le n
\]

\[
\rho(n,k)=T_{n-1}+k
\]

Chamber \(n\) contains exactly \(n\) cells, and the number of cells
through chamber \(N\) is \(T_N\). Reduction of \(k/n\) is not allowed to
erase the ambient denominator: if \(X/n\) reduces to \(p/q\), then

\[
q=\frac n{\gcd(X,n)}
\]

and the gcd is part of the unreduced record.

The native center is

\[
c_n=n-\frac12
\]

With

\[
\eta=k-\frac n2,\qquad j=2\eta=2k-n
\]

one has

\[
\boxed{j\equiv n\pmod2}
\]

Adjacent chambers therefore populate complementary parity sheets. Even
chambers contain the midpoint; odd chambers have the two nearest cells

\[
c_n\pm\frac1{2n}
\]

The parity rule matches the indexing grammar of the two-channel
finite-Weil construction. It does not by itself prove an operator
intertwiner, equal channel mass, or positivity.

\section{8. The centered rational
collar}\label{the-centered-rational-collar}

The centered displacement of a native cell inside its unit chamber is

\[
c_{n,k}=\left|\frac{k}{n}-\frac12\right|
=\frac{|2k-n|}{2n}
\]

Define

\[
\mathcal C_\triangle=
\left\{
\frac{|2k-n|}{2n}:
n\ge2,\ 1\le k\le n,\ 0<|2k-n|<n
\right\}
\]

\subsection{Proposition 8.1 - Exact rational
coverage}\label{proposition-8.1---exact-rational-coverage}

\[
\boxed{\mathcal C_\triangle=\mathbb Q\cap(0,1/2)}
\]

\textbf{Proof.} Every displayed cell displacement is rational and lies
in the required interval. Conversely, let \(c=r/s\) with \(0<r/s<1/2\).
Set

\[
n=2s,\qquad k=s+2r
\]

Then \(1\le k\le n\) and

\[
\frac{|2k-n|}{2n}=\frac{4r}{4s}=\frac rs
\qquad\square
\]

This exact coverage becomes load-bearing in the rational RH witness
theorem of Part V.

\begin{figure}
\centering
\includegraphics{figures/tfg_rational_collar.png}
\caption{Centered native cells through chamber 40 in the rational
collar.}
\end{figure}

\section{9. Cyclic orders and denominator
moments}\label{cyclic-orders-and-denominator-moments}

Let \(q_{n,k}\) be the reduced denominator of \(a(n,k)\). Then

\[
\boxed{
q_{n,k}=\frac n{\gcd(k,n)}
=\operatorname{ord}_{\mathbb Z/n\mathbb Z}([k])
}
\]

For \(j\ge0\), define

\[
M_j(n)=\sum_{d\mid n}d^j\varphi(d)
\]

The chamber/cyclic-order identity is

\[
\boxed{
M_j(n)=\sum_{k=1}^{n}q_{n,k}^{j}
=\sum_{g\in\mathbb Z/n\mathbb Z}\operatorname{ord}(g)^j
}
\]

The functions \(M_j\) are multiplicative and, for prime powers,

\[
M_j(p^a)=1+(p-1)p^j\frac{p^{(j+1)a}-1}{p^{j+1}-1}
\]

Their Dirichlet series is

\[
\boxed{
F_j(s)=\sum_{n\ge1}\frac{M_j(n)}{n^s}
=\zeta(s)\frac{\zeta(s-j-1)}{\zeta(s-j)}
\qquad \Re s>j+2
}
\]

For \(j\ge1\), a nontrivial zero \(\rho\) of \(\zeta\) appears as an
inherited pole at \(s=j+\rho\). This is an exact pole transport, not a
restriction on \(\rho\).

Triangularizing a reduced cell \(A/b\) gives

\[
\operatorname{den}(T_{A/b})=
\begin{cases}
b^2,&b\text{ odd}\\
2b^2,&b\text{ even}
\end{cases}
\]

The local factor at \(2\) is essential. Its resulting even-moment rail
selects ambient dimensions \(d=4m+1\).

\begin{figure}
\centering
\includegraphics{figures/cyclic_order_moments.png}
\caption{Cyclic orders across native chambers and their second
denominator moment.}
\end{figure}

\section{84. Reduced-denominator rays, the prime sieve, and the trace
law}\label{reduced-denominator-rays-the-prime-sieve-and-the-trace-law}

The one-based chamber distinguishes two exact partitions:
reduced-denominator rays and prime-divisibility layers. For native cell
\((n,k)\),

\[
q_{n,k}=\frac{n}{\gcd(n,k)}
\tag{84.1}
\]

is the reduced denominator of \(k/n\), and \(q_{dn,dk}=q_{n,k}\). Every
cell factors transversely as

\[
(n,k)=m(q,a),\qquad m=\gcd(n,k),\qquad \gcd(q,a)=1
\tag{84.2}
\]

The non-coprime cells form the prime-layer sieve

\[
\{(n,k):\gcd(n,k)>1\}
=\bigcup_{p\ \mathrm{prime}}\{(n,k):p\mid n,\ p\mid k\}
\tag{84.3}
\]

At \(N=32\) the exact counts are \(204\) non-coprime cells, \(176\) in
the \(p=2,3\) sieve, and \(28\) structured higher-prime residuals. At
\(N=64\) the corresponding counts are \(820\), \(704\), and \(116\). The
magenta family in the audit is the distinct constant ray \(q_{2m,m}=2\).

\begin{figure}
\centering
\includegraphics{figures/figure48_pp103_ray_sieve_audit.png}
\caption{One-based \(N=32\) chamber: red/gray marks the \(p=2,3\)
divisibility sieve within the drawn region; magenta marks the
reduced-denominator \(q=2\) ray.}
\end{figure}

For reduced \(a/q\), the triangular fold has exact denominator

\[
\operatorname{den}\!\left(\frac{a(a+q)}{2q^2}\right)
=\eta(q)q^2
\qquad
\eta(q)=\begin{cases}1,&q\ \text{odd}\\2,&q\ \text{even}
\end{cases}
\tag{84.4}
\]

Thus \(Q_{m+1}=\eta(Q_m)Q_m^2\). If \(q\) is odd, \(Q_m=q^{2^m}\); if
\(q=2^eu\) with \(u\) odd and \(e\ge1\), then

\[
Q_m=2^{2^m(e+1)-1}u^{2^m}
\tag{84.5}
\]

The separate one-based \(N=64\) sidecar has maximum \(2\cdot64^2=8192\).

\begin{figure}
\centering
\includegraphics{figures/figure49_pp103_triangularized_sidecar.png}
\caption{The one-based \(N=64\) triangularized-denominator chamber has
maximum \(2\cdot64^2=8192\).}
\end{figure}

The row identity

\[
\sum_{k=1}^{n}F(q_{n,k})
=\sum_{d\mid n}\varphi(d)F(d)
\tag{84.6}
\]

has a finite Vandermonde--Hankel owner and a rational Padé generating
model. These are exact arithmetic statements. The map from this finite
arithmetic owner into the Löwner matrix (53.4), the moment--Casimir
operator, or the source kernel remains open.

\section{94. Reciprocal-triangular telescope and the forced dyadic
ladder}\label{reciprocal-triangular-telescope-and-the-forced-dyadic-ladder}

Write

\[
\mathsf T_n=\frac{n(n+1)}2
\]

Then

\[
\boxed{
\frac1{\mathsf T_n}=2\left(\frac1n-\frac1{n+1}\right)
\qquad
\sum_{n=a}^{\infty}\frac1{\mathsf T_n}=\frac2a
}
\tag{94.1}
\]

More strongly, for every \(a\ge1\) and \(j\ge0\),

\[
\boxed{
\sum_{n=a2^j}^{a2^{j+1}-1}\frac1{\mathsf T_n}
=\frac1{a2^j}
}
\tag{94.2}
\]

Thus consecutive dyadic blocks of the reciprocal-triangular series are
exactly the terms of the reciprocal powers-of-two series. If block
endpoints \(b_j\) are required to have mass \(1/b_j\), the telescope
forces \(b_{j+1}=2b_j\). Doubling is also the unique dilation for which
the removed block and surviving tail have equal mass. Its endpoint pair
\((2m,m)\) is exactly the singleton \(q=2\) ray.

The ordinary generating function is

\[
\sum_{n\ge1}\frac{z^n}{\mathsf T_n}
=2\left[1+\frac{1-z}{z}\log(1-z)\right]
\qquad |z|<1
\tag{94.3}
\]

\section{\texorpdfstring{95. Centered \(c\)-geometry, spectral product,
and the rational \(\pi\)
collar}{95. Centered c-geometry, spectral product, and the rational \textbackslash pi collar}}\label{centered-c-geometry-spectral-product-and-the-rational-pi-collar}

The centered identity

\[
\boxed{8\mathsf T_n+1=(2n+1)^2}
\tag{95.1}
\]

gives

\[
\mathsf T_n-c
=\frac12\left[\left(n+\frac12\right)^2-\left(\frac14+2c\right)\right]
\tag{95.2}
\]

At \(c=1/8\) the centered roots are \(\pm1/\sqrt2\). At \(c=1/2\) the
factors are the golden reciprocal pair

\[
n(n+1)-1=(n-\phi^{-1})(n+\phi)
\tag{95.3}
\]

The genus-zero triangular product is

\[
\boxed{
\mathcal E(c)=\prod_{n\ge1}\left(1-\frac{c}{\mathsf T_n}\right)
=-\frac{\cos\!\left(\frac\pi2\sqrt{1+8c}\right)}{2\pi c}
}
\tag{95.4}
\]

with the value at zero filled by continuity. Its logarithmic derivative
is meromorphic away from \(c=\mathsf T_n\) and generates every
reciprocal-power sum.

Tying \(c\) to the sidecar value \(a(q)=\eta(q)q^2\) gives \(a(2)=8\),
\(c_2=1/8\), and the first resolvent term \(1/(8\mathsf T_1-1)=1/7\).

The visible \(1/8\)--\(1/7\) trap is the proved rational collar

\[
\boxed{
\frac18<\pi-3<\frac17
}
\quad\Longleftrightarrow\quad
\frac{25}{8}<\pi<\frac{22}{7}
\tag{95.5}
\]

It yields

\[
\boxed{
8-\frac17<11-\pi<8-\frac18<8
<8+\frac18<\pi+5<8+\frac17
}
\tag{95.6}
\]

Here \(56/7=8\) with \(56=2\mathsf T_7\), while \(64/8=8\) with
\(64=8^2\). The same upper bound is \(7\pi/8<11/4=2.75\). These are
scalar rational inequalities, not zeta-zero constraints.

\section{95A. Catalan normalization, Beta products, and the odd sine
slice}\label{a.-catalan-normalization-beta-products-and-the-odd-sine-slice}

Let \(C_m=\binom{2m}{m}/(m+1)\) for \(m\ge0\). These are integers
because \(C_m=\binom{2m}{m}-\binom{2m}{m+1}\). Since
\(T_{2m+1}=(2m+1)(m+1)\), cancellation gives the all-index identity

\[
\boxed{(2m+1)!=T_{2m+1}(m!)^2C_m,\qquad
\frac{T_{2m+1}}{(2m+1)!}=\frac1{(m!)^2C_m}.}
\tag{95A.1}
\]

In particular \(180=(3!)^2C_3\). The first denominators are
\(1,1,8,180,8064,604800,68428800\). Every odd index belongs to this
slice. The successor-prime numerator transition \(p\to1\) belongs
specifically to indices \(p-1\to p\) for an odd prime \(p\); it is not a
property of every odd index. At \(8\to9\) both reduced numerators are
\(1\).

Put \(D=z\,d/dz\). Since \(D(z^n)=nz^n\), triangular coefficient
weighting is \(T_D=D(D+1)/2\). Termwise differentiation of the entire
sine series gives

\[
\boxed{\sum_{m\ge0}\frac{(-1)^mz^{2m+1}}{(m!)^2C_m}
=T_D\sin z=z\cos z-\frac{z^2}{2}\sin z.}
\tag{95A.2}
\]

This coefficient operator is different from the interval operator in
Section 12A. Neither operation supplies a completed-source sign.

\textbf{Proposition 95A.1 (a specified product of moment measures).}
With the Beta filter \(b_m=4^m(m!)^2/(2m+1)!\) of Section 104A,

\[
\boxed{b_m\frac{C_m}{4^m}=\frac1{T_{2m+1}}.}
\tag{95A.3}
\]

The product of independent variables with densities
\(2\pi^{-1}\sqrt{(1-u)/u}\) and \(1/(2\sqrt{1-v})\) on \((0,1)\) has
density \((t^{-1/2}-1)dt\) on \((0,1)\).

\textbf{Proof.} On \(0<y<4\) define the probability measures

\[
 d\eta(y)=\frac{dy}{\pi\sqrt{y(4-y)}},\qquad
 d\nu_C(y)=\frac{4-y}{2}\,d\eta(y),\qquad
 d\omega(y)=\frac y2\,d\eta(y).
\tag{95A.4}
\]

The substitution \(y=4\sin^2(\theta/2)\) gives \(d\eta=d\theta/\pi\).
Even-sine integration yields \(\int y^j d\eta=\binom{2j}{j}\), and hence

\[
\int y^j d\nu_C
=\frac12\left[4\binom{2j}{j}-\binom{2j+2}{j+1}\right]=C_j.
\]

The reflected-polynomial orthogonality measure is \(\omega\), not
\(\nu_C\); they are interchanged by \(y\mapsto4-y\). Scaling \(y=4u\)
gives the first claimed density and moments \(C_m/4^m\). The second
variable has moments \(b_m\). Independence and (95A.1) give (95A.3).
Finally

\[
\int_0^1t^m(t^{-1/2}-1)dt
=\frac2{2m+1}-\frac1{m+1}=\frac1{T_{2m+1}}.
\tag{95A.5}
\]

Equality of all polynomial integrals determines finite measures on a
compact interval, proving the product law. Equivalently the last density
is the pushforward of \(2(1-v)dv\) under \(t=v^2\). \(\square\)

These are identities between specified comparison measures. They do not
identify any of them with the folded-zero functional or reverse a failed
finite-matrix Beta implication. Section 155C uses \(\nu_C\) in a
positive comparison whose arithmetic defect is kept separately.

\section{96. The reciprocal-triangular Hausdorff
benchmark}\label{the-reciprocal-triangular-hausdorff-benchmark}

The telescope normalizes the probability measure

\[
d\nu_\triangle(t)
=\frac12\sum_{n\ge1}\frac1{\mathsf T_n}\,
\delta_{1/\mathsf T_n}(dt)
\tag{96.1}
\]

Its moments \(m_k=\frac12\sum_n\mathsf T_n^{-(k+1)}\) form a Hausdorff
moment sequence, and every finite Hankel matrix
\(H_N^\triangle=[m_{i+j}]_{i,j=0}^{N-1}\) is strictly positive definite.
Cauchy--Binet gives

\[
\boxed{
\det H_N^\triangle
=2^{-N}\sum_{n_1<\cdots<n_N}
\frac{\prod_{r<s}(\mathsf T_{n_s}-\mathsf T_{n_r})^2}
{\prod_r\mathsf T_{n_r}^{2N-1}}>0
}
\tag{96.2}
\]

For comparison define the synthetic moment matrix

\[
G_N^{(a)}(0)=\sum_\nu a_\nu^{-1}
v_N(a_\nu^{-1})v_N(a_\nu^{-1})^{\mathsf T},
\qquad v_N(t)=(1,t,\ldots,t^{N-1})^{\mathsf T}
\]

where the sums converge. At \(a_\nu=\mathsf T_\nu\),

\[
H_N^\triangle=\frac12G_N^{(a_\nu=\mathsf T_\nu)}(0)
\tag{96.3}
\]

The determinant factor is therefore \(2^{-N}\) and the net atom exponent
is exactly \(2N-1\). Same Gram grammar does not identify the synthetic
atoms with the actual folded parameters \(q_\rho=\rho(1-\rho)\). Off RH
those parameters may be nonreal, so the actual folded atomic functional
is not known to be a positive real measure. The first missing line is an
RH-free positive kernel or uniform Gram domination transporting
\(2\nu_\triangle\) to that functional.

More generally, for distinct positive \(a_\nu\) with
\(\sum_\nu a_\nu^{-1}<\infty\), extend the synthetic matrix to \(x\ge0\)
by

\[
\begin{aligned}
G_N^{(a)}(x)&=\left[\sum_\nu\frac{a_\nu}{(x+a_\nu)^{i+j+2}}
\right]_{i,j=0}^{N-1},\\
\det G_N^{(a)}(x)&=
\sum_{\nu_1<\cdots<\nu_N}
\frac{\prod_r a_{\nu_r}\prod_{r<s}(a_{\nu_s}-a_{\nu_r})^2}
{\prod_r(x+a_{\nu_r})^{2N}}
\end{aligned}
\tag{96.4}
\]

Finite Cauchy--Binet gives the identity, and entry convergence together
with the increasing nonnegative determinant sums gives the infinite
limit. There is strict positivity whenever at least \(N\) distinct atoms
occur. At \(x=0\) the exponent reduces to \(2N-1\), recovering
(96.2)--(96.3). The positive-real-atom hypothesis is essential for this
comparison.

\section{96A. Euler's constant, the reciprocal-triangular zeta series,
and a simplex
ladder}\label{a.-eulers-constant-the-reciprocal-triangular-zeta-series-and-a-simplex-ladder}

Reader §R4A states these results, and this section proves them.
Throughout, \(T_k=k(k+1)/2\), \(H_n=\sum_{k\le n}1/k\),
\(\psi=\Gamma'/\Gamma\), and \(\gamma_{\mathrm E}\) is Euler's constant.
Every statement is \PROVED{}. Their substance is classical {[}95--109,
112{]}. The replay \texttt{qa/audit\_gamma\_bridge.py} checks every
display, and its receipt is \texttt{qa/GAMMA\_BRIDGE.json}.

\subsection{96A.1 The mirror lemma and Ramanujan's
coefficients}\label{a.1-the-mirror-lemma-and-ramanujans-coefficients}

\begin{quote}
\textbf{Lemma 96A.1 (mirror parity).} Put \(u=n+\tfrac12\). As
\(n\to\infty\), \[
\psi(n+1)-\tfrac12\log\bigl(n(n+1)\bigr)\sim\sum_{m\ge1}\frac{c_m}{u^{2m}},
\qquad
c_m=\frac{(1-2^{1-2m})B_{2m}+2^{-2m}}{2m}.
\tag{96A.1}
\] The formal expansion is invariant under \(n\mapsto-1-n\),
equivalently under \(u\mapsto-u\), and it is a power series in
\(1/T_n\).
\end{quote}

\emph{Proof.} Differentiating the generalized Stirling series {[}2,
§5.11{]} gives, for fixed \(h\),
\(\psi(z+h)\sim\log z-\sum_{k\ge1}(-1)^kB_k(h)/(kz^k)\). Take \(z=u\)
and \(h=\tfrac12\), so that \(z+h=n+1\). Since
\(B_k(1-x)=(-1)^kB_k(x)\), every odd-order \(B_k(\tfrac12)\) vanishes,
and \(B_{2m}(\tfrac12)=-(1-2^{1-2m})B_{2m}\) {[}2, §24.4{]}. Hence
\(\psi(n+1)\sim\log u+\sum_{m\ge1}(1-2^{1-2m})B_{2m}/(2m\,u^{2m})\).
Also \(n(n+1)=u^2-\tfrac14\), so
\(\tfrac12\log(n(n+1))=\log u-\sum_{m\ge1}(2m)^{-1}(4u^2)^{-m}\).
Subtracting the two expansions gives (96A.1). Both sides are even in
\(u\), and \(u^2=2T_n+\tfrac14\). \(\square\)

The uncentered subtraction \(\psi(n+1)-\log n\) keeps the odd term
\(1/(2n)\). The centered subtractions \(\log(n+\tfrac12)\) {[}100,
120{]} and \(\tfrac12\log(n(n+1))\) are both even in \(u\), and that is
why both converge faster.

\textbf{Corollary 96A.2 (Ramanujan's expansion).} Since
\(\psi(n+1)=H_n-\gamma_{\mathrm E}\), (96A.1) together with
\(u^{-2}=(2T_n)^{-1}\bigl(1+(8T_n)^{-1}\bigr)^{-1}\) gives

\[
H_n\sim\gamma_{\mathrm E}+\tfrac12\log(2T_n)+\sum_{k\ge1}\frac{\varrho_k}{T_n^{k}},
\tag{96A.1a}
\]

with the coefficients

\[
\begin{array}{c|ccccc}
k & 1 & 2 & 3 & 4 & 5\\ \hline
\varrho_k & \frac1{12} & -\frac1{120} & \frac1{630} & -\frac1{1680} & \frac1{2310}\\[6pt]
k & 6 & 7 & 8 & 9 & 10\\ \hline
\varrho_k & -\frac{191}{360360} & \frac{29}{30030} & -\frac{2833}{1166880}
 & \frac{140051}{17459442} & -\frac{6525613}{193993800}
\end{array}
\]

These agree with the values in the literature {[}97, 99, 119{]}.
Matching (96A.1a) against the \(1/n\) expansion is an overdetermined
linear system, with twenty equations for ten unknowns; it is consistent
precisely because of the parity in Lemma 96A.1. At \(n=5\), \(10\) and
\(40\), truncation after \(\varrho_6\) leaves an error smaller than
\(|\varrho_7|\,T_n^{-7}\). The coincidences \(\varrho_1=-\zeta(-1)\) and
\(\varrho_2=-\zeta(-3)\) do not continue: \(\varrho_3=\tfrac1{630}\),
while \(-\zeta(-5)=\tfrac1{252}\). \(\square\)

\subsection{96A.2 The reciprocal-triangular zeta
series}\label{a.2-the-reciprocal-triangular-zeta-series}

\begin{quote}
\textbf{Theorem 96A.3.} \[
\sum_{k\ge2}\frac{\zeta(k)}{T_k}=\log(2\pi)-\gamma_{\mathrm E}
=1.26066140150781262295\ldots
\tag{96A.2}
\]
\end{quote}

\emph{First proof (Raabe).} For \(|w|<1\),
\(\log\Gamma(1-w)=\gamma_{\mathrm E}w+\sum_{k\ge2}\zeta(k)w^k/k\) {[}2,
§5.7{]}. For \(0<w<1\) every term is positive, so monotone convergence
allows termwise integration over \([0,1]\). The left side is integrable
at \(w=1\), since \(\log\Gamma(1-w)\sim-\log(1-w)\). Use
\(\int_0^1w^k/k\,dw=1/(k(k+1))=1/(2T_k)\) and Raabe's integral
\(\int_0^1\log\Gamma(v)\,dv=\tfrac12\log(2\pi)\) {[}95{]}. They give
\(\tfrac12\log(2\pi)=\tfrac12\gamma_{\mathrm E}+\tfrac12\sum_{k\ge2}\zeta(k)/T_k\).
\(\square\)

\emph{Second proof (moment form).} By (4.5),
\(1/T_k=\int_0^1u^{k-1}\,2(1-u)\,du\), and
\(\sum_{k\ge2}\zeta(k)u^{k-1}=-\gamma_{\mathrm E}-\psi(1-u)\) for
\(0\le u<1\). All terms are nonnegative, and \(2(1-u)\psi(1-u)\) stays
bounded as \(u\to1\). So the sum equals
\(-\gamma_{\mathrm E}-2\int_0^1v\,\psi(v)\,dv\). Integrating by parts
against Raabe's integral gives
\(\int_0^1v\psi(v)\,dv=-\tfrac12\log(2\pi)\). \(\square\)

\emph{Third proof (triangular telescope and Stirling).} Evaluate (94.3)
at \(z=1/m\) and remove the \(n=1\) term. For \(m\ge2\) this gives

\[
\sum_{k\ge2}\frac{m^{-k}}{T_k}=2-\frac1m+2(m-1)\log\Bigl(1-\frac1m\Bigr),
\tag{96A.2a}
\]

and at \(m=1\) the value is \(1\), by Mengoli's telescope {[}91{]}. Sum
over \(m=1,\ldots,M\) and write \(a_m=m\log m\). The logarithmic terms
then telescope to \(2(\log M!-M\log M)\). For the truncated zeta
function \(\zeta_M(k)=\sum_{m\le M}m^{-k}\) this gives

\[
\boxed{\sum_{k\ge2}\frac{\zeta_M(k)}{T_k}=2M-H_M-2\log\frac{M^M}{M!}.}
\tag{96A.3}
\]

By Stirling's formula {[}94{]},
\(2\log M!-2M\log M=-2M+\log(2\pi M)+O(1/M)\). The right side of (96A.3)
is therefore \(\log(2\pi)-(H_M-\log M)+O(1/M)\), which tends to
\(\log(2\pi)-\gamma_{\mathrm E}\); the error is
\(-\frac1{3M}+\frac1{12M^2}+O(M^{-3})\). \(\square\)

\begin{quote}
\textbf{Corollary 96A.4 (unit-interval masses and the norm).} For
\(m\ge1\), \[
\sum_{k\ge2}\frac{m^{-k}}{T_k}=\int_{m-1}^{m}\Bigl(\frac{\{x\}}x\Bigr)^2dx,
\tag{96A.4}
\] and \[
\log(2\pi)-\gamma_{\mathrm E}
=\int_0^\infty\Bigl(\frac{\{x\}}x\Bigr)^2dx
=\int_0^\infty\{1/y\}^2\,dy
=\frac1{2\pi}\int_{-\infty}^{\infty}\frac{|\zeta(\tfrac12+it)|^2}{\tfrac14+t^2}\,dt.
\tag{96A.5}
\]
\end{quote}

\emph{Proof.} On \([m-1,m]\) we have \(\{x\}=x-(m-1)\), and

\[
\int_{m-1}^{m}\Bigl(1-\frac{m-1}{x}\Bigr)^2dx
=1-2(m-1)\log\frac m{m-1}+\frac{m-1}m,
\]

which equals (96A.2a). For \(m=1\) both sides equal \(1\). Summing over
\(m\) and using Theorem 96A.3 gives the first equality in (96A.5), and
the substitution \(y=1/x\) gives the second. For \(0<\sigma<1\),
\(\zeta(s)/s=-\int_0^\infty\{1/y\}\,y^{s-1}\,dy\) (from {[}3, (2.1.5){]}
after \(x=1/y\)), and \(\{1/y\}\in L^2(0,\infty)\). The
Mellin--Plancherel theorem {[}117{]} on \(\sigma=\tfrac12\) then gives
the last equality, because \(|s|^2=\tfrac14+t^2\) there. \(\square\)

The same constant is the squared Hardy norm (155G.6). Under \(z=1-1/s\),
the unit circle maps onto \(\sigma=\tfrac12\) with
\(|dz|=dt/(\tfrac14+t^2)\), and \(|\mathcal B(z)|=|\zeta(s)|\) on it. So
\(\sum_n(b_n^{\rm den})^2\) is exactly the last integral in (96A.5). The
fractional-part function in (96A.5) is the one in the Nyman--Beurling
closure problem {[}106--109{]}. A numerical evaluation of the
critical-line integral, up to \(|t|=200\) plus the second-moment tail,
agrees with \(\log(2\pi)-\gamma_{\mathrm E}\) to \(10^{-4}\). That check
is a sanity check, not a certificate; the proof above is the
certificate.

\subsection{96A.3 The simplex ladder}\label{a.3-the-simplex-ladder}

\begin{quote}
\textbf{Theorem 96A.5.} For \(d\ge2\), \[
\sum_{k\ge2}\frac{\zeta(k)}{\binom{k+d-1}d}
=d(d-1)\int_0^1u^{d-2}\log\Gamma(u)\,du-\gamma_{\mathrm E},
\tag{96A.6}
\] and, for \(m\ge1\), with \((m)_{j-1}=m(m-1)\cdots(m-j+2)\), \[
\int_0^1u^m\log\Gamma(u)\,du=\frac{\log(2\pi)}{2(m+1)}
-\sum_{j=1}^m\frac{(m)_{j-1}(-1)^j}{j!}\Bigl[\zeta'(-j)+H_j\,\zeta(-j)\Bigr].
\tag{96A.7}
\]
\end{quote}

\emph{Proof.} The Beta integral gives
\(1/\binom{k+d-1}d=d\,B(k,d)=d\int_0^1u^{k-1}(1-u)^{d-1}\,du\). As in
the second proof of Theorem 96A.3, and since
\(d\int_0^1(1-u)^{d-1}du=1\), the sum equals
\(-\gamma_{\mathrm E}-d\int_0^1v^{d-1}\psi(v)\,dv\). Integration by
parts gives
\(\int_0^1v^{d-1}\psi(v)\,dv=-(d-1)\int_0^1v^{d-2}\log\Gamma(v)\,dv\);
the boundary terms vanish for \(d\ge2\). This proves (96A.6).

For (96A.7), start from Lerch's formula
\(\log\Gamma(u)=\zeta'(0,u)+\tfrac12\log(2\pi)\) {[}2, §25.11{]}. Apply
\(\partial_u\zeta(s-1,u)=-(s-1)\zeta(s,u)\) repeatedly, integrating by
parts each time, and continue analytically. This gives the Hurwitz
moments
\(\int_0^1u^m\zeta(s,u)\,du=-\sum_{j=1}^m(m)_{j-1}\,\zeta(s-j)/\prod_{i=1}^{j}(s-i)\).
Differentiating at \(s=0\) gives (96A.7); compare {[}104{]}. \(\square\)

\textbf{Corollary 96A.6 (first rungs).} With
\(\log A=\tfrac1{12}-\zeta'(-1)\) {[}105{]} and
\(\zeta'(-2)=-\zeta(3)/(4\pi^2)\),

\[
\begin{aligned}
d=2:&\quad \log(2\pi)-\gamma_{\mathrm E},\\
d=3:&\quad \tfrac32\log(2\pi)-6\log A-\gamma_{\mathrm E},\\
d=4:&\quad 2\log(2\pi)-12\log A+\frac{3\zeta(3)}{\pi^2}-\gamma_{\mathrm E}.
\end{aligned}
\tag{96A.8}
\]

Written in the constants \(\log(2\pi)\), \(\gamma_{\mathrm E}\) and
\(\zeta'(-j)\), every rung has \(\gamma_{\mathrm E}\)-coefficient
\(-1\), which comes only from the term isolated in the proof of (96A.6),
and \(\log(2\pi)\)-coefficient \(d/2\), from (96A.7). Rung \(d\)
introduces \(\zeta'(2-d)\). At even arguments, differentiating the
functional equation {[}2, §25.4{]} at the trivial zero \(-2j\) gives
\(\zeta'(-2j)=(-1)^j(2j)!\,\zeta(2j+1)/\bigl(2(2\pi)^{2j}\bigr)\), so
odd zeta values enter exactly there. The replay evaluates
\(d=2,\ldots,9\) in three independent ways (the series, quadrature of
the moment, and (96A.7)), and they agree to 35 digits. With every
\(\zeta(k)\) replaced by \(1\), the left side of (96A.6) becomes the
reciprocal simplex telescope of §76, which sums to \(d/(d-1)\) from
\(k=1\). \(\square\)

\subsection{96A.4 Where Euler's constant meets the
zeros}\label{a.4-where-eulers-constant-meets-the-zeros}

\begin{quote}
\textbf{Corollary 96A.7.}
\(\lambda_1=\sum_{[\rho]}m_\rho/q_\rho=1+\tfrac12\gamma_{\mathrm E}-\tfrac12\log(4\pi)\),
and \[
\sum_{k\ge2}\frac{\zeta(k)}{T_k}+2\lambda_1=2-\log2=\sum_{k\ge1}\frac1{T_k}-\log2.
\tag{96A.9}
\]
\end{quote}

\emph{Proof.} Davenport's constant
\(B=-\tfrac12\gamma_{\mathrm E}-1+\tfrac12\log(4\pi)\) in the Hadamard
product of \(\xi\) satisfies \(B=-\sum_\rho\Re\rho^{-1}\) {[}110, 112,
Ch. 12{]}. The orbit pairing of (9.7) gives
\(\sum_{[\rho]}m_\rho/q_\rho=\sum_\rho\rho^{-1}\); equivalently,
\(\lambda_1=-\xi'(0)/\xi(0)=\xi'(1)/\xi(1)\). Add twice this to (96A.2).
\(\square\)

The replay confirms \(\xi'(1)/\xi(1)=\lambda_1\) to fifteen digits. It
also recovers \(\lambda_1\) to within \(10^{-5}\) from the first hundred
zeros and the smooth Riemann--von Mangoldt tail {[}111{]}. (96A.9) is
the triangular zero-sum formula for \(\gamma_{\mathrm E}\), written in
the folded coordinate. It is an identity, not a positivity statement.

\subsection{96A.5 Recomputed enclosures for
§R12}\label{a.5-recomputed-enclosures-for-r12}

These enclosures are computed in ball arithmetic: the zeta-zero routine
returns rigorous enclosures of the first two zeros, and every other
input is an exact constant. Every value below is \CERT{}.

\begin{longtable}[]{@{}
  >{\raggedright\arraybackslash}p{(\columnwidth - 2\tabcolsep) * \real{0.5000}}
  >{\raggedright\arraybackslash}p{(\columnwidth - 2\tabcolsep) * \real{0.5000}}@{}}
\toprule\noalign{}
\begin{minipage}[b]{\linewidth}\raggedright
quantity
\end{minipage} & \begin{minipage}[b]{\linewidth}\raggedright
enclosure
\end{minipage} \\
\midrule\noalign{}
\endhead
\bottomrule\noalign{}
\endlastfoot
\(\lambda_1\) &
\(0.0230957089661210338143102479065\pm4.8\times10^{-33}\) \\
\(q_1=\tfrac14+\gamma_1^2\) &
\(200.040454832386859463365185961\pm4.1\times10^{-28}\) \\
\(q_2=\tfrac14+\gamma_2^2\) &
\(442.176150574082490320460688768\pm2.1\times10^{-28}\) \\
\(q_1/q_2\) & \(0.45239991929160355477\pm2.7\times10^{-21}\) \\
\(C_0=q_2(\lambda_1-1/q_1)\) &
\(8.001938046160201229441097\pm9.7\times10^{-26}\) \\
\end{longtable}

The verified height \(H=3{,}000{,}175{,}332{,}800\) of §R13 enters
through (12.7a), which gives \(C_0\le C\le C_0+q_2\lambda_1/(2H^2)\).
Hence \(C=8.00193804616020122944\ldots\), and
\(C(q_1/q_2)^3=0.7409053672828\ldots<1\). These recomputed enclosures
replace the carried decimals of §R12.

An independent mpmath evaluation agrees with \(q_1\) and \(q_2\) to
twelve decimals. These enclosures were recomputed by the v5.0 build; the
receipt is \texttt{qa/GAMMA\_BRIDGE.json} and is bundled.

\section{97. Dyadic Gram amplification and the triangular exponent
law}\label{dyadic-gram-amplification-and-the-triangular-exponent-law}

Use \(v_N(t)=(1,t,\ldots,t^{N-1})^{\mathsf T}\). For a dyadic block
beginning at \(m\), set

\[
C_{m,N}=\sum_{n=m}^{2m-1}\frac1{2\mathsf T_n}
v_N(1/\mathsf T_n)v_N(1/\mathsf T_n)^{\mathsf T}
\tag{97.1}
\]

Its rank is \(\min(m,N)\), and it is positive definite exactly when
\(m\ge N\). After the exact congruence with
\(D_m=\operatorname{diag}(1,m^2,\ldots,m^{2(N-1)})\),

\[
C_{m,N}=\frac1{2m}D_m^{-1}G_{m,N}D_m^{-1}
\tag{97.2}
\]

Therefore

\[
\boxed{
\det C_{m,N}=2^{-N}m^{-\mathsf T_{2N-1}}\det G_{m,N}
\qquad
\mathsf T_{2N-1}=N(2N-1)
}
\tag{97.3}
\]

The ordered block eigenvalues have the exact scale ladder

\[
\boxed{
\lambda_{r+1}(C_{m,N})=\Theta_N(m^{-(4r+1)})
\qquad r=0,\ldots,N-1,\qquad m\ge2N
}
\tag{97.4}
\]

The exponents \(1,5,9,\ldots,4N-3\) sum to \(N(2N-1)=\mathsf T_{2N-1}\).
In particular, rank four has determinant exponent \(\mathsf T_7=28\);
under one asymptotic doubling its block determinant ratio tends to
\(2^{-28}\).

Define the normalized full tail

\[
R_{m,N}=mD_m\left[\sum_{n\ge m}\frac{1}{2\mathsf T_n}
v_N(1/\mathsf T_n)v_N(1/\mathsf T_n)^{\mathsf T}\right]D_m
\]

Splitting that tail at \(2m\) gives the exact renormalization

\[
R_{m,N}=\frac12G_{m,N}+\frac12ER_{2m,N}E
\qquad
E=\operatorname{diag}(1,4^{-1},\ldots,4^{-(N-1)})
\tag{97.5}
\]

The normalized cores converge at an explicit \(O_N(m^{-2})\) rate to the
compact Gram measure \(dy/\sqrt{2y}\) on \([1/2,2]\). This is the solved
synthetic analogue of the moving normalized source core. PP132 package
SHA-256:
\texttt{6861f7a7e4d2d695}\allowbreak \texttt{41d953ead56a078c}\allowbreak \texttt{5f0790787696221a}\allowbreak \texttt{5f59cf720875add6}.

\section{98. Decimal, cyclotomic, and divisor-ray
exponents}\label{decimal-cyclotomic-and-divisor-ray-exponents}

If

\[
\mathsf T_n=2^{\alpha_n}5^{\beta_n}m_n
\qquad(m_n,10)=1
\]

then the decimal preperiod of \(1/\mathsf T_n\) is
\(\max(\alpha_n,\beta_n)\) and its period is
\(\operatorname{ord}_{m_n}(10)\) when \(m_n>1\); for \(m_n=1\) it
terminates and the period is recorded as zero. The recurring remainders
in the nonterminating case form a coset \(r_n\langle10\rangle\) in the
primitive \(m_n\)-ray labels, covering all labels exactly when
\(\operatorname{ord}_{m_n}(10)=\varphi(m_n)\).

For \(n=7\),

\[
\frac1{28}=0.03\overline{571428}
\qquad(\text{prefix length},\text{period length})=(2,6)
\tag{98.1}
\]

Independently,

\[
\boxed{
2\mathsf T_n+1=\Phi_3(n)
\qquad
\operatorname{ord}_{2\mathsf T_n+1}(n)=3\quad(n>1)
}
\tag{98.2}
\]

At \(n=7\), \(57=111_7\) and \(1\mapsto7\mapsto49\mapsto1\pmod{57}\).
This order-three base-seven cycle is distinct from
\(\operatorname{ord}_{57}(10)=18\) and \(\operatorname{ord}_7(10)=6\). A
universal 57-prefix digit rule is false; \(n=6\) already gives
\(2\mathsf T_6+1=43\) but \(1/\mathsf T_6=0.\overline{047619}\).

There is nevertheless a complete radix exponent ladder. For every radix
\(b\ge2\), integer \(h\ge1\), and prime \(\ell\), put
\(Q=\Phi_\ell(b^h)\). Then

\[
\boxed{
\operatorname{ord}_{Q}(b)=\ell h
\qquad
\frac1Q
=0.\overline{
\underbrace{00\cdots0}_{(\ell-1)h}
\underbrace{(b-1)\cdots(b-1)}_{h}} _b
}
\tag{98.3}
\]

For \(\ell=3\), \(Q=2\mathsf T_{b^h}+1\) and the exact exponent sequence
is \(3h\). The complementary identity

\[
\Phi_3(n)=2\mathsf T_n+1=\Phi_6(n+1)
\]

gives orders \(3\) and \(6\) for the bases \(n\) and \(n+1\) modulo the
same number. At \(n=7\), the centered triad is \(55,56,57\); modulo
\(57\), the subgroups generated by \(7\), \(8\), and \(10\) have orders
\(3\), \(6\), and \(18\), with
\(\langle8\rangle\cap\langle10\rangle=\langle7\rangle\) and together the
larger two generate all \(36\) primitive labels. These exact arithmetic
orbits do not select the \(q=57\) determinant margin; the certificate
geometry does.

The exact reciprocal-triangular divisor trace is

\[
\sum_{k=1}^{n}\frac1{\mathsf T_{q_{n,k}}}
=\sum_{d\mid n}\frac{2\varphi(d)}{d(d+1)}
\tag{98.4}
\]

An elementary base-\(100\) seed illustrates this cyclic viewpoint: the
coefficients \(14,28,56,\ldots\) generate \(1/7\) before carrying, and
\(1/T_7=\tfrac14(1/7)\) transports its periodic tail. The multiplicative
order determines the period after the factors of the base are removed.

The arithmetic trace also explains the visible bands in the chamber
below. For each reduced denominator \(q\), the rays have slopes
\(k/n=a/q\) with \(\gcd(a,q)=1\). Each ray contains
\(\lfloor N/q\rfloor\) native cells, so

\[
\boxed{C_N(q)=\varphi(q)\lfloor N/q\rfloor}
\tag{98.5}
\]

Here \(C_N(q)\) counts cells in the whole \(q\)-family, including
\(q=1\) on the boundary \(k=n\). The low-denominator bands at \(N=256\)
are:

\begin{longtable}[]{@{}
  >{\raggedleft\arraybackslash}p{(\columnwidth - 6\tabcolsep) * \real{0.2667}}
  >{\raggedright\arraybackslash}p{(\columnwidth - 6\tabcolsep) * \real{0.2000}}
  >{\raggedleft\arraybackslash}p{(\columnwidth - 6\tabcolsep) * \real{0.2667}}
  >{\raggedleft\arraybackslash}p{(\columnwidth - 6\tabcolsep) * \real{0.2667}}@{}}
\toprule\noalign{}
\begin{minipage}[b]{\linewidth}\raggedleft
\(q\)
\end{minipage} & \begin{minipage}[b]{\linewidth}\raggedright
Reduced ray slopes \(k/n\)
\end{minipage} & \begin{minipage}[b]{\linewidth}\raggedleft
Colour value \(\eta(q)q^2\)
\end{minipage} & \begin{minipage}[b]{\linewidth}\raggedleft
Family cells
\end{minipage} \\
\midrule\noalign{}
\endhead
\bottomrule\noalign{}
\endlastfoot
1 & \(1\) (boundary diagonal) & 1 & 256 \\
2 & \(1/2\) & 8 & 128 \\
3 & \(1/3,2/3\) & 9 & 170 \\
4 & \(1/4,3/4\) & 32 & 128 \\
5 & \(1/5,2/5,3/5,4/5\) & 25 & 204 \\
6 & \(1/6,5/6\) & 72 & 84 \\
7 & \(1/7,2/7,\ldots,6/7\) & 49 & 216 \\
8 & \(1/8,3/8,5/8,7/8\) & 128 & 128 \\
\end{longtable}

The colour is constant along each listed ray. It is not monotone in
\(q\): the even denominator \(q=4\) has colour value \(32\), exceeding
the value \(25\) at \(q=5\). Horizontal and vertical crosshatching
additionally records the gcd/divisibility layers of (84.3); it is not a
second collection of constant-\(q\) rays.

\begin{figure}
\centering
\includegraphics{figures/figure51_ray_sieve_audit_N256_q2.png}
\caption{Extended one-based ray--sieve audit with the singleton \(q=2\)
family highlighted.}
\end{figure}

\begin{figure}
\centering
\includegraphics{figures/figure50_triangularized_denominator_N256.png}
\caption{Triangularized-denominator chamber at \(N=256\), coloured by
\(\eta(q)q^2\). The boundary has \(q=1\); the interior includes the
\(q=2,3,4,5,6,7,8\) ray families listed above and all higher reduced
denominators. Divisibility crosshatching overlays the rational-slope
structure. The exact maximum is \(2\cdot256^2=131072\).}
\end{figure}

\section{167. Denominator iteration, decimal periods, and useful
approximations}\label{denominator-iteration-decimal-periods-and-useful-approximations}

Triangular values at rational arguments can be tested on exact reduced
fractions. If \(\gcd(a,q)=1\), the reduced denominator of \(T_{a/q}\) is

\[
Q_1=\eta(q)q^2,\qquad
\eta(q)=\begin{cases}1,&q\text{ odd}\\2,&q\text{ even}
\end{cases}
\tag{167.1}
\]

Indeed both \(a\) and \(a+q\) are coprime to \(q\). For odd \(q\) their
product is even and cancels the factor \(2\); for even \(q\) it is odd.
Repeating the triangular map on the reduced value therefore gives

\[
Q_0=q,\qquad Q_{m+1}=\eta(Q_m)Q_m^2
\qquad
Q_m=\begin{cases}q^{2^m},&q\text{ odd}\\
(2q)^{2^m}/2,&q\text{ even}
\end{cases}
\tag{167.2}
\]

This is an exact sequence, not a fit to digit patterns. It has the
linear logarithmic reading

\[
\log_{10}Q_m=
\begin{cases}2^m\log_{10}q,&q\text{ odd}\\
2^m\log_{10}(2q)-\log_{10}2,&q\text{ even}
\end{cases}
\tag{167.3}
\]

For a reduced reciprocal with denominator \(Q=2^\alpha5^\beta d\),
\(\gcd(d,10)=1\), the nonrepeating prefix has length
\(\max\{\alpha,\beta\}\). If \(d>1\) the period is
\(\operatorname{ord}_d(10)\); if \(d=1\) the expansion terminates and
the period is recorded as zero. The following are exact integer results:

\begin{longtable}[]{@{}
  >{\raggedright\arraybackslash}p{(\columnwidth - 8\tabcolsep) * \real{0.2000}}
  >{\raggedright\arraybackslash}p{(\columnwidth - 8\tabcolsep) * \real{0.2000}}
  >{\raggedright\arraybackslash}p{(\columnwidth - 8\tabcolsep) * \real{0.2000}}
  >{\raggedright\arraybackslash}p{(\columnwidth - 8\tabcolsep) * \real{0.2000}}
  >{\raggedright\arraybackslash}p{(\columnwidth - 8\tabcolsep) * \real{0.2000}}@{}}
\toprule\noalign{}
\begin{minipage}[b]{\linewidth}\raggedright
initial \(q\)
\end{minipage} & \begin{minipage}[b]{\linewidth}\raggedright
iteration \(m\)
\end{minipage} & \begin{minipage}[b]{\linewidth}\raggedright
\(Q_m\)
\end{minipage} & \begin{minipage}[b]{\linewidth}\raggedright
prefix length of \(1/Q_m\)
\end{minipage} & \begin{minipage}[b]{\linewidth}\raggedright
period
\end{minipage} \\
\midrule\noalign{}
\endhead
\bottomrule\noalign{}
\endlastfoot
\(7\) & \(0,1,2\) & \(7,49,2401\) & \(0,0,0\) & \(6,42,2058\) \\
\(31\) & \(0,1,2\) & \(31,961,923521\) & \(0,0,0\) &
\(15,465,446865\) \\
\(44\) & \(0,1,2\) & \(44,3872,29984768\) & \(2,5,11\) &
\(2,22,2662\) \\
\end{longtable}

For these three rows the period laws are, respectively,
\(6\cdot7^{2^m-1}\), \(15\cdot31^{2^m-1}\), and \(2\cdot11^{2^m-1}\).
One proof checks the initial orders \(6,15,2\) and that \(p^2\) does not
divide \(10^\ell-1\) for \(p=7,31,11\) with the respective order
\(\ell\). Binomial expansion then lifts the order by a factor of \(p\)
at each higher prime power. For \(q=44\), the power of \(2\) is
\(3\cdot2^m-1\), giving the printed prefix lengths. The initial orders
and the non-Wieferich lifting checks are included in the exact
arithmetic receipt of this Dossier.

\subsection{Finite approximations can retain exact remainder
information}\label{finite-approximations-can-retain-exact-remainder-information}

The familiar \(1/7\) block expansion and the reciprocal of \(T_{44}\)
have particularly transparent tails. For \(J\ge1\), define

\[
\begin{aligned}
P_J&=\sum_{j=0}^{J-1}\frac{14\cdot2^j}{100^{j+1}}
&\frac17-P_J&=\frac1{7\cdot50^J}\\
Q_J^*&=\frac1{10}\sum_{j=1}^J100^{-j}
&\frac1{990}-Q_J^*&=\frac1{990\cdot100^J}
\end{aligned}
\tag{167.4}
\]

These are geometric-series identities. The logarithms of the relative
errors are \(-J\log_{10}50\) and \(-2J\). Thus a decimal approximation
can be retained with an exact convergence rate, rather than discarded or
mistaken for equality. In particular \(1/990=0.0\overline{01}\). The
symbol \(Q_J^*\) here is a partial sum, not the denominator sequence
\(Q_m\) of (167.2).

The nearby empirical marker \(11/4\) also has an exact triangular
address:

\[
\frac{11}{4}=\frac{T_4+1}{4},\qquad
T_{11/4}=\frac{165}{32},\qquad
\frac1{T_{11/4}}=\frac{32}{165}=\frac{192}{990}
=0.1\overline{93}
\tag{167.5}
\]

Thus \(T_{44}=192T_{11/4}\), and the denominator \(32=2\cdot4^2\) agrees
with (167.1). Its proximity to another marker has a precise error:

\[
\frac{11}{4}-\frac{7\pi}{8}
=\frac{22-7\pi}{8}
=\frac78\left(\frac{22}{7}-\pi\right)
\approx0.001106428109
\tag{167.6}
\]

This identifies the classical rational approximation to \(\pi\) inside
the marker comparison. Whether the marker predicts a gap statistic is
the separate empirical question in §168.

\section{168. Counting, unfolded gaps, and their triangular
coordinate}\label{counting-unfolded-gaps-and-their-triangular-coordinate}

The arithmetic identities above offer an exact change of coordinates for
zero-spacing measurements. They do not remove the fluctuating term in
zero counting. The smooth part of the Riemann--von Mangoldt count
{[}2{]} is

\[
N_0(t)=\frac{t}{2\pi}\left(\log\frac{t}{2\pi}-1\right)+\frac78,
\qquad N_0'(t)=\frac1{2\pi}\log\frac{t}{2\pi}
\tag{168.2}
\]

This normalization already connects Gamma, \(\pi\), and the logarithm.
With the logarithm of Gamma chosen continuously from \(t=0\),

\[
\begin{aligned}
\vartheta(t)&=\Im\log\Gamma\!\left(\frac14+\frac{it}{2}\right)
-\frac t2\log\pi,\\
\vartheta(t)&=\frac t2\log\frac{t}{2\pi}-\frac t2-\frac\pi8+\frac1{48t}+O(t^{-3}),
\qquad t\to+\infty
\end{aligned}
\tag{168.3}
\]

Thus \(1+\vartheta(t)/\pi=N_0(t)+1/(48\pi t)+O(t^{-3})\). The
coefficient follows from the shifted Stirling expansion {[}2{]}:
\(B_2(1/4)=(1/4)^2-1/4+1/6=-1/48\), and its imaginary contribution at
\(it/2\) is \(-B_2(1/4)/t\). This is a quarter-argument Gamma
coefficient, not a new estimate for the fluctuating zeta argument. The
actual zero count \(N(t)\) also contains the argument fluctuation of
zeta. The counting formula is unconditional: counting ordinates does not
determine the real parts of the zeros. Applying the polynomial \(T_x\)
gives

\[
\begin{aligned}
\frac{d}{dt}T_{N_0(t)}&=\left(N_0(t)+\frac12\right)N_0'(t),\\
T_{N(t)}-T_{N_0(t)}
&=\left(N_0(t)+\frac12\right)(N(t)-N_0(t))
+\frac12(N(t)-N_0(t))^2
\end{aligned}
\tag{168.4}
\]

The second equality is exact. In particular the unknown counting
fluctuation remains present after triangularization.

For consecutive ordinates above \(2\pi\), the usual smooth-unfolded gap
\(s_n=N_0(\gamma_{n+1})-N_0(\gamma_n)\) has the exact reading

\[
\boxed{s_n=
\frac{2\left(T_{N_0(\gamma_{n+1})}-T_{N_0(\gamma_n)}\right)}
{N_0(\gamma_{n+1})+N_0(\gamma_n)+1}}
\tag{168.5}
\]

This is the adjacent-product identity \(2(T_b-T_a)=(b-a)(a+b+1)\). The
denominator is positive in the stated range. It is a difference of
cumulative triangular coordinates; it is not \(T_{s_n}\). One may also
chart the positive gaps themselves by \(s\mapsto T_s\), which is
strictly increasing. The two operations have different meanings and
should not be interchanged.

\subsection{What the older enrichment experiment can
contribute}\label{what-the-older-enrichment-experiment-can-contribute}

The report {[}62{]} describes a preceding-window enrichment statistic
near small-gap valleys, with a broader follow-up scan from \(2.30\) to
\(3.20\) and lower controls. Its reported enrichment shelf is
approximately \(2.70\)--\(2.81\); \(11/4=2.75\) lies inside it. The
plotted ratios are observed to random-surrogate hit counts. They are not
average normalized zero spacing, and the reported record does not
establish \(2.75\) as a unique endpoint of the shelf.

The useful rational marker and the reported interval have exact
triangular coordinates:

\[
\begin{aligned}
T_{11/4}&=\frac{165}{32},\\
T_{[\,2.70,\,2.81\,]}&=
\left[\frac{999}{200},\frac{107061}{20000}\right],\\
T_{11/4\,\pm\,1/2000}
&=\frac{165}{32}\pm\frac{13}{8000}+\frac1{8000000}
\end{aligned}
\tag{168.6}
\]

The last line transports a symmetric original tolerance to its exact
slightly asymmetric triangular image. In general
\(T_{s+h}-T_s=(s+1/2)h+h^2/2\). To preserve a finite experiment,
transport all thresholds, windows, valley and merging rules, and support
conditions through the chosen coordinate map. Transforming only the
marker while leaving its window unchanged defines a different statistic.

Part E of this Dossier retains the historical tables and the earlier
five-marker figure, and reproduces the broader report pages that supply
their context. It also retains the modulus audit {[}61{]}. Those
empirical outcomes do not alter the exact gcd, totient, or decimal
identities in §§98 and 167. Their raw gap arrays and complete
randomization replay are not available here, so no new significance
level or confirmed shelf boundary is asserted.

For a further test, compare prespecified rational and \(\pi\)-based
markers on the same block support, retain the full scan correction, and
carry the entire statistic into the triangular chart. The existing
arithmetic gives that transport exactly. It does not yet supply a new
zero-location theorem.

\section{46. Native TFG nonlinear operator
route}\label{native-tfg-nonlinear-operator-route}

The scalar recognizers above provide no distribution estimate. A
proposed native operator should retain the unreduced grid and
distinguish actual arithmetic from an arbitrary prime mask. The
necessary proof obligations are:

\begin{enumerate}
\def\labelenumi{\arabic{enumi}.}
\tightlist
\item
  pairwise or higher interaction, not a weighted scalar prime sum;
\item
  explicit dependence on ambient denominator and gcd multiplicity;
\item
  failure under a generic support-preserving prime mask;
\item
  a proved intertwiner to one of the master positivity objects
  \(M_\xi\), \(Q_c\), \(J_c\), or \(\mathcal C_c\);
\item
  a quantifier-complete forcing theorem.
\end{enumerate}

The rational witness theorem gives the native TFG a rigorous indexing
role. It does not yet give it an arithmetic forcing role.

\Needspace{11\baselineskip}

\Needspace{11\baselineskip}

\section{169-172. Worked triangular
arithmetic}\label{worked-triangular-arithmetic}

The four sections that follow develop the reciprocal blocks, decimal
carries, gcd reduction, the Pascal/TFG address and the polar shells in
full. They are detail rather than argument: nothing in the Reading
Volume's continuous argument depends on reading them first. Their
identities return in §§9, 84, 98 and 167, in the Reading Volume at §R7
and §R17G, and in the source energy. Equation tags (4A.\emph{), (NR.}),
(RP.\emph{) and (AR.}) are unchanged.

\section{169. Reciprocal blocks, doubling, and decimal
carries}\label{reciprocal-blocks-doubling-and-decimal-carries}

The reciprocal law gives a particularly simple dyadic-index family:

\[
\frac1{T_{2^j-1}}=\frac{2^{1-j}}{2^j-1}\quad(j\ge1)
\qquad \frac1{T_7}=\frac1{28}=\frac14\frac17
\tag{4A.1}
\]

For every integer \(n\ge1\), the triangular-to-factorial conversion is
\(1/n!=[(n+1)/(2(n-1)!)]/T_n\). At \(n=7\) this gives the exact bridge

\[
\boxed{\frac17\ \xrightarrow{\ \div4\ }\ \frac1{T_7}=\frac1{28}
\ \xrightarrow{\ \div180\ }\ \frac1{7!}=\frac1{5040},\qquad
2T_7=56}
\tag{4A.1a}
\]

The repeated blocks in these reciprocals have an exact algebraic source.
Using base \(100\) to group decimal digits in pairs,

\[
\boxed{\frac17=\frac{14}{100-2}
=\sum_{j\ge0}\frac{14\cdot2^j}{100^{j+1}}}
\tag{4A.2}
\]

The raw coefficients begin \(14,28,56,112,224,\ldots\). They are not yet
base-\(100\) digits: the later coefficients create carries. After
removing the first two blocks, the remaining value, multiplied by
\(100^3\), is \(56/(1-2/100)=57+1/7\). This explains the carried block
\(57\) and closes the six-digit cycle exactly:

\[
\frac17=0.\overline{142857},\qquad
\boxed{\frac1{T_7}=T_4^{-4}\left(357+\frac17\right)
=0.0357\overline{142857},\quad T_4=10}
\tag{4A.3}
\]

The original \(1/7\) continues as the tail through successive halvings.
Each halving is exactly the midpoint of the preceding value and zero:

\[
\boxed{\begin{aligned}
\frac17&=0.\overline{142857}
&\frac1{14}&=0.07\overline{142857}\\
\frac1{T_7}=\frac1{28}&=0.0357\overline{142857}
&\frac1{2T_7}=\frac1{56}&=0.017857\overline{142857}
\end{aligned}}
\tag{4A.3a}
\]

The overbars in this display deliberately preserve the phase of the
original six-digit block. There is an exact rule for every halving step:

\[
\frac1{7\cdot2^r}=100^{-r}\left(A_r+\frac17\right),\qquad
A_r=\frac{50^r-1}{7},\qquad A_{r+1}=50A_r+7
\quad r\ge0
\tag{4A.3b}
\]

Indeed \(50\equiv1\pmod7\), so \(A_r\) is an integer and the fractional
part of \(100^r/(7\cdot2^r)\) is always \(1/7\). The prefix integers
\(0,7,357,17857,\ldots\) are placed in \(2r\) decimal positions, with
leading zeros included. Thus the repeated block is carried exactly; its
unchanged placement and the minimum length of a nonrepeating prefix are
two different readings of the same rational value.

\Needspace{12\baselineskip}

The following table uses the shortest prefixes. For example
\(0.0357\overline{142857}=0.03\overline{571428}\): moving the overbar
rotates the repeating block without changing the number.

\begin{longtable}[]{@{}
  >{\raggedright\arraybackslash}p{(\columnwidth - 6\tabcolsep) * \real{0.2143}}
  >{\raggedright\arraybackslash}p{(\columnwidth - 6\tabcolsep) * \real{0.2143}}
  >{\raggedleft\arraybackslash}p{(\columnwidth - 6\tabcolsep) * \real{0.2857}}
  >{\raggedleft\arraybackslash}p{(\columnwidth - 6\tabcolsep) * \real{0.2857}}@{}}
\toprule\noalign{}
\begin{minipage}[b]{\linewidth}\raggedright
Rational value
\end{minipage} & \begin{minipage}[b]{\linewidth}\raggedright
Decimal expansion
\end{minipage} & \begin{minipage}[b]{\linewidth}\raggedleft
Minimal prefix length
\end{minipage} & \begin{minipage}[b]{\linewidth}\raggedleft
Period length
\end{minipage} \\
\midrule\noalign{}
\endhead
\bottomrule\noalign{}
\endlastfoot
\(1/7\) & \(0.\overline{142857}\) & 0 & 6 \\
\(1/14\) & \(0.0\overline{714285}\) & 1 & 6 \\
\(1/T_7=1/28\) & \(0.03\overline{571428}\) & 2 & 6 \\
\(1/(2T_7)=1/56\) & \(0.017\overline{857142}\) & 3 & 6 \\
\(1/7!=1/5040\) & \(0.0001\overline{984126}\) & 4 & 6 \\
\(7!/990=56/11\) & \(5.\overline{09}\) & 0 & 2 \\
\end{longtable}

For a reduced denominator \(q=2^a5^b d\) with \(\gcd(d,10)=1\) and
\(d>1\), the minimal prefix length is \(\max(a,b)\) and the period
length is \(\operatorname{ord}_d(10)\). Here \(5040=2^4\cdot5\cdot63\),
with \(\operatorname{ord}_{63}(10)=6\), while \(10\equiv-1\pmod{11}\)
gives period two for \(56/11\). At the factorial scale,
\(1/5040=(1/720)(1/7)\) still retains the factor \(1/7\), but the
additional factor \(9\) in the reduced recurring denominator \(63\)
changes the block to \(984126\). That block is not a cyclic rotation of
\(142857\). The shared six-digit length has an order-theoretic
explanation even though the repeating blocks differ. Dossier §98 extends
this cyclic-order viewpoint; §76 connects the factorial and triangular
scales while retaining their different integer conditions.

The positional base and the logarithmic coordinate also fit together
exactly. This volume uses \(\log=\ln\) unless a base is displayed;
decimal plots use \(\log_{10}\) when labelled. The conversion is

\[
\boxed{v=\log_{10}n=\frac{\ln n}{\ln10},\qquad
n^{-s}=e^{-s\ln n}=10^{-sv}}
\tag{4A.4}
\]

Thus the source shift \(u=\ln n\) becomes \(u=(\ln10)v\). In an
integral, \(du=(\ln10)\,dv\); in a Fourier phase, \(tu=(t\ln10)v\).
Decimal cycles depend on the chosen base through multiplicative order.
The analytic identities remain the same after the coordinate, measure,
and frequency are converted together.

In the displayed family, \(\ln28=\ln7+2\ln2\), \(\ln56=\ln7+3\ln2\), and
\(\ln5040=\ln28+\ln180\). The same addition laws hold with
\(\log_{10}\). They express the reciprocal scalings in logarithmic
coordinates; the digit periods are determined separately by the
denominator orders above.

The same reciprocal scale also has an exact degree-to-radian rendering:

\[
\boxed{\left(\frac1{T_7}\right)^\circ
=\frac{\pi}{180T_7}\,\mathrm{rad}
=\frac{\pi}{7!}\,\mathrm{rad}
=\frac{\pi}{5040}\,\mathrm{rad}}
\tag{4A.5}
\]

Here the angle is assigned the degree measure \(1/T_7=1/28\), and the
conversion factor is \(\pi/180\). Its radian measure divided by \(\pi\)
is \(1/7!\). Grouping \(\pi\) with the angle-conversion factor makes the
triangular bridge explicit:

\[
\underbrace{\frac17\cdot\frac14}_{1/T_7}
\underbrace{\frac{\pi}{180}}_{\text{degree-to-radian factor}}
=\frac{\pi}{180T_7}=\frac{\pi}{7!}
\tag{4A.6}
\]

If the angle is \(\theta=(1/T_7)^\circ\), its radian measure is
\(\pi/5040\), while its fraction of a half-turn is
\(\theta/(180^\circ)=\theta/(\pi\,\mathrm{rad})=1/5040\). The \(\pi\)
cancels in this normalized fraction, not in the radian measure itself or
its decimal evaluation. A complex quarter-turn is multiplication by
\(i\); the displayed product changes an angle's scale and units.

\section{170. A triangle inside a triangle: 990, 496, and reduced
rays}\label{a-triangle-inside-a-triangle-990-496-and-reduced-rays}

The value \(990\) connects the reciprocal examples to an exact
composition of triangular coordinates. Start with \(T_2=3\), \(T_4=10\),
and \(T_{10}=55\):

\[
\boxed{T_{44}=T_2^2T_4(T_4+1)
=2T_2^2T_{T_4}=18T_{10}=990}
\tag{NR.1}
\]

Thus the inner product is precisely the triangular product at \(T_4\).
Its coefficient is \(T_2^2=9\); relative to the usual coefficient
\(1/2\), it multiplies the inner triangle by \(18\). This is a
composition identity, not an inference from the appearance of decimal
digits.

The general centered scaling law explains the neighboring perfect
number:

\[
\begin{aligned}
T_{az+(a-1)/2}&=a^2T_z+T_{(a-1)/2}\\
T_{31}&=T_{3\cdot10+1}=9T_{10}+1=496\\
2T_{31}&=T_{44}+2=992
\end{aligned}
\tag{NR.2}
\]

The first line is a polynomial identity for real or complex \(a,z\); odd
integer \(a\) preserves integer indices. Expanding either side, or using
\(8T_z+1=(2z+1)^2\), proves it. Another useful factorization is

\[
T_{ab}=a^2T_b-bT_{a-1},\qquad
T_{44}=15T_{11}=T_5T_{11}=10^3-10
\tag{NR.3}
\]

The divisors of \(496=16\cdot31\) come in two finite rows:

\begin{longtable}[]{@{}llllll@{}}
\toprule\noalign{}
power-of-two row & \(1\) & \(2\) & \(4\) & \(8\) & \(16\) \\
\midrule\noalign{}
\endhead
\bottomrule\noalign{}
\endlastfoot
multiply by \(31\) & \(31\) & \(62\) & \(124\) & \(248\) & \(496\) \\
distance to the next power of two & \(1\) & \(2\) & \(4\) & \(8\) &
\(16\) \\
\end{longtable}

Indeed \(31\cdot2^k=2^{k+5}-2^k\). Summing these divisors gives
\((1+2+4+8+16)(1+31)=31\cdot32=992=2\cdot496\), the classical
even-perfect construction {[}45{]}. The extension \(992\) continues the
doubling pattern but is not a divisor of \(496\).

For positive integers \(n,k\), the greatest common divisor
\(d=\gcd(n,k)\) removes their shared scale: \((n,k)=d(q,a)\) with
\(\gcd(q,a)=1\). Euler's totient \(\varphi(q)\) counts \(1\le a\le q\)
coprime to \(q\) (with \(\varphi(1)=1\)). Thus it counts primitive
rational directions at denominator \(q\), while gcd tells how often a
direction is repeated in larger triangular rows. The modulus comparison
has an exact common scale:

\[
\boxed{\varphi(496)=\varphi(990)=240=2T_{15},\qquad
\varphi(992)=480}
\tag{NR.4}
\]

This identifies equal counts of reduced residues; it does not identify
their arrangements or their decimal periods. Section 162 gives the
complete Pell family behind the offset \(2\), including the stricter
modular filter. The tested interval \(990\) through \(1009\) remains a
useful empirical search window. \EVID{} Section 168 and Part E of the
Dossier state exactly what the reported tests do and do not establish;
no result in this volume rests on them.

The distinction is visible in base ten. For \(496=2^4\cdot31\), removing
the factors of \(2\) and \(5\) leaves \(31\): \(1/496\) has a four-digit
nonrepeating prefix and period \(\operatorname{ord}_{31}(10)=15\). For
\(990=2\cdot5\cdot99\), the remaining denominator is \(99\):
\(1/990=0.0\overline{01}\) has a one-digit prefix and period two. The
same totient \(240\) therefore coexists with different decimal behavior.
Section 167 gives the general valuation and multiplicative-order rule.

\subsection{From gcd reduction to reciprocal triangular
weights}\label{from-gcd-reduction-to-reciprocal-triangular-weights}

For a native chamber address \((n,k)\), let \(d=\gcd(n,k)\), \(q=n/d\),
and \(a=k/d\). The rational direction is \(a/q\), and its denominator
weight is

\[
\boxed{\frac1{T_q}
=2\left(\frac1q-\frac1{q+1}\right)
=\frac{2d^2}{n(n+d)}}
\tag{NR.5}
\]

Thus the gcd and the reciprocal telescope describe the same weight in
two coordinates. Among all addresses \(1\le k\le n\le N\), the number
having reduced denominator \(q\) is

\[
C_N(q)=\varphi(q)\left\lfloor\frac Nq\right\rfloor
\qquad T_N=\sum_{q\le N}C_N(q)
\tag{NR.6}
\]

Consequently the complete weighted chamber sum is
\(\sum_{q\le N}C_N(q)/T_q\). This is the bridge to the TFG denominator
moments in the Technical Dossier. Here \(q\) is an integer denominator;
the later \(q(s)=s(1-s)\) is an analytic coordinate.

The distinction between a denominator weight and a folded rational value
is visible side by side:

\[
\frac1{T_q}=\frac2{q(q+1)},\qquad
T_{a/q}=\frac{a(a+q)}{2q^2},\qquad
\frac1{T_{a/q}}=\frac{2q^2}{a(a+q)}
\tag{NR.7}
\]

The last two formulas use the numerator as well. Section 167 follows
their reduced denominators under iteration, with exact decimal periods
and \(\log_{10}\) scales. These are useful arithmetic coordinates even
when a particular spacing experiment fails to detect a modulus effect.

\begin{figure}
\centering
\includegraphics{figures/v27_nested_triangles.png}
\caption{Nested triangular values organize 990 and 496 before either is
used as a modulus. The shared totient is a separate exact comparison.}
\end{figure}

\section{171. Reciprocal folding, Pascal coefficients, and the TFG
address}\label{reciprocal-folding-pascal-coefficients-and-the-tfg-address}

Pascal's triangle already contains the triangular numbers:
\(\binom{n+1}{2}=T_n\). The row \(1,5,10,10,5,1\) contains the repeated
\(10=T_4\) and pairs naturally when the two inputs are reciprocals. For
\(z\ne0\), define

\[
y=z+\frac1z,\qquad
\mathcal P_m(y):=z^m+z^{-m}
\tag{RP.1}
\]

Exchanging \(z\) and \(1/z\) leaves both \(y\) and \(\mathcal P_m\)
unchanged. The quadratic relation \(z^2-yz+1=0\) gives a recurrence
without a long binomial expansion:

\[
\boxed{\mathcal P_0=2,\quad\mathcal P_1=y,\quad
\mathcal P_m=y\mathcal P_{m-1}-\mathcal P_{m-2}\quad(m\ge2)}
\tag{RP.2}
\]

The first polynomials and their values at \(y=3\) are

\begin{longtable}[]{@{}lll@{}}
\toprule\noalign{}
\(m\) & reciprocal polynomial \(\mathcal P_m(y)\) & value at \(y=3\) \\
\midrule\noalign{}
\endhead
\bottomrule\noalign{}
\endlastfoot
\(0\) & \(2\) & \(2\) \\
\(1\) & \(y\) & \(3\) \\
\(2\) & \(y^2-2\) & \(7\) \\
\(3\) & \(y^3-3y\) & \(18\) \\
\(4\) & \(y^4-4y^2+2\) & \(47\) \\
\(5\) & \(y^5-5y^3+5y\) & \(123\) \\
\end{longtable}

Pascal's fifth row groups exactly as

\[
\left(z+\frac1z\right)^5
=\mathcal P_5(y)+5\mathcal P_3(y)+10\mathcal P_1(y)
\tag{RP.3}
\]

The unsigned coefficients \(1,5,10,10,5,1\) are the magnitudes of row 5
of the signed grid. Pairing reciprocal powers yields \(1,5,10\) in
(RP.3). The coefficients inside the reciprocal polynomials themselves
have a neighboring-cell formula, for \(m\ge1\):

\[
\mathcal P_m(y)=\sum_{j=0}^{\lfloor m/2\rfloor}
\bigl(c_{m-j,j}-c_{m-j-1,j-1}\bigr)y^{m-2j}.
\tag{RP.3a}
\]

It follows by substituting the signed Pascal recurrence into (RP.2). For
\(m=5\), the coefficients are \(1\), \(-4-1=-5\), and \(3-(-2)=5\). The
separate initial value \(\mathcal P_0=2\) counts the two reciprocal
powers; the coefficient-grid seed remains \(1\).

Thus \(y=3\) gives \(3^5=243=123+5\cdot18+10\cdot3\). The literal fifth
power of the sum is \(243\); the recovered sum of fifth powers is
\(123\). Distinguishing them is essential when using a binomial identity
as a computational shortcut. The two input branches are
\(z=(3\pm\sqrt5)/2\), and reciprocal inversion exchanges them.

This is the same quadratic recurrence mechanism as (4.4b). In general,
if \(A_m=\alpha^m+\beta^m\), then

\[
A_m=(\alpha+\beta)A_{m-1}-\alpha\beta A_{m-2}
\qquad A_0=2,\quad A_1=\alpha+\beta
\tag{RP.4}
\]

For adjacent reciprocals \(1/n,1/(n+1)\), its coefficients are
\(j_n/(2T_n)\) and \(1/(2T_n)\). For \(z,1/z\), they are \(y\) and
\(1\). The formulas therefore share an algebraic mechanism while keeping
their different products visible.

For a positive reduced TFG value \(z=a/q\), \(\gcd(a,q)=1\),

\[
y=\frac{a^2+q^2}{aq},\qquad
y^2-4=\frac{(a^2-q^2)^2}{a^2q^2}
\tag{RP.5}
\]

The first fraction is reduced: its numerator is coprime to both \(a\)
and \(q\). Its denominator is \(aq\), and numerator--denominator
exchange implements the reciprocal symmetry. This is an exact map of
rational values; it does not assert equality of different TFG aggregate
weights.

The analytic return appears in §R7. With \(z_C=s/(1-s)\) and
\(q(s)=s(1-s)\), that section gives

\[
y_C=z_C+\frac1{z_C}=\frac1{q(s)}-2
\tag{RP.6}
\]

Thus the same reciprocal coordinate is already present in the completed
fold. Pascal grouping, triangular coefficients, and the recurrence keep
the two inverse branches visible through a single invariant. On the unit
circle inversion is conjugation; away from it those operations differ.
The missing completed-source sign is still a further analytic condition.

This Pascal arithmetic returns in the source energy of §R17G. Its exact
matrix has entries \(\binom{r+s}{r}=(-1)^r c_{r+s,r}\) of the signed
grid. At \(r=2\) these follow the same grid's column \(j=2\):
\(c_{s+2,2}=T_{s+1}\), giving \(1,3,6,10,\ldots\). Dossier §163A
displays the matrix, proves its Gram factorization, and identifies the
polynomial reflection behind the change of basis.

\section{\texorpdfstring{172. Cyclic sums, polar shells, and the
\(1/12\)
marker}{172. Cyclic sums, polar shells, and the 1/12 marker}}\label{cyclic-sums-polar-shells-and-the-112-marker}

The first arithmetic bridge is built before any zero criterion. For a
positive integer \(q\), define the cyclic-order sum

\[
\Psi_{\rm cyc}(q)
=\sum_{a=1}^{q}\frac{q}{\gcd(a,q)}
=\sum_{d\mid q}d\varphi(d)
\tag{AR.1}
\]

The notation \(\Psi_{\rm cyc}\) is deliberate: it is an arithmetic
function of the integer \(q\) and is unrelated to the folded resolvent
\(S_\xi(x)\). Its Dirichlet series is

\[
\boxed{
F_{\rm cyc}(s)=\sum_{q\ge1}\frac{\Psi_{\rm cyc}(q)}{q^s}
=\frac{\zeta(s)\zeta(s-2)}{\zeta(s-1)}
\qquad \Re s>3}
\tag{AR.2}
\]

The pole at \(s=3\) has residue \(\zeta(3)/\zeta(2)\). First-Riesz
smoothing multiplies it by \(1/[s(s+1)]\), hence

\[
\boxed{
D
=\left.\frac1{s(s+1)}\right|_{s=3}
  \frac{\zeta(3)}{\zeta(2)}
=\frac1{12}\frac{\zeta(3)}{\zeta(2)}
=\frac1{2T_3}\frac{\zeta(3)}{\zeta(2)}}
\tag{AR.3}
\]

This is the exact reason for the positive \(1/12\) in the Riesz main
term: \(3\cdot4=2T_3=12\). The familiar continuation value

\[
\zeta(-1)=-\frac1{12}=-\frac1{2T_3}
\tag{AR.4}
\]

has the opposite sign but a different owner. Triangular reflection sends
\(3\mapsto-4\) and preserves \(T_3=T_{-4}\), so it preserves the
positive kernel factor \(1/12\); it does not create an automatic
reflected coefficient \(-D\).

The Polar Shell Renderer supplies a compatible picture, not the proof of
(AR.2)--(AR.4). Its reduced shell points are

\[
Z_{n,k}=\sqrt n\,e^{2\pi ik/n}
\qquad (k,n)=1
\tag{AR.5}
\]

and division by the shell radius gives primitive roots of unity. The
points \(k\) and \(n-k\) are conjugates, so a real shell sum never comes
from an unpaired complex phase. Combining the angular factor with the
Mellin weight gives

\[
U_{n,k}(s)=e^{2\pi ik/n}n^{1/2-s}
\qquad
S_n^{\rm shell}(s):=\sum_{(k,n)=1}U_{n,k}(s)
=\mu(n)n^{1/2-s}
\tag{AR.6}
\]

For squarefree \(n>1\),

\[
\left|S_n^{\rm shell}(s)\right|=n^{1/2-\sigma}=1
\iff \sigma=\frac12
\tag{AR.7}
\]

This is an exact critical-line detector, but no source identity here
forces a zeta zero to make its value one. At \(n=12\), the reduced shell
contains \(\varphi(12)=4\) primitive twelfth roots in two conjugate
pairs and sums to \(\mu(12)=0\); it is a clean conjugate-pair example,
not a modulus-one detector. The full cyclic/Riesz proof is in §§56--57,
and the renderer construction is cited separately in {[}55{]}.

\begin{figure}
\centering
\includegraphics{figures/figure83_v26_polar_shell_pairing.png}
\caption{Equal-area polar shells and the \(1/\sqrt n\) normalization.
The angular laws differ, but every complex phase must retain its
conjugate partner.}
\end{figure}

\section{173. The grid's own finite arithmetic: tracks, reflections,
collisions}\label{the-grids-own-finite-arithmetic-tracks-reflections-collisions}

These are properties of the signed coefficient grid itself, collected
out of the reading line. They are exact and they are finite. None of
them supplies an inequality, and none is used by any positivity argument
in either document.

The columns are read as tracks: the edge carries \(1\), the next counts
\(k\), the third is \(\binom{k}{2}=T_{k-1}\), and the fourth is
\(\binom{k}{3}=\sum_{r=1}^{k-2}T_r\) for \(k\ge3\). At \(k=6\) that
fourth entry is \(20=1+3+6+10\), beside the triangular \(15=T_5\).
Reflection supplies the other half:

\[
\binom{k}{k-j}=\binom{k}{j},\qquad
(-1)^{k-j}\binom{k}{k-j}=(-1)^k(-1)^j\binom{k}{j}
\tag{SW.4c}
\]

An even row has identical signed coefficients when reversed; an odd row
has opposite signed coefficients. This is a reflection of coefficients.
The attached terms also change: \(x^jZ_{k+j}\mapsto x^{k-j}Z_{2k-j}\).
They are not equal functions merely because their coefficient magnitudes
agree. Complex conjugation is separately retained by the folded zero sum
in (SW.2).

\textbf{Interior triangular collisions.} The third track is triangular
by construction. To ask a different, finite question, exclude the edge
and count tracks and retain one half of each reflected row. In the exact
domain \(1\le k\le20\), \(3\le j\le\lfloor k/2\rfloor\), the complete
list is

\[
\begin{aligned}
&\binom{10}{3}=T_{15},\qquad
\binom{10}{4}=T_{20},\qquad
\binom{14}{6}=T_{77},\\
&\binom{15}{5}=T_{77},\qquad
\binom{17}{8}=T_{220},\qquad
\binom{19}{5}=T_{152}
\end{aligned}
\tag{SW.4e}
\]

Reflection supplies the matching positions in the right halves. The test
is exact: a positive integer \(v\) is triangular if and only if \(8v+1\)
is an odd square. This finite census makes no classification claim
outside the stated domain and does not include the automatic \(j=2\)
track.

Row \(k=14\) gives

\[
\boxed{\binom{14}{6}=\binom{14}{8}=3003=T_{77}=\binom{78}{2}}
\tag{SW.4f}
\]

The eight positions of \(3003\) in Pascal's triangle (compare
Singmaster's problem {[}89{]}) are exactly \((14,6)\), \((14,8)\),
\((15,5)\), \((15,10)\), \((78,2)\), \((78,76)\), \((3003,1)\), and
\((3003,3002)\). To check exhaustiveness, reflect to \(j\le k/2\). For
\(j\ge7\) the smallest entry is at least \(\binom{14}{7}=3432>3003\).
For each \(1\le j\le6\), \(\binom{k}{j}\) increases strictly with
\(k\ge2j\); checking the neighboring values leaves precisely the four
left-half positions listed above. This is an arithmetic multiplicity
statement. It gives no source inequality and no additional positivity
for \(W_{14}\).

The complete row at \(k=12\) is also explicit:

\[
\begin{gathered}
(1,-12,66,-220,495,-792,924,\\
\hspace{12mm}-792,495,-220,66,-12,1),\\
\boxed{(-1)^9\binom{12}{9}=-220}
\end{gathered}
\tag{SW.4h}
\]

\textbf{The \(k=0\) seed and the reduced-rung boundary.} The coefficient
grid of the Reading Volume includes the valid coefficient row \([+1]\),
as used by \(F_{n,0}\) in (SW.3). The differential definition (R.3), the
normalized resolvent row (SW.4b), and the criterion (R.4) all require
\(k\ge1\). An antiderivative can be added by a separate convention. For
a chosen \(x_*>0\), define

\[
\mathcal A_0(x):=-\int_{x_*}^{x}S_\xi(t)\,dt,\qquad
\mathcal A_0'=-S_\xi,\qquad
-\mathcal A_0'-x\mathcal A_0''=W_1
\tag{SW.4g}
\]

This verifies the recurrence at the boundary for that convention. It
does not extend the rung definition to a negative derivative order or
assign a positivity condition to \(\mathcal A_0\). Adding a constant
changes the antiderivative but leaves its displayed recurrence
unchanged. The coefficient \(-1\) in this convention and the signed
Pascal seed \([+1]\) at \(k=0\) belong to different expressions.

The formal resolvent \(Z_0=\sum_{[\rho]}m_\rho\) diverges, and
\((2k-1)!=\Gamma(2k)\) is singular at \(k=0\). Continuing the isolated
coefficient ratio gives \(\Gamma(2k)/\Gamma(k)\to1/2\); with a sign
continuation tending to \(-1\), that isolated coefficient tends to
\(-1/2\). It does not define a canonical antiderivative or a rung.
Likewise \((-1)^j\binom{-1}{j}=1\) defines an infinite power-series row,
not another finite member of (SW.4b).

The normalization is itself an accumulated triangular product:

\[
(2k-1)!=\prod_{r=1}^{k-1}2T_{2r},\qquad
\frac{(2k+1)!}{(2k-1)!}=2T_{2k},\qquad
\frac{11!}{9!}=2T_{10}=110
\tag{SW.4d}
\]

The ratio is twice column \(j=2\) of row \(2k+1\) in the same grid:
\(2T_{2k}=2c_{2k+1,2}\). For \(k=5\), the entry \(c_{11,2}=55\)
therefore gives \(11!/9!=110\).

This is the odd-factorial denominator of the sine series (PR.6), and the
Bernstein location is \(B_{11,5}(a)=a^5W_6(a)/(5!6!)\) by (BR.3a).

\Needspace{11\baselineskip}

\chapter{Technical chapter III. Centered geometry and operator
models}\label{technical-chapter-iii.-centered-geometry-and-operator-models}

The centered product appears in Gamma products, Casimir spectra and
finite geometry. Each operator model retains its own space and
normalization.

\section{10. Odd-dimensional moment-Casimir
rail}\label{odd-dimensional-moment-casimir-rail}

Let

\[
u(z)=\frac{z(z-1)}2
\qquad
u_d(s)=\frac{(s-1)(s-d+1)}2
\]

For \(j\ge1\), the difference \(u_d(s)-u(s-j)\) is independent of \(s\)
if and only if

\[
\boxed{d=2j+1}
\]

At that dimension,

\[
\boxed{
u_{2j+1}(s)=u(s-j)-T_{j-1}
}
\]

At an inherited pole \(s=j+\rho\),

\[
u_{2j+1}(j+\rho)+T_{j-1}=u(\rho)
\]

The split-norm form is

\[
\boxed{
u_{2j+1}(j+\rho)
=\frac12\left[
\left(\rho-\frac12\right)^2
-\left(j-\frac12\right)^2
\right]
}
\]

The discrete spherical degree value has the parallel form

\[
\lambda_n^{(2j+1)}
=\frac12\left[
\left(n+j-\frac12\right)^2
-\left(j-\frac12\right)^2
\right]
\]

Both use the same fixed half-gap \(j-1/2\), but \(n\) and \(\rho\) live
in different spaces. The result is a common coordinate rail, not an
operator identification.

For every fixed \(j\ge1\),

\[
\mathrm{RH}
\iff
u_{2j+1}(j+\rho)\in\mathbb R_{<0}
\quad\text{for every nontrivial zero }\rho
\]

This is a coordinate equivalence: the imaginary part is
\(\gamma(\beta-1/2)\).

\section{11. Triangular canonical
product}\label{triangular-canonical-product}

Define

\[
\Pi_\triangle(u)=\prod_{n=1}^{\infty}
\left(1-\frac{u}{T_n}\right)
\]

Since \(\sum T_n^{-1}=2\), the product converges locally uniformly.
Under

\[
u(s)=\frac{s(s-1)}2
\]

the exact Gamma closure is

\[
\boxed{
\Pi_\triangle(u(s))
=\frac1{\Gamma(1+s)\Gamma(2-s)}
}
\]

At the fixed point \(s=1/2\),

\[
u(1/2)=-\frac18
\qquad
\boxed{\Pi_\triangle(-1/8)=\frac4\pi}
\]

The factors \(1/4\), \(1/2\), or \(2\) that turn this identity into
\(1/\pi\) or \(2/\pi\) are external normalizations and must be declared.
Repeated scalar contacts do not identify mechanisms.

\textbf{One of these contacts is not scalar, and the caution is
withdrawn for it alone.} The substitution \(u=s(s-1)/2\) \emph{is} the
fold of §16: with \(x=-2u=s(1-s)=q(s)\) and \(\nu=\sqrt{x-\tfrac14}\),
\(\Gamma\) reflection turns the Gamma closure above into

\[
\boxed{\ \prod_{n\ge1}\Bigl(1+\frac{x}{2T_n}\Bigr)=\frac{\sin\pi s}{\pi\,q(s)}
=\frac{\cosh\pi\nu}{\pi x},
\qquad
\sum_{n\ge1}\frac1{x+2T_n}=\frac{\pi\tanh\pi\nu}{2\nu}-\frac1x\ }
\tag{11.1a}
\]

the second being the logarithmic derivative of the first. So the
fixed-point value is not a numerical coincidence: \(u=-1/8\) is
\(x=q(1/2)=1/4\), the foot of the critical line, and
\(\Pi_\triangle=4/\pi\) is the value there. In this form the product is
visibly entire of order \(\tfrac12\) in \(x\) with zeros exactly at
\(x=-2T_n\), \(n\ge1\); \(x=0\) has removable value one, not a zero.
This model and the completed determinant have the same entire order, not
the same type; Section 12B proves the difference. Here the model
parameter is \(x=q(s)\), whereas the completed source parameter is
\(s(s-1)=-q(s)\). Reader Section R6A keeps these two evaluations
separate. The model nodes are real and positive. Section 12 supplies the
spherical degree interpretation; Section 12A realizes exactly the
multiplicity-one product by a Dirichlet--Neumann interval operator with
its zero mode removed. No identification with the completed spectrum is
proved.

\begin{figure}
\centering
\includegraphics{figures/triangular_product_convergence.png}
\caption{Convergence of the triangular canonical product at the fixed
point.}
\end{figure}

\section{12. Spherical Casimir and positive
degree}\label{spherical-casimir-and-positive-degree}

For \(d\ge2\), set

\[
\lambda_n^{(d)}=\frac{n(n+d-2)}2
\]

Let \(\Lambda_d\) be the Dirichlet-to-Neumann operator of the unit ball.
Harmonic extension sends a degree-\(n\) spherical harmonic to
\(r^nY_n\), so

\[
\Lambda_dY_n=nY_n
\]

Therefore

\[
\boxed{
-\frac12\Delta_{S^{d-1}}
=\frac12\Lambda_d\bigl(\Lambda_d+(d-2)I\bigr)
}
\]

Equivalently,

\[
\Lambda_d=
\sqrt{-\Delta_{S^{d-1}}+\frac{(d-2)^2}{4}I}
-\frac{d-2}{2}I
\]

This supplies a canonical positive branch of the reflection-completed
degree line. The centered square is

\[
\lambda_n^{(d)}+\frac{(d-2)^2}{8}
=\frac12\left(n+\frac{d-2}{2}\right)^2
\]

The dimension-\(3\) spectrum is exactly \(T_n\). The same shifted degree
is the Bessel order arising from radial separation.

The dimension-dependent determinant closes as

\[
\Pi_d(u_d(z))
=\frac{\Gamma(d-1)}{\Gamma(z)\Gamma(d-z)}
\qquad
u_d(z)=\frac{(z-1)(z-d+1)}2
\]

\begin{figure}
\centering
\includegraphics{figures/casimir_positive_branch.png}
\caption{Casimir spectra and recovery of the canonical positive degree
branch.}
\end{figure}

\section{12A. The triangular sine operator and its finite
projection}\label{a.-the-triangular-sine-operator-and-its-finite-projection}

The multiplicity-one triangular product has an explicit interval
realization. Use the reflected polynomials of Section 155A,
\(\mathcal P_0=1\), \(\mathcal P_1=3-y\) and
\(\mathcal P_{m+1}=(2-y)\mathcal P_m-\mathcal P_{m-1}\). Their
triangular expansion and finite geometric sum give, with \(n=2m+1\),

\[
\begin{aligned}
\mathcal P_m(y)&=(2m+1)\sum_{j=0}^m
 \frac{(-2y)^j}{(2j+1)!}\prod_{r=0}^{j-1}(T_m-T_r),\\
\mathcal P_m(4\sin^2(\theta/2))
 &=\sum_{h=-m}^{m}e^{ih\theta}
 =\frac{\sin(n\theta/2)}{\sin(\theta/2)}.
\end{aligned}
\tag{12A.1}
\]

These are the fourth-kind Chebyshev polynomials in the normalization of
{[}63, equation 18.5.4{]}. The quotient has its removable
interpretations.

\textbf{Proposition 12A.1 (triangular contraction and finite tent).}
Locally uniformly for complex \(z\),

\[
\boxed{\frac{\mathcal P_m(z^2/(2T_m))}{2m+1}
\longrightarrow\frac{\sin z}{z}.}
\tag{12A.2}
\]

Here \(m\ge1\) in the approximants and the value at \(z=0\) is one.
Indeed, the \(j\)th coefficient is \((-1)^j/(2j+1)!\) times
\(\prod_{r<j}(1-T_r/T_m)\). Each retained factor lies in \([0,1]\), the
product tends to one for fixed \(j\), and \(R^{2j}/(2j+1)!\) is a
summable majorant on \(|z|\le R\). This proves the uniform limit without
extrapolating finite computations. Also, directly from (12A.1),

\[
\frac1n\mathcal P_m\!\left(4\sin^2\frac{\pi u}{n}\right)
=\frac{\operatorname{sinc}_\pi(u)}{\operatorname{sinc}_\pi(u/n)},
\qquad \operatorname{sinc}_\pi(u)=\frac{\sin\pi u}{\pi u}.
\tag{12A.3}
\]

For bounded real \(u\) the error from \(\operatorname{sinc}_\pi(u)\) is
\(O(n^{-2})\). Squaring the real-angle sum gives the Fejer identity

\[
\frac{\mathcal P_m(4\sin^2(\theta/2))^2}{n}
=\sum_{h=-(n-1)}^{n-1}\left(1-\frac{|h|}{n}\right)e^{ih\theta}.
\tag{12A.4}
\]

There are exactly \(n-|h|\) pairs with difference \(h\). Positive and
negative lags each contribute \(T_{n-1}\) pairs; zero lag contributes
\(n\). Thus \(n+2T_{n-1}=n^2\) counts the whole square. For odd
\(n=2m+1\) this is also \(2T_{2m}+2m+1=(2m+1)^2=8T_m+1\). \(\square\)

The weighted sine identity of Section 95A is compatible with this limit:
\(g_m(z)=z\mathcal P_m(z^2/(2T_m))/(2m+1)\to\sin z\) locally uniformly,
so analytic differentiation gives \(T_Dg_m\to T_D\sin z\). This relates
two operations by a limit, not by equality of the finite polynomials.

\textbf{Theorem 12A.2 (a specified positive comparison operator).} On
\(L^2(0,\pi)\) let

\[
 A_\triangle=-\frac{d^2}{d\theta^2}-\frac14,\qquad
 \mathcal D(A_\triangle)=
 \{f\in H^2(0,\pi):f(0)=0,\ f'(\pi)=0\}.
\tag{12A.5}
\]

It is self-adjoint and nonnegative, with a complete orthonormal basis

\[
\psi_m(\theta)=\sqrt{\frac2\pi}\sin((m+\tfrac12)\theta),\qquad
A_\triangle\psi_m=2T_m\psi_m\quad(m\ge0).
\tag{12A.6}
\]

\textbf{Proof.} The separated Dirichlet--Neumann boundary conditions
give the regular self-adjoint second-derivative operator. Direct
differentiation produces \((m+1/2)^2-1/4=m(m+1)\). Orthonormality is
elementary sine integration. For completeness, reflect a function evenly
about \(\pi\) into \((0,2\pi)\); its Dirichlet sine expansion retains
exactly the odd modes. The resulting complete spectral resolution also
proves nonnegativity. With \(\omega\) as in (95A.4), the map

\[
(Uf)(\theta)=\sqrt{\frac2\pi}\sin(\theta/2)
 f(4\sin^2(\theta/2))
\tag{12A.7}
\]

is unitary from \(L^2(\omega)\) to \(L^2(0,\pi)\) and sends
\(\mathcal P_m\) to \(\psi_m\). Differentiation gives the polynomial
equation
\(y(4-y)\mathcal P_m''+(6-2y)\mathcal P_m'+2T_m\mathcal P_m=0\). The
unitary images, not the bare polynomials in \(\theta\), are the interval
eigenfunctions. \(\square\)

Remove the \(m=0\) zero mode and call the restriction \(A_+\). Since
\(\sum_{m\ge1}[m(m+1)]^{-1}=1\), its inverse is positive trace class and

\[
\boxed{\det(I+xA_+^{-1})
=\prod_{m\ge1}\left(1+\frac{x}{2T_m}\right)
=\frac{\cosh(\pi\sqrt{x-1/4})}{\pi x}.}
\tag{12A.8}
\]

The last equality is Section 11, including the entire continuation and
removable value one at zero. These factors have multiplicity one. The
full spherical Laplacian has additional angular multiplicities and is
not this determinant without a specified restriction.

\textbf{Theorem 12A.3 (projection and sine-kernel limit).} The
rank-\(N\) orthogonal projection onto \(\psi_0,\ldots,\psi_{N-1}\) has
kernel

\[
\boxed{\Pi_N(\theta,\phi)=\frac1{2\pi}
\left[\frac{\sin(N(\theta-\phi))}{\sin((\theta-\phi)/2)}
-\frac{\sin(N(\theta+\phi))}{\sin((\theta+\phi)/2)}\right].}
\tag{12A.9}
\]

For every fixed \(\theta_0\in(0,\pi)\), uniformly on compact real
\(u,v\) sets,

\[
\frac\pi N\Pi_N\!\left(\theta_0+\frac{\pi u}N,
 \theta_0+\frac{\pi v}N\right)
\longrightarrow\operatorname{sinc}_\pi(u-v).
\tag{12A.10}
\]

\textbf{Proof.} Product-to-sum and the finite cosine sum prove (12A.9),
with removable values on the diagonal. In its difference term the
denominator is its argument times \(1+O(N^{-2})\). In the reflected-sum
term it tends to \(\sin\theta_0>0\), while the numerator is bounded. The
normalization makes that second term \(O(N^{-1})\), proving the uniform
limit. At the two endpoints the reflected term instead survives:

\[
\begin{aligned}
\frac\pi N\Pi_N(\pi u/N,\pi v/N)
 &\longrightarrow\operatorname{sinc}_\pi(u-v)-\operatorname{sinc}_\pi(u+v),\\
\frac\pi N\Pi_N(\pi-\pi u/N,\pi-\pi v/N)
 &\longrightarrow\operatorname{sinc}_\pi(u-v)+\operatorname{sinc}_\pi(u+v).
\end{aligned}
\tag{12A.11}
\]

Both follow by substitution, uniformly on compact positive \(u,v\) sets.
Removing the zero mode changes the rescaled interior kernel by a term
tending to zero; its inclusion here does not change the convention in
(12A.8). \(\square\)

\textbf{A random-position comparison, not a zeta model.} Introduce the
density \(\det[\psi_{i-1}(\theta_j)]_{i,j=1}^N{}^2/N!\) on
\((0,\pi)^N\). Expansion of both determinants and orthonormality prove
that its integral is one. Integrating all but two variables gives
distinct-point intensity

\[
\rho_{2,N}(\theta,\phi)=\Pi_N(\theta,\theta)\Pi_N(\phi,\phi)
-\Pi_N(\theta,\phi)^2.
\tag{12A.12}
\]

After the scaling (12A.10) this tends, with the Jacobian factor
\((\pi/N)^2\), to \(1-\operatorname{sinc}_\pi(u-v)^2\). It is a
distribution of random positions, not of the deterministic triangular
eigenvalues. The latter unfold to the equally spaced sequence
\(\sqrt{2T_m+1/4}=m+1/2\). No identification with actual zeta zeros is
made. Section 61A supplies the precise ordinate-scale comparison;
Section 155C uses a weighted trace of this operator for a different
positive measure.

\section{12B. The comparison and completed determinants have different
types}\label{b.-the-comparison-and-completed-determinants-have-different-types}

Use separate evaluation parameters:
\(E_\triangle(z)=\prod_{n\ge1}(1+z/[n(n+1)])\) and
\(E_\xi(x)=X(-x)/X(0)=\xi(s_x)/\xi(1)\). The Gamma/reflection
substitution of Section 11 is \(z=q(s)\), whereas the completed source
uses \(x=-q(s)=s(s-1)\). Consequently

\[
\begin{aligned}
E_\triangle(1/4)&=4/\pi,\\
E_\xi(-1/4)&=X(1/4)/X(0)=\xi(1/2)/\xi(1),\\
E_\xi(1/4)&=X(-1/4)/X(0)
 =\xi((1+\sqrt2)/2)/\xi(1).
\end{aligned}
\tag{12B.1}
\]

The first two are centered evaluations in different parameter
conventions; the third is the completed determinant at the model's
positive numerical comparison parameter. They must not be identified.

\textbf{Proposition 12B.1 (order does not identify type).} As real
\(x\to\infty\),

\[
\begin{aligned}
\log E_\triangle(x)&=\pi\sqrt x-\log(2\pi x)+O(x^{-1/2}),\\
\log E_\xi(x)&=\frac{\sqrt x}{4}\log x
 -\frac{1+\log(2\pi)}2\sqrt x+\frac78\log x+C_\xi
 +O\!\left(\frac{\log x}{\sqrt x}\right).
\end{aligned}
\tag{12B.2}
\]

Both functions have order one-half. The triangular determinant has
finite type \(\pi\) at this order; the completed determinant has
infinite type, without assuming RH.

\textbf{Proof.} Expand the closed quotient
\(\cosh(\pi\sqrt{x-1/4})/(\pi x)\) for the first formula. For the
second, \(\frac{d}{dx}\log E_\xi=S_\xi\), and integration of (OP.1)
gives \(\sqrt x\log x/4-(1+\log(2\pi))\sqrt x/2+(7/8)\log x\). Its
remainder is integrable at infinity and leaves a constant plus
\(O(\log x/\sqrt x)\). The closed quotient bounds the triangular circle
maximum by an algebraic factor times \(\exp(\pi\sqrt{r+1/4})\), and its
positive axis attains type \(\pi\). The completed positive-axis growth
already makes \(\log M(r)/\sqrt r\) unbounded; the folded order-one-half
statement is established in Section 16. \(\square\)

RH concerns the negative-real-zero set of \(E_\xi\), not equality of its
entire-function type or spectral density with the triangular model. The
operator and sine-kernel results of Section 12A remain comparison
results with their stated multiplicities.

\section{13. Reciprocal and Cayley
uniformization}\label{reciprocal-and-cayley-uniformization}

For \(d>2\), put \(D=d-2\) and

\[
p=\frac{z-1}{D}
\qquad
r=\frac{p}{1-p}
\]

Then reflection \(z\mapsto d-z\) becomes

\[
p\mapsto1-p
\qquad
r\mapsto r^{-1}
\]

For real \(0<p<1\) (equivalently \(r>0\)), the bounded invariant is

\[
\boxed{
\tau=p(1-p)=\frac{r}{(1+r)^2}
=\frac14\operatorname{sech}^2 y
\qquad r=e^{2y}
}
\]

The reciprocal \((1+r)^2/r\) is not the bounded variable \(\tau\). On
the zeta-centered right sheet, \(p=1/2+c\) and hence

\[
\boxed{
\tau=\frac14-c^2
\qquad
r=\frac{1+2c}{1-2c}
\qquad
y=\operatorname{artanh}(2c)
}
\]

This is the exact collar coordinate used in the flux synthesis below.

\subsection{Proposition 13.1 - The self-scaling conic
pencil}\label{proposition-13.1---the-self-scaling-conic-pencil}

The triangular-circle relation from the geometric branch can be written

\[
T_x
=\frac12y\left(c+\frac12\right)
-\frac12y^2-\frac18
\]

With \(X=x+1/2\), this is exactly

\[
\boxed{
X^2+y^2-\left(c+\frac12\right)y=0
}
\]

For fixed \(c\), let

\[
r_c=\frac{c+1/2}{2}
\]

The locus is the circle

\[
X^2+(y-r_c)^2=r_c^2
\]

For \(r_c\ne0\) it has center \(C_c=(0,r_c)\), radius \(|r_c|\), passes
through \(F=(0,0)\), and is tangent there to \(L:y=0\). Its centers
satisfy

\[
|C_c-F|=\operatorname{dist}(C_c,L)
\]

Because the focus lies on the directrix, the center locus is the
degenerate focus-directrix parabola \(X=0\). In the original coordinate
it is \(x=-1/2\).

Now impose a linear selector

\[
c=\lambda X+\mu
\]

Eliminating \(c\) gives

\[
\boxed{
X^2-\lambda Xy+y^2-\left(\mu+\frac12\right)y=0
}
\]

Its conic discriminant is

\[
\Delta=\lambda^2-4
\]

When \(\mu\ne-1/2\), the locus is a nondegenerate ellipse, parabola, or
hyperbola according as \(|\lambda|<2\), \(|\lambda|=2\), or
\(|\lambda|>2\). The exceptional value \(\mu=-1/2\) removes the linear
term and produces the corresponding homogeneous degeneration.

For the selector \(c=X\) (\(\lambda=1,\mu=0\)), the original coordinates
obey

\[
x^2-xy+y^2+x-y+\frac14=0
\]

The center is \((-1/3,1/3)\). With

\[
\xi=x+\frac13,\qquad
\eta=y-\frac13,\qquad
u=\frac{\xi+\eta}{\sqrt2},\qquad
v=\frac{\xi-\eta}{\sqrt2}
\]

the ellipse is

\[
\boxed{
\frac{u^2}{1/6}+\frac{v^2}{1/18}=1
}
\]

Its foci are

\[
\left(-\frac13\pm\frac1{3\sqrt2},
\frac13\pm\frac1{3\sqrt2}\right)
\]

with matching signs. The selected circle has real intersections
precisely for \(-1/6\le c\le1/2\).

For \(c=2X\) one obtains the focus-directrix parabola

\[
\boxed{\left(x+\frac12-y\right)^2=\frac y2}
\]

with focus and directrix

\[
F=\left(-\frac9{16},\frac1{16}\right)
\qquad
L:\ x+y=-\frac58
\]

For \(c=3X\) the eliminated curve is the hyperbola

\[
\left(x+\frac12\right)^2
-3\left(x+\frac12\right)y+y^2-\frac12y=0
\]

Finally, the PP74 level circles

\[
\left(U+1+\frac1{2k}\right)^2+V^2=\frac1{4k^2}
\]

form the same focus-directrix center construction after a Euclidean
similarity: their centers lie on \(V=0\), every circle passes through
\((-1,0)\), and each is tangent there to \(U=-1\). This upgrades the old
visual comparison to an exact conic-pencil bridge. It remains geometric;
it supplies no zero-discrepancy estimate by itself.

\begin{figure}
\centering
\includegraphics{figures/conic_pencil_diagnostic.png}
\caption{Exact circle-pencil and selector-conic diagnostic.}
\end{figure}

\section{14. Retained exact operator
families}\label{retained-exact-operator-families}

The following families connect the centered coordinate to operator
models. Their scalar identities and operator domains have different
scopes.

\begin{longtable}[]{@{}
  >{\raggedright\arraybackslash}p{(\columnwidth - 4\tabcolsep) * \real{0.3333}}
  >{\raggedright\arraybackslash}p{(\columnwidth - 4\tabcolsep) * \real{0.3333}}
  >{\raggedright\arraybackslash}p{(\columnwidth - 4\tabcolsep) * \real{0.3333}}@{}}
\toprule\noalign{}
\begin{minipage}[b]{\linewidth}\raggedright
Family
\end{minipage} & \begin{minipage}[b]{\linewidth}\raggedright
Exact content
\end{minipage} & \begin{minipage}[b]{\linewidth}\raggedright
Boundary
\end{minipage} \\
\midrule\noalign{}
\endhead
\bottomrule\noalign{}
\endlastfoot
Higher-simplex products & Root-symmetric Gamma closure; order \(1/p\) &
Entire-type convention merits specialist review \\
Continuous Pascal row & Gamma interpolation, dimension raising, central
splitting & Scalar identities do not create an arithmetic sieve \\
Cube-sphere lift & Degree/Casimir and Boolean-noise/Poisson-semigroup
intertwiners & Different Hilbert spaces; norm multiplier retained \\
Gaussian-Laguerre & Explicit Fourier eigenmodes in harmonic sectors & No
eventual-sign theorem \\
Mellin-Fourier & Degree-resolved Gamma multiplier and fixed-line
unitarity & Maximal unbounded-operator domains not developed \\
Sphere/Bessel & Coordinate characteristic functions and Gram positivity
& Radial sign observables remain nonlinear \\
Log-Gamma/Laplace & Determinant amplitude as a characteristic function &
Probability law does not imply zeta-zero location \\
Phase-complement collar & Exact complement, Cayley, rapidity, and
\(\tau\) bookkeeping & Physical-source operators remain independent \\
\end{longtable}

\section{\texorpdfstring{77. Cube roots, \(A_3=D_3\), and the twelve-ray
theorem}{77. Cube roots, A\_3=D\_3, and the twelve-ray theorem}}\label{cube-roots-a_3d_3-and-the-twelve-ray-theorem}

The primary six-cardinal construction fixes the cone-figure cube

\[
Q_\pi=\{(\pm\pi/2,\pm\pi/2,\pm\pi/2)\}
\tag{77.1}
\]

Its twelve edge midpoints are

\[
C_\pi=
\{(\pm\pi/2,\pm\pi/2,0)\ \text{and permutations}\}
\tag{77.2}
\]

each of norm \(\pi/\sqrt2\). This is the cone cube, not the distinct
inward \(\sqrt3\) cube retained elsewhere in the geometry corpus.

Partition \(Q_\pi\) by the sign of \(xyz\). Each four-point class is a
regular tetrahedron, and the two classes are negatives of one another.
If \(v_1,\ldots,v_4\) are the vertices of either class, then

\[
\left\{\frac{v_j-v_i}{2}:i\ne j\right\}=C_\pi
\tag{77.3}
\]

After normalization by \(\pi/\sqrt2\), (77.2) becomes

\[
\boxed{
\mathcal E=\frac1{\sqrt2}
\{(\pm1,\pm1,0),(\pm1,0,\pm1),(0,\pm1,\pm1)\}
=\Phi(A_3)=\Phi(D_3)}
\tag{77.4}
\]

Thus the cube contains the twelve directed tetrahedral edge directions
and the cuboctahedral root system. Both parity tetrahedra generate the
identical set.

Choose the threefold axis and registered plane basis

\[
u=\frac{(1,1,1)}{\sqrt3},\qquad
e_1=\frac{(1,-1,0)}{\sqrt2},\qquad
e_2=\frac{(1,1,-2)}{\sqrt6}
\tag{77.5}
\]

Let \(P_u\) be orthogonal projection to \(u^\perp\), expressed in this
basis. Here \(u\) is the perspective or viewing ray. The symbol \(i\)
enters only after the image plane is identified with \(\mathbb C\) by
\(z(v)=(v\cdot e_1)+i(v\cdot e_2)\): multiplication by \(i\) is the
oriented quarter-turn inside the image, not the viewing ray itself.
Replacing \(i\) by \(-i\) reverses the handedness of the registered
view. Direct substitution in the twelve coordinates gives

\[
\boxed{
P_u(\mathcal E)
=\mu_6\ \sqcup\
\frac1{\sqrt3}\zeta_{12}\mu_6}
\tag{77.6}
\]

The six roots orthogonal to \(u\) form the outer radius-one hexagon. The
six remaining roots form the staggered radius-\(1/\sqrt3\) hexagon.
Their twelve rays have exactly the arguments

\[
0^\circ,30^\circ,\ldots,330^\circ
\tag{77.7}
\]

Hence radial normalization of the nonzero projected points gives
\(\mu_{12}\). The projection itself does not.

\begin{quote}
\textbf{Theorem 77.1 (radial impossibility).} No orthogonal projection
of \(\mathcal E\) onto a plane places all twelve projected points on one
circle.
\end{quote}

\textbf{Proof.} The normalized roots form the tight frame

\[
\sum_{e\in\mathcal E}ee^{\mathsf T}=4I
\tag{77.8}
\]

If projection along a unit normal \(n\) gave equal radii, then
\(1-|e\cdot n|^2\) would be constant. Taking the trace against
\(nn^{\mathsf T}\) in (77.8) forces

\[
|e\cdot n|^2=\frac13
\quad\text{for every }e
\tag{77.9}
\]

Comparing \((1,1,0)/\sqrt2\) with \((1,-1,0)/\sqrt2\) forces
\(n_1n_2=0\). The analogous pairs force \(n_1n_3=n_2n_3=0\). Thus \(n\)
is a coordinate axis. For a coordinate axis, however, \(|e\cdot n|^2\)
takes the two values \(0\) and \(1/2\), contradicting (77.9).
\(\square\)

There is also a group-theoretic obstruction. The full root automorphism
group is the octahedral group \(O_h\) of order \(48\), isomorphic to a
central \(C_2\) extension of \(S_4\). Its element orders divide
\(1,2,3,4,\) or \(6\); it has no element of order \(12\). Therefore the
geometry induces no intrinsic \(C_{12}\) action.

\begin{figure}
\centering
\includegraphics{figures/cube_a3_twelve_ray_projection_v1_2.png}
\caption{The cone cube yields the \(A_3\) roots and twelve exact rays,
but equal radii appear only after a separate radial normalization.}
\end{figure}

The strongest honest classification is therefore

\[
\boxed{\begin{gathered}
\text{geometric }A_3\text{ skeleton}
+\text{ exact twelve-ray angles}\\
+\text{ radial and cyclic no-go}
\end{gathered}}
\tag{77.10}
\]

The shared value \(1/\sqrt3\) in the inward cube cascade and in (77.6)
is a shared three-coordinate normalization. No intertwiner between the
two operators is known. The normalized cube-corner value \(\pm1/8\) and
the triangular completed-square correction \(-1/8\) have cubic and
quadratic homogeneity respectively; they remain an exact normalization
contact only.

\section{\texorpdfstring{78. Balanced stabilizers, factorial mid-shells,
and the intrinsic
\(D_8\)}{78. Balanced stabilizers, factorial mid-shells, and the intrinsic D\_8}}\label{balanced-stabilizers-factorial-mid-shells-and-the-intrinsic-d_8}

Let \(\mathcal F_{24}\) be the complete ordered tetrahedral frames.
Choosing a reference ordering identifies it with the regular
\(S_4\)-set. For the directed edge \((1,2)\) and the unoriented edge
\(\{1,2\}\), the stabilizers are

\[
H=\{1,(34)\}\cong C_2
\qquad
K=\langle(12),(34)\rangle\cong V_4
\tag{78.1}
\]

Therefore

\[
\boxed{
\mathcal F_{24}\cong S_4
\longrightarrow
E_{12}^{\to}\cong S_4/H
\longrightarrow
E_6\cong S_4/K}
\tag{78.2}
\]

is a balanced \(2:1\) then \(2:1\) tower of sets. It follows that

\[
\boxed{12^2=24\cdot6,\qquad12=\sqrt{3!4!}}
\tag{78.3}
\]

The owner of (78.3) is the stabilizer chain, not a factorial-induced
map.

More generally, if \(1<H<K<G\) are finite and \(|H|=[K:H]=r\), then
\(|K|=r^2\) and

\[
\boxed{|G/H|^2=|G|\,|G/K|}
\tag{78.4}
\]

For an \(n\)-simplex, the complete frames, directed edges, and edges
have sizes

\[
(n+1)!,\qquad n(n+1),\qquad\frac{n(n+1)}2
\]

The first fiber has size \((n-1)!\) and the second has size \(2\). The
tower is balanced binary only when \((n-1)!=2\), hence uniquely for
\(n=3\). This is the structural reason that the tetrahedral
\(24\to12\to6\) packet is exceptional.

\begin{figure}
\centering
\includegraphics{figures/balanced_stabilizer_tower_v1_2.png}
\caption{The middle value 12 is the exact cardinal geometric mean
because the tetrahedral frame tower has two equal stabilizer fibers.}
\end{figure}

The number \(12\) is not, however, the unique factorial mid-shell. From

\[
q^2=n!(n+1)!=(n!)^2(n+1)
\]

one obtains

\[
\boxed{q\in\mathbb Z\iff n+1=a^2,\qquad q=a\,n!}
\tag{78.5}
\]

This is an infinite family. The value \(12\) is uniquely the
\emph{doubling} member: \(q=2n!\) and \((n+1)!=2q\) force \(n+1=4\), so
\(n=3\) and \(q=12\). The Pell identity \(17^2-2\cdot12^2=1\) is exact,
but no recurrence intertwines the Pell orbit with (78.5); it remains a
one-hit shared scale.

\subsection{78.1 The four-state fiber}\label{the-four-state-fiber}

Over each unoriented edge sit four complete frames. Two intrinsic
involutions act on the fiber:

\[
A(i,j;k,l)=(j,i;k,l)
\qquad
B(i,j;k,l)=(i,j;l,k)
\tag{78.6}
\]

They commute, are fixed-point-free, and generate a regular
\(V_4=\langle A,B\rangle\).

The registered cyclotomic bijection \(\Psi:\mathcal F_{24}\to\mu_{24}\)
transports a quarter-turn \(L\) with

\[
L^2=B,\qquad L^4=1,\qquad ALA=L^{-1}
\tag{78.7}
\]

Although choosing \(L\) rather than \(L^{-1}\) requires an oriented
frame, the group it generates with \(A\) does not:

\[
\boxed{
\langle A,L\rangle
=C_{\operatorname{Sym}(\mathcal Q_{\bar e})}(B)
\cong D_8}
\tag{78.8}
\]

Indeed the centralizer of a fixed-point-free double transposition in
\(S_4\) has order eight. Both possible quarter-turn orientations
generate this same centralizer. Consequently the envelope \(D_8\) is
intrinsic, while the regular \(C_4=\langle L\rangle\) and its label
\(+i\) are registration-relative. On the same four points,

\[
V_4\cap C_4=\langle B\rangle\cong C_2
\qquad
\langle V_4,C_4\rangle=D_8
\tag{78.9}
\]

\begin{figure}
\centering
\includegraphics{figures/intrinsic_d8_envelope_v1_2.png}
\caption{The intrinsic \(V_4\) and transported \(C_4\) sit in one
intrinsic dihedral envelope; only the orientation of the quarter-turn is
registered.}
\end{figure}

With the determinant registration, write a top exponent as \(K=r+6q\),
\(q\in\mathbb Z/4\). Then

\[
L:q\mapsto q+1
\qquad
B:q\mapsto q+2
\qquad
A:q\mapsto-q-1
\tag{78.10}
\]

On roots, where \(y=z^4=\zeta_6^r\), endpoint reversal is the literal
fiber reflection

\[
\boxed{A(z)=-i\,\zeta_{12}^{r}\,\overline z}
\tag{78.11}
\]

The cyclotomic order is constant around the \(C_4\) cycle exactly when
\(r\) is odd. Geometrically these are the three apex edges incident to
the registered threefold-axis vertex; the even-\(r\) fibers are the
three base-face edges. The primitive eight states are exactly the two
complete fibers \(r=1,5\). This is cyclotomic primitivity, not integer
primality.

The tetrahedral tower remains a coset tower of sets. Since \(H\) is not
normal in \(S_4\), \(S_4/H\) is not a quotient group. Likewise
\(E_6\simeq A_4/\langle(12)(34)\rangle\) is a homogeneous space, not a
quotient group. The root-of-unity tower

\[
\mu_{24}\xrightarrow{z^2}\mu_{12}\xrightarrow{z^2}\mu_6
\]

is a genuine group tower. Their maximal unregistered common structure is
a free \(C_2\)-set; stagewise registered maps do not make the ambient
groups isomorphic.

The finite-rank statements below specify their certificate geometry and
mathematical grade. Sections 92--93 give the global result through order
eight; §138 gives the distinct rank-twelve result at one fixed
basepoint.

\Needspace{11\baselineskip}

\chapter{Technical chapter IV. Completion, folded support, and the Gamma
source}\label{technical-chapter-iv.-completion-folded-support-and-the-gamma-source}

Completion selects the particular resolvent. Moment and matrix criteria
apply to that function; the auxiliary Gamma scaffold is then compared
with it.

\section{15. Recognition is not
distribution}\label{recognition-is-not-distribution}

Wilson's theorem and the successor gate give exact finite prime
recognizers. To enter analytic number theory one must pass to prime
powers with the von Mangoldt measure

\[
\nu=\sum_{m\ge2}\Lambda(m)\,\delta_{\log m}
\]

whose Laplace transform is

\[
\mathcal L\nu(s)=-\frac{\zeta'(s)}{\zeta(s)}
\qquad \Re s=\sigma>1
\]

In the finite-visibility convention used here this means

\[
s=\sigma+it,\qquad \Re s=\sigma>1
\]

This closes the finite-to-prime adapter. It does not provide the
cancellation or positivity needed beyond \(\Re s=1\).

The polar/Ramanujan bridge is likewise exact:

\[
c_n(m)=\sum_{(k,n)=1}e^{2\pi i mk/n}
\qquad
\sum_{n\ge1}\frac{c_n(m)}{n^{s}}
=\frac{\sigma_{1-s}(m)}{\zeta(s)}
\]

The series identity holds for each fixed positive integer \(m\) and
\(\Re s>1\). For \(m=1\) it reaches \(1/\zeta(s)\) directly. It does not
show analyticity in \(\Re s=\sigma>1/2\).

\section{16. Completion and folded
coordinate}\label{completion-and-folded-coordinate}

Use the completed entire function

\[
\xi(s)=\frac12s(s-1)
\pi^{-s/2}\Gamma(s/2)\zeta(s)
\qquad
\xi(s)=\xi(1-s)
\]

Here and in every completion formula,

\[
\boxed{
s=\sigma+it,\qquad
1-s=(1-\sigma)-it,\qquad
\bar s=\sigma-it}
\tag{16.0}
\]

Put

\[
q(s)=s(1-s)
\]

Functional-equation invariance gives a unique entire function \(X\) such
that

\[
\xi(s)=X(q)
\]

For a zero \(\rho=\beta+i\gamma\),

\[
q_\rho=\rho(1-\rho)
=\frac14+\gamma^2-(\beta-\tfrac12)^2
-2i\gamma(\beta-\tfrac12)
\]

All folded zeros satisfy \(\Re q_\rho>0\) unconditionally. RH is the
stronger statement

\[
\boxed{q_\rho\in\mathbb R_{>0}\quad\text{for every zero orbit}}
\]

Define the folded resolvent

\[
M_\xi(q)=-\frac{X'(q)}{X(q)}
\]

On the negative real axis, let

\[
s(x)=\frac{1+\sqrt{1+4x}}2>1
\qquad
S_\xi(x)=M_\xi(-x)
\]

Thus the source evaluation is the specialization

\[
s(x)=\sigma_x+i0,\qquad
\sigma_x=\frac{1+\sqrt{1+4x}}2>1,\qquad
x=s(x)(s(x)-1)=2T_{s(x)-1}
\tag{16.1a}
\]

Then

\[
S_\xi(x)=\frac1{\sqrt{1+4x}}
\frac{\xi'(s(x))}{\xi(s(x))}
\]

The \(s>1\) prime/archimedean formula is

\[
\boxed{
S_\xi(x)=\frac1{\sqrt{1+4x}}
\left[
\frac1s+\frac1{s-1}-\frac12\log\pi
+\frac12\psi(s/2)
-\sum_{n\ge2}\frac{\Lambda(n)}{n^s}
\right]
}
\]

\section{\texorpdfstring{152. Full \(s=\sigma+it\) expansion and the
triangular
center}{152. Full s=\textbackslash sigma+it expansion and the triangular center}}\label{full-ssigmait-expansion-and-the-triangular-center}

Starting from \(s=\sigma+it\), direct multiplication gives the entire
coordinate chart:

\begin{longtable}[]{@{}
  >{\raggedright\arraybackslash}p{(\columnwidth - 4\tabcolsep) * \real{0.3333}}
  >{\raggedright\arraybackslash}p{(\columnwidth - 4\tabcolsep) * \real{0.3333}}
  >{\raggedright\arraybackslash}p{(\columnwidth - 4\tabcolsep) * \real{0.3333}}@{}}
\toprule\noalign{}
\begin{minipage}[b]{\linewidth}\raggedright
object
\end{minipage} & \begin{minipage}[b]{\linewidth}\raggedright
Cartesian expansion
\end{minipage} & \begin{minipage}[b]{\linewidth}\raggedright
critical-line specialization \(\sigma=1/2\)
\end{minipage} \\
\midrule\noalign{}
\endhead
\bottomrule\noalign{}
\endlastfoot
\(s\) & \(\sigma+it\) & \(1/2+it\) \\
\(1-s\) & \((1-\sigma)-it\) & \(1/2-it=\overline{s}\) \\
\(\overline{s}\) & \(\sigma-it\) & \(1/2-it=1-s\) \\
\(s^{\,2}\) & \((\sigma^2-t^2)+2i\sigma t\) & \((1/4-t^2)+it\) \\
\(q=s(1-s)\) & \(\sigma(1-\sigma)+t^2+i\,t(1-2\sigma)\) &
\(1/4+t^2\in\mathbb R_{>0}\) \\
\(x=s(s-1)=-q\) & \(\sigma^2-\sigma-t^2+i\,t(2\sigma-1)\) &
\(-(1/4+t^2)\) \\
\(T_{s-1}=s(s-1)/2\) & \((\sigma^2-\sigma-t^2)/2+i\,t(\sigma-1/2)\) &
\(-1/8-t^2/2\) \\
centered coordinate \(j=2s-1\) & \((2\sigma-1)+2it\) & \(2it\) \\
\end{longtable}

Thus

\[
\boxed{
\begin{aligned}
q(s)
&=(\sigma+it)(1-\sigma-it)\\
&=\sigma(1-\sigma)+t^2+i\,t(1-2\sigma)\\[1mm]
T_{s-1}
&=\frac{\sigma^2-\sigma-t^2}{2}
+i\,t\!\left(\sigma-\frac12\right)
\end{aligned}}
\tag{152.1}
\]

Two identities hold without placing \(s\) on any line:

\[
\boxed{
4q(s)=1-(2s-1)^2,\qquad
8T_{s-1}+1=(2s-1)^2}
\tag{152.2}
\]

The same centered coordinate therefore owns the fold and the triangular
completion. On the real safe axis put \(s_x=n+1\) with \(n\ge0\). Then

\[
\boxed{\begin{aligned}
s_x&=n+1=\sigma_x+i0\\
x&=s_x(s_x-1)=n(n+1)=2T_n\\
R&=2s_x-1=2n+1=j_n=\sqrt{1+4x}
\end{aligned}}
\tag{152.3}
\]

This is the analytic version of \(8T_n+1=(2n+1)^2\). It does
\textbf{not} put a nontrivial zeta zero at a triangular point: on the
critical line \(x=-q<0\), whereas the source half-line uses \(x>0\) and
\(t=0\). The identity is a shared coordinate chart, not a zero-location
theorem.

For a zero \(\rho=\beta+i\gamma\), (152.1) becomes

\[
\boxed{
q_\rho=\rho(1-\rho)
=\frac14+\gamma^2-\left(\beta-\frac12\right)^2
-2i\gamma\left(\beta-\frac12\right)}
\tag{152.4}
\]

Consequently \(q_\rho\) is real for a nonreal zero exactly when
\(\beta=1/2\). The omitted \(\gamma=0\) branch is the reason the
informal statement \(\Im q=0\iff\sigma=1/2\) is false for arbitrary
\(s\); the exact real-axis distinction is

\[
\Im q=t(1-2\sigma)=0
\iff t=0\ \text{or}\ \sigma=\frac12
\tag{152.5}
\]

The coordinate atlas in Reading Volume §R6 displays this same fold.

\section{\texorpdfstring{153. The self-dual Mellin weight
\(n^{-1/2}=1/\sqrt n\)}{153. The self-dual Mellin weight n\^{}\{-1/2\}=1/\textbackslash sqrt n}}\label{the-self-dual-mellin-weight-n-121sqrt-n}

For \(n>0\) and \(s=\sigma+it\),

\[
\boxed{
n^{-s}
=n^{-(\sigma+it)}
=n^{-\sigma}e^{-it\log n}
=\frac{\cos(t\log n)-i\sin(t\log n)}{n^\sigma}}
\tag{153.1}
\]

The modulus and phase separate:

\[
|n^{-s}|=n^{-\sigma},\qquad
|n^{-it}|=1,\qquad
\arg(n^{-s})=-t\log n\pmod{2\pi}
\tag{153.2}
\]

Reflection and conjugation give two different products:

\begin{longtable}[]{@{}
  >{\raggedright\arraybackslash}p{(\columnwidth - 4\tabcolsep) * \real{0.3333}}
  >{\raggedright\arraybackslash}p{(\columnwidth - 4\tabcolsep) * \real{0.3333}}
  >{\raggedright\arraybackslash}p{(\columnwidth - 4\tabcolsep) * \real{0.3333}}@{}}
\toprule\noalign{}
\begin{minipage}[b]{\linewidth}\raggedright
pairing
\end{minipage} & \begin{minipage}[b]{\linewidth}\raggedright
expanded calculation
\end{minipage} & \begin{minipage}[b]{\linewidth}\raggedright
result
\end{minipage} \\
\midrule\noalign{}
\endhead
\bottomrule\noalign{}
\endlastfoot
functional-equation reflection &
\(n^{-s}n^{-(1-s)}=n^{-(\sigma+it)}n^{-((1-\sigma)-it)}\) &
\(n^{-1}=1/n\) for every \(\sigma,t\) \\
Hermitian conjugation &
\(n^{-s}\overline{n^{-s}}=n^{-(\sigma+it)}n^{-(\sigma-it)}\) &
\(n^{-2\sigma}\) \\
\end{longtable}

They agree for every \(n>1\) exactly when \(2\sigma=1\). Equivalently,

\[
\boxed{
|n^{-s}|=|n^{-(1-s)}|
\iff \sigma=\frac12
\qquad n^{-1/2}=\frac1{\sqrt n}}
\tag{153.3}
\]

Thus \(1/\sqrt n\) is not an arbitrary normalization in the explicit
formula. It is the unique pointwise magnitude at which the two
functional-equation partners carry the same length. On the critical
line,

\[
n^{-(1/2+it)}=\frac{e^{-it\log n}}{\sqrt n},\qquad
n^{-(1/2-it)}=\frac{e^{it\log n}}{\sqrt n}
=\overline{n^{-(1/2+it)}}
\tag{153.4}
\]

The resolvent carrier makes the same equality visible. With
\(a_x=\sqrt{x+1/4}\) and
\(f_x(u)=e^{-a_x|u|}(1/(4a_x)+\tfrac12\operatorname{sgn}u)\), declare
\(\widehat f_x(t)=\int_{\mathbb R}f_x(u)e^{-itu}\,du\). The sign of the
transform is part of the definition. For real \(t\),

\[
\boxed{
\widehat f_x(t)
=\frac{1/2-it}{t^2+a_x^2}
=\left.\frac{1-s}{x+q(s)}
\right|_{s=1/2+it},\qquad
|\widehat f_x(t)|^2=\frac{q}{(x+q)^2}}
\tag{153.5}
\]

The opposite convention gives the conjugate numerator:

\[
\int_{\mathbb R}f_x(u)e^{+itu}\,du
=\frac{1/2+it}{t^2+a_x^2}
=\overline{\widehat f_x(t)}
\tag{153.5a}
\]

The modulus-square equality in (153.5) uses \(1-s=\overline{s}\) and
\(q=s(1-s)=s\overline{s}\) on the line. Off the line, reflection and
conjugation separate, so replacing the holomorphic square by a modulus
square would assume the desired zero placement.

The rotation diagram in Reading Volume §R8 displays the conjugate
phases.

The Gauss-circle connection is exact only at the level of a coordinate
and a weight. For a finite disk and
\(r_2(n)=\#\{(a,b)\in\mathbb Z^2:a^2+b^2=n\}\),

\[
\sum_{0<a^2+b^2\le R^2}(a^2+b^2)^{-w}
=\sum_{1\le n\le R^2}r_2(n)n^{-w}
\tag{153.6}
\]

At \(w=1/2\), each lattice vector has weight
\(n^{-1/2}=1/\sqrt{a^2+b^2}=1/r\). This is a useful visual and Mellin
comparison, but the multiplicity \(r_2(n)\) changes the arithmetic
object. No Gauss-circle estimate, Dirichlet-divisor estimate, or RH
implication is transported by (153.6).

\begin{figure}
\centering
\includegraphics{figures/figure79_sqrt_weight_and_gauss_circle_comparison.png}
\caption{The square-root weight and its lattice-radius reading. The
right panel is a comparison model, not an identification of proof
mechanisms.}
\end{figure}

\section{17. Stieltjes, heat, and moment
equivalences}\label{stieltjes-heat-and-moment-equivalences}

The folded moment chain is

\[
\begin{aligned}
\mathrm{RH}
&\iff (\mu_k)\text{ is a Hamburger moment sequence}\\
&\iff [\mu_{i+j}]_{i,j=0}^N\succeq0\quad\forall N\\
&\iff M_\xi\text{ is Herglotz-Pick}\\
&\iff S_\xi\text{ is Stieltjes}\\
&\iff \Theta_\xi\text{ is completely monotone}
\end{aligned}
\]

The folded heat trace is

\[
\Theta_\xi(u)=\sum_{[\rho]}m_\rho e^{-u q_\rho}
\qquad
S_\xi(x)=\int_0^\infty e^{-xu}\Theta_\xi(u)\,du
\]

The central-Pascal change of basis between folded moments and Li
coefficients is (155.3).

Using Gaussian subordination, the prime contribution to the heat trace
becomes the explicit log-Gaussian comb

\[
-\frac{e^{-u/4}}{2\sqrt{\pi u}}
\sum_{n\ge2}\frac{\Lambda(n)}{\sqrt n}
\exp\!\left[-\frac{(\log n)^2}{4u}\right]
\]

balanced by an explicit Gamma term.

The full load-bearing PP1 theorem remains:

\begin{quote}
Starting only from the \(s>1\) prime/archimedean formula, prove that
\(S_\xi\) is Stieltjes, that \(\Theta_\xi\) is completely monotone, or
that every prime-side Loewner matrix is positive semidefinite.
\end{quote}

The first six reduced-Widder diagonal inequalities are source-certified
in §§35--38 and 74--74A; \(W_1\) is analytic and \(W_2\)--\(W_6\) are
A-CAT from one bundled certificate. Reading Volume §R12 states the
controlling rung status. This is genuine finite-order progress inside
the criterion, but it does not establish the complete hierarchy required
for the Stieltjes class. Ordinary positivity of \(S_\xi(x)\) is
unconditional and is not the target.

\begin{figure}
\centering
\includegraphics{figures/stieltjes_heat_model.png}
\caption{Finite positive folded spectrum illustrating Stieltjes
derivative signs and heat decay.}
\end{figure}

\section{18. Li/transverse defect
identity}\label{litransverse-defect-identity}

For \(s=1/2+u+iv\),

\[
s(1-s)=\frac14+v^2-u^2-2iuv
\]

Thus the adjacent triangular polarization \(P(u,v)=2uv\) is exactly the
imaginary folded-zero defect. If \(C_\xi e_\rho=q_\rho^{-1}e_\rho\),
then

\[
\mathcal E_\xi
=\frac14\|C_\xi-C_\xi^*\|_{\mathrm{HS}}^2
=\sum_{[\rho]}m_\rho
\frac{P(\beta-\tfrac12,\gamma)^2}{|q_\rho|^4}
\]

Consequently

\[
\boxed{
\mathrm{RH}
\iff
\sum_{[\rho]}\frac{m_\rho}{|\rho(1-\rho)|^2}
=4\lambda_1-\lambda_2
}
\]

This identifies the missing object as a nonholomorphic absolute-square
norm. The arithmetic construction of that norm remains open.

\begin{figure}
\centering
\includegraphics{figures/transverse_defect_geometry.png}
\caption{The folded imaginary coordinate is exactly the transverse
critical-line defect.}
\end{figure}

\section{155. Li positivity, PP97, and the Löwner source
gate}\label{li-positivity-pp97-and-the-luxf6wner-source-gate}

To keep the analytic variable visible, define the Li coefficients with a
dummy differentiation variable \(z\):

\[
\boxed{
\lambda_n=
\left.\frac1{(n-1)!}\frac{d^n}{dz^n}
\left[z^{n-1}\log\xi(z)\right]\right|_{z=1}}
\tag{155.1}
\]

Li's criterion is

\[
\boxed{
\mathrm{RH}\iff \lambda_n\ge0\quad\text{for every }n\ge1}
\tag{155.2}
\]

This is an all-order criterion. A long positive finite prefix is not the
quantified statement in (155.2). The log-frequency notation
\(L_n=\log n\) in the Montgomery--Vaughan calculation is unrelated to
\(\lambda_n\).

The folded moments and Li coefficients are joined by the exact
central-Pascal transform

\[
\boxed{
\mu_{k-1}=\sum_{j=1}^{k}(-1)^{j+1}
\binom{2k}{k-j}\lambda_j}
\tag{155.3}
\]

The first rows, printed rather than hidden inside the binomial, are

\[
\begin{aligned}
\mu_0&=\lambda_1\\
\mu_1&=4\lambda_1-\lambda_2\\
\mu_2&=15\lambda_1-6\lambda_2+\lambda_3\\
\mu_3&=56\lambda_1-28\lambda_2+8\lambda_3-\lambda_4\\
\mu_4&=210\lambda_1-120\lambda_2+45\lambda_3-10\lambda_4+\lambda_5
\end{aligned}
\tag{155.4}
\]

The second row is the particularly important bridge:

\[
\boxed{
\mu_1=4\lambda_1-\lambda_2}
\tag{155.5}
\]

On the folded-zero side this is holomorphic data in
\(q_\rho=\rho(1-\rho)\). The exact RH-equivalent equality

\[
\boxed{
\sum_{[\rho]}\frac{m_\rho}{|\rho(1-\rho)|^2}
=4\lambda_1-\lambda_2}
\tag{155.6}
\]

compares a nonholomorphic left side with the holomorphic folded moment
on the right. A proof of (155.6) must create conjugation through a
two-point, sesquilinear, or variational mechanism; scalar holomorphic
bookkeeping cannot do so by itself.

Now define

\[
S_\xi(x)=\sum_{[\rho]}\frac{m_\rho}{x+q_\rho},\qquad
h(x)=xS_\xi(x),\qquad x>0
\tag{155.7}
\]

Under RH, \(q_\rho>0\) and \(S_\xi\) is Stieltjes, so \(h\) is complete
Bernstein and operator monotone. Conversely, the Pick continuation of
\(h\) excludes nonreal folded poles and forces RH. For positive nodes
\(x_i\) define

\[
[\mathcal L_h]_{ij}=
\begin{cases}
\dfrac{h(x_i)-h(x_j)}{x_i-x_j},&i\ne j\\[2mm]
h'(x_i),&i=j
\end{cases}
\tag{155.8}
\]

\begin{figure}
\centering
\includegraphics{figures/v26c_two_node_compatibility.png}
\caption{A two-node Löwner matrix tests off-diagonal compatibility as
well as scalar signs. With \(K=[h(x)-h(y)]/(x-y)\), positive
semidefiniteness requires \(K^2\le W_1(x)W_1(y)\). The global
order-eight theorem already includes this two-node test. The all-size
quantifier below is the remaining requirement.}
\end{figure}

Then

\[
\boxed{
\mathrm{RH}\iff \mathcal L_h(x_1,\ldots,x_N)\succeq0
\quad\text{for every }N\text{ and every positive node set}}
\tag{155.9}
\]

On the safe source sheet

\[
s_x=\sigma_x+i0=\frac{1+\sqrt{1+4x}}2>1,\qquad
R=2s_x-1=\sqrt{1+4x}
\tag{155.10}
\]

the explicit formula splits

\[
S_\xi(x)=\frac1x+G(x)-P(x)
\tag{155.11}
\]

with

\[
G(x)=\frac{\psi((1+R)/4)-\log\pi}{2R},\qquad
P(x)=\frac1R\sum_{n\ge2}\frac{\Lambda_{\mathrm{vM}}(n)}{n^{s_x}}
\tag{155.12}
\]

Because constants disappear from divided differences,

\[
\boxed{\mathcal L_h=L_\Gamma-L_P}
\tag{155.13}
\]

This is PP97's exact source gate. The \(1\times1\) diagonal is

\[
h'(x)=S_\xi(x)+xS_\xi'(x)=W_1(x)>0
\tag{155.14}
\]

but the full statement is

\[
\boxed{L_\Gamma-L_P\succeq0
\quad\text{for every rank and every positive node set}}
\tag{155.15}
\]

Equation (155.15) is equivalent to RH and remains unproved. Its
diagonal, finite-rank, and finite-rung shadows must be kept separate:

\begin{longtable}[]{@{}
  >{\raggedright\arraybackslash}p{(\columnwidth - 4\tabcolsep) * \real{0.3333}}
  >{\raggedright\arraybackslash}p{(\columnwidth - 4\tabcolsep) * \real{0.3333}}
  >{\raggedright\arraybackslash}p{(\columnwidth - 4\tabcolsep) * \real{0.3333}}@{}}
\toprule\noalign{}
\begin{minipage}[b]{\linewidth}\raggedright
layer
\end{minipage} & \begin{minipage}[b]{\linewidth}\raggedright
exact content
\end{minipage} & \begin{minipage}[b]{\linewidth}\raggedright
present authority
\end{minipage} \\
\midrule\noalign{}
\endhead
\bottomrule\noalign{}
\endlastfoot
Li sequence & \(\lambda_n\ge0\) for all \(n\) & RH-equivalent; all-order
sign open \\
folded moment & \([\mu_{i+j}]\succeq0\) at all sizes & RH-equivalent;
finite minors do not close it \\
Widder scalar & \(W_k(x)\ge0\) for every \(k,x\) & \(W_1\) analytic;
\(W_2\)--\(W_6\) certified from the source (A-CAT); source-side
\(W_7\rightsquigarrow F^{(14)}\) and the hierarchy open \\
local Dobsch & derivative matrices at one \(x_0\) and all ranks &
all-rank statement RH-equivalent \\
global Löwner & all positive node sets & ranks through eight certified;
rank nine/all ranks open \\
source split & \(\mathcal L_h=L_\Gamma-L_P\) & identity exact; dominance
(155.15) open \\
\end{longtable}

Reader §§R15--R16 display the Li--moment--matrix dictionary. The proof
map in Reading Volume §R24 places its all-size positivity condition
alongside the other RH criteria.

The global owner is the completed matrix \(\mathcal L_h\). Its
source-side Gamma and prime pieces are separately signed and nearly
cancel. Positive translation-defect representations retain a negative
scalar anchor, while compression to the scalar symbol in §159 loses
off-line reflection data. These are concrete restrictions on a source
proof, not a change of target.

\section{155A. Reflected Li forms and explicit positive-node
witnesses}\label{a.-reflected-li-forms-and-explicit-positive-node-witnesses}

Put \(y=q^{-1}\) and \(w=1-1/\rho\). Reflection gives
\(w_{1-\rho}=w_\rho^{-1}\) and \(y=2-w-w^{-1}\); conjugation is a
different operation off the critical line. Define the polynomials
\(Q_n(y)=2-w^n-w^{-n}\) and \(\mathcal P_m(y)=\sum_{r=-m}^{m}w^r\). They
satisfy

\[
\mathcal P_0=1,\quad \mathcal P_1=3-y,\quad
\mathcal P_{m+1}=(2-y)\mathcal P_m-\mathcal P_{m-1}.
\tag{155A.1}
\]

The calligraphic symbol distinguishes these polynomials from the
reciprocal polynomials of Section 171 and the source rows of Section
163.

\textbf{Proposition 155A.1 (triangular coefficients and reflected
squares).} For \(m\ge0\),

\[
\begin{aligned}
\mathcal P_m(y)&=\sum_{j=0}^m d_{m,j}y^j,\\
d_{m,j}&=(-1)^j(2m+1)\frac{(m+j)!}{(m-j)!(2j+1)!},\\
\mathcal P_m(y)&=(2m+1)\sum_{j=0}^m
 \frac{(-2y)^j}{(2j+1)!}\prod_{r=0}^{j-1}(T_m-T_r),\\
Q_{2m+1}(y)&=y\mathcal P_m(y)^2,\\
y\mathcal P_r(y)\mathcal P_s(y)&=Q_{r+s+1}(y)-Q_{|r-s|}(y).
\end{aligned}
\tag{155A.2}
\]

\textbf{Proof.} The coefficient formula satisfies (155A.1) and its two
initial conditions. The product formula follows from
\((m+j)!/(m-j)!=\prod_{r<j}(m-r)(m+r+1)=2^j\prod_{r<j}(T_m-T_r)\).
Multiply \(1-w^{2m+1}=(1-w)w^m\mathcal P_m\) by its reciprocal-reflected
version to get the square. For the cross identity, put
\(w=e^{i\theta}\): \(\mathcal P_m=\sin((m+1/2)\theta)/\sin(\theta/2)\),
and use the sine product identity. Polynomial identity extends the
result from \(0<y<4\) to all complex \(y\). \(\square\)

For the actual folded functional \(\mathcal L[y^j]=\mu_j\), the
fixed-degree sums converge absolutely because \(\sum_qm_q/|q|<\infty\).
Set \(H_N=[\mu_{i+j}]_{0\le i,j<N}\) and
\(\mathsf C_N=[d_{i,j}]_{0\le i,j<N}\), with zero entries above the
diagonal. Then

\[
\boxed{K_{r,s}=\lambda_{r+s+1}-\lambda_{|r-s|},\qquad
 K_N=\mathsf C_NH_N\mathsf C_N^{\mathsf T},\qquad
 K_{m,m}=\lambda_{2m+1}.}
\tag{155A.3}
\]

Here \(\lambda_0=0\) and \(\det\mathsf C_N=(-1)^{N(N-1)/2}\), so the
leading determinants and inertias agree. In particular
\(\lambda_3=9\mu_0-6\mu_1+\mu_2\) and
\(\det K_2=\lambda_1\lambda_3-(\lambda_2-\lambda_1)^2
=\mu_0\mu_2-\mu_1^2\). The diagonal \(\lambda_{2m+1}\) already mixes the
original off-diagonal data. These are holomorphic squares, not absolute
squares; positive individual moments do not establish their signs.

\textbf{The even counterpart.} Put
\(\mathcal B_r(y)=\sum_{j=0}^r(-1)^j\binom{r+j+1}{2j+1}y^j
=U_r(1-y/2)\). The same sine calculation gives

\[
\begin{aligned}
y(4-y)\mathcal B_r(y)\mathcal B_s(y)
 &=Q_{r+s+2}(y)-Q_{|r-s|}(y),\\
J_{r,s}&=\lambda_{r+s+2}-\lambda_{|r-s|},\\
J_N&=\mathsf B_N[4\mu_{i+j}-\mu_{i+j+1}]\mathsf B_N^{\mathsf T},
\qquad J_{r,r}=\lambda_{2r+2}.
\end{aligned}
\tag{155A.4}
\]

The coefficient matrix \(\mathsf B_N\) is triangular with diagonal
\((-1)^r\). The extra weight \(4-y\) distinguishes the even and odd
forms.

\textbf{The diagonal cases.} Putting \(r=s=m\) in (155A.2) and
\(r=s=m-1\) in (155A.4), and using \(Q_0=0\), prints the square form of
every reflected rung in one place:

\[
\boxed{Q_{2m+1}(y)=y\,\mathcal P_m(y)^2,
\qquad
Q_{2m}(y)=y(4-y)\,\mathcal B_{m-1}(y)^2}
\tag{155A.4a}
\]

The odd case is the fourth line of (155A.2); the even case is a
specialization of (155A.4) and not an independent result. Both weights
are nonnegative exactly on \(0\le y\le4\), the interval that Reading
Volume (10C.6) identifies with the critical-line condition for the orbit
carrying \(y\). The first rows are \(Q_1=y\), \(Q_2=y(4-y)\),
\(Q_3=y(3-y)^2\), \(Q_4=y(4-y)(2-y)^2\) and \(Q_5=y(y^2-5y+5)^2\).
Reading Volume §R10C states the same pair as the endpoint of the
pole-to-Li dictionary; nothing here constrains where the zeros lie,
since \(\lambda_n\) sums \(Q_n\) over all orbits.

\textbf{Theorem 155A.2 (the actual odd-subsequence criterion).} For the
actual Riemann function, RH is equivalent to eventual nonnegativity of
\(\lambda_{2m+1}\), and hence also to nonnegativity at every odd index.

\textbf{Proof.} Write \(F(z)=\xi(1/(1-z))/\xi(1)\) and \(D=F'/F\). The
parity projection
\(A(z)=(D(z)+D(-z))/2=\sum_{m\ge0}\lambda_{2m+1}z^{2m}\) is analytic
near every real \(0\le z<1\), since \(\xi\) has no real zeros. Eventual
nonnegativity, after subtracting a polynomial, and Pringsheim's theorem
force its Taylor radius to be at least one. If RH fails, a zero with
\(\Re\rho>1/2\) gives a pole at \(z_\rho=1-1/\rho\) inside the disk. Its
cancellation by \(D(-z)\) would require a zero at \(\rho/(2\rho-1)\),
but

\[
\left|\Im\frac{\rho}{2\rho-1}\right|
=\frac{|\Im\rho|}{(2\Re\rho-1)^2+4(\Im\rho)^2}
\le\frac1{4|\Im\rho|}<\frac14.
\]

This contradicts the low-height zero exclusion already implicit in the
volume's first-zero input: every nontrivial zero has \(|\Im\rho|>1\).
Thus the pole survives. Under RH, (155A.2) gives the converse term by
term. No new low-height certificate is being claimed. \(\square\)

A stronger version will be useful for a source bound. If RH fails, let
\(r_0<1\) be the least pole modulus of \(D\). For a fixed \(R>r_0\)
before the next poles, subtraction of the finitely many poles \(a\) on
\(|a|=r_0\) gives

\[
\lambda_n=-\sum_{|a|=r_0}m_a a^{-n}+O(R^{-n}).
\tag{155A.5}
\]

All these actual poles have \(\Re a>0\), by the same low-height
exclusion. Simultaneous recurrence of the finitely many phases
\(2m\arg a\) to zero on their compact torus gives arbitrarily large
\(m\) with \(\sum m_a\cos((2m+1)\arg a)\) bounded below by a positive
constant. (Recurrence follows by the pigeonhole principle on equal small
torus boxes.) Consequently, for some \(c>0\),
\(\lambda_{2m+1}\le-c r_0^{-(2m+1)}\) infinitely often. Tied dominant
poles have been retained, not replaced by one preferred zero.

\textbf{Proposition 155A.3 (a constructive source witness).} For fixed
\(m\), set \(x_i=(i+1)\varepsilon\), \(0\le i\le m\), and

\[
c_i(\varepsilon)=\frac{(-1)^i}{i!}
 \sum_{r=i}^m\frac{d_{m,r}}{(r-i)!\,\varepsilon^r}.
\qquad
\boxed{c(\varepsilon)^{\mathsf T}(L_\Gamma-L_P)c(\varepsilon)
 \longrightarrow\lambda_{2m+1}.}
\tag{155A.6}
\]

\textbf{Proof.} The elementary partial-fraction identity gives

\[
\sum_{i=0}^m\frac{c_i(\varepsilon)}{q+x_i}
=\sum_{r=0}^m\frac{d_{m,r}}{\prod_{j=0}^r(q+(j+1)\varepsilon)}
\longrightarrow q^{-1}\mathcal P_m(q^{-1}).
\]

The exact divided difference in (155.8) is
\(\sum_qm_qq/((x_i+q)(x_j+q))\). Thus its quadratic form tends to
\(\sum_qm_qq^{-1}\mathcal P_m(q^{-1})^2\). Since \(\Re q>0\),
\(|q+x_i|\ge|q|\); the summands for fixed \(m\) are bounded by
\(C_m m_q/|q|\), using the positive minimum of \(|q|\). Dominated
convergence proves (155A.6). \(\square\)

A negative odd value therefore gives negative positive-node tests inside
every neighborhood of zero. This is a necessary consequence, not a
contradiction supplied by the source. At \(m=8\) it tests one rank-nine
quadratic form involving moments through \(\mu_{16}\), not every
rank-nine minor. Large coalescing coefficients require their own
numerical error control.

\section{155B. The regularized Li source and its triangular
curvature}\label{b.-the-regularized-li-source-and-its-triangular-curvature}

The preceding identities specify a source inequality at the boundary.
Its regularization can be stated explicitly. Put

\[
\Psi(X)=\sum_{\ell\le X}\Lambda_{\mathrm{vM}}(\ell),\quad
M(u)=\Psi(e^u)-e^u+1,\quad
-\frac{\zeta'}{\zeta}(1+v)=\frac1v+\sum_{j\ge0}\eta_jv^j.
\tag{155B.1}
\]

In particular \(\eta_0=-\gamma_E\). These are coefficients of the local
regular part, not an unregularized Euler sum at \(s=1\).

\textbf{Theorem 155B.1 (reserve and signed prime remainder).} For every
\(n\ge1\),

\[
\begin{aligned}
\mathcal A_n&=n\lambda_1+
 \sum_{\substack{\ell\ge3\\\ell\ {\rm odd}}}
 \left[(1-1/\ell)^n-1+n/\ell\right],\\
\mathcal E_n&=\sum_{j=2}^n\binom nj\eta_{j-1}
 =\int_0^\infty e^{-u}M(u)\left[L_{n-1}^{(2)}(u)-n\right]du,\\
\lambda_n&=\mathcal A_n-\mathcal E_n,\qquad
\mathcal A_n\ge n\lambda_1>0,\qquad
\mathcal A_n=\tfrac12n\log n+O(n).
\end{aligned}
\tag{155B.2}
\]

The symbols \(\mathcal A_n,\mathcal E_n\) here do not denote the
cutoff-row norm \(A_N\) or the energy \(\mathscr E_N\). Nor is this
renormalized scalar reserve the separately signed matrix \(L_\Gamma\).

\textbf{Proof of the split and reserve.} For analytic \(g\) at one the
Li operation is \(\sum_{j=1}^n\binom nj g^{(j)}(1)/(j-1)!\). Apply it to
\(\log s+\log((s-1)\zeta(s))-s\log\pi/2+\log\Gamma(s/2)\). The polygamma
series and the regular part in (155B.1) give

\[
\lambda_n=1-\frac n2(\gamma_E+\log(4\pi))
 +\sum_{j=2}^n(-1)^j\binom nj(1-2^{-j})\zeta(j)
 -\sum_{j=1}^n\binom nj\eta_{j-1}.
\tag{155B.3}
\]

Separate \(\eta_0\) and expand \((1-2^{-j})\zeta(j)\) over odd positive
integers. The \(\ell=1\) sum is \(n-1\); the binomial theorem gives
(155B.2) for the remaining denominators. For \(0\le z\le1\) and
\(n\ge2\),

\[
0\le(1-z)^n-1+nz
=n(n-1)\int_0^z(z-v)(1-v)^{n-2}dv
\le\min\{nz,n(n-1)z^2/2\}.
\]

This proves convergence and positivity. Splitting at \(\ell=n\) and
using the odd harmonic sum yields
\(n\lambda_1\le\mathcal A_n\le n\lambda_1+(n/2)\log n+(n-1)/2\). The
lower bound \(n/\ell-1\) on the initial segment supplies the matching
asymptotic. The case \(n=1\) is immediate.

\textbf{Why the boundary integral exists.} The classical
prime-number-theorem error {[}3{]} supplies constants \(C,a>0\) with
\(|M(u)|\le C e^{u-a\sqrt u}\), after enlargement on a compact interval.
This is an unconditional existence bound, not a numerical interval
certificate. It justifies every fixed-degree integral and endpoint limit
below. For \(v>0\),

\[
\int_0^\infty e^{-(1+v)u}dM(u)
=-\frac{\zeta'}{\zeta}(1+v)-\frac1v.
\]

Use \(L_{n-1}^{(1)}(u)=\sum_{j=1}^n\binom nj(-u)^{j-1}/(j-1)!\) and
\((e^{-u}L_{n-1}^{(1)}(u))'=-e^{-u}L_{n-1}^{(2)}(u)\). Integration by
parts, followed by right differentiation at \(v=0\), gives

\[
\begin{aligned}
\lim_{T\to\infty}\left[
 \sum_{\ell\le e^T}\frac{\Lambda_{\mathrm{vM}}(\ell)}{\ell}
 L_{n-1}^{(1)}(\log\ell)-\int_0^T L_{n-1}^{(1)}(u)du\right]
 &=\sum_{j=1}^n\binom nj\eta_{j-1}\\
 &=\int_0^\infty e^{-u}M(u)L_{n-1}^{(2)}(u)du.
\end{aligned}
\tag{155B.4}
\]

The boundary term is \(e^{-T}M(T)L_{n-1}^{(1)}(T)\to0\) and \(M(0)=0\).
At \(n=1\) the integral equals \(\eta_0\); subtracting \(n\eta_0\)
proves the \(-n\) subtraction in (155B.2). Omitting either subtraction
changes the functional. \(\square\)

\textbf{Corollary 155B.2 (the accumulated curvature target).} Define
\(\mathcal A_0=\mathcal E_0=\lambda_0=0\) and
\(\Delta^2 f_n=f_{n+2}-2f_{n+1}+f_n\). Then

\[
\begin{aligned}
G_n:=\Delta^2\mathcal A_n
 &=\sum_{\substack{\ell\ge3\\\ell\ {\rm odd}}}
       \frac{(1-1/\ell)^n}{\ell^2}>0,\\
I_n:=\Delta^2\mathcal E_n
 &=\int_0^\infty e^{-u}M(u)L_{n+1}^{(0)}(u)du,\\
\lambda_N&=N\lambda_1+
 \sum_{n=0}^{N-2}(N-1-n)(G_n-I_n).
\end{aligned}
\tag{155B.5}
\]

The first equality is a finite difference of the convergent reserve. The
Laguerre recurrence
\(L_r^{(\alpha)}-L_{r-1}^{(\alpha)}=L_r^{(\alpha-1)}\) proves the
second; \(n=0\) follows from \(L_1^{(0)}=1-u\). Twice summing finite
differences proves the last line. The positive weights sum to
\(T_{N-1}\).

The unproved sign target is \(\mathcal E_{2m+1}\le\mathcal A_{2m+1}\)
for every \(m\), equivalently the nonnegativity of (155B.5) at every odd
\(N\). An upper bound

\[
\forall\epsilon>0\ \exists C_\epsilon<\infty\ \forall m\ge0:
\quad\mathcal E_{2m+1}\le C_\epsilon e^{\epsilon(2m+1)}
\tag{155B.6}
\]

already suffices, by (155A.5) and the growth of \(\mathcal A_n\).
Conversely RH implies \(\mathcal E_n\le\mathcal A_n\), so (155B.6) is
RH-equivalent. The constants and the all-rate quantifier matter.

The coefficientwise absolute PNT estimate does not prove (155B.6). It
gives

\[
|\mathcal E_n|\le
2C\left[\frac n{a^2}+\sum_{j=0}^{n-1}
 \binom{n+1}{j+2}\frac{(2j+1)!}{j!a^{2j+2}}\right].
\tag{155B.7}
\]

Its last summand has \(n\)th root asymptotic to \(4n/(ea^2)\), so this
particular majorant is superexponential. That is not a lower bound on
\(|\mathcal E_n|\). At \(n=3\) the actual kernel is \(u^2/2-4u+3\), of
changing sign. Earlier nonnegative Li coefficients do not control the
next triangular weighted sum; cancellation, not the reserve constant, is
the missing source estimate.

\section{155C. A compensated positive reserve and the arithmetic
defect}\label{c.-a-compensated-positive-reserve-and-the-arithmetic-defect}

Section 155B gives a positive scalar reserve
\(\mathcal A_n\ge n\lambda_1>0\). That does not make its reflected
matrix positive. The obstruction and its exact repair can both be stated
in the same polynomial basis as the actual completed matrix \(K_N\).

\textbf{Proposition 155C.1 (the original reserve is indefinite).} Put
\(g_j=\sum_{\ell\ge3,\ \ell\ {\rm odd}}\ell^{-j}\) for \(j=2,3\). The
first reserve block is

\[
\mathsf A_2=
\begin{pmatrix}
\lambda_1&\lambda_1+g_2\\
\lambda_1+g_2&3\lambda_1+3g_2-g_3
\end{pmatrix},\qquad
(3,-1)\mathsf A_2(3,-1)^{\mathsf T}
=6\lambda_1-3g_2-g_3<-\frac{47}{600}.
\tag{155C.1}
\]

\textbf{Proof.} Expanding the first three reserves in (155B.2) gives the
matrix. For an exact sign bound, the positive atanh series gives
\(\log2>842/1215\) and \(\log3>263/240\). The decreasing sequence
\(H_n-\log n\) yields \(\gamma_E<H_{32}-5(842/1215)<3/5\). Also
\(\pi>3\) and \(\log(4\pi)>\log12>2(842/1215)+263/240>99/40\). The
formula \(\lambda_1=1+\gamma_E/2-\log(4\pi)/2\) therefore gives
\(\lambda_1<1/16\). Since \(g_2>1/9+1/25\) and \(g_3>0\),

\[
6\lambda_1-3g_2-g_3<\frac38-3\left(\frac19+\frac1{25}\right)
=-\frac{47}{600}<0.
\]

For example \(\pi>3\) itself follows by integrating the alternating
lower polynomial \(\sum_{k=0}^7(-1)^kt^{2k}\) for \((1+t^2)^{-1}\) on
\([0,1]\). Every comparison used above is rational after these series
bounds. Since \(3\mathcal P_0-\mathcal P_1=y\), the direction is the
\(y^2\) moment of the reserve functional. It is not a counterexample to
the actual completed matrix. \(\square\)

\textbf{Theorem 155C.2 (sharp linear anchor in the reserve family).} For
real \(c\) and \(n\ge0\) define

\[
\begin{aligned}
\mathcal A_n^{[c]}&=cn+
 \sum_{\substack{\ell\ge3\\\ell\ {\rm odd}}}
 \bigl[(1-1/\ell)^n-1+n/\ell\bigr],\\
\mathsf R_N^{[c]}&=
 [\mathcal A_{r+s+1}^{[c]}-\mathcal A_{|r-s|}^{[c]}]_{0\le r,s<N}.
\end{aligned}
\tag{155C.2}
\]

Put

\[
\boxed{c_*=
\sum_{\substack{\ell\ge3\\\ell\ {\rm odd}}}\frac1{\ell(2\ell-1)}
=\frac\pi4+\frac{\log2}{2}-1.}
\tag{155C.3}
\]

Then

\[
\boxed{\mathsf R_N^{[c]}\succeq0\quad\hbox{for every }N\ge1
\quad\Longleftrightarrow\quad c\ge c_*.}
\tag{155C.4}
\]

For \(c\ge c_*\) every finite block is positive definite. At \(c=c_*\)
its Gram measure is the following Catalan measure times a resolvent
weight:

\[
\boxed{d\nu_+(y)=W(y)d\nu_C(y),\qquad
W(y)=\sum_{\substack{\ell\ge3\\\ell\ {\rm odd}}}
\frac1{(2\ell-1)[1+\ell(\ell-1)y]},\quad0<y<4.}
\tag{155C.5}
\]

Here \(\nu_C\) is the probability measure (95A.4), while the mass of
\(\nu_+\) is \(c_*\). The sharpness assertion is within the family
(155C.2), not an optimization over every possible comparison form.

\textbf{Reconciliation from the Gamma pole lattice.} The
quarter-centered repair of the Gamma pole expansion has positive
measures \(d\sigma_m\). Under the reciprocal pushforward
\(y=1/t\in(0,4)\), the sum over \(m\ge1\) is exactly the resolvent
measure in (155C.5), while the \(m=0\) term is the Catalan measure
itself. The total baseline cost is \(\pi/4+(\log2)/2\), so after
removing the \(m=0\) unit mass one recovers exactly \(c_*\) of (155C.3).
Thus the pole-derived reserve is the existing compensated reserve in
another coordinate, not a new positive owner.

\textbf{Proof of the positive representation.} Set \(L=r+s+1\),
\(d=|r-s|\) and \(h=L-d=2\min(r,s)+1\). Under \(y=4\sin^2(\theta/2)\),
the measure \(\eta\) of (95A.4) is \(d\theta/\pi\). The cosine
representation \(\mathcal P_r=1+2\sum_{j=1}^r\cos(j\theta)\) gives
\(\int\mathcal P_r\mathcal P_s\,d\eta=h\). For \(0<t<1\) the convergent
geometric Poisson kernel is

\[
p_t(\theta)=\frac{1-t^2}{1-2t\cos\theta+t^2}
=1+2\sum_{j\ge1}t^j\cos(j\theta).
\]

The cross identity
\(y\mathcal P_r\mathcal P_s=2\cos(d\theta)-2\cos(L\theta)\) therefore
gives \(\int\mathcal P_r\mathcal P_s(y/2)p_t\,d\eta=t^d-t^L\). For
\(t=1-1/\ell\) and \(a_\ell=\ell(\ell-1)\), direct algebra yields

\[
\frac2{2\ell-1}-\frac y2
\frac{1-t^2}{(1-t)^2+ty}
=\frac{4-y}{2(2\ell-1)(1+a_\ell y)}.
\tag{155C.6}
\]

Multiplying by \(\mathcal P_r\mathcal P_s\) and integrating gives

\[
\boxed{\int_0^4\frac{\mathcal P_r(y)\mathcal P_s(y)}
{(2\ell-1)(1+a_\ell y)}\,d\nu_C(y)
=\frac{2h}{2\ell-1}-(t^d-t^L).}
\tag{155C.7}
\]

The right side equals the reflected difference of the single-\(\ell\)
term in (155C.2), plus \(h/[\ell(2\ell-1)]\). At \(r=s=0\) it shows that
this positive component measure has mass \(1/[\ell(2\ell-1)]\). Tonelli
and the summability of these masses justify the sum of measures;
boundedness of every fixed polynomial on \([0,4]\) justifies every Gram
entry. Consequently

\[
[\mathsf R_N^{[c_*]}]_{r,s}
=\int_0^4\mathcal P_r(y)\mathcal P_s(y)d\nu_+(y),
\qquad \nu_+([0,4])=c_*.
\tag{155C.8}
\]

The density is positive on \((0,4)\), so every nonzero polynomial has
strictly positive squared integral. This proves positive definiteness.
For general \(c\) the representing finite signed measure is
\(\nu_c=\nu_++(c-c_*)\eta\), which is positive if \(c\ge c_*\).

\textbf{Sharpness.} Near \(y=4\) the terms of \(W\) have a uniformly
summable \(O(\ell^{-3})\) bound, so \(W\) is continuous and finite
there. The density of \(\nu_c\) relative to \(\eta\) is
\((4-y)W(y)/2+c-c_*\) and tends to \(c-c_*\). If \(c<c_*\), it is
negative on an interval next to \(4\). Choose a nonzero continuous
squared test supported inside that interval and approximate its square
root uniformly by polynomials. Finite total variation passes its
negative integral to some polynomial square. Because the
\(\mathcal P_r\) span the polynomials, some finite block then has a
negative direction. This proves the reverse implication in (155C.4).

Finally the absolutely convergent sum in (155C.3) equals

\[
c_*=\frac12\int_0^1\frac{x^{1/4}-x^{1/2}}{1-x}dx
=2\int_0^1\frac{t^4}{(1+t)(1+t^2)}dt
=\frac\pi4+\frac{\log2}{2}-1.
\tag{155C.9}
\]

The substitution is \(x=t^4\); polynomial division and partial fractions
evaluate the last integral. \(\square\)

\textbf{A second finite-entry calculation.} The Catalan transform is

\[
\int_0^4\frac{d\nu_C(y)}{1+ay}
=\frac2{1+\sqrt{1+4a}},\qquad a\ge0.
\tag{155C.10}
\]

The substitution of (95A.4) and the elementary integral of
\((1+2a-2a\cos\theta)^{-1}\) prove it, with continuity at \(a=0\). At
\(a=a_\ell\) its value is \(1/\ell\), independently checking the
component mass in (155C.7). If
\(M_{j,\ell}=\int y^j/(1+a_\ell y)d\nu_C\), then \(M_{0,\ell}=1/\ell\)
and \(M_{j,\ell}=(C_{j-1}-M_{j-1,\ell})/a_\ell\) for \(j\ge1\). These
exact rational moments give a separate route to every finite entry.

\textbf{The same weight as an operator trace.} Restrict \(A_\triangle\)
of Section 12A to the closed span of \(\psi_{2r}\), \(r\ge1\), and
denote the restriction by \(A_{\rm e}\). Its eigenvalues and odd
coordinates are \(a_r=2T_{2r}=(2r)(2r+1)\) and \(\sqrt{1+4a_r}=4r+1\).
Substitution \(\ell=2r+1\) in (155C.5) proves

\[
\boxed{W(y)=\operatorname{Tr}
\bigl[(I+4A_{\rm e})^{-1/2}(I+yA_{\rm e})^{-1}\bigr],\qquad y>0.}
\tag{155C.11}
\]

The terms are \(O_y(r^{-3})\). At \(y=0\) this trace diverges, although
\(W(y)d\nu_C(y)\) has the finite mass already proved. It is a weighted
resolvent trace, not the unweighted resolvent or a bounded endpoint
multiplier. It uses exactly the even-triangular indices of Section 133A,
but does not identify them with the actual folded zeros.

\textbf{Corollary 155C.3 (positive reserve minus the complete defect).}
Put \(\delta=c_*-\lambda_1>0\), and shift both source channels:

\[
\mathcal A_n^+=\mathcal A_n+\delta n,\qquad
\mathcal E_n^+=\mathcal E_n+\delta n,\qquad
\lambda_n=\mathcal A_n^+-\mathcal E_n^+.
\tag{155C.12}
\]

The positivity of \(\delta\) follows from \(\lambda_1<1/16\) and the
first four positive terms of \(c_*\), whose sum exceeds \(1/10\). For
orientation, \(c_*=0.1319717536774209643\ldots\) and
\(\delta=0.1088760447112999305\ldots\); these decimals are not
certificates. Define, for \(0\le r,s<N\),

\[
\begin{aligned}
[\mathsf R_N]_{r,s}&=\mathcal A_{r+s+1}^+-\mathcal A_{|r-s|}^+,
 &\mathsf R_N&\succ0,\\
[\mathsf D_N]_{r,s}&=\mathcal E_{r+s+1}-\mathcal E_{|r-s|}
 +\delta(2\min(r,s)+1),
 &K_N&=\mathsf R_N-\mathsf D_N.
\end{aligned}
\tag{155C.13}
\]

The reserve is exactly the Gram matrix (155C.8). For an explicit prime
expression set \(\Phi_0(u)=0\) and \(\Phi_n(u)=L_{n-1}^{(2)}(u)-n\) for
\(n\ge1\). Then

\[
[\mathsf D_N]_{r,s}=\delta(2\min(r,s)+1)
+\int_0^\infty e^{-u}M(u)
 [\Phi_{r+s+1}(u)-\Phi_{|r-s|}(u)]du.
\tag{155C.14}
\]

Here \(M(u)=\Psi(e^u)-e^u+1\) is exactly (155B.1). All fixed-degree
integrals converge by Theorem 155B.1. The \(-n\) subtraction is
retained, and \(\Phi_0\) is defined separately rather than using
\(L_{-1}^{(2)}\). Equations (155C.12)--(155C.14) follow by subtraction;
omitting the compensation \(\delta(2\min(r,s)+1)\) would change the
source. \(\square\)

The remaining comparison is

\[
\boxed{\mathsf D_N\preceq\mathsf R_N\quad\hbox{for every }N\ge1.}
\tag{155C.15}
\]

It is equivalent to requiring the largest eigenvalue of
\(\mathsf R_N^{-1/2}\mathsf D_N\mathsf R_N^{-1/2}\) to be at most one at
every rank. This is a one-sided form bound: it does not require
\(\mathsf D_N\) positive or its absolute operator norm bounded by one.
By Section 155A's congruence, it is RH-equivalent and remains unproved.
Its diagonal condition cancels the adjustment and is precisely
\(\mathcal E_{2m+1}\le\mathcal A_{2m+1}\), the existing odd-Li target.
Thus the positive reserve has not supplied a hidden sign proof.

This is a boundary comparison. No isometric identification with the full
Bernstein--Laguerre row at \(a=2\) is asserted. Section 166F defines and
approximates a separate equivalent target directly at scale two. Section
166G reconstructs the actual boundary block from a finite Taylor
polynomial of the full cutoff source at two, with a uniform normalized
error. Its coupled degree and cutoff are essential; it does not evaluate
the raw cutoff at zero or prove the source upper bound.

\section{155D. Reserve growth and the one-sided source
criterion}\label{d.-reserve-growth-and-the-one-sided-source-criterion}

This section keeps the completed reserve/defect problem in the reflected
polynomial basis while suppressing the longer conditioning derivation
retained in earlier editions. For nonzero \(p\) of degree \(<N\) let
\(\mathcal R[p]\) be the positive reserve form and \(\mathcal D[p]\) the
actual completed defect. The extremal generalized eigenvalues are

\[
\beta_N^+=\sup_{p\ne0}\frac{\mathcal D[p]}{\mathcal R[p]},\qquad
\beta_N^-=\inf_{p\ne0}\frac{\mathcal D[p]}{\mathcal R[p]}.
\tag{155D.2}
\]

\subsection{Polynomial conditioning of the
reserve}\label{polynomial-conditioning-of-the-reserve}

With \(y=2-2\cos\theta\), the third-kind Chebyshev basis

\[
V_j(y)=\frac{\cos((j+\tfrac12)\theta)}{\cos(\theta/2)},
\qquad d\nu_C=\frac2\pi\cos^2(\theta/2)d\theta
\tag{155D.4}
\]

is orthonormal for the Catalan measure. The reflected basis has Gram
matrix

\[
(G_N)_{rs}=\begin{cases}4r+1&r=s,\\4\min(r,s)+2&r\ne s,
\end{cases}
\tag{155D.5}
\]

and the explicit inverse gives the two-sided conditioning estimate

\[
\boxed{\frac{I}{125(2N^2-1)}\preceq\mathsf R_N
\preceq\frac{c_*N(4N^2-1)}3I.}
\tag{155D.6}
\]

Thus the reserve costs only polynomial factors in \(N\).

\subsection{Spectral growth and the exact
criterion}\label{spectral-growth-and-the-exact-criterion}

For \(y\notin[0,4]\) write \(w+w^{-1}=2-y\), \(|w|>1\), and put
\(r(y)=|w|\); set \(r(y)=1\) on \([0,4]\). At actual folded zeros this
is the reciprocal Li coordinate,

\[
\chi_\rho=1-\rho^{-1},\qquad
\mathfrak L=\sup_\rho|\chi_\rho|=\sup_{[\rho]}r(y_\rho)\ge1.
\tag{155D.8}
\]

The reserve bounds and the Chebyshev expansion imply

\[
\|\mathsf Z_N\|_{\rm op}
\le1+125M_1N(2N-1)^2\mathfrak L^{2N-2}.
\tag{155D.10}
\]

The converse exponential rate is obtained from one maximizing nonreal
folded orbit. Hence the actual one-sided source problem is exactly

\[
\boxed{\begin{aligned}
\mathrm{RH}&\iff
\forall\epsilon>0\ \exists C_\epsilon\ \forall N:\
\mathsf D_N\preceq C_\epsilon e^{\epsilon N}\mathsf R_N\\
&\iff\exists C,A\ge0\ \forall N:\
\mathsf D_N\preceq C(1+N)^A\mathsf R_N\\
&\iff\|\mathsf Z_N\|_{\rm op}=O(N)
\iff\|\mathsf Z_N\|_{\rm op}\text{ is polynomially bounded}.
\end{aligned}}
\tag{155D.11}
\]

No inequality in this display is proved for the actual source; the
theorem identifies the strength required.

The linear compensation remains controlled. With
\(d\eta=dy/[\pi\sqrt{y(4-y)}]\) and \(\mathcal H[p]=\int p^2d\eta\),

\[
0\le\delta\mathcal H[p]\le125\delta N^2\mathcal R[p],
\qquad\deg p<N.
\tag{155D.13}
\]

Thus the open term is the signed arithmetic remainder after this
explicit reserve and compensation are retained. The longer basis
inversion and extremal-orbit proof are preserved in the edition-4.6
package of the Zenodo version chain.

\section{155E. Curvature, progression filters, and cosine
coordinates}\label{e.-curvature-progression-filters-and-cosine-coordinates}

Let \(M(u)=\Psi(e^u)-e^u+1\) and

\[
I_n=\int_0^\infty e^{-u}M(u)L_{n+1}^{(0)}(u)\,du.
\]

For the actual boundary source, \(\mathcal E_0=\mathcal E_1=0\) and

\[
\mathcal E_N=\sum_{n=0}^{N-2}(N-1-n)I_n.
\tag{155E.1}
\]

\textbf{Theorem 155E.1 (one-sided scalar criterion).} The following are
equivalent:

\[
\boxed{\begin{aligned}
\mathrm{RH}&\iff\sup_{n\ge0}I_n<\infty\\
&\iff\exists C,A\ge0\ \forall n:\ I_n\le C(1+n)^A\\
&\iff\forall\epsilon>0\ \exists C_\epsilon\ \forall n:\
I_n\le C_\epsilon e^{\epsilon n}.
\end{aligned}}
\tag{155E.2}
\]

The upper bounds may hold only eventually. Under RH one has the stronger
uniform estimate

\[
|I_n|\le\frac{\pi^2}{8}-1+2\lambda_1.
\tag{155E.3}
\]

The proof is short: under RH, second differences of the absolutely
grouped Li sum are bounded trigonometric moments; conversely (155E.1)
turns any all-rate one-sided upper bound into the odd-index source
criterion of §155B.

\subsection{The generating function and its
poles}\label{the-generating-function-and-its-poles}

With \(H_{\rm reg}(s)=-\zeta'(s)/\zeta(s)-1/(s-1)\),

\[
\mathcal I(z)=\sum_{n\ge0}I_nz^n
=\frac{H_{\rm reg}(1/(1-z))+\gamma_E}{z}.
\tag{155E.4}
\]

A zero \(\rho\) contributes the genuine pole

\[
z_\rho=1-\rho^{-1},\qquad
\operatorname*{Res}_{z=z_\rho}\mathcal I(z)
=-\frac{m_\rho}{z_\rho\rho^2}\ne0.
\tag{155E.5}
\]

Pringsheim's theorem applied separately to the positive and negative
coefficient parts gives the same exponential obstruction on either sign.
Admissible arithmetic progressions retain it when the pole angles cannot
collide; arbitrary sparse subsequences do not inherit that conclusion.

\subsection{Cosine-coordinate matrix}\label{cosine-coordinate-matrix}

After the adjacent-difference basis change the arithmetic defect has the
Toeplitz-plus-Hankel form

\[
\begin{gathered}
\widetilde{\mathsf D}_d=\delta I_d+\widetilde{\mathsf E}_d,\\
\widetilde{\mathsf E}_{00}=0,\quad
\widetilde{\mathsf E}_{0r}=J_r/\sqrt2,\quad
\widetilde{\mathsf E}_{rs}=(J_{r+s}+J_{|r-s|})/2.
\end{gathered}
\tag{155E.14}
\]

The reserve obeys

\[
\widetilde{\mathsf R}_d\succeq
\frac{1-\cos(\pi/(2d))}{125}I_d.
\tag{155E.15}
\]

Therefore a polynomial one-sided bound on either parity of \(I_n\) is an
explicit sufficient arithmetic target for RH. The detailed progression
angle estimate and cosine-basis diagonalization remain in earlier
editions.

\section{155F. Direct curvature cutoffs and the envelope
obstruction}\label{f.-direct-curvature-cutoffs-and-the-envelope-obstruction}

Section 155E isolates the actual curvature \(I_{k-1}=\int_0^\infty
 e^{-u}M(u)L_k^{(0)}(u)du\), with \(M(u)=\Psi(e^u)-e^u+1\). A finite
approximation must retain its regularization. This section gives one
directly on the real axis; Section 166G gives a different boundary
approximation from completed Taylor data at two. Neither proves the
source upper bound.

For \(U>0\) fill the removable quotient at \(v=0\) and define

\[
\begin{aligned}
\widehat H_U(v)&=\sum_{2\le n\le e^U}\frac{\Lambda(n)}{n^{1+v}}
                 -\frac{1-e^{-vU}}v,\\
\widehat I_{k-1}(U)&=\sum_{2\le n\le e^U}\frac{\Lambda(n)}n
 L_k^{(-1)}(\log n)+L_{k+1}^{(-2)}(U),\qquad k\ge1.
\end{aligned}
\tag{155F.1}
\]

The first function is entire in \(v\) and approximates the regular part
\(H_{\rm reg}(1+v)\), not the unregularized Euler sum. Negative Laguerre
parameters here denote finite polynomial continuations; they do not
assert positive-weight orthogonality.

\textbf{Proposition 155F.1 (finite regularization).} For \(U>0\) and
\(k\ge1\),

\[
\boxed{\int_0^U e^{-u}M(u)L_k^{(0)}(u)du
=\widehat I_{k-1}(U)-e^{-U}M(U)L_k^{(-1)}(U).}
\tag{155F.2}
\]

\textbf{Proof.} The polynomial identities
\((e^{-u}L_k^{(-1)}(u))'=-e^{-u}L_k^{(0)}(u)\) and
\((L_{k+1}^{(-2)}(u))'=-L_k^{(-1)}(u)\) have the required zero endpoint
values at \(u=0\). Integrate by parts against
\(dM=\sum_{n\ge2}\Lambda(n)\delta_{\log n}-e^u du\). The atoms give the
finite prime sum; the continuous term gives the polynomial counterterm;
the remaining endpoint has the displayed minus sign. At a prime-power
cutoff the atom and endpoint jumps cancel. \(\square\)

For example \(\widehat I_0(U)=U^2/2-
\sum_{n\le e^U}\Lambda(n)(\log n)/n\). Omitting its \(U^2/2\) does not
approximate \(I_0\). Omitting the endpoint from (155F.2) is also
incorrect; it can be omitted from an approximant only with its error
included.

\textbf{Theorem 155F.2 (uniform direct-boundary remainder).} Retain the
existence constants \(C,c>0\) in the PNT estimate
\(|M(u)|\le Ce^{u-c\sqrt u}\), enlarged on a compact interval as in
Section 155B. For \(K\ge1\), \(U\ge K\) and \(v=c\sqrt U\), set

\[
 e_K(U)=\frac{Ce^{K^2/U}}{K!}
 \left[U^Ke^{-v}+\frac2{c^{2K+2}}\Gamma(2K+2,v)\right].
 \quad
 \max_{1\le k\le K}|I_{k-1}-\widehat I_{k-1}(U)|\le e_K(U).
\tag{155F.3}
\]

Here \(\Gamma(s,v)=\int_v^\infty t^{s-1}e^{-t}dt\) is the upper
incomplete Gamma function. These are not newly certified numerical
constants for the primes.

\textbf{Proof.} Reversing the finite coefficients gives
\(|L_k^{(0)}(u)|\le (u^k/k!)\sum_{r=0}^k k^{2r}/(r!u^r)
\le u^ke^{k^2/u}/k!\). The absolute coefficients of \(L_k^{(-1)}\) are
no larger. For \(u\ge U\ge K\) and \(k\le K\) both are bounded by
\(u^Ke^{K^2/U}/K!\). Apply the PNT estimate to the endpoint and the
remaining integral in (155F.2), using
\(\int_U^\infty u^Ke^{-c\sqrt u}du=2c^{-2K-2}\Gamma(2K+2,c\sqrt U)\).
\(\square\)

For any fixed \(\eta>0\), a sufficient schedule is

\[
\begin{aligned}
U_K&=\max\left\{K,\left[\frac{(1+\eta)K\log(K+2)}c\right]^2\right\},\\
\log e_K(U_K)&=-\eta K\log K+2K\log\log K+O_{C,c,\eta}(K).
\end{aligned}
\tag{155F.4}
\]

Indeed, for \(v>2K+1\) the incomplete Gamma term is between
\(v^{2K+1}e^{-v}\) and that quantity divided by \(1-(2K+1)/v\).
Substitute \(v=(1+\eta)K\log(K+2)\) and use Stirling's formula. The
error is smaller than every fixed exponential in degree, but the prime
cutoff is \(X=e^{U_K}\), much larger than the exponential-in-rank Taylor
cutoff of Section 166G. No necessity or optimality is asserted.

In the cosine coordinates (155E.14), replace each \(I_n\) by its
approximant, keep \(J_0=0\), and retain the entire compensation
\(\delta I_d\). For \(d\ge2\) the highest index is \(2d-3\). The row-sum
bound and (155E.15) give

\[
\left\|\widetilde R_d^{-1/2}
 (\widetilde D_d-\widehat{\widetilde D}_{d,U})
 \widetilde R_d^{-1/2}\right\|_{\rm op}
\le\frac{125d}{1-\cos(\pi/(2d))}\,e_{2d-2}(U).
\tag{155F.5}
\]

Its prefactor is \(O(d^3)\), so (155F.4) controls the full mixed block.
This is a second boundary approximant, not an identification with the
Taylor polynomial or a sign theorem for either matrix.

\subsection{What the absolute PNT envelope
loses}\label{what-the-absolute-pnt-envelope-loses}

Put \(F_k(c)=2(2k+1)!/(k!c^{2k+2})\) and
\(Q_k(c)=\int_0^\infty e^{-c\sqrt u}L_k^{(0)}(u)du\). Then

\[
\begin{aligned}
Q_k(c)&=(-1)^kF_k(c)(e^{-c^2/4}+o(1)),\\
\left(\int_0^\infty e^{-c\sqrt u}|L_k^{(0)}(u)|du\right)^{1/k}
&\sim\frac{4k}{ec^2}.
\end{aligned}
\tag{155F.6}
\]

To prove both assertions, integrate the finite expansion and reverse its
index. The ratio \((-1)^kQ_k/F_k\) is \(\sum_{r=0}^k(-1)^rA_{k,r}\),
where

\[
 A_{k,r}=\frac1{r!}\prod_{j=0}^{r-1}
 \frac{c^2(k-j)}{2(2(k-j)+1)}
 \longrightarrow\frac{(c^2/4)^r}{r!},\qquad
 0\le A_{k,r}\le\frac{(c^2/4)^r}{r!}.
\tag{155F.7}
\]

Dominated convergence gives the signed limit. The absolute integral lies
between \(|Q_k|\) and \(F_ke^{c^2/4}\); factorial asymptotics give its
root limit. Thus preserving the Laguerre oscillation alone does not
repair an estimate that has already replaced the actual remainder by
this envelope. It is not a lower bound on the actual \(I_n\).

One fixed positive counting measure shows that the obstruction is also
one-sided. Let \(F(X)=(X-1)(1-e^{-\sqrt{\log X}})\) for \(X\ge1\) and
\(\widetilde\Psi_s=F\). It is increasing and nonnegative, with remainder
\(\widetilde M_s(u)=-(e^u-1)e^{-\sqrt u}\). At \(c=1\) the alternating
sum in (155F.7) is at least \(3/4\). Since
\(\int_0^\infty e^{-u}|L_k(u)|du\le2^k\), odd \(k=2m+1\) give

\[
 \widetilde I_{2m,s}\ge
 \frac32\frac{(4m+3)!}{(2m+1)!}-2^{2m+1}.
\tag{155F.8}
\]

The integer-supported version
\(\widetilde\Psi_d(X)=F(\lfloor X\rfloor)\) has positive weights
\(F(n)-F(n-1)\) at \(n\ge2\). Differentiation gives \(0\le F'\le3/2\)
(use \((1-e^{-t})/\sqrt t\le1\)). Hence

\[
\begin{aligned}
|\widetilde M_d(u)|&\le(1+\tfrac32e^{1/4})e^{u-\sqrt u},\\
\widetilde I_{2m,d}&\ge
\frac32\frac{(4m+3)!}{(2m+1)!}-\frac52\,2^{2m+1}.
\end{aligned}
\tag{155F.9}
\]

These are artificial measures, not the von Mangoldt measure or a
Gamma-completed zeta function. They show that a PNT-quality envelope,
monotonicity, nonnegative weights and integer support do not alone force
the curvature upper bound. The next section keeps prime-power support as
well, and identifies the arithmetic it still does not retain.

\section{155G. The actual denominator and the limits of positive prime
support}\label{g.-the-actual-denominator-and-the-limits-of-positive-prime-support}

Define the disk function

\[
\mathcal B(z)=z\zeta\!\left(\frac1{1-z}\right)
=\sum_{n\ge0}b_n^{\rm den}z^n,
\qquad \mathcal B(0)=1.
\tag{155G.1}
\]

Its zeros in \(|z|<1\) are exactly the right-half-strip images
\(1-1/\rho\). Combining \(\mathcal B\) with the curvature generating
function of §155E gives the exact signed recurrence

\[
\boxed{\sum_{j=0}^n b_j^{\rm den}I_{n-j}
=-(n+2)b_{n+2}^{\rm den}+(2n+\gamma_E+1)b_{n+1}^{\rm den}
 +(1-n)b_n^{\rm den}.}
\tag{155G.3}
\]

It is a convolution identity, not a positivity recurrence. Already
\(b_1^{\rm den}=\gamma_E-1<0\), from the Laurent expansion of \(\zeta\)
at \(s=1\). Numerically \(b_2^{\rm den},\ldots,b_{16}^{\rm den}\) are
positive, and a directed Euler--Maclaurin certificate gives
\(b_{17}^{\rm den}<0\), where the sign changes again. Neither sign is a
negative Li coefficient or a negative Widder rung.

The coefficients are square summable with the classical normalization

\[
\boxed{\sum_{n\ge0}(b_n^{\rm den})^2=\log(2\pi)-\gamma_E.}
\tag{155G.6}
\]

but not absolutely summable. Thus an unweighted Neumann/Rouche argument
is unavailable; finite-radius control additionally needs a lower bound
for the truncated denominator. The constant in (155G.6) is the one of
Theorem 96A.3: by (96A.5) it is also \(\sum_{k\ge2}\zeta(k)/T_k\) and a
weighted mean square of \(\zeta\) on the critical line.

The prime-supported countermodel retained from earlier editions is
deliberately stronger than a cardinality toy model: it keeps positive
prime-power weights, a PNT-size envelope, a meromorphic Euler product
and an arbitrarily long exact prime prefix, yet introduces a conjugate
pair of interior poles. Its role is only to show that those reduced
hypotheses do not imply the required one-sided curvature bound. It is
not a model of the completed functional equation and is not a
counterexample to RH.

The next section therefore restores the exact divisor arithmetic. Full
coefficient enclosures, the countermodel Euler product and its pole
separation are retained in earlier editions.

\section{155H. Exact divisor arithmetic and a Mellin squared-error
target}\label{h.-exact-divisor-arithmetic-and-a-mellin-squared-error-target}

Here \(\mu(d)\) is the Möbius function, not a boundary moment \(\mu_j\).
The actual primitive factors satisfy \(\Lambda*1=\log\),
\(\mu*1=\delta_1\) and \(\Lambda=\mu*\log\). In particular, for
\(N=\lfloor X\rfloor\) and \(H_N=\sum_{k\le N}1/k\),

\[
\boxed{\sum_{k=1}^N M\!\left(\log\frac Xk\right)
 =\log(N!)-XH_N+N.}
\tag{155H.1}
\]

Finite reordering gives
\(\sum_{k\le N}\Psi(X/k)=\sum_{d\le N}\Lambda(d)\lfloor X/d\rfloor\).
Regrouping by the product gives
\(\sum_{n\le N}\sum_{d\mid n}\Lambda(d)=\log(N!)\). Both subtractions in
\(M\) are retained. The perturbed weights in Section 155G add the
discrepancy
\(\sum_{n\le N}(\sum_{d\mid n}\widetilde\Lambda_P(d)-\log n)\). Thus
(155H.1) uses information absent from that model, but alone supplies no
sign for the delayed, correlated Laguerre transforms.

Take finitely supported real coefficients \(c_d\) with
\(\sum_dc_d/d=0\), and define

\[
\begin{aligned}
D_c(s)&=\sum_dc_dd^{-s},&
 h_c(x)&=1-\sum_dc_d\lfloor x/d\rfloor,\\
\mathscr N(c)&=\int_1^\infty h_c(x)^2\frac{dx}{x^2}
 &=\sum_{n\ge1}\frac{(1-\sum_dc_d\lfloor n/d\rfloor)^2}{n(n+1)}.
\end{aligned}
\tag{155H.2}
\]

The scalar \(D_c\) is not the defect matrix \(\mathsf D_d\); this Mellin
error is not a Bernstein row energy. Balance gives
\(h_c=1+\sum_dc_d\{x/d\}\) and \(|h_c|\le1+\sum|c_d|\), so the norm is
finite. Jumps occur at integers, proving the series. The weight
\(1/[n(n+1)]=1/(2T_n)\) is exactly the integral of \(x^{-2}\) over that
unit interval. The square follows the signed divisor sum.

\textbf{Proposition 155H.1 (Mellin identity).} For \(\Re s>0\),

\[
\begin{aligned}
E_c(s)&:=\int_1^\infty h_c(x)x^{-s-1}dx
 =\frac{1-\zeta(s)D_c(s)}s,\\
\mathscr N(c)&=\frac1{2\pi}\int_{\mathbb R}
 |E_c(1/2+it)|^2dt.
\end{aligned}
\tag{155H.3}
\]

\textbf{Proof.} For \(\Re s>1\), each floor integral equals
\(\zeta(s)d^{-s}/s\). Summing proves the identity there. The bounded
error makes its integral holomorphic for \(\Re s>0\); \(D_c(1)=0\)
cancels the zeta pole, so continuation proves it throughout. At
\(x=e^u\), \(h_c(e^u)e^{-u/2}\) lies in \(L^1\cap L^2\) on \(u\ge0\);
Fourier Plancherel gives the factor \(1/(2\pi)\) and the norm.
\(\square\)

If \(L\) is a common multiple of the support indices, balance also gives
\(h_c(n+L)=h_c(n)\). Summing through \(T\) has the rigorous tail budget

\[
0\le\mathscr N(c)-\sum_{n=1}^T\frac{h_c(n)^2}{n(n+1)}
\le\frac{(1+\sum_d|c_d|)^2}{T+1}.
\tag{155H.4}
\]

\subsection{Exact-prefix cancellation and zero
detection}\label{exact-prefix-cancellation-and-zero-detection}

Require additionally \(c_d=\mu(d)\) for \(1\le d\le K\). The
inverse-divisor identity gives
\(\sum_{d\le n}\mu(d)\lfloor n/d\rfloor=1\) for every \(n\ge1\);
coefficients beyond \(K\) cannot affect the first \(K\) intervals. Thus
\(h_c(x)=0\) for \(1\le x<K+1\).

\textbf{Theorem 155H.2.} Every actual zero \(\rho=\sigma+i\gamma\) with
\(\sigma>1/2\) forces, for every balanced exact-prefix family,

\[
\boxed{\mathscr N(c)\ge
\frac{2\sigma-1}{|\rho|^2}(K+1)^{2\sigma-1}.}
\tag{155H.5}
\]

\textbf{Proof.} At a zero, \(E_c(\rho)=1/\rho\). By the vanishing prefix
and Cauchy--Schwarz,
\(|\rho|^{-2}\le\mathscr N(c)\int_{K+1}^\infty x^{-2\sigma}dx
=\mathscr N(c)(K+1)^{1-2\sigma}/(2\sigma-1)\). No phase selection or
verified height is used. \(\square\)

Consequently the following is sufficient for RH, even along one fixed
unbounded sequence of actual prefix lengths:

\[
\boxed{\forall\epsilon>0\ \exists C_\epsilon<\infty\ \forall K:
\quad\mathscr N(c^{(K)})\le C_\epsilon(K+1)^\epsilon.}
\tag{155H.6}
\]

The family must be specified and preserve balance and the exact prefix.
No converse is asserted for an arbitrary choice of extensions. A fixed
\(O(K^a)\) bound, \(0\le a<1\), excludes only \(\Re\rho>(1+a)/2\), and
reflection gives the corresponding symmetric strip. Unlike a polynomial
matrix-rank allowance, the prefix-length allowance here must be below
every positive power.

One explicit balanced extension is

\[
 a_K=\sum_{d\le K}\frac{\mu(d)}d,\quad q=K+1,\qquad
 D_K(s)=\sum_{d\le K}\mu(d)d^{-s}-q\,a_Kq^{-s},\qquad
 \mathscr N(c^{(K)})\le4(K+1).
\tag{155H.7}
\]

Indeed \(Ka_K=1+\sum_{d\le K}\mu(d)\{K/d\}\) gives \(|a_K|\le1\). Then
\(|h_c|\le2(K+1)\) and its vanishing prefix gives the displayed linear
bound. For \(\sigma=\Re s>1\) the same family satisfies

\[
\left|D_K(s)-\frac1{\zeta(s)}\right|
\le\frac{K^{1-\sigma}}{\sigma-1}+(K+1)^{1-\sigma}\longrightarrow0.
\tag{155H.8}
\]

This is safe-half-plane convergence, not the missing subpower norm
estimate. At \(K=1\), \(D_1(s)=1-2^{1-s}\) and \(h_c(n)\) is one for
even \(n\), zero for odd \(n\), giving

\[
 \mathscr N(c^{(1)})=\sum_{j\ge1}\frac1{2j(2j+1)}=1-\log2.
\tag{155H.9}
\]

\subsection{Why a natural balanced primitive product
fails}\label{why-a-natural-balanced-primitive-product-fails}

For integer \(P\ge2\) set \(Q_P=\prod_{p\le P}p\) and

\[
 A_P(s)=\prod_{p\le P}(1-p^{-s})=\sum_{d\mid Q_P}\mu(d)d^{-s},
 \quad\alpha_P=A_P(1),\quad D_P^{\rm E}(s)=A_P(s)-\alpha_P.
\tag{155H.10}
\]

This uses actual inverse Euler factors. The subtraction enforces
balance, but changes the coefficient at one: it is not an exact-prefix
family. Although \(A_P\to1/\zeta\) for \(\Re s>1\) and
\(\alpha_P^{-1}\ge H_P\to\infty\), the norm satisfies

\[
\boxed{\mathscr N(c^{\rm E,P})
\ge\alpha_P^2(P+1-H_{P+1})
\ge\frac{P^{1/3}}{18e^2}\longrightarrow\infty\quad(P\ge5).}
\tag{155H.11}
\]

\textbf{Proof.} Inclusion--exclusion makes
\(\sum_{d\mid Q_P}\mu(d)\lfloor n/d\rfloor\) count the integers at most
\(n\) coprime to \(Q_P\). Only one survives for \(1\le n\le P\). The
balancing correction therefore makes \(h_{\rm E,P}(n)=\alpha_P n\).
Their contribution to (155H.2) is the first lower bound. For an
elementary product bound, adding all integers coprime to six gives
\(\alpha_P\ge\tfrac13\prod_{5\le n\le P,(n,6)=1}(1-1/n)\). Use
\(-\log(1-1/n)\le1/n+1/[2n(n-1)]\) and pair \(6j-1,6j+1\). Their
reciprocal sum through \(J=\lceil(P+1)/6\rceil\le P\) is at most
\((\log P+1)/3+1/70\), since
\(1/(6j-1)+1/(6j+1)=1/(3j)+1/[3j(36j^2-1)]\). The remaining logarithm
error is at most \(1/8\), and \(1/3+1/70+1/8<1\). Thus
\(\alpha_P\ge e^{-1}/(3P^{1/3})\). Finally \(P+1-H_{P+1}\ge P/2\).
\(\square\)

Even division by \(1-\alpha_P\) to restore the first coefficient leaves
the ramp \(\alpha_P(n-1)/(1-\alpha_P)\) and a positive multiple of the
same \(P^{1/3}\) lower bound. This excludes the displayed balancing
scheme, not all Euler products or mollifiers. The successful family
would need an upper bound such as (155H.6); neither tested construction
provides it. The framework is classical Mellin/fractional-part
approximation, with no claim of literature priority.

\section{155I. One denominator residual, Mellin norm, and Laguerre
expansion}\label{i.-one-denominator-residual-mellin-norm-and-laguerre-expansion}

The denominator of Section 155G and the balanced arithmetic approximant
of Section 155H are connected by an exact norm identity. Keep the same
\(c_d\), \(D_c\), \(h_c\) and \(\mathscr N(c)\); put \(s(z)=1/(1-z)\)
and

\[
\begin{aligned}
\mathcal C_c(z)&=D_c(s(z))/z,\qquad
 \mathcal C_c(0)=D_c'(1),\\
\mathcal F_c(z)&=1-\mathcal B(z)\mathcal C_c(z)
 =1-\zeta(s(z))D_c(s(z)).
\end{aligned}
\tag{155I.1}
\]

Balance removes the apparent singularity in \(\mathcal C_c\). The
residual at zero is \(1-D_c'(1)\), not \(1\): treating
\(\zeta(1)D_c(1)\) as zero would lose the removable product's limit.

\textbf{Theorem 155I.1 (exact isometry).} With the normalized Hardy
coefficient norm on \(\mathbb D\), the residual belongs to
\(H^2(\mathbb D)\) and

\[
\boxed{\begin{aligned}
\mathcal F_c(z)&=\sum_{n\ge0}r_n(c)z^n,\qquad
 r_n(c)=\int_0^\infty e^{-u}h_c(e^u)L_n^{(0)}(u)du,\\
\|\mathcal F_c\|_{H^2}^2&=\sum_{n\ge0}|r_n(c)|^2=\mathscr N(c).
\end{aligned}}
\tag{155I.2}
\]

\textbf{Proof.} The exact Mellin identity (155H.3), with
\(g(u)=h_c(e^u)\), becomes

\[
\mathcal F_c(z)=\frac1{1-z}\int_0^\infty e^{-u}g(u)
          \exp\!\left(-\frac{uz}{1-z}\right)du.
\tag{155I.3}
\]

The kernel is the ordinary Laguerre generating function. On a small
closed disk \(|z|\le r<1/2\), its absolute coefficient majorant is
\((1-r)^{-1}\exp(ur/(1-r))\). Since \(g\) is bounded, termwise
integration is justified locally. The Laguerre polynomials are
orthonormal in \(L^2(e^{-u}du)\) by Rodrigues and integration by parts.
They are complete: if \(f\) is orthogonal to all polynomials, then
\(A(t)=\int_0^\infty f(u)e^{-u}e^{tu}du\) is analytic for \(\Re t<1/2\)
by Cauchy--Schwarz and all its derivatives at zero vanish. Analytic
continuation and Fourier uniqueness on the imaginary axis give \(f=0\).
Parseval now proves (155I.2), since \(\int|g|^2e^{-u}du=\mathscr N(c)\).
The square-summable series defines an \(H^2\) function on the whole
disk; identity with (155I.1) follows from local equality. This does not
infer a Hardy bound from pointwise boundary values alone. \(\square\)

For a normalization check, \(s(e^{i\theta})=1/2+(i/2)\cot(\theta/2)\)
and \(|d\theta|=dt/(1/4+t^2)\). The normalized circle norm is therefore
exactly (155H.3), with its single factor \(1/(2\pi)\).

The finite Dirichlet polynomial gives explicit multiplier coefficients:

\[
\begin{aligned}
\mathcal C_c(z)&=\sum_{n\ge0}t_n(c)z^n,
 &t_n(c)&=\sum_d\frac{c_d}{d}L_{n+1}^{(-1)}(\log d),\\
r_n(c)&=\mathbf1_{n=0}-\sum_{j=0}^n b_j^{\rm den}t_{n-j}(c).
\end{aligned}
\tag{155I.4}
\]

Indeed \(D_c(s(z))=\sum_d(c_d/d)\exp[-(\log d)z/(1-z)]\) and
\(\exp[-az/(1-z)]=\sum_{k\ge0}L_k^{(-1)}(a)z^k\); its constant term
cancels by balance. Each convolution is finite. Thus the same error is

\[
\mathscr N(c)=\sum_{n\ge0}\left|
\mathbf1_{n=0}-\sum_{j=0}^n b_j^{\rm den}
  \sum_d\frac{c_d}{d}L_{n-j+1}^{(-1)}(\log d)\right|^2.
\tag{155I.5}
\]

Signs remain inside the square; absolute summability of
\(b_j^{\rm den}\) is neither used nor true. Schwarz's lemma applied to
the numerator of \(\mathcal C_c\) gives
\(\|\mathcal C_c\|_\infty\le\sum|c_d|/\sqrt d\). This and (155G.6) do
not establish subpower growth for chosen families. For \(c_1=1,c_2=-2\),
\(r_0=1-\log2\) and \(\mathscr N(c)=1-\log2\); the constant coefficient
contributes only \((1-\log2)^2\) to that full squared norm.

\subsection{Prefix support is the additional analytic
information}\label{prefix-support-is-the-additional-analytic-information}

For \(c_d=\mu(d)\) through \(K\), set \(T=K+1\). The support
cancellation \(h_c=0\) on \([1,T)\) gives

\[
\mathcal G_{c,T}(z)=T^{s(z)-1/2}\mathcal F_c(z)\in H^2,
\qquad \|\mathcal G_{c,T}\|_{H^2}^2=\mathscr N(c).
\tag{155I.6}
\]

To prove this, apply the same isometry to \(x\mapsto T^{-1/2}h_c(Tx)\).
Its norm is \(\int_T^\infty h_c(v)^2v^{-2}dv=\mathscr N(c)\) and its
Mellin transform is \(T^{s-1/2}E_c(s)\). The prefix is essential:
\(T^{s(z)-1/2}\) is not a bounded multiplier on arbitrary \(H^2\)
functions.

At an actual zero \(\rho=\sigma+i\gamma\) with \(\sigma>1/2\),
\(\mathcal F_c(z_\rho)=1\) and \(1-|z_\rho|^2=(2\sigma-1)/|\rho|^2\).
The coefficient evaluation inequality
\(|\mathcal G(z)|^2\le\|\mathcal G\|_{H^2}^2/(1-|z|^2)\) recovers
exactly (155H.5), with the same constant and exponent.

This joins the analytic denominator and exact-divisor error in one
residual, not an improved zero lower bound. The remaining sufficient
estimate is (155H.6). A finite prefix of coefficients in (155I.5) gives
a lower bound on the full norm, not an upper tail certificate. These
infinite Laguerre coefficients are not the curvature \(I_n\), the finite
Bernstein--Laguerre row at two, or the reserve-normalized boundary
matrix. The shared polynomial basis does not identify different inputs
or norms.

\section{155J. Fixed plateau exponent, scalar logarithmic target, and
proof
firewalls}\label{j.-fixed-plateau-exponent-scalar-logarithmic-target-and-proof-firewalls}

Sections 155H--155I gave a sufficient exact-prefix construction target.
For one specified family the exponent can be identified exactly. Define

\[
V_N(s)=\sum_{d\le N}\mu(d)\left(1-\frac dN\right)d^{-s},
\qquad
U_N=\int_0^\infty\left|\mathbf1_{[1,\infty)}(x)
+\sum_{d\le N}\mu(d)\left(1-\frac dN\right)\{x/d\}
\right|^2\frac{dx}{x^2}.
\tag{155J.1}
\]

For \(K<L\), put

\[
A_{K,L}=\frac{LV_L-KV_K}{L-K},
\qquad
C_{K,L}(s)=A_{K,L}(s)-(K+1)A_{K,L}(1)(K+1)^{-s}.
\tag{155J.2}
\]

The coefficients of \(A_{K,L}\) equal \(\mu(d)\) through \(K\) and then
taper linearly to zero. The final term balances at the first index after
the exact prefix, so it changes no coefficient through \(K\). With
\(E=\log(2\pi)-\gamma_E\), Minkowski and the exact initial cancellation
give

\[
\boxed{\mathscr N(C_{K,L})\le E
\left(\frac{L\sqrt{U_L}+K\sqrt{U_K}}{L-K}\right)^2},
\qquad
\mathscr N(C_{K,2K})\le9E\max(U_K,U_{2K}).
\tag{155J.3}
\]

The constant is independent of the prefix length; this is not a
constant-cost repair theorem for arbitrary input coefficients.

\textbf{Theorem 155J.1 (fixed-family exponent).} Let
\(\Theta_\zeta=\sup_{\rho}\Re\rho\), over nontrivial zeros. Then

\[
\boxed{\lim_{N\to\infty}\frac{\log(1+U_N)}{\log N}
=2\Theta_\zeta-1,
\qquad
\lim_{K\to\infty}\frac{\log(1+\mathscr N(C_{K,2K}))}{\log K}
=2\Theta_\zeta-1.}
\tag{155J.4}
\]

\textbf{Proof.} The upper exponent is as before: fix
\(\sigma_0>\Theta_\zeta\), use the zero-free-half-plane Möbius estimate
quoted in Báez-Duarte {[}25, Lemma 2.1{]}, and remove the damping by
finite Abel summation. This gives
\(U_N\ll_{\sigma_0,\epsilon}N^{2\sigma_0-1+\epsilon}\). The lower
exponent for the \emph{raw taper} needs a bootstrap because \(V_N\) does
not have an exact Möbius prefix. For a zero \(\rho=\sigma+i\gamma\) put
\(\delta=2\sigma-1>0\) and \(b_\rho=\sqrt{\delta}/(|\rho|\sqrt E)\). For
\(K<N\) the exact-prefix plateau relation implies \[
\sqrt{U_N}\ge\left(1-\frac KN\right)b_\rho(K+1)^{\delta/2}
-\frac KN\sqrt{U_K}.
\tag{155J.4a}
\] If \(U_K\le DK^a\) with \(a>\delta\), choose
\(K=\lfloor(rN)^{2/(2+a-\delta)}\rfloor\) with
\(0<r\le b_\rho/(4\sqrt D)\). Then, for all sufficiently large \(N\), \[
U_N\ge c_{\rho,a,D}N^{2\delta/(2+a-\delta)}.
\tag{155J.4b}
\] Use the global upper exponents \(a>2\Theta_\zeta-1\), then let the
zero real parts approach \(\Theta_\zeta\) and \(a\) decrease to that
value. The endpoint \(\Theta_\zeta=1\) uses the elementary linear bound;
at \(1/2\) nonnegativity of \(\log(1+U_N)\) supplies the lower limit.
Equation (155J.3) transfers the exponent to the balanced dyadic plateau.
\(\square\)

Thus, for this specified family,

\[
\boxed{\mathrm{RH}\iff
\forall\epsilon>0:\ U_N=O_\epsilon(N^\epsilon),}
\tag{155J.5}
\]

with the estimate allowed on one fixed unbounded set of indices. This is
a fixed-family converse; it does not assert boundedness under RH and
does not apply automatically to arbitrary exact-prefix extensions.

Write

\[
\mathcal X_N=2\sum_{d<e\le N}\mu(d)\mu(e)
\left(1-\frac dN\right)\left(1-\frac eN\right)G(d,e),
\qquad
G(d,e)=\int_0^\infty\{x/d\}\{x/e\}\frac{dx}{x^2}.
\tag{155J.6}
\]

The diagonal and linear pieces can be evaluated separately and give

\[
\boxed{U_N-\mathcal X_N=\kappa\log N+\kappa_0+o(1),
\qquad\kappa=\frac{\log(2\pi)-\gamma_E}{\zeta(2)}.}
\tag{155J.7}
\]

Hence the remaining norm target may be written one-sidedly as
\(\mathcal X_N\le C_\epsilon N^\epsilon\) for every \(\epsilon>0\) on
one fixed unbounded index set. The signs \(\mu(d)\mu(e)\) remain outside
the positive Gram kernel; positivity of \(G\) does not supply this upper
bound.

A second, different scalar target is

\[
B_{\log}(x)=1-\sum_{d\le x}\frac{\mu(d)}d\log(x/d).
\tag{155J.8}
\]

Its Mellin transform is \(1/s-1/[s^2\zeta(1+s)]\). The same
pole/positive-envelope argument gives

\[
\boxed{\mathrm{RH}\iff
\forall\epsilon>0:\ B_{\log}(n)\le
C_\epsilon n^{-1/2+\epsilon}\quad(n\gg_\epsilon1).}
\tag{155J.9}
\]

No positivity of \(B_{\log}\) is assumed. The eventual-integer form can,
however, be replaced by two structured meshes. Put
\(\mathfrak M(x)=\sum_{d\le x}\mu(d)\) and \[
w(t)=\log t+\gamma_E-H_{\lfloor t\rfloor}
-(\gamma_E-1)\frac{\{t\}}t.
\tag{155J.9a}
\] Finite divisor convolutions give \[
\boxed{B_{\log}(x)=\frac{1+(\gamma_E-1)\log x}{x}
-\frac1x\int_1^x\mathfrak M(x/t)w(t)dt.}
\tag{155J.9b}
\] Here \(|w(t)|<1/t\), and the rational interval certificate gives
\(\int_1^\infty|w(t)|t^{-1/2}dt<0.950760<1\). Hence
\(|\mathfrak M(y)|\le Ay^\theta\) through \(X\), \(\theta\ge1/2\),
implies \(|B_{\log}(x)|\le(1+A)x^{\theta-1}\) on the same finite range.
This is one-way smoothing, not a contraction equation for the same
unknown.

In logarithmic coordinates, the exact Dirichlet Green formula yields \[
B_{\log}(x)\le\max\{B_{\log}(a),B_{\log}(b)\}
+\frac14\log^2(b/a),\qquad a\le x\le b.
\tag{155J.9c}
\] Consequently (155J.9) is equivalent to its restriction to all
sufficiently large fourth powers \(x=k^4\), and also to the squared
triangular mesh \(x=T_k^2\), with the same arbitrary positive slack. The
interpolation costs are \(O(x^{-1/2})\). An arbitrary sparse sequence is
still insufficient.

The original balanced canonical prefix family has the same sharp power
exponent. If \(J_K\) is the raw untapered norm and \(N_K\) the balanced
canonical norm, then \(N_K\le EJ_K\) and \[
\boxed{\lim_{K\to\infty}\frac{\log(1+N_K)}{\log K}
=\lim_{K\to\infty}\frac{\log(1+J_K)}{\log K}
=2\Theta_\zeta-1.}
\tag{155J.9d}
\] The accepted unconditional Mertens estimate also transfers to the
fixed taper: \[
\boxed{U_N\ll N\exp(-c\sqrt{\log N}),\qquad
\mathcal X_N\ll N\exp(-c\sqrt{\log N})}
\tag{155J.9e}
\] for some \(c>0\); the second statement is one-sided. These improve
the finite allowance while leaving power exponent one, so neither proves
RH.

The arithmetic norm also has the exact sine realization

\[
P_c(x)=\frac{\sin(2\pi x)}\pi+\sum_d\frac{c_d}{d}\{dx\},
\qquad
\boxed{\mathscr N(c)=\int_0^\infty P_c(x)^2\frac{dx}{x^2}
=\frac{\pi^2}{2}\int_0^1\frac{P_c(x)^2}{\sin^2(\pi x)}dx.}
\tag{155J.10}
\]

This is an isometry for the same signed arithmetic residual, not a
replacement by the positive sine-comparison measure of Section 12A. It
gives no missing upper estimate.

\subsection{Determinant and topology
checks}\label{determinant-and-topology-checks}

A Vandermonde on \(n\) nodes has scaling degree \(T_{n-1}\). For
confluent multiplicities \(m_i\), total \(n\), the degree is
\(T_{n-1}-\sum_iT_{m_i-1}=\sum_{i<j}m_im_j\); raw derivative columns
also carry \(\prod_i\prod_{j=0}^{m_i-1}j!\), removed by Taylor
normalization. This explains recurring triangular determinant exponents
but identifies no new source operator.

For finite coefficient optimization, if \(H>0\), \(Cp=u\), and \(C\) has
full row rank, the constrained minimum is

\[
p_*=H^{-1}C^*(CH^{-1}C^*)^{-1}u,
\qquad
\min_{Cp=u}p^*Hp=u^*(CH^{-1}C^*)^{-1}u.
\tag{155J.11}
\]

It is a useful exact solver for prefix/balance constraints, but a
minimum over adjustable coefficients is not an upper bound on the fixed
taper above. A positive kernel update likewise need not decrease a
normalized defect: the reviewed rational counterexample has
\(K_4-K_3\succeq0\) while \(25/99<4225/15136\). Finally, completion
arguments must retain continuity of the zero observation. Our
exact-prefix Mellin route does so through

\[
\left|\int_q^\infty h_c(x)x^{-s-1}dx\right|
\le\sqrt{\mathscr N(c)}\,
\frac{q^{1/2-\sigma}}{\sqrt{2\sigma-1}},
\qquad q=K+1,\ \sigma>1/2.
\tag{155J.12}
\]

These checks close tempting shortcuts; they do not improve the open
upper bound in (155J.7) or (155J.9).

\section{53. Complete-Bernstein and operator-monotone
owner}\label{complete-bernstein-and-operator-monotone-owner}

The Stieltjes criterion in PP1 has a natural increasing form.

\begin{quote}
\textbf{Theorem 53.1.} With \(h(x)=xS_\xi(x)\), \[
\boxed{\begin{aligned}
\mathrm{RH}&\iff S_\xi\text{ is Stieltjes}\\
&\iff h\text{ is a complete Bernstein function}\\
&\iff h\text{ is operator monotone on }(0,\infty)
\end{aligned}}
\tag{53.1}
\]
\end{quote}

The middle equivalence is the standard Stieltjes--complete-Bernstein
correspondence: a nonnegative \(f\) is complete Bernstein precisely when
\(f(x)/x\) is Stieltjes.

A direct proof is instructive. Under RH,

\[
h(x)=\sum_{[\rho]}m_\rho\frac{x}{x+q_\rho}
\tag{53.2}
\]

For \(q>0\), the map

\[
f_q(X)=X(X+qI)^{-1}=I-q(X+qI)^{-1}
\tag{53.3}
\]

is operator monotone increasing because inversion reverses
positive-operator order. The series converges locally in operator norm
by (48.1), hence \(h\) is operator monotone.

Conversely, Löwner's theorem extends an operator-monotone function on
the positive half-line to a Pick function holomorphic in the upper
half-plane. The genus-zero continuation of (53.2) has poles at
\(-q_\rho\). If a folded zero were nonreal, conjugation symmetry would
place one of \(-q_\rho,-\overline{q_\rho}\) in the upper half-plane.
Common-sign multiplicities prevent cancellation, contradicting
holomorphy. Thus every \(q_\rho\) is positive real and RH follows from
(50.12).

For positive nodes \(x_1,\ldots,x_N\), define the Löwner matrix

\[
\mathcal L_h(x_1,\ldots,x_N)
=\left[\ell_{ij}\right]_{i,j=1}^N
\qquad
\ell_{ij}=
\begin{cases}
\dfrac{h(x_i)-h(x_j)}{x_i-x_j},&i\ne j\\[2mm]
h'(x_i),&i=j
\end{cases}
\tag{53.4}
\]

Then

\[
\boxed{
\mathrm{RH}
\iff
\mathcal L_h(x_1,\ldots,x_N)\succeq0
\quad\text{for every finite positive node set}}
\tag{53.5}
\]

The \(1\times1\) condition is exactly

\[
h'(x)=S_\xi(x)+xS_\xi'(x)=W_1(x)>0
\tag{53.6}
\]

Thus the analytic \(W_1\) theorem is the first diagonal Löwner
condition. The higher scalar Widder rungs are local derivative shadows
of the full matrix problem; no finite collection of shadows replaces
(53.5). Section 48C gives an exact rational rank-two witness showing why
the centered Gamma pole terms cannot themselves be split into positive
Löwner sublattices; this closes a proposed decomposition route without
changing (53.5).

\section{54. Exact Gamma--prime source
matrix}\label{exact-gammaprime-source-matrix}

On the safe real half-line \(s>1\), put

\[
R=\sqrt{1+4x},\qquad
s=\sigma+i0=\frac{1+R}{2}>1
\tag{54.1}
\]

The explicit source formula separates exactly as

\[
\boxed{
S_\xi(x)=\frac1x+G(x)-P(x)}
\tag{54.2}
\]

where

\[
G(x)=
\frac{\psi((1+R)/4)-\log\pi}{2R}
\qquad
P(x)=\frac1R\sum_{n\ge2}\frac{\Lambda(n)}{n^s}
\tag{54.3}
\]

Indeed,

\[
\frac1R\left(\frac1s+\frac1{s-1}\right)=\frac1x
\tag{54.4}
\]

Therefore

\[
h(x)=1+xG(x)-xP(x)
\tag{54.5}
\]

The constant disappears from divided differences. Define

\[
L_\Gamma=
\left[
\frac{x_iG(x_i)-x_jG(x_j)}{x_i-x_j}
\right]_{i,j}
\qquad
L_P=
\left[
\frac{x_iP(x_i)-x_jP(x_j)}{x_i-x_j}
\right]_{i,j}
\tag{54.6}
\]

with diagonal entries obtained by differentiation. Then

\[
\boxed{\mathcal L_h=L_\Gamma-L_P}
\tag{54.7}
\]

This identifies the principal missing theorem in source-native form.

\begin{quote}
\textbf{Open Source Forcing Theorem.} Prove directly from the explicit
Gamma and von-Mangoldt terms that \[
\boxed{L_\Gamma-L_P\succeq0}
\tag{54.8}
\] for every finite positive node set.
\end{quote}

By Theorem 53.1, (54.8) is equivalent to RH. Unlike a zero-side
restatement, it keeps all arithmetic on the absolutely convergent real
Dirichlet half-line. It is stronger and more coherent than proving
\(W_4,W_5,\ldots\) one at a time, although finite-rung certificates
remain useful local information.

The proof map in Reading Volume §R24 displays this all-node source
obligation with its quantifiers. Reader §R15 supplies the
moment-to-matrix dictionary.

\section{131. Gamma-triangular analytic
scaffold}\label{gamma-triangular-analytic-scaffold}

The triangular fold is the exact invariant-coordinate identity

\[
q=s(1-s)=-2\mathsf T_{s-1}
\tag{131.1}
\]

To turn this coordinate identity into an analytic statement about the
Gamma factor, antisymmetrize its logarithmic derivative across the two
sheets.

Put

\[
A_\Gamma(s)
=\frac12\psi(s/2)-\frac12\log\pi
\qquad
R=2s-1
\qquad
q=s(1-s)
\tag{131.2}
\]

and define folded antisymmetrization by

\[
\mathfrak F[A](q)
=\frac{A(s)-A(1-s)}{2R}
\tag{131.3}
\]

The classical digamma partial fraction gives

\[
A_\Gamma(s)
=-\frac{\gamma_E+\log\pi}{2}
+\sum_{m=0}^{\infty}
\left(
\frac1{2(m+1)}-\frac1{s+2m}
\right)
\tag{131.4}
\]

For every \(m\ge0\),

\[
(s+2m)(1-s+2m)
=q+2m(2m+1)
=q+2\mathsf T_{2m}
\tag{131.5}
\]

whereas

\[
-\frac1{s+2m}+\frac1{1-s+2m}
=\frac{R}{q+2\mathsf T_{2m}}
\tag{131.6}
\]

The regularizing terms in (131.4) cancel under (131.3). Local normal
convergence away from the pole set follows from
\(2\mathsf T_{2m}=4m^2+2m\).

\begin{quote}
\textbf{Theorem 131.1 (Gamma-triangular scaffold).} Meromorphically in
\(q\), \[
\boxed{
\mathfrak G_\triangle(q)
:=\mathfrak F[A_\Gamma](q)
=\frac12\sum_{m=0}^{\infty}
\frac1{q+2\mathsf T_{2m}}
}
\tag{131.7}
\] Hence the nonpositive even Gamma lattice folds exactly to \[
q=-2\mathsf T_{2m}=0,-6,-20,-42,-72,\ldots
\tag{131.8}
\]
\end{quote}

This is the first explicit analytic intertwiner from the Gamma factor to
the triangular lattice in this volume. It is stronger than (131.1), but
its weights and spectral orientation remain distinct from the synthetic
measure of Section 96.

\section{132. All-order Löwner geometry of the
scaffold}\label{all-order-luxf6wner-geometry-of-the-scaffold}

For \(q>0\), (131.7) is a Stieltjes transform with measure

\[
d\mu_\Gamma(a)
=\frac12\sum_{m=0}^{\infty}
\delta_{2\mathsf T_{2m}}(da)
\tag{132.1}
\]

The Stieltjes integrability condition holds because the positive atoms
grow quadratically. Therefore, for every \(k\ge0\),

\[
(-1)^k\mathfrak G_\triangle^{(k)}(q)
=\frac{k!}{2}\sum_{m=0}^{\infty}
\frac1{(q+2\mathsf T_{2m})^{k+1}}>0
\tag{132.2}
\]

Define

\[
g_\Gamma(q)=q\mathfrak G_\triangle(q)
=\frac12+\frac12\sum_{m=1}^{\infty}
\frac{q}{q+a_m}
\qquad
a_m=2\mathsf T_{2m}
\tag{132.3}
\]

For distinct positive nodes \(q_1,\ldots,q_N\), direct divided
differences give

\[
\boxed{
\mathcal L_{g_\Gamma}(q_1,\ldots,q_N)_{ij}
=\frac12\sum_{m=1}^{\infty}
\frac{a_m}{(q_i+a_m)(q_j+a_m)}
}
\tag{132.4}
\]

Every summand is a rank-one Gram atom. Strictness follows because a null
vector would make the rational function \(\sum_i c_i/(q_i+a)\) vanish at
every \(a=a_m\); its numerator has degree at most \(N-1\), so infinitely
many distinct zeros force every \(c_i\) to vanish.

\begin{quote}
\textbf{Corollary 132.1 (all-order scaffold Löwner theorem).} The
function \(g_\Gamma\) is complete Bernstein on \((0,\infty)\), and every
finite Löwner matrix in (132.4) is positive definite at distinct
positive nodes.
\end{quote}

This solves an all-order triangular Gamma model. It does not solve the
source matrix (54.7), because \(g_\Gamma\) uses both sheets through
\(A_\Gamma(s)-A_\Gamma(1-s)\).

\section{133. Completion cancels the
scaffold}\label{completion-cancels-the-scaffold}

Write

\[
\frac{\xi'(s)}{\xi(s)}
=A_0(s)+A_\Gamma(s)+A_\zeta(s)
\qquad
A_0(s)=\frac1s+\frac1{s-1}
\qquad
A_\zeta(s)=\frac{\zeta'(s)}{\zeta(s)}
\tag{133.1}
\]

The completed logarithmic derivative is odd under \(s\mapsto1-s\). Hence

\[
M_\xi(q)
=\mathfrak F[A_0](q)
+\mathfrak G_\triangle(q)
+\mathfrak F[A_\zeta](q)
\qquad
\mathfrak F[A_0](q)=-\frac1q
\tag{133.2}
\]

For \(m\ge1\), the Gamma term has residue \(+1/2\) at
\(q=-2\mathsf T_{2m}\). The simple trivial zero of \(\zeta\) at
\(s=-2m\) contributes residue \(-1/2\) through the reflected zeta term.
At \(q=0\), the elementary, Gamma, and zeta residues are respectively

\[
-1,\qquad\frac12,\qquad\frac12
\tag{133.3}
\]

\begin{quote}
\textbf{Proposition 133.1 (exact trivial-scaffold cancellation).} Every
pole \(q=-2\mathsf T_{2m}\) of the Gamma scaffold is removable in the
completed folded resolvent. None is a zero or spectral atom of \(\xi\).
\end{quote}

Binet's formula supplies a complementary positive-kernel description.
For \(z>0\),

\[
\log\Gamma(z)
=\left(z-\frac12\right)\log z-z+\frac12\log(2\pi)
+\int_0^\infty e^{-zt}K_B(t)\,\frac{dt}{t}
\tag{133.4}
\]

where

\[
K_B(t)
=\frac1{e^t-1}-\frac1t+\frac12
=\frac12\coth(t/2)-\frac1t>0
\tag{133.5}
\]

Indeed, with \(u=t/2\), positivity is equivalent to
\(u\cosh u-\sinh u>0\); its derivative is \(u\sinh u>0\).
Differentiation gives

\[
A_\Gamma(s)
=\frac12\log\frac{s}{2\pi}-\frac1{2s}
-\frac12\int_0^\infty e^{-st/2}K_B(t)\,dt
\tag{133.6}
\]

Thus the Binet remainder is completely monotone, but it enters
\(A_\Gamma\) with a minus sign. The safe-half-line block cannot be
replaced by the positive scaffold. In fact, with \(G\) from (54.3) and
\(h_\Gamma(x)=xG(x)\),

\[
\boxed{
\lim_{x\downarrow0}h_\Gamma'(x)
=\frac{\psi(1/2)-\log\pi}{2}
=-\frac{\gamma_E+\log(4\pi)}2<0
}
\tag{133.7}
\]

\begin{quote}
\textbf{Proposition 133.2 (Gamma-only source no-go).} The actual source
block \(L_\Gamma\) is not positive semidefinite even at order one for
all positive nodes. Completion-level Gamma-prime cancellation is
necessary.
\end{quote}

The scaffold figure in Reading Volume §R17 summarizes these exact
cancellations.

\section{133A. Quarter-Gamma products and a positive triangular
comparison}\label{a.-quarter-gamma-products-and-a-positive-triangular-comparison}

The half-sized Gamma argument selects the even triangular ladder. For
\(a=s/2\), \(b=(1-s)/2\), and \(q=s(1-s)\),

\[
4(m+a)(m+b)=2T_{2m}+q.
\tag{133A.1}
\]

On \(s=1/2+it\), the two arguments are conjugates. Set
\(C(q)=\Gamma(s/2)\Gamma((1-s)/2)\), independent of the selected root of
\(s(1-s)=q\). The convergent Gamma product gives

\[
\begin{aligned}
\frac{\Gamma(1/4)^2}{|\Gamma(1/4+it/2)|^2}
 &=\prod_{m=0}^\infty\left(1+\frac{t^2}{2T_{2m}+1/4}\right),\\
-\frac{d}{dq}\log C(q)
 &=\sum_{m=0}^\infty\frac1{q+2T_{2m}}=2\mathfrak G_\triangle(q),\\
E(q):=\frac{2\sqrt\pi}{qC(q)}
 &=\frac{\Gamma(3/2)}{\Gamma(1+s/2)\Gamma(1+(1-s)/2)}
 =\prod_{m=1}^\infty\left(1+\frac q{2T_{2m}}\right).
\end{aligned}
\tag{133A.2}
\]

Here the logarithmic derivative is for \(q>0\), and \(E(0)=1\) by
removal of the zero-point singularity. The normalization follows from
\(C(q)\sim2\sqrt\pi/q\). Alternatively the finite product equals
\(\Gamma(N+1+a)\Gamma(N+1+b)\Gamma(3/2)\) divided by
\(\Gamma(1+a)\Gamma(1+b)\Gamma(N+1)\Gamma(N+3/2)\); the \(N\)-dependent
ratio tends to one because \(a+b=1/2\). This proves the product and its
normal convergence directly.

Consequently \(E'(0)=\sum_{m\ge1}(2T_{2m})^{-1}=1-\log2\), recovering
(TR.3). Euler's Beta substitution \(v=e^x/(1+e^x)\) also gives

\[
C(1/4+t^2)=2\sqrt{2\pi}\int_0^\infty
 \frac{\cos(ty)}{\sqrt{\cosh y}}dy.
\tag{133A.3}
\]

The reciprocal sum is the negative logarithmic derivative of this
integral, not the integral itself. Differentiating (133A.2) gives
\((\log C)''(q)=\sum_{m\ge0}(q+2T_{2m})^{-2}>0\) and
\((\log E)''(q)=-\sum_{m\ge1}(q+2T_{2m})^{-2}<0\). The signed cosine
integrand is not a positive measure in the parameter \(t\); the product
proves these \(q\)-curvature statements.

The reciprocal telescope has a compatible Gaussian difference formula:

\[
\frac1{T_x}=\frac2{\sqrt\pi}\int_0^\infty
 \frac{e^{-x^2u}-e^{-(x+1)^2u}}{\sqrt u}du\quad(x>0),\qquad
\sum_{n=1}^N\frac1{T_n}=2\left(1-\frac1{N+1}\right).
\tag{133A.4}
\]

Scaling Euler's Gamma integral proves the first identity; finite
telescoping and monotone convergence give the infinite sum two. With
\(f_n(t)=n^{-1/2-it}\), the full adjacent product retains the phase:

\[
(f_n-f_{n+1})\overline{(f_n+f_{n+1})}
=\frac1{2T_n}+\frac{2i}{\sqrt{2T_n}}
 \sin\left(t\log\frac{n+1}{n}\right).
\tag{133A.5}
\]

Its real parts sum to one; the complex series converges absolutely for
fixed \(t\). This is an exact connection to the Mellin weights of
Section 153, not a zero-location assertion.

\textbf{Proposition 133A.1 (all-order Li positivity of the comparison).}
Define \(\mathcal M(s)=E(s(s-1))\), explicitly \(E(-q(s))\) rather than
\(E(q(s))\). For \(a_m=2T_{2m}\), let \(\cos\theta_m=1-1/(2a_m)\) and
define \(\log\mathcal M(1/(1-z))=\sum_{n\ge1}\ell_nz^n/n\). Then

\[
\boxed{\ell_n=2\sum_{m\ge1}[1-\cos(n\theta_m)],\qquad
0<\ell_n\le n^2(1-\log2).}
\tag{133A.6}
\]

\textbf{Proof.} Under \(s=1/(1-z)\) each factor of the product becomes
\((1-e^{i\theta_m}z)(1-e^{-i\theta_m}z)/(1-z)^2\). Taking its local
logarithm gives the coefficient in (133A.6). The bound
\(4\sin^2(n\theta_m/2)\le n^2/a_m\) proves convergence; for fixed \(n\),
sufficiently small positive \(\theta_m\) give strictly positive
summands. \(\square\)

Thus \(\ell_1=1-\log2\), whereas the actual first Li value is
\(\lambda_1=(1-\log2)+(\gamma_E-\log\pi)/2\). The first correction is
already negative. The remaining local logarithm
\(\log(\xi(s)/(\xi(1)\mathcal M(s)))\) needs a separate source estimate.
The cancellation in Section 133 still applies; no comparison spectrum is
identified with the actual completed spectrum.

For comparison, scalar log-convexity alone is insufficient. The positive
real-axis Laplace transform
\(\mathcal F(s)=(2+\cosh(4(s-1/2)))/(2+\cosh2)\) is reflection symmetric
and has
\((\log\mathcal F)''=16(2\cosh(4(s-1/2))+1)/(2+\cosh(4(s-1/2)))^2>0\) on
the real axis. Its zeros nevertheless have real parts
\(1/2\pm\operatorname{arcosh}(2)/4\) and imaginary parts
\((2k+1)\pi/4\). Positivity of the comparison in (133A.6) uses its
particular product, not merely scalar convexity.

\Needspace{11\baselineskip}

\chapter{Technical chapter V. Flux geometry and rational
witnesses}\label{technical-chapter-v.-flux-geometry-and-rational-witnesses}

One energy gives the collar flux, the parabolic pullback and the
rational witness criterion. Its zero-height slice leads to the
source-rung calculations.

\section{20. Normalization}\label{normalization}

Define

\[
\Xi(t)=\xi\!\left(\frac12+it\right)
\]

and for \(c>0\),

\[
E_c(t)=\xi\!\left(\frac12+c-it\right)=A_c(t)-iB_c(t)
\]

The Wronskian is

\[
\Omega_c(t)=A_cB_c'-A_c'B_c
\]

Cauchy-Riemann gives the exact identities

\[
\boxed{
\Omega_c(t)=\frac12\partial_c|E_c(t)|^2
=|E_c(t)|^2\Re\frac{\xi'}{\xi}
\!\left(\frac12+c-it\right)
}
\]

The canonical normalization in this volume is

\[
\boxed{J_c(t)=\frac{\Omega_c(t)}c}
\]

The division by \(c\) makes \(J_c\) the derivative with respect to the
square coordinate \(c^2\), as proved in §25.

\section{21. Generalized Laguerre
hierarchy}\label{generalized-laguerre-hierarchy}

For a real entire function \(f\), define

\[
L_n[f](t)=\frac1{(2n)!}
\sum_{k=0}^{2n}(-1)^{n+k}\binom{2n}{k}
f^{(k)}(t)f^{(2n-k)}(t)
\]

In particular,

\[
L_1[f]=f'^2-ff''
\]

The exact generating laws are

\[
\boxed{
|E_c(t)|^2=\sum_{n=0}^{\infty}L_n[\Xi](t)c^{2n}
}
\]

and

\[
\boxed{
J_c(t)=\sum_{n=1}^{\infty}nL_n[\Xi](t)c^{2n-2}
}
\]

At the fold,

\[
J_0(t)=L_1[\Xi](t)=\Xi'(t)^2-\Xi(t)\Xi''(t)
\]

This series already suggests the correct differentiation variable:
\(c^2\), not \(c\).

\begin{figure}
\centering
\includegraphics{figures/laguerre_taylor_xi.png}
\caption{Direct xi evaluation and its even generalized-Laguerre Taylor
reconstruction at t equals 10.}
\end{figure}

\section{22. Energy and flux kernels}\label{energy-and-flux-kernels}

Let \(\Phi\) be the even theta kernel with

\[
\Xi(t)=\widehat\Phi(t)
\]

If \(f_c(x)=e^{-cx}\Phi(x)\), then the energy kernel

\[
\mathcal H_c=f_c*\widetilde f_c
\]

satisfies

\[
\widehat{\mathcal H_c}(t)=|E_c(t)|^2\ge0
\]

Differentiating in the collar direction gives the flux kernel

\[
\boxed{
\mathcal C_c(x)=\frac1{2c}
\int_{\mathbb R}(x-2u)\sinh(c(x-2u))
\Phi(u)\Phi(x-u)\,du
}
\]

It obeys

\[
\boxed{J_c(t)=\widehat{\mathcal C_c}(t)}
\]

Pointwise \(\mathcal C_c(x)\ge0\) because \(y\sinh(cy)\ge0\). However,

\[
\mathcal C_c(x)\ge0
\centernot\Longrightarrow
\widehat{\mathcal C_c}(t)\ge0
\]

The missing statement is positive definiteness of \(\mathcal C_c\).

\begin{figure}
\centering
\includegraphics{figures/fourier_positivity_firewall.png}
\caption{A nonnegative function with a sign-changing Fourier transform:
the positive-definiteness firewall.}
\end{figure}

\section{22A. Every rung test has essential Fourier
cancellation}\label{a.-every-rung-test-has-essential-fourier-cancellation}

The firewall closing §22 is not a gap in the argument; it is a theorem,
and it holds at every order. Fix \(x>0\), put \(b=\tfrac12\) and
\(a=\sqrt{x+\tfrac14}>b\), and let

\[
f_{k,x}(t)=\left[\frac{4x(t^2+b^2)}{(t^2+a^2)^2}\right]^{k},
\qquad
\widehat f_{k,x}(u)=\int_{\mathbb R}f_{k,x}(t)e^{-itu}\,dt
\tag{22A.1}
\]

Here \(t\) is an independent real frequency; it is not a substitution
for the imaginary part of an off-line zero, and (22A.1) is a scalar
test, not the completed Weil form.

\begin{quote}
\textbf{Theorem 22A.1.} For every \(x>0\) and every integer \(k\ge1\),
the real even transform \(\widehat f_{k,x}\) has at least \(k\) sign
changes on \(u>0\). Consequently \(f_{k,x}\) is not positive definite.
For odd \(k\) the transform is eventually negative; for even \(k\) it is
eventually positive and negative on some nonempty interval.
\end{quote}

\emph{Proof.} Parameter differentiation of the elementary rational
transform gives \(\Phi_m(u;a)=\int_{\mathbb R}(t^2+a^2)^{-m}e^{-itu}dt
=\frac{(-1)^{m-1}}{(m-1)!}\partial_{a^2}^{\,m-1}\bigl(\tfrac\pi ae^{-a|u|}\bigr)\).
Expanding \(t^2+b^2=(t^2+a^2)-x\),

\[
\widehat f_{k,x}(u)=(4x)^k\sum_{j=0}^k\binom kj(-x)^j\Phi_{k+j}(u;a)
\tag{22A.2}
\]

which for \(u\ge0\) is \(e^{-au}\) times a real polynomial of degree
\(2k-1\) whose leading coefficient is
\((-1)^k2\pi x^{2k}/\bigl[(2k-1)!\,a^{2k}\bigr]\). So the transform has
finitely many sign changes, exponential moments for \(|y|<a\), and
\(\widehat f_{k,x}(0)=\int f_{k,x}>0\). Analytic inversion in the strip
\(|\Im t|<a\) gives

\[
\frac1{2\pi}\int_{\mathbb R}e^{yu}\widehat f_{k,x}(u)\,du
=f_{k,x}(-iy)
=\left[\frac{4x(b^2-y^2)}{(a^2-y^2)^2}\right]^{k}
\tag{22A.3}
\]

whose right side has a zero of order exactly \(k\) at each of
\(y=\pm b\), both strictly inside the strip since \(x>0\).

Now the variation-diminishing step. A continuous real \(g\not\equiv0\)
with \(M\) sign changes and finite exponential polynomial moments has a
bilateral Laplace transform with at most \(M\) real zeros, counted with
multiplicity, inside the interval of convergence: for \(M=0\) the
transform has one strict sign; for the induction step, multiplying \(g\)
by \((u-u_0)\) at a sign change removes one sign change and replaces
\(F\) by \(F'-u_0F=e^{u_0y}\bigl(e^{-u_0y}F\bigr)'\), and Rolle's
theorem with multiplicities closes the induction on each compact
subinterval. Applied to (22A.3), the \(2k\) zeros force at least \(2k\)
sign changes on the whole axis; evenness and \(\widehat f_{k,x}(0)>0\)
put at least \(k\) of them on \(u>0\). The leading coefficient gives the
eventual signs. \(\square\)

The obstruction is already visible at weight zero: the identity

\[
\boxed{\ \int_{\mathbb R}\cosh(u/2)\,\widehat f_{k,x}(u)\,du=0\ }
\tag{22A.4}
\]

rules out an everywhere nonnegative transform, since a nonzero
nonnegative density would have strictly positive weighted integral.

\textbf{Scope, and the \(k=1\) consistency check.} This closes, at all
orders, the route that asserts the rung sign by claiming these scalar
tests have nonnegative transforms. It does \emph{not} exclude a separate
Weil-functional representation, a two-variable positivity construction,
or an argument using more of the arithmetic source. For \(k=1\) the
transform's unique positive root is exactly the crossover
\(u_x=(1+\tfrac1{2x})/a\) recorded at §R11 --- at \(x=1,3,10\) it is
\(1.34164079\), \(0.64715023\), \(0.32796490\) --- so Theorem 22A.1
contains the volume's exact first-rung statement as its first case
rather than competing with it. A residue computation of (22A.2) for
\(k\le4\) at \(x=1,3,10\) returns degree \(2k-1\), leading sign
\((-1)^k\), and exactly \(k\) positive roots in all twelve cases;
whether the count is exactly \(k\) in general is \textbf{open and is not
used}.

\section{23. Signed zero blocks}\label{signed-zero-blocks}

Let an off-axis zero quartet have horizontal displacement \(a>0\). Its
paired contribution to the normalized horizontal log-derivative is

\[
\boxed{
H_{a,c}(u)=\frac1c\left[
\frac{c-a}{(c-a)^2+u^2}
+\frac{c+a}{(c+a)^2+u^2}
\right]
}
\]

Equivalently,

\[
H_{a,c}(u)=
\frac{2(c^2+u^2-a^2)}
{((c-a)^2+u^2)((c+a)^2+u^2)}
\]

If \(a<c\), the block is positive for every \(u\). If \(a>c\), it is
negative exactly when

\[
|u|<\sqrt{a^2-c^2}
\]

Thus

\[
b=\sqrt{a^2-c^2}
\]

is the exact negative-core half-width. It is not an optimal smoothing
radius.

At a fixed \(c\), a dangerous block can be hidden. If \(k\)
critical-line blocks are placed at the same ordinate, complete
compensation occurs exactly when

\[
k(a^2-c^2)\ge2c^2
\]

This firewall is correct, but it concerns a fixed horizontal line. It
does not survive the limit in which \(c\) approaches the target zero
itself; that limit is the key to Theorem 27.

\begin{figure}
\centering
\includegraphics{figures/signed_zero_block_phase.png}
\caption{Exact sign region of a synthetic off-axis zero block.}
\end{figure}

\section{24. Master energy and collar
variables}\label{master-energy-and-collar-variables}

Put

\[
\mathcal M(c,t)
=\left|\xi\!\left(\frac12+c-it\right)\right|^2
\qquad c>0
\]

Functional equation and conjugation imply

\[
\mathcal M(-c,t)=\mathcal M(c,t)
\]

Hence \(\mathcal M\) is analytic in the square variable

\[
x=c^2
\]

On the complementary collar \(0<c<1/2\), define

\[
p=\frac12+c
\qquad
\tau=p(1-p)=\frac14-c^2
\qquad
r=\frac{1+2c}{1-2c}
\]

The fold is \(c=0\), \(p=1/2\), \(\tau=1/4\), \(r=1\). Moving outward on
either sheet increases \(|c|\), decreases \(\tau\), and sends \(r\) away
from its fixed radius.

\section{25. Theorem A - Cayley-Laguerre flux
identity}\label{theorem-a---cayley-laguerre-flux-identity}

For every \(c>0\) and every real \(t\) at which the displayed
logarithmic derivative is finite,

\[
\boxed{
\frac{\partial}{\partial(c^2)}\mathcal M(c,t)=J_c(t)
}
\]

and

\[
\boxed{
\frac{\partial}{\partial(c^2)}\log\mathcal M(c,t)
=Q_c(t)
:=\frac1c\Re\frac{\xi'}{\xi}
\!\left(\frac12+c-it\right)
}
\]

Equivalently,

\[
\boxed{
J_c(t)=-\frac{\partial}{\partial\tau}\mathcal M(\tau,t)
\qquad
Q_c(t)=-\frac{\partial}{\partial\tau}\log\mathcal M(\tau,t)
}
\]

\textbf{Proof.} By the Wronskian identity,

\[
\partial_c\mathcal M(c,t)=2\Omega_c(t)=2cJ_c(t)
\]

Since \(\partial_{c^2}=(2c)^{-1}\partial_c\), the first identity
follows. Dividing by \(\mathcal M\) gives the logarithmic identity.
Finally \(\tau=1/4-c^2\), so \(\partial_\tau=-\partial_{c^2}\).
\(\square\)

The phase-complement product \(\tau\), the generalized Laguerre flux
\(J_c\), and the horizontal log-derivative \(Q_c\) are linked by these
derivatives of one energy. The next theorem identifies its
folded-resolvent coordinate.

\section{26. Theorem B - Folded-resolvent parabolic
pullback}\label{theorem-b---folded-resolvent-parabolic-pullback}

Let

\[
q=s(1-s),\qquad \xi(s)=X(q),\qquad
M_\xi(q)=-\frac{X'(q)}{X(q)}
\]

For

\[
s=\frac12+c-it
\]

the folded coordinate is

\[
\boxed{
q(c,t)=\tau+t^2+2ict
=\frac14-c^2+t^2+2ict
}
\]

The logarithmic derivative factors as

\[
\boxed{
\frac{\xi'}{\xi}(s)=2(c-it)M_\xi(q(c,t))
}
\]

Consequently

\[
\boxed{
Q_c(t)=\frac2c
\Re\bigl[(c-it)M_\xi(q(c,t))\bigr]
}
\]

\textbf{Proof.} Since \(dq/ds=1-2s=-2(c-it)\) and \(X'/X=-M_\xi\),

\[
\frac{\xi'}{\xi}(s)
=\frac{X'}X(q)\frac{dq}{ds}
=2(c-it)M_\xi(q)
\]

Take the real part and divide by \(c\). \(\square\)

\subsection{Corollary 26.1 - Fixed-focus parabolic
foliation}\label{corollary-26.1---fixed-focus-parabolic-foliation}

Write

\[
q(c,t)=u+iv
\]

For fixed \(c>0\),

\[
u=\tau+t^2,\qquad v=2ct
\]

and elimination of \(t\) gives

\[
\boxed{
v^2=4c^2(u-\tau)
\qquad
\tau=\frac14-c^2
}
\]

Thus every fixed-\(c\) pullback path is a focus-directrix parabola. Its
vertex, focus, and directrix are

\[
\boxed{
V_c=\left(\frac14-c^2,0\right),\qquad
F_*=\left(\frac14,0\right),\qquad
L_c:\ u=\frac14-2c^2
}
\]

The focus \(F_*\) is independent of \(c\), the Cayley complement
\(\tau\) is the vertex coordinate, and \(c^2\) is the focal length.
Directly,

\[
\left|q(c,t)-\frac14\right|
=c^2+t^2
=\operatorname{dist}(q(c,t),L_c)
\]

At \(c=0\) the leaf degenerates to the real focal ray \([1/4,\infty)\).
Hence the folded zero statement admits the exact conic form

\[
\boxed{
\mathrm{RH}
\iff
\text{every folded zero lies on the degenerate focal ray,}
}
\]

whereas a zero \(\rho=1/2+a+i\gamma\) with \(a\ne0\) lies on the genuine
confocal leaf \(c=|a|\). This is the analytic focus-directrix bridge
described by the shared centered square; it is a reformulation of zero
location, not a forcing theorem.

\begin{figure}
\centering
\includegraphics{figures/folded_zero_geometry.png}
\caption{Synthetic reflection orbits in the folded q-plane; critical
orbits map to the real ray.}
\end{figure}

\subsection{Consequence}\label{consequence}

PP1 and the signed zero-Poisson program are two coordinate views of one
resolvent. PP1 studies \(M_\xi\) on the negative real \(q\)-axis, where
the prime Dirichlet series converges after the \(s(x)>1\) substitution.
The HB/Poisson program studies directional real parts of the same
function along the parabolas \(q(c,t)\).

This is an exact intertwiner, not a new estimate.

\subsection{Prior-art boundary}\label{prior-art-boundary}

The horizontal-modulus core of the next theorem is not claimed as new.
Sondow and Dumitrescu proved that \(|\xi|\) is horizontally increasing
in a zero-free right half-plane and extracted the RH reformulation;
Matiyasevich, Saidak, and Zvengrowski gave a related treatment and
traced the idea through Jensen/Polya and Lagarias {[}12-14{]}. The
contribution of this volume is the proof-preserving identification of
that known monotonicity criterion with \(J_c\), \(\tau\), the folded
resolvent, the flux kernel, and the native rational TFG witness system.
The proof is repeated to make every normalization and converse
quantifier explicit.

\section{27. Theorem C - Outward-modulus criterion in flux
coordinates}\label{theorem-c---outward-modulus-criterion-in-flux-coordinates}

Let \(F(z)=\xi(1/2+z)\). The following are equivalent.

\begin{enumerate}
\def\labelenumi{\arabic{enumi}.}
\tightlist
\item
  The Riemann hypothesis is true.
\item
  For every \(c>0\) and \(t\in\mathbb R\), \[
  \Re\frac{F'}F(c-it)>0
  \]
\item
  For every \(c>0\) and \(t\in\mathbb R\), \[
  J_c(t)>0
  \]
\item
  For every fixed \(t\), the function \[
  c\longmapsto\mathcal M(c,t)
  \] is strictly increasing on \((0,\infty)\).
\item
  On \(0<c<1/2\), the same energy is strictly decreasing as a function
  of \(\tau=1/4-c^2\).
\item
  For every \(c>0\), the flux kernel \(\mathcal C_c\) is positive
  definite.
\end{enumerate}

The strict inequalities may be replaced by nonnegative/nondecreasing
inequalities without changing the equivalence.

\subsection{Proof that RH implies outward
positivity}\label{proof-that-rh-implies-outward-positivity}

Under RH the zeros of \(F\) are \(\pm i\gamma\), with multiplicity. The
paired Hadamard product is

\[
F(z)=F(0)\prod_{\gamma>0}
\left(1+\frac{z^2}{\gamma^2}\right)^{m_\gamma}
\]

Because \(\sum\gamma^{-2}<\infty\), its logarithmic derivative has the
absolutely convergent real-part expansion

\[
\frac{F'}F(z)=
\sum_{\gamma>0}m_\gamma
\left(\frac1{z-i\gamma}+\frac1{z+i\gamma}\right)
\]

At \(z=c-it\),

\[
\Re\frac{F'}F(c-it)
=\sum_{\gamma>0}m_\gamma\left[
\frac{c}{c^2+(t+\gamma)^2}
+\frac{c}{c^2+(t-\gamma)^2}
\right]>0
\]

Theorem A then gives strict outward monotonicity and \(J_c>0\).

\subsection{Proof that outward nonnegativity implies
RH}\label{proof-that-outward-nonnegativity-implies-rh}

If RH fails, functional-equation symmetry supplies a zero

\[
z_0=a-i\gamma
\]

of \(F\) with \(a>0\). Let its multiplicity be \(m\). In a disk
containing no other zero,

\[
\frac{F'}F(z)=\frac m{z-z_0}+h(z)
\]

where \(h\) is analytic. For \(d>0\) small, put

\[
w=z_0-d=(a-d)-i\gamma
\]

Then

\[
\Re\frac{F'}F(w)=-\frac md+O(1)<0
\]

Thus \(J_{a-d}(\gamma)<0\) for all sufficiently small \(d\). This
contradicts nonnegativity. Hence all zeros lie on the critical line.

\subsection{Positive-definiteness
equivalence}\label{positive-definiteness-equivalence}

The kernel \(\mathcal C_c\) is continuous and rapidly decreasing, and
\(\widehat{\mathcal C_c}=J_c\). Bochner's theorem therefore identifies
\(J_c\ge0\) with positive definiteness of \(\mathcal C_c\). This
completes the equivalence. \(\square\)

\section{28. Quantitative local witness
theorem}\label{quantitative-local-witness-theorem}

The preceding converse has a useful quantitative form.

Let \(z_0=a-i\gamma\) be an off-axis zero of multiplicity \(m\). Choose
\(r>0\) so that \(z_0\) is the only zero in \(|z-z_0|\le r\), after
multiplicity is absorbed into the target factor. Put

\[
h(z)=\frac{F'}F(z)-\frac m{z-z_0}
\qquad
B=\sup_{|z-z_0|\le r/2}|\Re h(z)|
\]

For

\[
0<d<\min\left\{a,\frac r2,\frac m{2B}\right\}
\]

with the last bound omitted if \(B=0\), one has

\[
\boxed{
d\Re\frac{F'}F(z_0-d)
\le-\frac m2
}
\]

Indeed,

\[
d\Re\frac{F'}F(z_0-d)
=-m+d\Re h(z_0-d)
\le-m+dB\le-\frac m2
\]

Therefore every off-axis zero creates an \textbf{open} negative flux
region in the \((c,t)\) plane. Fixed-\(c\) compensation cannot remove
that region because the target pole diverges as \(c\uparrow a\) while
all non-target contributions remain bounded locally.

\begin{figure}
\centering
\includegraphics{figures/local_witness_scaling.png}
\caption{Synthetic local-pole scaling: the normalized witness tends to
minus the target multiplicity.}
\end{figure}

\section{29. Theorem D - Native TFG rational witness
criterion}\label{theorem-d---native-tfg-rational-witness-criterion}

Let

\[
\mathcal C_\triangle=\mathbb Q\cap(0,1/2)
\]

be represented by centered native cells as in Proposition 8.1.
Independently, address rational heights by unreduced integer pairs:

\[
\mathcal T_{\mathbb Q}
=\left\{\frac pq:p\in\mathbb Z,\ q\in\mathbb N\right\}
=\mathbb Q
\]

Then RH is equivalent to the countable family of inequalities

\[
\boxed{
\mathcal M(c_2,t)\ge\mathcal M(c_1,t)
\quad
\text{for all }
t\in\mathcal T_{\mathbb Q}
\quad
c_1,c_2\in\mathcal C_\triangle
\quad
c_1<c_2
}
\]

Equivalently, for all finite indices satisfying

\[
c_i=\frac{|2k_i-n_i|}{2n_i}
\qquad
0<|2k_i-n_i|<n_i
\]

and all rational heights \(t=p/q\), the outward finite difference is
nonnegative.

\textbf{Proof.} RH implies strict monotonicity for every real \(c,t\) by
Theorem C, so it implies the displayed rational inequalities.
Conversely, if RH fails, the quantitative local witness theorem produces
an open rectangle on which \(\partial_c\mathcal M<0\). Choose rational
\(t\) in its height projection and rational \(c_1<c_2\) in its
horizontal projection. Then

\[
\mathcal M(c_2,t)<\mathcal M(c_1,t)
\]

and Proposition 8.1 supplies finite native-cell indices for both
\(c_i\). \(\square\)

\subsection{Interpretation}\label{interpretation}

The exact rational coverage in §8 turns any open negative-flux region
into a finite, rationally addressed witness cell. No literature-priority
claim is made. A finite witness can refute the full criterion; proving
it still requires a uniform mechanism that orders every addressed cell.

\section{30. Synthetic diagnostic}\label{synthetic-diagnostic}

The following figure compares two exact finite zero models. It is a
diagnostic of the theorem, not data about unknown zeta zeros.

\begin{figure}
\centering
\includegraphics{figures/cayley_laguerre_flux_diagnostic.png}
\caption{Synthetic critical-line and off-axis flux fields, together with
the Cayley collar coordinate.}
\end{figure}

The left panel uses a polynomial with a critical-line pair
\(\pm i\gamma\); the outward logarithmic flux is positive. The middle
panel uses the symmetric off-axis quartet \(\pm a\pm i\gamma\); a blue
negative cell opens immediately to the left of the target. The right
panel shows that outward motion in \(c\) decreases the complement
product \(\tau\) and increases log odds.

\section{31. Theorem E - Axial flux-Widder
identity}\label{theorem-e---axial-flux-widder-identity}

The parabolic pullback contains the negative-axis folded resolvent as
its zero-height slice. Thus the reduced-Widder hierarchy is the axial
derivative ladder of the same master flux object, not a second analytic
program.

For \(x>0\), set

\[
c(x)=\sqrt{x+\frac14}
\qquad
s(x)=\frac12+c(x)=\frac{1+\sqrt{1+4x}}2
\]

Then \(c>1/2\), and at \(t=0\) the folded coordinate of Theorem B is

\[
q(c,0)=\frac14-c^2=-x
\]

Consequently,

\[
\boxed{
Q_{c(x)}(0)=2M_\xi(-x)=2S_\xi(x)
}
\tag{31.1}
\]

Since \(x=c^2-1/4\), differentiation in \(x\) agrees with
differentiation in \(c^2\) along this slice. Theorem A therefore also
gives

\[
\boxed{
\frac{d}{dx}
\log\left|
\xi\!\left(\frac{1+\sqrt{1+4x}}2\right)
\right|^2
=2S_\xi(x)
}
\tag{31.2}
\]

Define the reduced-Widder diagonal

\[
W_k(x)=(-1)^{k-1}D_x^{2k-1}[x^kS_\xi(x)]
\qquad k\ge1
\]

Substitution of (31.1) gives the exact axial-flux formula

\[
\boxed{
W_k(x)=\frac{(-1)^{k-1}}2
D_x^{2k-1}
\left[x^kQ_{\sqrt{x+1/4}}(0)\right]
}
\tag{31.3}
\]

Thus the hierarchy measures alternating higher differential curvature of
the outward logarithmic flux, with the collar parameter and the folded
negative-axis parameter related by \(x=c^2-1/4\).

\begin{figure}
\centering
\includegraphics{figures/axial_flux_widder_bridge.png}
\caption{The negative folded axis is the zero-height slice of the
confocal parabolic flux, and its axial flux generates every
reduced-Widder rung.}
\end{figure}

\section{42. One master object, several
renderings}\label{one-master-object-several-renderings}

The exact coordinate relationships are

\[
\boxed{
\begin{array}{c}
X(q)\\[2pt]
\downarrow\ -\partial_q\log\\[2pt]
M_\xi(q)
\end{array}
}
\,\xrightarrow{\ q=s(1-s)\ }\,
\boxed{
Q_c(t)=\frac1c\Re\frac{\xi'}{\xi}(s)
}
\,\xrightarrow{\ \times\mathcal M\ }\,
\boxed{J_c(t)}
\,\xleftrightarrow{\ \mathcal F\ }\,
\boxed{\mathcal C_c(x)}
\]

The coordinate collar

\[
c^2=\frac14-\tau
\qquad
\tau=p(1-p)=\frac{r}{(1+r)^2}
\]

turns \(J_c\) into the negative \(\tau\)-derivative of the energy. The
Laguerre hierarchy is its Taylor expansion at the fold. The signed
Poisson blocks are its zero-side partial fractions. The PP1 Stieltjes
problem is its folded analytic-class statement. Each rendering retains
the same completed resolvent; a sign theorem still requires its full
domain and quantifiers.

\section{43. Mesoscopic localization and rational witness
size}\label{mesoscopic-localization-and-rational-witness-size}

For a hypothetical target \(z_0=a-i\gamma\), choose \(0<c<a\) and write

\[
d=a-c
\qquad
D=a+c
\qquad
b=\sqrt{dD}
\qquad
w=c-i\gamma
\]

The target contributes exactly

\[
d\Re\frac{m_0}{w-z_0}=-m_0
\]

The PP76 mesoscopic owner is

\[
M_{\mathrm{meso}}
=d\Re\sum_{Cd<|\lambda-z_0|<b}
\frac{m_\lambda}{w-\lambda}
\]

Current every-target counting gives a crowding scale
\(C\asymp\log\gamma\). Proving \(C=o(\log\gamma)\) would be genuine
progress; an \(O(\sqrt{\log\gamma})\) result would be major.

The local witness theorem fixes the role of this mesoscopic estimate:

\begin{enumerate}
\def\labelenumi{\arabic{enumi}.}
\tightlist
\item
  \textbf{It is not needed to prove the equivalence.} Near a fixed
  off-axis zero, the target pole dominates automatically.
\item
  \textbf{It remains useful for fixed-line propagation.} If positivity
  is known at selected \(c\) values and one wants to force a neighboring
  zero or propagate control between scales, the signed discrepancy is
  the least-lossy owner.
\item
  \textbf{A sublog bound alone is not RH.} The local sign hypothesis,
  global coverage in \(c\), and an iteration or termination theorem are
  still needed.
\item
  \textbf{The signed field is primary.} Counts, quadratic energy, second
  harmonics, and fixed finite-Weil compressions lose directional
  information.
\end{enumerate}

The recommended successor problem is therefore two-pronged:

\begin{quote}
\textbf{Source question --- collar monotonicity.} Prove
\(-\partial_\tau\log\mathcal M\ge0\) from the explicit prime/Gamma
source, equivalently prove the PP1 Stieltjes/Pick/Loewner property.
\end{quote}

\begin{quote}
\textbf{Quantitative question --- witness localization.} Develop
validated bounds for the size of the negative rational witness box
around a hypothetical off-axis zero and determine what zero-spacing
input would make a finite mesh certificate possible up to height \(T\).
\end{quote}

The source question is RH-equivalent. The quantitative localization
question controls witness size; it is not by itself a proof of RH.

\Needspace{11\baselineskip}

\chapter{Technical chapter VI. Widder source proofs and their asymptotic
limits}\label{technical-chapter-vi.-widder-source-proofs-and-their-asymptotic-limits}

The full derivative identity fixes the normalization. The first four
source proofs appear together, followed by fixed-rung tails and the
limits of phase-based detection.

\section{32. Theorem F - Full Widder-zero-heat
identity}\label{theorem-f---full-widder-zero-heat-identity}

For integers \(n,k\ge0\), define the full Widder functional

\[
\mathcal F_{n,k}[S](x)
=(-1)^nD_x^{n+k}[x^kS(x)]
\]

The notation \(\mathcal F_{n,k}\) here is the same functional denoted by
\(F_{n,k}\) in reader (13.2). The reduced diagonal is

\[
W_k=\mathcal F_{k-1,k}
\]

Because \(X\) has order \(1/2\), its folded canonical product has genus
zero:

\[
X(q)=X(0)
\prod_{[\rho]}
\left(1-\frac q{q_\rho}\right)^{m_\rho}
\qquad
q_\rho=\rho(1-\rho)
\]

There is one factor per functional-equation orbit. An off-line quartet
folds to a conjugate pair \(q_\rho,\overline{q_\rho}\); it must not be
counted four times. Logarithmic differentiation on \(q=-x\) yields

\[
S_\xi(x)=\sum_{[\rho]}\frac{m_\rho}{x+q_\rho}
\tag{32.1}
\]

Polynomial division, or \(n+k\) direct differentiations, gives the
one-pole identity

\[
(-1)^nD_x^{n+k}\frac{x^k}{x+q}
=(n+k)!\frac{q^k}{(x+q)^{n+k+1}}
\]

Hence

\[
\boxed{
\mathcal F_{n,k}[S_\xi](x)
=(n+k)!
\sum_{[\rho]}m_\rho
\frac{q_\rho^k}{(x+q_\rho)^{n+k+1}}
}
\tag{32.2}
\]

Now introduce the folded heat trace

\[
\Theta_\xi(u)=\sum_{[\rho]}m_\rho e^{-uq_\rho}
\qquad u>0
\]

For fixed \(u>0\), this series and all of its \(u\)-derivatives converge
absolutely. The zero-counting estimate and \(q_\rho\asymp\gamma^2\) also
give absolute convergence after the weighted Laplace integration below.
Since

\[
\frac{(n+k)!}{(x+q)^{n+k+1}}
=\int_0^\infty e^{-(x+q)u}u^{n+k}\,du
\]

Tonelli-Fubini applied to absolute values gives

\[
\boxed{
\mathcal F_{n,k}[S_\xi](x)
=\int_0^\infty
e^{-xu}u^{n+k}(-1)^k\Theta_\xi^{(k)}(u)\,du
}
\tag{32.3}
\]

In particular,

\[
\boxed{
W_k(x)
=(2k-1)!
\sum_{[\rho]}m_\rho
\frac{q_\rho^k}{(x+q_\rho)^{2k}}
=\int_0^\infty e^{-xu}u^{2k-1}
(-1)^k\Theta_\xi^{(k)}(u)\,du
}
\tag{32.4}
\]

This proves the heat bridge directly from the spectral expansion. No
integration by parts is used, so there are no unverified endpoint terms.
It also explains the universal completion cancellation. The Laplace
transform of the completion term \(1/x\) is the constant heat mode
\(1\); every \(k\ge1\) heat derivative annihilates it.

\subsection{Folded multiplicity and preserved
information}\label{folded-multiplicity-and-preserved-information}

Here \([\rho]\) is the reflection pair \(\{\rho,1-\rho\}\), and
\(m_\rho\) is its common zero multiplicity. For an off-line quartet, the
two factors in the product are \((1-q/q_\rho)^{m_\rho}\) and
\((1-q/\bar q_\rho)^{m_\rho}\). On the critical line there is one such
factor, not two copies of it. No assumption of simplicity is made.

The factorization \(q(s_1)-q(s_2)=(s_1-s_2)(1-s_1-s_2)\) proves that the
only identification made by the fold is \(s\sim1-s\). Since nontrivial
zeros have nonzero imaginary part, \(q_\rho\) is real exactly when
\(\Re\rho=1/2\); in that case \(q_\rho=1/4+(\Im\rho)^2>1/4\). Reciprocal
and sign changes of \(q_\rho\) are invertible, whereas a real-part
observation is a further map. Any probability construction made from
these data must declare its own atoms and multiplicity convention; none
is needed for Theorem F.

\section{32A. The Widder array as a Laguerre transform of the heat
trace}\label{a.-the-widder-array-as-a-laguerre-transform-of-the-heat-trace}

Theorem F places the derivatives on the heat trace. Moving them onto the
kernel instead turns the whole array \(F_{n,k}\) into a single integral
transform against \(\Theta_\xi\), with a scalar kernel whose sign
pattern is explicit at every index.

\begin{quote}
\textbf{Theorem 32A.1.} For all integers \(n,k\ge0\) and every \(x>0\),
\[
\boxed{\ F_{n,k}[S_\xi](x)=k!\int_0^\infty e^{-xu}u^{n}L_k^{(n)}(xu)\,\Theta_\xi(u)\,du\ }
\tag{32.14}
\] with \(L_k^{(n)}(t)=\sum_{j=0}^k(-1)^j\binom{n+k}{k-j}t^j/j!\). In
particular, for every \(k\ge1\), \[
\boxed{\ W_k(x)=k!\int_0^\infty e^{-xu}u^{k-1}L_k^{(k-1)}(xu)\,\Theta_\xi(u)\,du\ }
\tag{32.15}
\] and for \(k\ge1\) the scalar kernel in (32.14) has exactly \(k\)
simple sign changes on \(u>0\).
\end{quote}

\emph{Proof.} Leibniz's rule gives the kernel identity directly:

\[
(-1)^nD_x^{\,n+k}\bigl[x^ke^{-xu}\bigr]=k!\,e^{-xu}u^nL_k^{(n)}(xu)
\tag{32.16}
\]

Differentiating the absolutely convergent heat representation under the
integral is justified term by term, since after expansion each
contribution is dominated by
\(\sum_{[\rho]}m_\rho\int_0^\infty e^{-(x+\Re q_\rho)u}u^{n+j}du
=(n+j)!\sum_{[\rho]}m_\rho(x+\Re q_\rho)^{-(n+j+1)}<\infty\), including
the least-decaying case \(n=j=0\), because \(\Re q_\rho\asymp\gamma^2\)
against a zero count \(O(T\log T)\). No sign for \(\Theta_\xi\) and no
hypothesis is used. Setting \(n=k-1\) gives (32.15), which requires
\(k\ge1\); this is consistent with \(W_k=F_{k-1,k}\) and introduces no
reduced \(W_0\).

For the sign count, Rodrigues' formula makes \(L_k^{(n)}\) orthogonal to
every polynomial of degree below \(k\) against the positive weight
\(t^ne^{-t}dt\); the \(k\) integrations by parts have vanishing boundary
terms because \(n\ge0\). If there were fewer than \(k\) positive sign
changes, multiplying by the polynomial with those points as roots would
give a nonzero one-sign integrand of degree below \(k\), contradicting
orthogonality. Degree \(k\) then forces exactly \(k\) simple positive
roots, and every other kernel factor is positive for \(u>0\).
\(\square\)

\textbf{What the two placements identify, and what neither creates.}
Nonnegativity of an undifferentiated heat trace cannot force these
signs: at any prescribed \(x\), a smooth nonnegative function supported
where \(L_k^{(n)}(xu)<0\) gives a negative functional. That is an
abstract countermodel to the inference, not a statement about the actual
Riemann heat trace. Keeping the derivatives on the heat trace instead,
as in Theorem F, requires exactly the complete-monotonicity signs
already equivalent to RH (§33). \textbf{The two exact placements of the
derivatives name the missing information from either side; neither
supplies it.} This is the same firewall as §22A in the Laguerre basis
rather than the Fourier one.

\section{33. The reduced-Widder criterion and its
firewall}\label{the-reduced-widder-criterion-and-its-firewall}

Reader (12.1)--(12.4) prints the first six differential rows and their
common recurrence. Here the full family supplies the criterion.

Widder's real-variable characterization, in the concise form restated
and proved by Sokal, says that a smooth real function \(S\) on
\((0,\infty)\) is Stieltjes if and only if

\[
S(x)\ge0
\quad\text{and}\quad
\mathcal F_{k-1,k}[S](x)\ge0
\quad(k\ge1, x>0)
\]

For \(S_\xi\), the baseline \(S_\xi(x)>0\) is unconditional: the
positive theta-kernel representation makes \(\xi(1/2+y)\) increasing for
\(y>0\). Combining Widder's theorem with the folded Stieltjes
equivalence gives

\[
\boxed{
\mathrm{RH}
\iff
W_k(x)\ge0
\quad\text{for every }k\ge1\text{ and every }x>0
}
\tag{33.1}
\]

This is an exact RH criterion. Its quantifiers are load-bearing.
Positivity of \(W_1,W_2,W_3\), or of any other currently verified finite
prefix, does not supply the complete diagonal hypothesis in Widder's
theorem. Nothing in the present work promotes a finite prefix to the
full Stieltjes class.

\section{34. Theorem G - General source-jet
ladder}\label{theorem-g---general-source-jet-ladder}

On the negative axis write

\[
\begin{gathered}
x=s(s-1),\qquad s=\sigma+i0>1,\qquad R=2s-1=\sqrt{1+4x}\\
L(s)=\frac{\xi'(s)}{\xi(s)},\qquad S_\xi(x)=\frac{L(s)}R
\end{gathered}
\]

Then \(D_x=R^{-1}D_s\). The elementary Jacobian satisfies

\[
D_xR^{-p}=-2pR^{-p-2},\qquad
D_x^rR^{-1}=(-1)^r2^r(2r-1)!!R^{-2r-1}
\quad(r\ge0)
\tag{34.0a}
\]

where \((-1)!!=1\). Each differentiation of the inverse Jacobian
advances its denominator exponent by two. At the Widder derivative order
\(r=2k-1\), this gives \(2r+1=4k-1\): the successive clearing powers are
\(R^3,R^7,R^{11},R^{15},R^{19}\). The coordinate value \(R=2n+1\) and
the derivative index \(r\) have separate roles, joined by this explicit
differential law. For every \(k\ge1\) there are polynomials
\(C_{k,j}(x,R)\) such that

\[
\boxed{
R^{4k-1}W_k
=\sum_{j=0}^{2k-1}C_{k,j}(x,R)L^{(j)}(s)
}
\tag{34.1}
\]

The top coefficient is

\[
\boxed{
C_{k,2k-1}(x,R)
=(-1)^{k-1}x^kR^{2k-1}
}
\tag{34.2}
\]

These source-jet coefficients come from the factorially weighted
signed-grid rows in (SW.6), followed by \(S=L/R\) and \(D_x=R^{-1}D_s\).
The product rule introduces the additional Jacobian factors in (34.0a).
In particular the top coefficient reads the right-hand grid edge:
\(C_{k,2k-1}=-c_{k,k}x^kR^{2k-1}\). The lower polynomial coefficients
include derivatives of the Jacobian and are not unweighted grid entries.

Indeed, before differentiation the coefficient of \(L\) in \(x^kL/R\)
has one power of \(R\) in the denominator. Each application of
\(D_x=R^{-1}D_s\) can increase the denominator exponent by at most two,
so \(R^{4k-1}\) clears the \(2k-1\) differentiations. The term
\(L^{(2k-1)}\) arises only when every \(D_s\) differentiates \(L\);
before clearing it has coefficient \((-1)^{k-1}x^k/R^{2k}\), proving
(34.2).

Regularize the zeta pole by

\[
F(s)=(s-1)\zeta(s)>0
\qquad
g(s)=\frac{F'(s)}{F(s)}
\]

Then

\[
L(s)=\frac1s-\frac12\log\pi
+\frac12\psi(s/2)+g(s)
\]

which is regular at \(s=1\). Define the polynomialized owner

\[
\mathfrak N_k=F^{2k}R^{4k-1}W_k
\]

Because \(g^{(2k-1)}\) contains \(F^{(2k)}/F\) with coefficient \(1\),
the coefficient of the highest regularized source jet in
\(\mathfrak N_k\) is

\[
\boxed{
(-1)^{k-1}x^kR^{2k-1}F^{2k-1}
}
\tag{34.3}
\]

It never vanishes for \(x>0\). Thus

\[
\boxed{W_k\rightsquigarrow F^{(2k)}}
\tag{34.4}
\]

is a genuine source-control ladder, not merely bookkeeping notation.

\begin{longtable}[]{@{}
  >{\raggedright\arraybackslash}p{(\columnwidth - 10\tabcolsep) * \real{0.1429}}
  >{\raggedright\arraybackslash}p{(\columnwidth - 10\tabcolsep) * \real{0.1429}}
  >{\raggedleft\arraybackslash}p{(\columnwidth - 10\tabcolsep) * \real{0.1905}}
  >{\raggedleft\arraybackslash}p{(\columnwidth - 10\tabcolsep) * \real{0.1905}}
  >{\raggedleft\arraybackslash}p{(\columnwidth - 10\tabcolsep) * \real{0.1905}}
  >{\raggedright\arraybackslash}p{(\columnwidth - 10\tabcolsep) * \real{0.1429}}@{}}
\toprule\noalign{}
\begin{minipage}[b]{\linewidth}\raggedright
Rung
\end{minipage} & \begin{minipage}[b]{\linewidth}\raggedright
Cleared owner
\end{minipage} & \begin{minipage}[b]{\linewidth}\raggedleft
Highest \(L\) derivative
\end{minipage} & \begin{minipage}[b]{\linewidth}\raggedleft
Highest prime moment
\end{minipage} & \begin{minipage}[b]{\linewidth}\raggedleft
Highest \(F\) jet
\end{minipage} & \begin{minipage}[b]{\linewidth}\raggedright
Status
\end{minipage} \\
\midrule\noalign{}
\endhead
\bottomrule\noalign{}
\endlastfoot
\(W_1\) & \(R^3W_1\) & \(L'\) & \(P_1\) & \(F^{(2)}\) & Analytic,
global \\
\(W_2\) & \(R^7W_2\) & \(L'''\) & \(P_3\) & \(F^{(4)}\) & A-CAT,
global \\
\(W_3\) & \(R^{11}W_3\) & \(L^{(5)}\) & \(P_5\) & \(F^{(6)}\) & A-CAT,
global \\
\(W_4\) & \(R^{15}W_4\) & \(L^{(7)}\) & \(P_7\) & \(F^{(8)}\) & A-CAT,
global \\
\(W_5\) & \(R^{19}W_5\) & \(L^{(9)}\) & \(P_9\) & \(F^{(10)}\) & A-CAT,
global \\
\(W_6\) & \(R^{23}W_6\) & \(L^{(11)}\) & \(P_{11}\) & \(F^{(12)}\) &
A-CAT, global \\
\end{longtable}

\begin{figure}
\centering
\includegraphics{figures/v26c_source_depth.png}
\caption{The same source-depth dictionary as Reading Volume §R12.
Independent Gamma--prime source certificates now reach \(W_6\) and
\(F^{(12)}\); the next source rung is \(W_7\rightsquigarrow F^{(14)}\).}
\end{figure}

\section{35. The first rung - an analytic theta
proof}\label{the-first-rung---an-analytic-theta-proof}

Direct differentiation gives

\[
\boxed{R^3W_1=(1+2x)L+xRL'}
\tag{35.1}
\]

There is a more revealing proof of its sign. Put

\[
y=s-\frac12=\frac R2>\frac12
\qquad
f(y)=\xi\!\left(\frac12+y\right)
\]

The classical theta representation can be normalized as

\[
f(y)=\int_0^\infty\Phi(u)\cosh(yu)\,du
\qquad
\Phi(u)>0
\]

Let \(\ell=f'/f\). Then \(\ell(y)>0\). After writing the cosh integral
as a bilateral exponential integral against a positive symmetric
measure, \((\log f)''=\ell'\) is the variance of \(u\) under the
exponentially tilted probability measure. Hence \(\ell'(y)\ge0\).

Since \(x=y^2-1/4\) and \(S=\ell/(2y)\),

\[
xS(x)=\left(\frac y2-\frac1{8y}\right)\ell(y)=a(y)\ell(y)
\]

where \(a(y)>0\) and \(a'(y)>0\) for \(y>1/2\). Therefore

\[
W_1(x)=\frac1{2y}\left[a'(y)\ell(y)+a(y)\ell'(y)\right]>0
\]

Thus

\[
\boxed{W_1(x)>0\quad(x>0)}
\tag{35.2}
\]

This proof is unconditional, source-theta-side, and uses no
nontrivial-zero location data.

\section{36. Theorem H - Global positivity of the second
rung}\label{theorem-h---global-positivity-of-the-second-rung}

Write

\[
b=1+2x
\qquad
c=2x^2+2x+1
\]

Three applications of \(D_x=R^{-1}D_s\) give the exact owner

\[
\boxed{
R^7W_2
=12cL-6RcL'-6xR^2bL''-x^2R^3L'''
}
\tag{36.1}
\]

The completion part \(L_{\rm pole}=R/x\) gives \(S_{\rm pole}=1/x\), so
\(-D_x^3[x^2S_{\rm pole}]=0\). For

\[
P_j(s)=\sum_{n\ge2}\Lambda(n)(\log n)^j n^{-s}
\]

the arithmetic block is

\[
T_2=-12cP_0-6RcP_1+6xR^2bP_2-x^2R^3P_3
\tag{36.2}
\]

The negative coefficient of \(P_3\) is retained; no termwise positivity
is asserted. With \(g=F'/F\), singular cancellation gives the regular
owner

\[
T_2=6(b+2)+12cg-6Rcg'-6xR^2bg''-x^2R^3g'''
\tag{36.3}
\]

Multiplication by \(F^4\) produces the polynomialized owner
\(\mathfrak N_2=F^4R^7W_2\) and retains the load-bearing term
\(-x^2R^3F^{(4)}F^3\).

\begin{quote}
\textbf{Computer-assisted theorem H.} For every \(x>0\), \[
\boxed{W_2(x)>0}
\]
\end{quote}

\subsection{Finite continuum
certificate}\label{finite-continuum-certificate}

The primary MPFR implementation certifies \(1\le s\le10\) with 44
complete continuum cells, not a point grid. It uses a common-\(\delta\)
Taylor model for the correlated \(F\)-jet, 256-bit directed rounding,
Euler--Maclaurin cutoff \(N=30\), order \(M=12\), and Taylor degree
\(6\). The weakest standard cell is

\[
[1,1.05]
\qquad
\mathfrak N_2>1.7722913290995005\times10^{-4}
\]

The adversarial replay changes simultaneously to 512 bits, \(N=60\),
\(M=16\), and degree \(8\), and gives

\[
\mathfrak N_2([1,1.05])
>1.7722971519681892\times10^{-4}
\]

A materially separate directed-Decimal implementation covers
\(1\le s\le40\) with 40 differently partitioned cells. Its weakest cell
is

\[
[1,1.0625]
\qquad
\mathfrak N_2>
1.450410501387284617086325089\times10^{-4}
\]

It uses 70 decimal digits in the standard run and 120 digits in its
adversarial run. The independent replay reproduced every standard cell,
adversarial cell, and exact analytic-tail field in the published Decimal
receipt.

\subsection{Analytic outer theorem}\label{analytic-outer-theorem}

For \(s\ge10\), one-sided positive-real digamma and polygamma
asymptotics give

\[
A_{\Gamma,2}(s)\ge A_L(s)
\qquad
A_L(10)>72975.29282199671
\]

After replacing the positive logarithmic coefficient by
\(\log(5/\pi)>23/50\), differentiation yields

\[
A_L'(s)\ge\frac{(2s-1)Q_2(s)}{525s^{10}}
\]

where

\[
\begin{aligned}
Q_2(s)={}&13146s^{12}-19446s^{11}+9723s^{10}-2100s^9+525s^8\\
&+13440s^7+17360s^6-325280s^5+2515160s^4\\
&-6609920s^3+7797440s^2-4596480s+1088640
\end{aligned}
\]

Every coefficient of \(Q_2(10+t)\) is positive, so \(A_L\) is
increasing. On the prime side, \(\Lambda(n)\le\log n\), integral
comparison, and \(\log2<7/10\) give

\[
|T_2(s)|
<\frac{15826993}{2460375}s^7 2^{-s}<62986
\]

Thus \(R^7W_2>72975-62986>9989\) for \(s\ge10\), overlapping the finite
certificate.

\subsection{Artifact classification}\label{artifact-classification}

The unmodified MPFR source hash is

\begin{Shaded}
\begin{Highlighting}[]
\NormalTok{923ead8447dd2ba48841b3e8a8325397a62f8261dd4b8dbb60c6018f67f6ba40}
\end{Highlighting}
\end{Shaded}

and the independent Decimal source hash is

\begin{Shaded}
\begin{Highlighting}[]
\NormalTok{6beed41c4644be8ab539fba3995bc6e85f0e174073d08e57c6df9e897fdec7c3}
\end{Highlighting}
\end{Shaded}

The theorem is unconditional in its mathematical assumptions,
computer-assisted on a finite interval, Gamma/prime analytic outside it,
independent of zero data, and classified \textbf{A-CAT}. Section 74A
proves the same sign again on the uniform partition and tail it shares
with \(W_3,\ldots,W_6\), in the certificate bundled with this edition.

\begin{figure}
\centering
\includegraphics{figures/v26c_certificate_cover.png}
\caption{A global source certificate covers the entire \(s>1\) branch:
directed finite intervals meet proved analytic tails. The handoffs are
\(s=10\) for \(W_2\) and \(s=20\) for \(W_3\); the drawing compresses
the unbounded axis.}
\end{figure}

\section{37. Theorem I - Global positivity of the third
rung}\label{theorem-i---global-positivity-of-the-third-rung}

Five differentiations give

\[
\boxed{R^{11}W_3=\sum_{j=0}^{5}C_jL^{(j)}}
\tag{37.1}
\]

where

\[
\begin{aligned}
C_0&=720(x+1)(2x^2+x+1)\\
C_1&=-360R(x+1)(2x^2+x+1)\\
C_2&=-60R^2(4x^3+3x^2-1)\\
C_3&=60xR^3(3x^2+3x+1)\\
C_4&=15x^2R^4(2x+1)\\
C_5&=x^3R^5
\end{aligned}
\]

Again the completion contribution is annihilated exactly. The source
decomposition is

\[
R^{11}W_3=A_{\Gamma,3}+T_3
\]

with

\[
\begin{aligned}
A_{\Gamma,3}={}&
\frac{C_0}{2}[\psi(s/2)-\log\pi]
+\frac{C_1}{4}\psi_1(s/2)
+\frac{C_2}{8}\psi_2(s/2)\\
&+\frac{C_3}{16}\psi_3(s/2)
+\frac{C_4}{32}\psi_4(s/2)
+\frac{C_5}{64}\psi_5(s/2)
\end{aligned}
\]

and

\[
T_3=-C_0P_0+C_1P_1-C_2P_2+C_3P_3-C_4P_4+C_5P_5
\tag{37.2}
\]

The mixed signs are retained. Singular regularization gives

\[
\boxed{
T_3=240(3x^2+5x+5)+\sum_{j=0}^{5}C_jg^{(j)}
}
\tag{37.3}
\]

If \(G_0=F'\) and

\[
G_{j+1}=FG_j'-(j+1)F'G_j
\]

then \(g^{(j)}=G_j/F^{j+1}\) and

\[
\boxed{
\mathfrak N_3:=F^6R^{11}W_3
=F^6A_{\Gamma,3}
+240(3x^2+5x+5)F^6
+\sum_{j=0}^{5}C_jG_jF^{5-j}
}
\tag{37.4}
\]

The coefficient of \(F^{(6)}\) is \(C_5F^5=x^3R^5F^5>0\), so the sixth
source jet is genuinely load-bearing.

\begin{quote}
\textbf{Computer-assisted theorem I.} For every \(x>0\), \[
\boxed{W_3(x)>0}
\]
\end{quote}

\subsection{Finite continuum
certificate}\label{finite-continuum-certificate-1}

The directed MPFR artifact certifies \(1\le s\le20\) in 135 continuum
cells at 256 bits. It uses a common-variable degree-6 Taylor model, with
center Euler--Maclaurin parameters \((N,M)=(24,9)\) and
uniform-remainder parameters \((12,7)\). The weakest cell is

\[
[1,1.01]
\qquad
\mathfrak N_3>1.6950218693828748\times10^{-5}
\]

The adversarial replay changes to 512 bits, Taylor degree \(8\), center
parameters \((40,12)\), and remainder parameters \((24,10)\). It gives

\[
1.6951786515869995\times10^{-5}
<\mathfrak N_3([1,1.01])
<2.2185472840400943\times10^{-5}
\]

\subsection{Analytic outer theorem}\label{analytic-outer-theorem-1}

For \(s\ge20\), elementary polygamma series bounds and
\(\log(s/(2\pi))>23/20\) yield

\[
A_{\Gamma,3}>\frac{3Q_3(s)}{s^3}
\]

where

\[
\begin{aligned}
Q_3(s)={}&324s^9-5652s^8+23266s^7-47632s^6+58430s^5\\
&-46350s^4+24055s^3-8052s^2+1519s-120
\end{aligned}
\]

Every coefficient of \((Q_3(s)-93s^9)|_{s=20+t}\) is positive;
consequently

\[
A_{\Gamma,3}>279s^6
\]

On the arithmetic side, \(\Lambda(n)\le\log n\) and exact log-moment
integration give

\[
|T_3|<14s^{11}2^{-s}
\]

Since

\[
14\frac{20^5}{2^{20}}=\frac{21875}{512}<279
\]

and \(s^5/2^s\) decreases for \(s\ge20\),

\[
|T_3|<279s^6<A_{\Gamma,3}
\]

This closes the analytic tail with no zero-location input.

\begin{figure}
\centering
\includegraphics{figures/widder_certificate_margins.png}
\caption{The finite \(W_2\) and \(W_3\) continuum certificates have
positive lower bounds on every band; the two panels show different
polynomialized owners and should not be compared in absolute scale.}
\end{figure}

\begin{figure}
\centering
\includegraphics{figures/widder_tail_domination.png}
\caption{The explicit analytic tails: an Archimedean lower bound
dominates the complete absolute von-Mangoldt block beyond the stated
handoff.}
\end{figure}

\subsection{Reproducible source
certificate}\label{reproducible-source-certificate}

The certificate uses exact rational Bernoulli data and outward-rounded
interval arithmetic on 135 continuum cells, and its adversarial replay
reproduced every cell and every analytic-tail constant. Together with
the tail theorem this proves the global result with grade A-CAT. Section
74A proves the same sign again on the uniform partition and tail it
shares with the other certified rungs, in the certificate bundled with
this edition.

\section{38. The fourth owner and the non-raising
recurrence}\label{the-fourth-owner-and-the-non-raising-recurrence}

The next reduced rung is

\[
W_4(x)=-D_x^7[x^4S_\xi(x)]
\]

Exact differentiation gives

\[
\boxed{R^{15}W_4=\sum_{j=0}^{7}D_jL^{(j)}}
\tag{38.1}
\]

where, with

\[
A(x)=10x^4+20x^3+18x^2+10x+5
\]

the coefficients are

\[
\begin{aligned}
D_0&=20160A(x)\\
D_1&=-10080RA(x)\\
D_2&=-10080R^2(x+1)(2x^3-1)\\
D_3&=840R^3(22x^4+32x^3+18x^2+4x-1)\\
D_4&=-840xR^4(x+1)(2x^2+2x+1)\\
D_5&=-84x^2R^5(10x^2+10x+3)\\
D_6&=-28x^3R^6(2x+1)\\
D_7&=-x^4R^7
\end{aligned}
\]

The array \(D_0,\ldots,D_7\) is row \(k=4\) of the signed-grid
construction after the factorial and Jacobian operations in (SW.6) and
§34. Its top entry \(D_7=-x^4R^7\) uses \(c_{4,4}=+1\). An exact
rational differentiator independently reproduces every one of these
coefficients and verifies the general top coefficient through \(k=8\).
The highest prime moment is \(P_7\) and the highest regularized source
jet is \(F^{(8)}\).

There is an exact recurrence between successive rungs:

\[
\boxed{
W_{k+1}=-j_kW_k'-xW_k''
=-x^{-2k}D_x[x^{2k+1}W_k'],\qquad j_k=2k+1
}
\tag{38.2}
\]

To prove it, write \(x^{k+1}S=x(x^kS)\) and use
\(D_x^m(xf)=xD_x^mf+mD_x^{m-1}f\) with \(m=2k+1\). The recurrence is not
a positivity-raising operator: knowing only \(W_k\ge0\) gives no control
over the two derivatives on the right. An independent source-side
continuum certificate and analytic tail close the sign, as proved in
Theorem 74.1:

\[
\boxed{W_4(x)>0\qquad(x>0)}
\]

\section{38A. Positive descent, not positive
raising}\label{a.-positive-descent-not-positive-raising}

For the actual completed rungs the non-raising recurrence (38.2) has a
positive inverse in the downward direction:

\[
\boxed{W_k(x)=\frac1{2k}\left[
 x^{-2k}\int_0^x t^{2k}W_{k+1}(t)dt
 +\int_x^\infty W_{k+1}(t)dt\right],\qquad k\ge1.}
\tag{38A.1}
\]

\textbf{Proof.} The recurrence is \((x^{2k+1}W_k')'=-x^{2k}W_{k+1}\).
Analyticity at zero removes the first integration constant. Each fixed
rung tends to zero at infinity by its folded expansion and dominated
convergence. Integrating twice, and interchanging the order, gives
(38A.1). This interchange is absolute: for \(\Re q>0\),
\(|t+q|^2\ge(t+|q|)^2/2\), so the integral of the absolute value of one
\((k+1)\)-st rung is at most \(C_k|q|^{-k}\), summable over the folded
zeros. \(\square\)

Both kernels are nonnegative. Thus \(W_{k+1}\ge0\) on the full positive
axis implies \(W_k\ge0\) there. Positivity on that axis at every index
in any unbounded set implies the whole Widder criterion; odd indices are
one example. This does not raise a finite prefix, and values at \(x=0\)
alone do not satisfy its hypotheses. The separate odd Li criterion is
Theorem 155A.2.

\section{74. Global fourth-rung
theorem}\label{global-fourth-rung-theorem}

This section proves the fourth scalar rung directly from the source. It
is a source-side theorem: the finite certificate evaluates the
regularized zeta/Gamma source on the real half-line \(s>1\), and the
analytic tail uses only polygamma inequalities and the absolutely
convergent von-Mangoldt series. It does not use a zeta zero table, a
verified RH height, or RH.

Put

\[
x=s(s-1),\qquad R=2s-1,\qquad s>1
\]

and

\[
L(s)=\frac{\xi'(s)}{\xi(s)}
\]

Section 38 proved the exact owner

\[
R^{15}W_4=\sum_{j=0}^{7}D_jL^{(j)}
\tag{74.1}
\]

with the displayed polynomial coefficients \(D_0,\ldots,D_7\).
Regularize the pole of \(\zeta\) at \(1\) by

\[
F(s)=(s-1)\zeta(s)>0
\]

and define

\[
\boxed{\mathfrak N_4(s)=F(s)^8R^{15}W_4(x(s))}
\tag{74.2}
\]

Since \(F\) and \(R\) are positive on \(s>1\), \(\mathfrak N_4\) and
\(W_4\) have the same sign.

\begin{quote}
\textbf{Theorem 74.1 (global fourth source rung).} For every \(x>0\), \[
\boxed{W_4(x)>0}
\tag{74.3}
\] The accepted classification is \textbf{A-CAT / SOURCE-SIDE}, as
recorded in §R12.

This is the source-side A-CAT fourth-rung theorem, established by a
finite continuum certificate and a separate analytic tail. The section
keeps the original proof, with its own handoff at \(s=40\) and its own
constants; §74A proves the same sign again inside the bundled
certificate that covers \(W_2,\ldots,W_6\) on one partition. The
independent verified-height theorem also includes this scalar sign. No
finite positive prefix establishes the all-order RH criterion.
\end{quote}

\subsection{\texorpdfstring{Finite continuum
\(1\le s\le40\)}{Finite continuum 1\textbackslash le s\textbackslash le40}}\label{finite-continuum-1le-sle40}

The regular source identity is

\[
L(s)=\frac1s-\frac12\log\pi+\frac12\psi\!\left(\frac s2\right)
+\frac{F'(s)}{F(s)}
\tag{74.4}
\]

Substituting (74.4) in (74.1), multiplying by \(F^8\), and performing
every operation in one common Taylor variable preserves the correlations
among \(F,F',\ldots\). This point matters: independent interval
substitution for the jets is far too wide near \(s=1\).

The standard replay uses GNU MPFR directed rounding at 256 bits. It
covers \([1,40]\) by the following exact schedule:

\begin{longtable}[]{@{}rrr@{}}
\toprule\noalign{}
band & cell width & cells \\
\midrule\noalign{}
\endhead
\bottomrule\noalign{}
\endlastfoot
\([1,1.1]\) & \(0.01\) & 10 \\
\([1.1,1.3]\) & \(0.02\) & 10 \\
\([1.3,1.5]\) & \(0.025\) & 8 \\
\([1.5,2]\) & \(0.02\) & 25 \\
\([2,3]\) & \(0.05\) & 20 \\
\([3,5]\) & \(0.1\) & 20 \\
\([5,10]\) & \(0.25\) & 20 \\
\([10,20]\) & \(0.5\) & 20 \\
\([20,40]\) & \(1\) & 20 \\
\end{longtable}

Thus 153 complete continuum cells are certified. Taylor degree eight,
center Euler--Maclaurin parameters \((N,M)=(24,9)\), and remainder
parameters \((12,7)\) require correlated center jets through
\(F^{(16)}\) and a remainder through \(F^{(17)}\). Bernoulli numbers are
exact rationals; every arithmetic operation, including the
Euler--Maclaurin remainder, is rounded outward.

The weakest standard cell is the endpoint cell:

\[
s\in[1,1.01],\qquad
\mathfrak N_4(s)>
3.2148256503423557\times10^{-6}
\tag{74.5}
\]

An independent adversarial replay changes all material numerical
parameters at once: 512 bits, Taylor degree ten, center
\((N,M)=(40,12)\), and remainder \((24,10)\). It gives

\[
\mathfrak N_4([1,1.01])>
3.2149708080858507\times10^{-6}
\tag{74.6}
\]

The two lower endpoints need not coincide; each is a rigorous enclosure
from a different Taylor/Euler--Maclaurin model. Their agreement at the
displayed scale comes from an adversarial replay, not from an
assumption.

\subsection{\texorpdfstring{Analytic tail
\(s\ge40\)}{Analytic tail s\textbackslash ge40}}\label{analytic-tail-sge40}

Completion annihilation separates

\[
R^{15}W_4=A_{\Gamma,4}+T_4
\tag{74.7}
\]

where

\[
A_{\Gamma,4}
=\frac{D_0}{2}\left[\psi(s/2)-\log\pi\right]
+\sum_{j=1}^{7}\frac{D_j}{2^{j+1}}\psi^{(j)}(s/2)
\tag{74.8}
\]

and \(T_4\) is the von-Mangoldt source. Integral comparison in the
classical polygamma series gives, for \(m\ge1\) and \(z>0\),

\[
(m-1)!z^{-m}
\le (-1)^{m+1}\psi^{(m)}(z)
\le (m-1)!z^{-m}+m!z^{-m-1}
\tag{74.9}
\]

together with \(\psi(z)\ge\log z-z^{-1}\). Substitution with the signs
of the \(D_j\) yields a rational lower function \(B_4(s)\). After
clearing \(s^8\),

\[
s^8\bigl(B_4(s)-25000s^8\bigr)=P_4(s)
\tag{74.10}
\]

and every coefficient of \(P_4(40+t)\) is a strictly positive integer.
Consequently

\[
\boxed{A_{\Gamma,4}>25000s^8\qquad(s\ge40)}
\tag{74.11}
\]

For the prime source, \(x<s^2\), \(R<2s\), and \(\Lambda(n)\le\log n\)
give

\[
|T_4|<K s^{15}2^{-s}
\qquad
K=\frac{14727251778604209636316819}
{593716227296025600000000}<25
\tag{74.12}
\]

Since

\[
\frac{40^7}{2^{40}}=\frac{78125}{524288}<1
\]

and \(s^7/2^s\) decreases for \(s\ge40\),

\[
\boxed{|T_4|<25s^8\qquad(s\ge40)}
\tag{74.13}
\]

Equations (74.11)--(74.13) imply

\[
R^{15}W_4>24975s^8>0\qquad(s\ge40)
\tag{74.14}
\]

The finite continuum and the analytic tail overlap at \(s=40\), proving
Theorem 74.1. This is the original fourth-rung argument, in its own
coordinates: a directed finite certificate on \(1\le s\le40\) and the
analytic tail (74.11)--(74.14) beyond it. Section 74A proves \(W_4>0\)
again, on the uniform partition and tail shared by \(W_2,\ldots,W_6\),
with a bundled certificate that the v5.0 build re-executed; the two
arguments are independent of each other in partition, threshold and
handoff point. Independently of both, the verified-height theorem in
Reading Volume §R13 includes \(W_4>0\) in its finite zero-assisted
prefix.

\begin{figure}
\centering
\includegraphics{figures/w4_global_certificate_v1_2.png}
\caption{Every cell lower bound is positive; the separate analytic tail
leaves a factor-\(999\) source margin.}
\end{figure}

\subsection{What the fourth rung does and does not
establish}\label{what-the-fourth-rung-does-and-does-not-establish}

The fourth-rung theorem itself established the inherited source prefix

\[
\boxed{W_1,W_2,W_3,W_4>0\quad\text{globally}}
\tag{74.15}
\]

Section 74A extends the independent source prefix to
\(W_1,\ldots,W_6>0\), proving the four certified rungs above the first
on one partition and one tail. The non-raising recurrence (38.2) is
unchanged, and the next independent owner is
\(W_7\rightsquigarrow F^{(14)}\). A finite source prefix is not a
zero-location theorem and does not imply the full Widder hierarchy.

\section{\texorpdfstring{74A. One certificate for the source rungs
\(W_2,\ldots,W_6\)}{74A. One certificate for the source rungs W\_2,\textbackslash ldots,W\_6}}\label{a.-one-certificate-for-the-source-rungs-w_2ldotsw_6}

The five certified source rungs are proved together, by one certificate:
one partition of a finite \(s\)-interval, one exact analytic tail beyond
it, and two programs that share no arithmetic. No zero table, verified
height or RH assumption enters anywhere.

\begin{quote}
\textbf{Theorem 74A.1 (global source rungs).} For every \(x>0\), \[
\boxed{W_k(x)>0\qquad(k=2,3,4,5,6).}
\tag{74A.1}
\] The classification is \textbf{A-CAT / SOURCE-SIDE}. The owners reach
\(F^{(4)}\) at the second rung and \(F^{(12)}\) at the sixth, for
\(F(s)=(s-1)\zeta(s)\).
\end{quote}

Use \(x=s(s-1)\), \(R=2s-1\), \(D_x=R^{-1}D_s\) and \(L=\xi'/\xi\), so
that \(W_k(x)=(-1)^{k-1}D_x^{2k-1}[x^kL/R]\). On \(1\le s\le128\) the
regularization \[
F(s)=1+(s-1)\left(\zeta(s)-\frac1{s-1}\right),\qquad
\log\xi(s)=\log s+\log F(s)+\log\Gamma(s/2)
-\frac s2\log\pi-\log2
\tag{74A.2}
\] removes the pole at one. The certified quantity is \[
\mathcal W_k(s)=\frac{s^{2k-1}}{(2k-1)!}W_k(s(s-1)),
\tag{74A.3}
\] which is positive exactly where \(W_k\) is, because
\(s^{2k-1}/(2k-1)!>0\) and \(s\mapsto s(s-1)\) carries \([1,\infty)\)
onto \([0,\infty)\).

\subsection{\texorpdfstring{The finite continuum,
\(1\le s\le128\)}{The finite continuum, 1\textbackslash le s\textbackslash le128}}\label{the-finite-continuum-1le-sle128}

A complete rational partition into continuum cells proves \[
\boxed{\mathcal W_k(s)>\theta_k\quad(1\le s\le128)},\qquad
\begin{array}{c|ccccc}
k&2&3&4&5&6\\\hline
\theta_k&10^{-5}&10^{-8}&10^{-10}&10^{-12}&10^{-15}
\end{array}
\tag{74A.4}
\] Each cell is accepted from a bound that holds on the whole interval,
never from a point grid. With \(m\) the midpoint of a cell and \(r\) its
radius, \(p_j\) an enclosure of \(\mathcal W_k^{(j)}(m)/j!\), and
\(Q_D\) an enclosure of \(\mathcal W_k^{(D)}(\sigma)/D!\) for every
\(\sigma\) in the cell, Taylor's theorem with Lagrange remainder gives
\[
\mathcal W_k(s)\ \ge\
p_0-\sum_{j=1}^{D-1}|p_j|r^j-\sup|Q_D|\,r^D
\qquad(|s-m|\le r).
\tag{74A.5}
\] Every operation is outward rounded, the endpoints are exact dyadic
rationals, adjacency and both ends of \([1,128]\) are verified exactly,
and a cell whose bound does not clear \(\theta_k\) is bisected and
retried.

The thresholds in (74A.4) are deliberately far below the certified
values. The smallest enclosures occur at the left endpoint, where
(12.4a) identifies \(\mathcal W_k(1)=\mu_{k-1}\), the folded boundary
moment:

\begin{longtable}[]{@{}
  >{\raggedright\arraybackslash}p{(\columnwidth - 10\tabcolsep) * \real{0.1667}}
  >{\raggedright\arraybackslash}p{(\columnwidth - 10\tabcolsep) * \real{0.1667}}
  >{\raggedright\arraybackslash}p{(\columnwidth - 10\tabcolsep) * \real{0.1667}}
  >{\raggedright\arraybackslash}p{(\columnwidth - 10\tabcolsep) * \real{0.1667}}
  >{\raggedright\arraybackslash}p{(\columnwidth - 10\tabcolsep) * \real{0.1667}}
  >{\raggedright\arraybackslash}p{(\columnwidth - 10\tabcolsep) * \real{0.1667}}@{}}
\toprule\noalign{}
\begin{minipage}[b]{\linewidth}\raggedright
\(k\)
\end{minipage} & \begin{minipage}[b]{\linewidth}\raggedright
2
\end{minipage} & \begin{minipage}[b]{\linewidth}\raggedright
3
\end{minipage} & \begin{minipage}[b]{\linewidth}\raggedright
4
\end{minipage} & \begin{minipage}[b]{\linewidth}\raggedright
5
\end{minipage} & \begin{minipage}[b]{\linewidth}\raggedright
6
\end{minipage} \\
\midrule\noalign{}
\endhead
\bottomrule\noalign{}
\endlastfoot
\(\mu_{k-1}=\mathcal W_k(1)\) & \(3.71006\times10^{-5}\) &
\(1.43678\times10^{-7}\) & \(6.59828\times10^{-10}\) &
\(3.19389\times10^{-12}\) & \(1.57589\times10^{-14}\) \\
\(\theta_k\) & \(10^{-5}\) & \(10^{-8}\) & \(10^{-10}\) & \(10^{-12}\) &
\(10^{-15}\) \\
\end{longtable}

\subsection{\texorpdfstring{The exact tail,
\(s\ge128\)}{The exact tail, s\textbackslash ge128}}\label{the-exact-tail-sge128}

Generate the source polynomials recursively by \[
P_j^{(r+1)}=R(P_j^{(r)})'-2(2r+1)P_j^{(r)}+RP_{j-1}^{(r)},
\qquad P_0^{(0)}=x^k,
\tag{74A.6}
\] which is the product rule for \(R^{-1}D_s\) applied to
\(P(s)L^{(j)}(s)R^{-m}\) at \(m=1,3,\ldots,4k-3\). With
\(C_{k,j}=(-1)^{k-1}P_j^{(2k-1)}\) this gives the exact identity \[
R^{4k-1}W_k=\sum_{j=0}^{2k-1}C_{k,j}(s)L^{(j)}(s),
\qquad
C_{k,2k-1}=(-1)^{k-1}x^kR^{2k-1}.
\tag{74A.7}
\] Completion annihilates the elementary \(1/s+1/(s-1)\) terms, which is
checked as a polynomial identity, and the recurrence is checked against
the closed-form \(x\)-Leibniz values of \(D_x^{2k-1}[x^{k+m}]\).
Separating the Gamma and von-Mangoldt pieces and using
\(0<(1-e^{-t})^{-1}-t^{-1}-1/2<t/12\) gives an exact rational lower
function \(B_k\) for the Gamma part; with
\(\alpha_k=\tfrac14\operatorname{lead}(B_k)\) and
\(\Lambda(n)\le\log n\) with integral comparison for the complete prime
part, \[
A_{\Gamma,k}(s)>\alpha_ks^{2k},
\qquad
|T_k(s)|<\tfrac12\alpha_ks^{2k}
\qquad(s\ge128),
\tag{74A.8}
\] both by exact rational checks: every coefficient of
\(s^{2k+1}\bigl(B_k(s)-\alpha_ks^{2k}\bigr)\), expanded at \(s=128+t\),
is nonnegative with a positive constant term, and the endpoint ratio of
the prime bound to the Gamma margin is a rational number below \(1/2\),
after which it decreases because \(s^{2k-1}2^{-s}\) does. Hence \[
\boxed{W_k(s(s-1))>\frac{\alpha_k}{2}\cdot\frac{s^{2k}}{(2s-1)^{4k-1}}
\qquad(s\ge128)}
\tag{74A.9}
\] with the exact constants

\begin{longtable}[]{@{}
  >{\raggedright\arraybackslash}p{(\columnwidth - 10\tabcolsep) * \real{0.1667}}
  >{\raggedright\arraybackslash}p{(\columnwidth - 10\tabcolsep) * \real{0.1667}}
  >{\raggedright\arraybackslash}p{(\columnwidth - 10\tabcolsep) * \real{0.1667}}
  >{\raggedright\arraybackslash}p{(\columnwidth - 10\tabcolsep) * \real{0.1667}}
  >{\raggedright\arraybackslash}p{(\columnwidth - 10\tabcolsep) * \real{0.1667}}
  >{\raggedright\arraybackslash}p{(\columnwidth - 10\tabcolsep) * \real{0.1667}}@{}}
\toprule\noalign{}
\begin{minipage}[b]{\linewidth}\raggedright
\(k\)
\end{minipage} & \begin{minipage}[b]{\linewidth}\raggedright
2
\end{minipage} & \begin{minipage}[b]{\linewidth}\raggedright
3
\end{minipage} & \begin{minipage}[b]{\linewidth}\raggedright
4
\end{minipage} & \begin{minipage}[b]{\linewidth}\raggedright
5
\end{minipage} & \begin{minipage}[b]{\linewidth}\raggedright
6
\end{minipage} \\
\midrule\noalign{}
\endhead
\bottomrule\noalign{}
\endlastfoot
\(\alpha_k\) & \(7\) & \(396\) & \(54{,}000\) & \(13{,}406{,}400\) &
\(5{,}258{,}131{,}200\) \\
prime/Gamma ratio at \(s=128\) & \(1.9\times10^{-33}\) &
\(1.2\times10^{-30}\) & \(3.3\times10^{-28}\) & \(5.1\times10^{-26}\) &
\(5.2\times10^{-24}\) \\
\end{longtable}

The finite continuum and the exact tail overlap at \(s=128\), so (74A.4)
and (74A.9) prove Theorem 74A.1.

\subsection{Two arithmetics, one
statement}\label{two-arithmetics-one-statement}

The certificate is executed twice, by programs that share no arithmetic,
no special-function implementation and no computer-algebra system:

\begin{itemize}
\tightlist
\item
  \texttt{qa/source\_rungs/source\_rungs\_arb.py} evaluates the jets in
  FLINT/Arb ball arithmetic through python-flint, with SymPy for the
  exact tail. It generalizes the fifth- and sixth-rung certificate of
  the 4.5 edition to all five rungs; the evaluator, the cell test and
  the tail argument are the same, with \(\theta_k\) and \(\alpha_k\) as
  parameters.
\item
  \texttt{qa/source\_rungs/source\_rungs\_decimal.py} uses the Python
  standard library alone. Arithmetic is \texttt{decimal} (libmpdec) with
  directed rounding on every operation; the \(\zeta\) jet is
  Euler--Maclaurin (DLMF 25.2.9) with Lehmer's bound on the periodic
  Bernoulli function and Cauchy estimates for the Taylor coefficients of
  the remainder; the \(\psi\) jet is the recurrence with the Stirling
  series (DLMF 5.11.2) and the classical remainder bound
  \(|R_m(w)|\le|B_{2m}|/(2m(\operatorname{Re}w)^{2m})\); the source
  polynomials and the whole tail are exact \texttt{Fraction} arithmetic.
  It imports neither FLINT/Arb nor MPFR, SymPy or mpmath.
\end{itemize}

Each program also checks the repeated \(s\)-operator against a second
formulation of the same quantity: the \(x\)-Leibniz formula in the
first, the polynomial identity (74A.7) in the second. Both are run a
second time with every material parameter changed at once -- precision,
Taylor order and partition, and in the standard-library program the
Euler--Maclaurin and Stirling truncations as well.
\texttt{qa/source\_rungs/cross\_backend\_check.py} compares the runs:
the exact tail data agree digit for digit, including the source
polynomials, the signs, \(\alpha_k\), the prime bound and the endpoint
ratio, and the sampled enclosures of \(\mathcal W_k\) overlap. The
receipts record every accepted cell, and the build re-executes both
programs.

Thus the source-side prefix is \(W_1,\ldots,W_6>0\): the first rung
analytically (§35), the next five by this certificate. The recurrence
(38.2) remains non-raising. The next independent source sign is
\(W_7=-x^{-12}(x^{13}W_6')'>0\), whose owner reaches \(F^{(14)}\), and
it is open. No finite prefix establishes the all-order Widder criterion.

\subsection{Pole-sublattice calibration at the next
rung}\label{pole-sublattice-calibration-at-the-next-rung}

A Gamma-pole analysis supplies an auxiliary NO-GO, not a completed
\(W_7\) sign. For a centered pole term \(f_m\) let \(d\sigma_{f_m}\) be
the signed measure of \(f_m/x\). Its density changes sign once, at
\(u=m+1/2\), and

\[
\int u^{-k}\,d\sigma_{f_m}(u)<0\qquad(k\ge2).
\tag{74A.10}
\]

The publication label is important: because \(h=xE_\xi'/E_\xi\) already
contains the factor \(x\),

\[
\boxed{\frac{W_k[f_m](0)}{(2k-1)!}
=\int u^{-k}\,d\sigma_{f_m}(u),}
\tag{74A.11}
\]

not \(W_k[f_m/x]\). For \(m=0\),

\[
\frac{W_k[f_0](0)}{(2k-1)!}
=-\binom{2k-2}{k-2},\qquad k\ge2,
\tag{74A.12}
\]

so the complementary centered pole channel contributes at most \(-792\)
to the normalized seventh rung at the boundary. This is an obstruction
to treating the unused Gamma poles as a hidden positive reserve. It is
\textbf{not} a negative statement about completed \(W_7\), whose Gamma
and prime pieces nearly cancel and whose all-\(x\) source sign remains
open.

\section{55. Fourth-rung theorem and fixed-rung
asymptotics}\label{fourth-rung-theorem-and-fixed-rung-asymptotics}

The fourth reduced-Widder owner is

\[
W_4(x)
=-840S_\xi'''(x)-840xS_\xi^{(4)}(x)
-252x^2S_\xi^{(5)}(x)
-28x^3S_\xi^{(6)}(x)-x^4S_\xi^{(7)}(x)
\tag{55.1}
\]

and

\[
W_4=-7W_3'-xW_3''
\tag{55.2}
\]

Its denominator-cleared source owner reaches \(F^{(8)}\) for
\(F(s)=(s-1)\zeta(s)\). Direct source-side \(W_4(x)>0\) for all \(x>0\)
is an A-CAT theorem. Section 74A carries the same source-side
classification for \(W_2,\ldots,W_6\) in one bundled certificate,
reaching \(F^{(12)}\) at the sixth rung. The original fourth-rung proof
remains in §74.

A complementary unconditional outer theorem controls every fixed rung.

\begin{quote}
\textbf{Theorem 55.1 (eventual positivity of every fixed rung).} For
every fixed integer \(k\ge1\), \[
\boxed{
W_k(x)
=\kappa_kx^{1/2-k}\log x
+O_k(x^{1/2-k})}
\tag{55.3}
\] where \[
\boxed{
\kappa_k=
\frac{(2k-1)!!(2k-3)!!}{2^{2k+2}}>0}
\tag{55.4}
\] Hence \(W_k(x)>0\) for all sufficiently large \(x\), with the
threshold allowed to depend on \(k\).
\end{quote}

\textbf{Proof.} The explicit source formula and the differentiated
Stirling--polygamma expansion imply, for each fixed \(m\ge0\),

\[
D_x^mS_\xi(x)
=D_x^m\!\left(\frac{\log x}{8\sqrt x}\right)
+O_m(x^{-m-1/2})
\tag{55.5}
\]

To justify differentiation, note that \(s\asymp\sqrt x\) and
\(D_xs=1/R=O(x^{-1/2})\); every fixed derivative of the von-Mangoldt
series is exponentially small in \(s\), while the differentiated
polygamma remainder is a classical symbol of the required order. Thus
(55.5), not merely an undifferentiated big-\(O\) estimate, is available.

Now

\[
x^kS_\xi(x)
=\frac18x^{k-1/2}\log x+O_k(x^{k-1/2})
\tag{55.6}
\]

with the same differentiated-symbol control through order \(2k-1\). The
coefficient of \(\log x\) after \(2k-1\) derivatives is

\[
\frac18
\left(k-\frac12\right)_{\underline{2k-1}}
x^{1/2-k}
\tag{55.7}
\]

The falling factorial has sign \((-1)^{k-1}\) and magnitude

\[
\left|
\left(k-\frac12\right)_{\underline{2k-1}}
\right|
=\frac{(2k-1)!!(2k-3)!!}{2^{2k-1}}
\tag{55.8}
\]

Multiplication by the Widder sign gives (55.3)--(55.4). \(\square\)

For \(k=4\),

\[
\boxed{
W_4(x)
=\frac{1575}{1024}x^{-7/2}\log x
+O(x^{-7/2})}
\tag{55.9}
\]

This reduces any fixed-rung proof to a compact interval plus an
effective tail. It does not solve RH because the constants and
thresholds are not uniform in \(k\).

\begin{figure}
\centering
\includegraphics{figures/fixed_rung_asymptotic_tails.png}
\caption{The positive fixed-rung asymptotic constants and the decay of
the first four leading tails.}
\end{figure}

\section{75. Exact reserve phase transition and all-rung shifted
handoff}\label{exact-reserve-phase-transition-and-all-rung-shifted-handoff}

This section audits the triangular coefficient proposed for logarithmic
jets of the same regular source \(F(s)=(s-1)\zeta(s)\). For \(m\ge1\)
put

\[
P_{m-1}(s)=\sum_{n\ge2}\frac{\Lambda(n)(\log n)^{m-1}}{n^s}
\]

\[
\rho_m(s)=\frac{(s-1)^mP_{m-1}(s)}{(m-1)!}
\qquad
\mathcal R_m(s)=1-\rho_m(s)
\tag{75.1}
\]

Termwise differentiation of the Euler product on \(s>1\) gives the exact
identity

\[
\boxed{
\mathcal R_m(s)
=(-1)^{m-1}(\log F)^{(m)}(s)
\frac{(s-1)^m}{(m-1)!}}
\tag{75.2}
\]

Fix a rung \(k\) and the minimal termwise-positive handoff

\[
s_k=1+\frac{2k}{\log2}
\tag{75.3}
\]

If \(1\le m\le2k\) and \(s\ge s_k\), every source atom has
\(y=(s-1)\log n\ge m\). Since \(y^me^{-y}\) decreases for \(y\ge m\),
\(\rho_m(s)\) decreases with \(s\). At fixed \(s_k\), the atom ratio
from order \(m\) to order \(m+1\) equals \(y/m\ge1\), so \(\rho_m(s_k)\)
increases with \(m\). Both monotonicity reductions are therefore exact,
and

\[
\boxed{
\inf_{\substack{s\ge s_k\\1\le m\le2k}}\mathcal R_m(s)
=\mathcal R_{2k}(s_k)=1-A_k}
\tag{75.4}
\]

Writing \(Y=2k+\Delta\) and \(\lambda=\log_2(p^a)\) gives the
prime-power owner

\[
A_k(\Delta)
=\frac1{(2k-1)!}
\sum_{p^a}
\frac{(Y\lambda)^{2k}e^{-Y\lambda}}{a p^a}
\tag{75.5}
\]

The proposed triangular margin is

\[
c_k=\frac1{kT_k}=\frac{2}{k^2(k+1)}
\tag{75.6}
\]

It holds on the minimal shell precisely when \(A_k(0)<1-c_k\).

\begin{quote}
\textbf{Theorem 75.1 (exact minimal-shell ladder).} The inequality \[
(-1)^{m-1}(\log F)^{(m)}(s)
>c_k\frac{(m-1)!}{(s-1)^m}
\tag{75.7}
\] for every \(s\ge s_k\) and \(1\le m\le2k\) holds exactly for \[
\boxed{2\le k\le11}
\tag{75.8}
\] It fails at \(k=1\) and at every \(k\ge12\).
\end{quote}

The finite values \(1\le k\le12\) are enclosed by directed MPFR
evaluation of (75.5) through prime powers \(p^a\le20000\), followed by
an explicit incomplete-gamma tail. At \(k=12\), the two atoms \(n=2,3\)
already give

\[
A_{12}(0)>1.00660782373417>1
\tag{75.9}
\]

so the first failure is independent of the rest of the prime source. For
\(k\ge13\), the \(n=2\) atom

\[
B_k=\frac{(2k)^{2k}e^{-2k}}{2(2k-1)!}
\tag{75.10}
\]

is increasing and \(B_{13}>1.01385265607324\). Thus one atom proves
every later failure. Stirling explains the transition:

\[
B_k=\sqrt{\frac{k}{4\pi}}\,e^{-\theta_{2k}}
\qquad 0<\theta_{2k}<\frac1{24k}
\tag{75.11}
\]

so the obstruction scale is \(4\pi\), not a triangular number.

At the fourth rung, the sharp reserve is

\[
0.28687152864220360<C_4^*=1-A_4
<0.28687152864220367
\tag{75.12}
\]

and therefore

\[
\frac27<C_4^*<\frac3{10}
\]

The rational value \(1/40\) is also a valid lower margin; it is not the
optimal constant.

\begin{quote}
\textbf{Theorem 75.2 (all-rung shifted handoff).} For every integer
\(k\ge2\), if \[
s\ge1+\frac{2k+\sqrt{2k\log k}}{\log2}
\qquad 1\le m\le2k
\tag{75.13}
\] then \[
\boxed{
(-1)^{m-1}(\log F)^{(m)}(s)
>\frac25\frac{(m-1)!}{(s-1)^m}}
\tag{75.14}
\]
\end{quote}

For completeness, let \(B_k(\Delta)\) be the \(n=2\) atom and let
\(E_k\) be the relative contribution of all remaining atoms. Atom-ratio
monotonicity gives

\[
A_k(\Delta)\le(1+E_2)B_k(\Delta)
\qquad 1+E_2<\frac85
\tag{75.15}
\]

With \(\Delta_k=\sqrt{2k\log k}\), Stirling and
\(\delta-\log(1+\delta)\ge\delta^2/2-\delta^3/3\) give

\[
B_k(\Delta_k)
\le\frac1{2\sqrt\pi}
\exp\!\left(\frac{(\log k)^{3/2}}{3\sqrt{2k}}\right)
\tag{75.16}
\]

The function \((\log k)^{3/2}/\sqrt k\) is maximized at \(k=e^3\). Its
explicit Stirling majorant has maximum
\(U_*=(2\sqrt\pi)^{-1}\exp(\sqrt{3/(2e^3)})\). Numerical evaluation
gives

\[
B_k(\Delta_k)\le U_*<\frac38,\qquad
U_*\approx0.37074729122661467
\tag{75.17}
\]

Therefore

\[
A_k(\Delta_k)<\frac85\frac38=\frac35
\]

which proves (75.14). No zeros and no RH enter either reserve theorem.

\begin{figure}
\centering
\includegraphics{figures/triangular_reserve_phase_v1_2.png}
\caption{The minimal triangular margin survives exactly through rung 11;
moving the handoff by \(\sqrt{2k\log k}\) restores an absolute reserve
for every rung.}
\end{figure}

\section{39. Theorem J - Conjugate-block phase and the blind
radius}\label{theorem-j---conjugate-block-phase-and-the-blind-radius}

Let \(q=a+ib\) with \(b\ne0\) be a nonreal folded zero, paired with
\(\bar q\). Its contribution to the reduced rung is

\[
\boxed{
W_k^{\{q,\bar q\}}(x)
=2(2k-1)!
\left(\frac{|q|}{|x+q|^2}\right)^k
\cos\!\left(k\,[\arg q-2\arg(x+q)]\right)
}
\tag{39.1}
\]

This follows immediately from (32.4) by adding a term to its conjugate.
Define

\[
\phi_q(x)=\arg q-2\arg(x+q)
\]

Direct calculation gives

\[
\operatorname{Im}\frac{q}{(x+q)^2}
=\frac{b(x^2-|q|^2)}{|x+q|^4}
\tag{39.2}
\]

Hence \(x=|q|\) is the unique positive \textbf{blind radius} at which
\(\phi_q(x)=0\) and the isolated conjugate block is positive for every
\(k\). For every \(x\ne|q|\), the phase is nonzero modulo \(2\pi\). If
\(\phi_q(x)/(2\pi)\) is rational, its finite rotation orbit contains
negative cosines and repeats them forever; if it is irrational, the
orbit is dense. Therefore the isolated block is negative for infinitely
many Widder orders.

A useful conditional detection statement follows. Put

\[
A_q(x)=\frac{|q|}{|x+q|^2}
\]

If at some \(x\ne|q|\) a nonreal pair has a strict spectral gap

\[
A_q(x)>\sup_{p\ne q,\bar q}A_p(x)
\]

then the remaining absolutely summable blocks are exponentially smaller
along a subsequence for which \(\cos(k\phi_q(x))\le-1/2\). Consequently
\(W_k(x)<0\) for infinitely many \(k\). This is an exact conditional
high-order detector. The strict-gap hypothesis is sufficient for this
orbitwise argument. Theorem 49.1 supplies a stronger general conclusion
without it: an RH failure gives infinitely many negative rungs on an
open interval of scales. That converse is proved; the independent source
sign remains open.

\begin{figure}
\centering
\includegraphics{figures/folded_block_phase.png}
\caption{For a synthetic nonreal folded pair, the normalized Widder
contribution oscillates in the rung index away from the unique blind
radius.}
\end{figure}

\subsection{Theorem J.1 - Secondary Cayley spectrum and the critical
barrier}\label{theorem-j.1---secondary-cayley-spectrum-and-the-critical-barrier}

The order index \(k\) itself carries an exact folded spectrum. Define

\[
A_x(q)=\frac{q}{(x+q)^2}
\qquad
a_k(x)=\frac{W_k(x)}{(2k-1)!}
\]

Equation (32.4) becomes the power-sum identity

\[
\boxed{
a_k(x)=\sum_{[\rho]}m_\rho A_x(q_\rho)^k
}
\tag{39.3}
\]

The scaled eigenvalue is precisely the Cayley complement already present
in the triangular collar. With \(r=q/x\) and \(p=r/(1+r)\),

\[
\boxed{
\alpha_x(q):=xA_x(q)
=\frac{r}{(1+r)^2}=p(1-p)
}
\tag{39.4}
\]

Consequently

\[
\boxed{
A_x(q)=A_x\!\left(\frac{x^2}{q}\right)
}
\tag{39.5}
\]

Thus the same reciprocal-complement involution that organized the
original Cayley geometry also organizes the Widder-order spectrum. At
the blind radius \(x=|q|\), the reciprocal image \(x^2/q\) is exactly
\(\bar q\); the conjugate pair therefore collapses to one positive
secondary eigenvalue.

Because \(\sum_{[\rho]}m_\rho|A_x(q_\rho)|<\infty\), the genus-zero
product

\[
\mathcal D_x(z)=
\prod_{[\rho]}\left(1-zA_x(q_\rho)\right)^{m_\rho}
\]

converges locally uniformly. Its logarithmic derivative is the
generating resolvent of the complete reduced-Widder ladder:

\[
\boxed{
\mathcal G_x(z):=
\sum_{k\ge1}\frac{W_k(x)}{(2k-1)!}z^{k-1}
=\sum_{[\rho]}m_\rho
\frac{A_x(q_\rho)}{1-zA_x(q_\rho)}
=-\frac{\mathcal D_x'(z)}{\mathcal D_x(z)}
}
\tag{39.6}
\]

The poles are

\[
Z_x(q)=\frac1{A_x(q)}
=\frac{(x+q)^2}{q}
=2x+q+\frac{x^2}{q}
\tag{39.7}
\]

If \(q>0\), then the arithmetic-geometric mean inequality gives

\[
0<\alpha_x(q)\le\frac14
\qquad
Z_x(q)\ge4x
\]

with equality exactly at \(q=x\). If \(q=|q|e^{i\theta}\) is nonreal,
then \(\Re q>0\) for a folded zeta zero, so \(0<|\theta|<\pi/2\). At its
blind radius \(x=|q|\),

\[
\boxed{
\alpha_x(q)
=\frac1{2(1+\cos\theta)}
=\frac1{4\cos^2(\theta/2)}>\frac14,\quad
Z_x(q)=2x(1+\cos\theta)\in(2x,4x)
}
\tag{39.8}
\]

The blind radius is therefore blind only to phase: it is
\textbf{supercritical in amplitude}. Every off-line folded zero forces a
positive real pole of \(\mathcal G_x\) strictly below the critical
barrier \(4x\) at its own scale. Conversely, critical-line folded zeros
give only poles in \([4x,\infty)\). Hence the moving-scale
secondary-pole criterion

\[
\boxed{
\mathrm{RH}
\iff
\operatorname{Pole}(\mathcal G_x)\subset[4x,\infty)
\quad\text{for every }x>0
}
\tag{39.9}
\]

is exact. It is a zero-side reorganization, not a source-side proof, but
it sharpens the research target: the oscillatory detector away from the
blind circle and the strict amplitude-barrier crossing on the blind
circle are complementary signatures of one nonreal folded orbit.

\begin{figure}
\centering
\includegraphics{figures/secondary_cayley_spectrum.png}
\caption{The normalized Widder spectrum uses the same reciprocal Cayley
complement: the critical folded ray stays below the one-quarter
amplitude barrier, while an off-axis blind pole enters the forbidden
interval.}
\end{figure}

\section{40. Source competition in heat
variables}\label{source-competition-in-heat-variables}

The source-side heat decomposition is

\[
\Theta_\xi(u)=1+\Theta_\Gamma(u)-\Theta_P(u)
\]

where the prime comb is the pointwise-positive function

\[
\Theta_P(u)=
\frac{e^{-u/4}}{2\sqrt{\pi u}}
\sum_{n\ge2}\frac{\Lambda(n)}{\sqrt n}
\exp\!\left[-\frac{(\log n)^2}{4u}\right]
\]

It enters with a minus sign. For \(k\ge1\), Theorem F gives

\[
W_k(x)=
\int_0^\infty e^{-xu}u^{2k-1}
(-1)^k\left[\Theta_\Gamma^{(k)}(u)-\Theta_P^{(k)}(u)\right]du
\tag{40.1}
\]

Thus the exact cancellation of \(1/x\) is the heat-side annihilation of
the constant completion mode. The \(W_2\) and \(W_3\) theorems establish
positivity of the weighted global Gamma/prime assembly; they do not
establish pointwise positivity of \((-1)^k\Theta_\xi^{(k)}\). The latter
for every \(k\) would be complete monotonicity of the heat trace and
hence RH.

The recurrence (38.2) is also the Laplace shadow of heat
differentiation. Writing \(H_k=(-1)^k\Theta_\xi^{(k)}\) and integrating
the derivative of \(e^{-xu}u^{2k+1}\), with the boundary terms vanishing
termwise by the folded spectral sum, gives

\[
\int e^{-xu}u^{2k+1}H_{k+1}(u)\,du
=-(2k+1)W_k'(x)-xW_k''(x).
\]

\Needspace{11\baselineskip}

\chapter{Technical chapter VII. One-scale moment and derivative
criteria}\label{technical-chapter-vii.-one-scale-moment-and-derivative-criteria}

Noncancelling poles give negative-rung witnesses. Hausdorff and
Beta--Hankel arguments reduce the full criterion to one prescribed
scale, still at every order.

\section{48. Meromorphic completion of the secondary Cayley
spectrum}\label{meromorphic-completion-of-the-secondary-cayley-spectrum}

Let \(q_\rho=\rho(1-\rho)\), one representative per folded orbit. The
zero count gives the genus-zero summability

\[
\sum_{[\rho]}\frac{m_\rho}{|q_\rho|}<\infty.
\tag{48.1}
\]

For fixed \(x>0\) put \(A_x(q)=q/(x+q)^2\) and
\(a_k(x)=\sum m_\rho A_x(q_\rho)^k\). Then

\[
\mathcal G_x(z)=\sum_{k\ge1}a_k(x)z^{k-1}
=\sum_{[\rho]}m_\rho\frac{A_x(q_\rho)}{1-zA_x(q_\rho)},
\tag{48.4}
\]

with secondary poles

\[
z_x(q)=A_x(q)^{-1}=2x+q+\frac{x^2}{q}.
\tag{48.5}
\]

For nonreal \(q=a+ib\), \(r=|q|\), the phase of \(A_x(q)\) changes sign
only at \(x=r\), where

\[
A_r(q)=\frac1{4r\cos^2(\arg q/2)}>\frac1{4r}.
\tag{48.13}
\]

Thus every nonreal folded point has an open supercritical interval
around its blind radius. The full derivation of its reciprocal endpoints
and the shifted Joukowski form of \(z_x(q)\) is retained in earlier
editions; only the consequences used later are kept here.

\subsection{48A. Exact radius deficit}\label{a.-exact-radius-deficit}

For an off-line zero write \(d=2\beta-1\), \(v=\gamma^2\),
\(a=v+(1-d^2)/4\), \(r=|q|\) and \(w=r-a\). The exact remainder is

\[
\boxed{d^2-D(q)=\frac{w[w^2+c(3r-a)]}{vr}>0,}
\tag{48.29}
\]

and therefore

\[
\boxed{\frac{4\gamma^2}{4\gamma^2+1}d^2<2w<D(q)<d^2.}
\tag{48.30}
\]

For fixed \(x\) the secondary-pole radius is attained. With
\(\Delta=\sup_x(4x-R_x)\) and \(d_*=\sup_\rho|2\Re\rho-1|\),

\[
0\le\Delta\le d_*^2\le1,
\tag{48.31}
\]

\[
\boxed{\Delta<1\iff d_*<1,\qquad \Delta=0\iff\mathrm{RH}.}
\tag{48.32}
\]

The high-scale statistic is weaker:

\[
\boxed{\limsup_{x\to\infty}(4x-R_x)
=\limsup_{|\Im\rho|\to\infty}(2\Re\rho-1)^2.}
\tag{48.33}
\]

It cannot detect finitely many off-line exceptions. At verified height
\(H\), the corresponding two-sided comparison differs by at most
\((4H^2+1)^{-1}\); this is a conditional error estimate, not a source
proof. The whole rung ladder at fixed \(x\) is one angular profile: with
\(w=z/(4x)\) and \(\cos\varphi=1-2w\),

\[
4xw\,\mathcal G_x(4xw)
=2\tan\frac\varphi2\,\Im h(xe^{i\varphi}),
\qquad h(z)=zS_\xi(z).
\tag{48.28}
\]

This is the exact bridge from the scalar \(W_k\) coefficients to the
Pick/Löwner owner. Product rewrites do not change the sign obligation:
Euler products, finite prime-power products, and the completed zero
product encode different cutoffs, and none by itself implies that the
folded poles are positive real.

\subsection{48C. Gamma pole-index sublattices do not absorb the prime
source}\label{c.-gamma-pole-index-sublattices-do-not-absorb-the-prime-source}

An adversarial Gamma-pole analysis tested the idea that the regular
Gamma pole lattice might dominate the prime-power sublattice. Use the
completion-centered terms

\[
f_m(x)=\frac{x}{R_x}\left(\frac1{2m+1}-\frac1{2m+s_x}\right),
\qquad m\ge0.
\tag{48.40}
\]

At the two positive nodes \(x_1=5/16\), \(x_2=3/4\) and \(v=(1,-1)^T\),
every single term has the same negative Löwner direction. With
\(a=4m+1\),

\[
\boxed{v^TL_{f_m}v
=-\frac{388a^3+2140a^2+2473a+624}
{1512(a+1)(a+2)^2(2a+3)^2}<0.}
\tag{48.41}
\]

For \(m=0\) alone,

\[
L_{f_0}=\begin{pmatrix}5/27&4/21\\4/21&3/16\end{pmatrix},
\qquad \det L_{f_0}=-\frac{11}{7056}.
\tag{48.42}
\]

So neither the prime-power-indexed nor complementary centered Gamma
sublattice is a positive matrix reserve. After the exact baseline repair
\(f_m=g_m-c_mb\), the \(g_m\) have positive Stieltjes measures; §166B
identifies that repaired reserve with the already-published Catalan
reserve. For a set \(A\) of retained indices, let \(B\) be its
complement and put \(F_{A,c}=cb+\sum_{m\in A}g_m-h_P\). On the upper
bank \(x=-u+i0\), with \(u>1/4\) and \(y=\sqrt{4u-1}\), the continuation
gives

\[
\boxed{\Im F_{A,c}(-u+i0)=\frac{u}{y}\left[c-a_\Gamma-
\sum_{m\in B}\frac{2y^2}{(4m+1)((4m+1)^2+y^2)}\right].}
\tag{48.43}
\]

If \(\sum_{m\in B}(4m+1)^{-1}=\infty\), monotone convergence makes the
bracket negative for all sufficiently large \(y\). Löwner's theorem
therefore rules out positive-node Löwner matrices of every rank for any
finite real \(c\). This applies both when \(A\) is the prime-power
sublattice and when \(A\) is its complement. It is an all-rank NO-GO for
the proposed absorber, not a counterexample to the completed source. The
first unsupported line remains the signed prime-error estimate.

\section{49. Pringsheim no-dominance
theorem}\label{pringsheim-no-dominance-theorem}

\begin{quote}
\textbf{Theorem 49.1 (open-interval negative-rung theorem).} If RH is
false, there is a nonempty open interval \(J\subset(0,\infty)\) such
that, for every \(x\in J\), \[
W_k(x)<0
\quad\text{for infinitely many }k
\tag{49.1}
\]
\end{quote}

\textbf{Proof.} Choose a nonreal folded zero \(q_0\). Its supercritical
interval \(I_{q_0}\) is nonempty. Remove the locally finite set of blind
radii of all nonreal folded zeros and choose a nonempty open subinterval
\(J\subset I_{q_0}\) disjoint from that set.

Fix \(x\in J\). A pole associated with \(q_0\) has modulus below \(4x\),
so the radius of convergence \(R_x\) of (48.4) satisfies

\[
R_x<4x
\tag{49.2}
\]

The minimum pole modulus is attained: by (48.5), pole moduli tend to
infinity with \(|q|\), while only finitely many folded zeros lie in any
bounded region. The common-sign residue lemma prevents cancellation, so
this minimum is the actual power-series radius.

There is no positive-real singularity in \((0,4x)\). Indeed, for a
positive folded zero \(p\),

\[
z_x(p)=2x+p+\frac{x^2}{p}\ge4x
\tag{49.3}
\]

and for a nonreal folded zero, \(A_x(q)\) is nonreal because \(x\) is
not a blind radius. Normal convergence away from the locally finite pole
set shows that there are no other singularities.

Suppose \(a_k(x)\ge0\) for every sufficiently large \(k\). Removing
finitely many initial terms does not change \(R_x\). The remaining power
series has nonnegative real coefficients and finite radius \(R_x\). The
Vivanti--Pringsheim theorem forces a singularity at the positive point
\(z=R_x\). This contradicts the absence of positive singularities below
\(4x\). Hence the real sequence \(a_k(x)\) is negative infinitely often.
Since \((2k-1)!>0\), the same is true of \(W_k(x)\). \(\square\)

This strictly strengthens the conditional dominant-orbit detector. It
uses neither a largest unique amplitude nor a gap between secondary
eigenvalues.

\begin{quote}
\textbf{Corollary 49.2 (dense-countable eventual positivity).} For every
dense set \(D\subset(0,\infty)\), \[
\boxed{
\mathrm{RH}
\iff
\forall x\in D\ \exists K_x\ \forall k\ge K_x:
W_k(x)\ge0}
\tag{49.4}
\]
\end{quote}

Under RH every rung is nonnegative. If RH fails, Theorem 49.1 supplies
an open interval \(J\); density gives \(D\cap J\ne\varnothing\), where
eventual positivity fails. In particular, \(D=\mathbb Q_{>0}\) is
sufficient.

\section{50. One prescribed Hausdorff
sequence}\label{one-prescribed-hausdorff-sequence}

The preceding criterion still uses a dense family of scales. A stronger
normalization compresses the full problem to any one scale chosen in
advance. Fix \(x_0>0\) and define

\[
\boxed{
h_n(x_0)
=(4x_0)^n\frac{W_{n+1}(x_0)}{(2n+1)!}}
\tag{50.1}
\]

Set

\[
t_x(q)=4xA_x(q)=\frac{4xq}{(x+q)^2}
\tag{50.2}
\]

Then

\[
h_n(x)
=\sum_{[\rho]}m_\rho A_x(q_\rho)t_x(q_\rho)^n
\tag{50.3}
\]

\begin{quote}
\textbf{Theorem 50.1 (fixed-scale Hausdorff criterion).} For every
prescribed \(x_0>0\), \[
\boxed{
\mathrm{RH}
\iff
(h_n(x_0))_{n\ge0}
\text{ is a Hausdorff moment sequence on }[0,1]}
\tag{50.4}
\]
\end{quote}

\textbf{RH implies the moment property.} Under RH, every \(q_\rho\) is
positive. The arithmetic--geometric mean inequality gives

\[
0<t_x(q_\rho)\le1
\tag{50.5}
\]

The finite positive measure

\[
d\mu_x(t)
=\sum_{[\rho]}m_\rho A_x(q_\rho)
\delta_{t_x(q_\rho)}(dt)
\tag{50.6}
\]

has finite mass by (48.1), and

\[
h_n(x)=\int_0^1t^n\,d\mu_x(t)
\tag{50.7}
\]

\textbf{The moment property implies RH.} Suppose (50.7) holds at the
prescribed \(x=x_0\) for some finite positive measure on \([0,1]\). Its
Markov transform

\[
H_x(z):=\sum_{n\ge0}h_n(x)z^n
=\int_0^1\frac{d\mu_x(t)}{1-tz}
\tag{50.8}
\]

is holomorphic on \(\mathbb C\setminus[1,\infty)\). On a neighborhood of
the origin, the secondary expansion gives

\[
H_x(z)
=\sum_{[\rho]}m_\rho
\frac{A_x(q_\rho)}{1-4xzA_x(q_\rho)}
\tag{50.9}
\]

Normal convergence gives a meromorphic continuation. Every candidate
pole

\[
z_\rho=\frac{1}{4xA_x(q_\rho)}
\tag{50.10}
\]

has residue \(-m_\rho/(4x)\); coincident poles cannot cancel. Uniqueness
of analytic continuation therefore forces every \(z_\rho\) onto
\([1,\infty)\). Equivalently,

\[
A_x(q_\rho)\in\left(0,\frac1{4x}\right]
\tag{50.11}
\]

If \(q_\rho\) were nonreal, (48.12) would force \(x=|q_\rho|\), while
(48.13) would give the strict reverse inequality. Hence all folded zeros
are positive real. For \(\rho=\beta+i\gamma\),

\[
\Im\bigl(\rho(1-\rho)\bigr)=\gamma(1-2\beta)
\tag{50.12}
\]

Nontrivial zeros have \(\gamma\ne0\), so \(\beta=1/2\). This is RH.
\(\square\)

The scale is not selected after inspecting zeros. The theorem holds for
every arbitrary scale, including \(x_0=1\):

\[
\boxed{
\mathrm{RH}
\iff
\left(
4^n\frac{W_{n+1}(1)}{(2n+1)!}
\right)_{n\ge0}
\text{ is Hausdorff}}
\tag{50.13}
\]

\begin{figure}
\centering
\includegraphics{figures/fixed_scale_hausdorff_map.png}
\caption{At one prescribed scale, the critical folded ray maps into the
compact Hausdorff support, while a nonreal folded zero cannot remain
there.}
\end{figure}

\section{51. Finite witnesses and a positive
contraction}\label{finite-witnesses-and-a-positive-contraction}

Hausdorff's theorem states that a real sequence is a moment sequence on
\([0,1]\) exactly when all forward differences are completely monotone.
Thus Theorem 50.1 is equivalent to

\[
\boxed{
(-1)^r\Delta^rh_n(x_0)\ge0
\qquad(n,r\ge0)}
\tag{51.1}
\]

In the original rungs this becomes

\[
\boxed{
\sum_{j=0}^{r}(-1)^j\binom rj
(4x_0)^{n+j}
\frac{W_{n+j+1}(x_0)}{(2n+2j+1)!}
\ge0}
\tag{51.2}
\]

Consequently the negation has the strong quantifier order

\[
\boxed{
\neg\mathrm{RH}
\Longrightarrow
\forall x_0>0\ \exists n,r\ge0:
(-1)^r\Delta^rh_n(x_0)<0}
\tag{51.3}
\]

At every fixed scale, an off-line zero therefore has a finite Hausdorff
witness, even if no individual late rung at that same scale has yet been
isolated.

The same theorem has a canonical operator realization.

\begin{quote}
\textbf{Theorem 51.1 (positive-contraction criterion).} For every fixed
\(x_0>0\), RH is equivalent to the existence of a Hilbert space
\(\mathcal H\), a vector \(v\in\mathcal H\), and a positive contraction
\(0\le T\le I\) such that \[
h_n(x_0)=\langle v,T^nv\rangle
\quad(n\ge0)
\tag{51.4}
\]
\end{quote}

Under RH, take \(\mathcal H=L^2([0,1],\mu_{x_0})\), let \(T\) be
multiplication by \(t\), and let \(v=1\). Conversely, the spectral
theorem for a positive contraction supplies a positive measure supported
on \([0,1]\), and Theorem 50.1 applies.

\begin{figure}
\centering
\includegraphics{figures/hausdorff_finite_witnesses.png}
\caption{Finite Hausdorff differences in a positive synthetic spectrum
and a synthetic off-axis spectrum; the latter exhibits a finite negative
witness.}
\end{figure}

\section{52. Cayley--Bernstein and Markov divided
differences}\label{cayleybernstein-and-markov-divided-differences}

The normalized secondary eigenvalue has a square-complement form. Define

\[
c_x(q)=\frac{q-x}{q+x}
\tag{52.1}
\]

Then

\[
\boxed{1-t_x(q)=c_x(q)^2}
\tag{52.2}
\]

Under RH, \(c_x(q)\in(-1,1)\) and \(t_x(q)=1-c_x(q)^2\in(0,1]\). The
finite Hausdorff differences have the exact zero-side form

\[
(-1)^r\Delta^rh_n(x)
=\sum_{[\rho]}m_\rho A_x(q_\rho)
t_x(q_\rho)^n\bigl(1-t_x(q_\rho)\bigr)^r
\tag{52.3}
\]

This is a compact Cayley--Bernstein basis: \(t^n(1-t)^r\) is nonnegative
on the critical support.

The Markov function is also a reciprocal divided difference of the
single owner

\[
h(w)=wS_\xi(w)
\tag{52.4}
\]

Choose \(u,v\) so that

\[
uv=x^2,\qquad u+v=2x(1-2z)
\tag{52.5}
\]

Then

\[
(q+x)^2-4xzq=(q+u)(q+v)
\tag{52.6}
\]

and termwise

\[
\frac{u/(u+q)-v/(v+q)}{u-v}
=\frac{q}{(u+q)(v+q)}
\tag{52.7}
\]

Summing the normally convergent folded expansion yields

\[
\boxed{
H_x(z)=\frac{h(u)-h(v)}{u-v}}
\tag{52.8}
\]

For \(z=\sin^2(\theta/2)\) one may take \(u=xe^{i\theta}\) and
\(v=xe^{-i\theta}\), so

\[
\boxed{
H_x\!\left(\sin^2\frac\theta2\right)
=\frac{\Im h(xe^{i\theta})}{x\sin\theta}}
\tag{52.9}
\]

The fixed-scale moment problem is therefore a circular Pick-sign problem
for the folded resolvent.

Under RH, (50.8) also gives the positive Pick kernel

\[
K_x(z,w)
=\frac{H_x(z)-\overline{H_x(w)}}{z-\overline w}
=\int_0^1
\frac{t\,d\mu_x(t)}
{(1-tz)(1-t\overline w)}
\tag{52.10}
\]

Every finite matrix \([K_x(z_i,z_j)]\) is positive semidefinite.

\begin{figure}
\centering
\includegraphics{figures/markov_divided_difference_circle.png}
\caption{The reciprocal pair \(u,v\) runs on a circle, and the Markov
parameter is \(z=\sin^2(\theta/2)\).}
\end{figure}

\section{104. The one-point Beta--Hankel and Hausdorff
theorem}\label{the-one-point-betahankel-and-hausdorff-theorem}

Put

\[
k(u)=-\Theta_\xi'(u)
\]

The congruenced local Dobsch matrix is

\[
\widehat H_N(x)_{ij}
=\frac1{(i+j+1)!}
\int_0^\infty e^{-xu}u^{i+j+1}k(u)\,du
\qquad 0\le i,j<N
\tag{104.1}
\]

The beta integral

\[
\frac1{(i+j+1)!}
=\frac1{i!\,j!}\int_0^1t^i(1-t)^j\,dt
\tag{104.2}
\]

turns its quadratic form, for \(P_c(z)=\sum_{j<N}c_jz^j/j!\), into

\[
\begin{aligned}
c^*\widehat H_N(x)c
={}&
\int_0^\infty e^{-xu}u\,k(u)
\int_0^1
\overline{P_c(ut)}P_c(u(1-t))\,dt\,du \\
={}&
\int_{a,b>0}
e^{-x(a+b)}k(a+b)
\overline{P_c(a)}P_c(b)\,da\,db
\end{aligned}
\tag{104.3}
\]

This exposes the additive Hankel kernel

\[
K_x(a,b)=e^{-x(a+b)}k(a+b)
\tag{104.4}
\]

\begin{quote}
\textbf{Theorem 104.1 (one prescribed basepoint).} For every prescribed
\(x_0>0\), \[
\widehat H_N(x_0)\succeq0\quad\hbox{for every }N
\iff
-\Theta_\xi'\hbox{ is completely monotone}
\tag{104.5}
\]
\end{quote}

\textbf{Proof.} Complete monotonicity gives
\(k(u)=\int_0^\infty e^{-\lambda u}\,d\mu(\lambda)\) with \(d\mu\ge0\).
Substitution into (104.3) produces

\[
\int_0^\infty
\left|
\int_0^\infty e^{-(x_0+\lambda)a}P_c(a)\,da
\right|^2d\mu(\lambda)\ge0
\tag{104.6}
\]

Conversely, the actual completed source obeys, for every
\(3/2<\beta<2\),

\[
|k(a+b)|
\le C_\beta
\bigl(1+a^{-\beta/2}\bigr)
\bigl(1+b^{-\beta/2}\bigr)
\tag{104.7}
\]

The exponential-polynomial core \(e^{-x_0a}\mathbb C[a]\) is dense in
the corresponding weighted \(L^1\) space. Positivity on every finite
polynomial core therefore extends to compactly supported tests.
Approximating point masses gives \([k(a_r+a_s)]\succeq0\) for every
finite positive node set. Thus \(k\) is exponentially convex and has a
bilateral Laplace representation by a positive measure. The
unconditional decay \(k(u)\to0\) rules out support on \((-\infty,0]\),
leaving a Bernstein--Widder measure on \((0,\infty)\). Therefore \(k\)
is completely monotone. \(\square\)

Since \(\Theta_\xi(u)\to0\),

\[
-\Theta_\xi'\hbox{ completely monotone}
\iff
\Theta_\xi\hbox{ completely monotone}
\tag{104.8}
\]

Combining this with the source criterion above gives an exact
reorganization, not an RH proof:

\[
\boxed{
\mathrm{RH}
\iff
\widehat H_N(x_0)\succeq0\quad\hbox{for every }N
}
\qquad(x_0>0\hbox{ prescribed})
\tag{104.9}
\]

Define

\[
a_n(x_0)
=\frac{(-1)^n}{(n+1)!}h^{(n+1)}(x_0)
\qquad
b_n(x_0)=x_0^na_n(x_0)
\tag{104.10}
\]

Then \(\widehat H_N(x_0)=[a_{i+j}(x_0)]\), and under the positive
Bernstein measure,

\[
b_n(x_0)=\int_0^1r^n\,d\rho_{x_0}(r)
\qquad
r=\frac{x_0}{x_0+\lambda}
\tag{104.11}
\]

Hausdorff's theorem now gives

\[
\boxed{
\mathrm{RH}
\iff
G_{n,k}(x_0):=(-1)^k\Delta^kb_n(x_0)\ge0
\quad\hbox{for every }n,k\ge0
}
\tag{104.12}
\]

The grid recursion is

\[
G_{n,k+1}=G_{n,k}-G_{n+1,k}
\tag{104.13}
\]

On the zero side it has the exact closed form

\[
\boxed{
G_{n,k}(x)
=x^n\sum_{[\rho]}m_\rho
\frac{q_\rho^{\,k+1}}{(x+q_\rho)^{n+k+2}}
}
\tag{104.14}
\]

Under RH, the representing measure is explicitly

\[
\rho_{x_0}
=\sum_{[\rho]}m_\rho
\frac{q_\rho}{(x_0+q_\rho)^2}
\delta_{x_0/(x_0+q_\rho)}
\tag{104.15}
\]

The qd/S-fraction object in PP145 and the canonical moments in PP152 use
the same algorithm on different moment sequences. The first is zero-side
Stieltjes data; the second is source-side Hausdorff data. They are not
one sequence; their different moment measures must be retained.

The one-point theorem is exact only at infinite order. Every usable
certificate is finite, and the PP153 scans show that the one-point
prefix can lose spectral reach relative to optimizing \(x\). The
fixed-scale bound must retain every order, as in (104.12).

\begin{figure}
\centering
\includegraphics{figures/v26c_sector_array.png}
\caption{The height-to-sector theorem certifies a finite rectangle at
every \(x>0\). Its diagonal \(W_j=F_{j-1,j}\) shares the same certified
height input. The open continuation lies beyond a finite index boundary,
not at the fifth rung.}
\end{figure}

\section{104A. Fixed-depth filtering, weighted recovery, and triangular
damping}\label{a.-fixed-depth-filtering-weighted-recovery-and-triangular-damping}

This section averages the boundary moments before forming a determinant
or Li form. It is different from the positive-scale moments in (104.10)
and from multiplying an already computed Li coefficient by a positive
scalar. Put

\[
b_j=\frac{4^j(j!)^2}{(2j+1)!}
 =\int_0^1u^j\frac{du}{2\sqrt{1-u}},\qquad
H_N^{[k]}=[b_{i+j}^{\,k}\mu_{i+j}],\quad D_N(k)=\det H_N^{[k]}.
\tag{104A.1}
\]

The measure is a probability measure, and depth \(k\) is multiplication
of the support coordinate by \(k\) independent variables of this law.
The matrix size, filter, and underlying moment array stay fixed; the
depth of each observation is identified when outputs are compared.

\textbf{The exact rank-two test.} With \(A=\mu_0\mu_2\) and
\(B=\mu_1^2\),

\[
\begin{aligned}
D_2(k)&=(8/15)^k A-(4/9)^k B,\\
D_2(k+2)&=\tfrac{44}{45}D_2(k+1)-\tfrac{32}{135}D_2(k),\\
D_2(0)&=\tfrac{33}{8}D_2(1)-\tfrac{135}{32}D_2(2).
\end{aligned}
\tag{104A.2}
\]

For \((\mu_0,\mu_1,\mu_2)=(1,4/5,3/5)\) the three values are
\(-1/25,8/225,448/10125\). Thus two positive observations recover a
negative unfiltered determinant. If each observation has absolute error
at most \(\eta\), the recovered error is at most \(267\eta/32\). A
single observation does not suffice: \((1,0,1/15)\) has the same
\(D_2(1)\) and a positive \(D_2(0)\).

\textbf{The triangular curvature law.} The ratio
\(b_{j+1}/b_j=(2j+2)/(2j+3)\) gives

\[
\begin{aligned}
\frac{b_{j+1}^2}{b_jb_{j+2}}&=1-\frac1{T_{2j+3}},\qquad
\Delta^2\log b_j=-\log\left(1-\frac1{T_{2j+3}}\right)>0,\\
\prod_{j=0}^{n-1}\left(1-\frac1{T_{2j+3}}\right)
 &=\frac{2n+3}{3(n+1)},\\
\prod_{i=0}^{\infty}\left(1-\frac1{T_{4i+3}}\right)
 &=\frac{4\Gamma(3/4)}{\sqrt\pi\,\Gamma(1/4)}.
\end{aligned}
\tag{104A.3}
\]

The last finite product is \((1/2)_n(5/4)_n/((3/4)_n(1)_n)\), so the
Gamma-ratio limit proves the infinite product. In particular, the
squared normalized adjacent entry at \((i,i+1)\) is multiplied by
\((1-1/T_{4i+3})^k\). These are specified triangular indices, not prime
Euler factors.

\textbf{Theorem 104A.1 (fixed-rank smoothing).} For any real array with
\(\mu_{2i}>0\), \(0\le i<N\), the matrix \(H_N^{[k]}\) is positive
definite for all sufficiently large integer \(k\), even if \(H_N\) is
indefinite.

\textbf{Proof.} Diagonal normalization leaves off-diagonal entry

\[
\frac{\mu_{i+j}}{\sqrt{\mu_{2i}\mu_{2j}}}
 \left(\frac{b_{i+j}}{\sqrt{b_{2i}b_{2j}}}\right)^k.
\]

Strict Cauchy--Schwarz in the nondegenerate Beta measure makes the ratio
strictly between zero and one for \(i\ne j\). For fixed \(N\) the
normalized matrix tends to \(I_N\); a maximum off-diagonal absolute row
sum below one already proves positivity. The normalization is a
nonsingular real congruence. \(\square\)

\textbf{Theorem 104A.2 (fixed depth retains the complete criterion).}
For the actual folded moments, and each one fixed finite integer
\(k\ge0\),

\[
\boxed{\mathrm{RH}\iff H_N^{[k]}\succeq0\quad\hbox{for every }N\ge1.}
\tag{104A.4}
\]

\textbf{Proof.} Under RH, push the positive measure
\(\sum_q(m_q/q)\delta_{1/q}\) through multiplication by the \(k\) Beta
variables. Conversely, all Hankel matrices positive semidefinite give a
positive real representing measure \(\nu\) by the Hamburger moment
theorem. Put \(C=\sum_qm_q/|q|\) and \(Y=\max_q|q|^{-1}\). Since
\(b_j\le1\), \(|b_j^k\mu_j|\le CY^j\). The even moments then give
\(\nu(|t|>Y+\epsilon)(Y+\epsilon)^{2n}\le CY^{2n}\), so
\(\operatorname{supp}\nu\subset[-Y,Y]\). The function
\(G_k(z)=\int(1-zt)^{-1}d\nu(t)\) is therefore holomorphic at every
nonreal \(z\).

Let \(V\) be the product of the \(k\) Beta variables, with \(V=1\) at
depth zero. Then

\[
B_k(z)=\mathbb E(1-zV)^{-1}=\sum_{j\ge0}b_j^kz^j,
\qquad G_k(z)=\sum_q\frac{m_q}{q}B_k(z/q)
\tag{104A.5}
\]

near zero. \(B_k\) is holomorphic off \([1,\infty)\) and singular at
one. For \(k=0\) it has a pole; for \(k\ge1\), the asymptotic
\(b_j^k\sim(\sqrt\pi/2)^kj^{-k/2}\) makes its derivative of order
\(\lceil k/2\rceil\) diverge along \(z\uparrow1\).

If a nonreal folded point exists, choose one, \(q_*\), of least modulus
among such points and approach it along \(z=tq_*\), \(0<t<1\). Its term
becomes nonanalytic there. No distinct cut \(q[1,\infty)\) contains
\(q_*\): a real cut cannot, and a nonreal one would require a smaller
modulus. Coincident values have already been combined into positive
multiplicities. All other finite terms are analytic near \(q_*\); the
infinite tail is normally convergent there, with majorant a constant
times \(m_q/|q|\). They cannot cancel the distinguished singularity.
This contradicts the holomorphy of \(G_k\) at \(q_*\). Thus every folded
point is real, and \(\Im q_\rho=\Im\rho(1-2\Re\rho)\), with
\(\Im\rho\ne0\), gives RH. \(\square\)

Fixing depth before imposing all ranks is essential. Theorem 104A.1 can
make \(\forall N\ \exists k(N)\) positivity automatic; it does not
establish (104A.4).

\textbf{Proposition 104A.3 (signed finite-mode recovery).} For known
distinct positive rates \(\mathcal R\), \(M=|\mathcal R|\), define
weights by

\[
1-\prod_{r\in\mathcal R}(1-z/r)=\sum_{\ell=1}^M w_\ell z^\ell.
\quad
F(k)=\sum_{r\in\mathcal R}A_rr^k
\ \Longrightarrow\ F(0)=\sum_{\ell=1}^M w_\ell F(\ell).
\tag{104A.6}
\]

Evaluation at each \(z=r\) proves the result. For independent errors
bounded by \(\eta\), the exact worst-case amplifier is
\(\kappa=\sum|w_\ell|=\prod_{r\in\mathcal R}(1+1/r)-1\). The weights
alternate in sign.

For determinants, expand over permutations and group equal nonzero
monomial/rate contributions:
\(D_N(k)=\sum_\pi\operatorname{sgn}(\pi)\prod_i\mu_{i+\pi(i)}
(\prod_i b_{i+\pi(i)})^k\). The multiplicities \(\nu_j\) in each
monomial obey \(\sum\nu_j=N\) and \(\sum j\nu_j=N(N-1)=2T_{N-1}\). The
exact enumerated rate counts for ranks \(1\) through \(9\) are
\(1,2,5,16,58,225,980,4520,22244\). This is finite enumeration, not an
all-rank count theorem. A geometric filter would give a single common
factor; the strict curvature in (104A.3) separates the rates.

For one odd Li form the problem is much smaller. Put
\(c_{m,j}=[y^j]\mathcal P_m(y)^2\), \(J=2m\), and
\(\Lambda_m(k)=\sum_{j=0}^Jc_{m,j}b_j^k\mu_j\). Use rates
\(b_0,\ldots,b_J\) in (104A.6). Then
\(\lambda_{2m+1}=\sum_{\ell=1}^{2m+1}w_\ell\Lambda_m(\ell)\):
\(\lambda_{17}\) uses 17 observations, not 22,244 determinant modes. The
independent-error amplifier satisfies
\(\log\kappa_m=\tfrac12(2m+1)\log(2m+1)+O(m)\).

\textbf{A mandatory negative control.} Take \(q_0=17/4\) of multiplicity
two and \(q_\pm=147/16\pm3i/2\) of multiplicity one. Every model moment
\(\mu_j^*=2q_0^{-j-1}+2\Re(q_+^{-j-1})\) is positive, since
\(|q_+|>q_0\). Exact rational evaluation gives 37 positive observations
\(\Lambda_{18}^*(1),\ldots,\Lambda_{18}^*(37)\) but recovers
\(-0.598<\lambda_{37}^*<-0.597\). The separate directed-error ledger
reproduces this negative sign. This is a polynomial comparison model,
not an actual-zeta counterexample. The actual source still has to prove
the signed comparison between the odd- and even-depth terms in (104A.6).
Section 166E treats common, rather than independent, errors.

\section{136. The completed-source Pascal
owner}\label{the-completed-source-pascal-owner}

Let

\[
h(x)=\frac{x}{R}\Phi(s),\qquad x=s(s-1),\qquad R=2s-1
\]

and let \(G_{n,r}\) be the completed Widder/Pascal cells defined in
(104.14). They satisfy the exact addition law

\[
\boxed{G_{n,r}=G_{n+1,r}+G_{n,r+1}}
\tag{136.1}
\]

Whenever \(G_{n,r}>0\), define the branch ratios

\[
p_{n,r}=\frac{G_{n+1,r}}{G_{n,r}},\qquad
q_{n,r}=\frac{G_{n,r+1}}{G_{n,r}}
\]

Then \(p_{n,r}+q_{n,r}=1\). Differentiating (104.14) gives
\(x\partial_xG_{n,r}=nG_{n,r}-(n+r+2)G_{n+1,r}\). Where both adjacent
cells \(G_{n+1,r}\) and \(G_{n,r+1}\) are nonnegative, the division-free
flow corridor is

\[
\boxed{
-(r+2)G_{n,r}
\le x\partial_xG_{n,r}
\le nG_{n,r}}
\tag{136.2}
\]

Equation (136.2) requires the two adjacent signs, but no division by
\(G_{n,r}\). Where the denominator is positive, those signs also put
both branch ratios in \([0,1]\). PP186 certifies all \(81\) cells with
\(0\le n,r\le8\) at \(x=56\) by directed rational intervals. PP187
certifies the directed confluent rank-eleven block

\[
\boxed{\mathfrak B_{11}^{\mathrm{conf}}(56)\succ0}
\tag{136.3}
\]

These are finite source certificates. By the finite-rank camouflage
theorem, no finite prefix proves the first unsupported one-point line
without additional completed-source structure:

\[
\boxed{G_{n,r}(56)\ge0\qquad(n,r\ge0)}
\tag{136.4}
\]

\Needspace{11\baselineskip}

\chapter{Technical chapter VIII. Finite Löwner certificates and spectral
reach}\label{technical-chapter-viii.-finite-luxf6wner-certificates-and-spectral-reach}

The normalized determinant and quartet signature precede the certificate
details. Global matrix order and rank at a single basepoint remain
distinct.

\section{85. Finite determinant gates through rank
six}\label{finite-determinant-gates-through-rank-six}

Throughout the finite-gate sections, put \(h=xS_\xi\) and, for
\(r\ge2\), \[
M_r(x)=\frac{(-1)^r h^{(r-1)}(x)}{(r-1)!},\qquad
\Delta_N^{\mathrm L}(x)=\det[M_{i+j}(x)]_{i,j=1}^N
\] The alternating-sign diagonal congruence identifies this determinant
with that of the normalized confluent Löwner matrix. In particular
\(M_2=h'=W_1\). This fixes the normalization used in (85.1) and (93.1).

The PP98--PP111 certificate chain establishes, without assuming RH,

\[
\boxed{\Delta_N^{\mathrm L}(x)>0
\quad(x>0),\qquad N=1,2,3,4,5,6}
\tag{85.1}
\]

The ranks \(4,5,6\) are classified \textbf{A-CAT}: exact symbolic owners
are joined to directed interval or exact zero-box compact certificates
and analytic source tails. The finite-zero lanes use authenticated zero
boxes only in their declared ranges; the source tails use no zero
ordinates. Equation (85.1) is a finite leading-minor statement, not the
all-order source theorem.

\section{102. Off-line quartet signature, determinant factorization, and
reach}\label{off-line-quartet-signature-determinant-factorization-and-reach}

Let an off-line zero be

\[
\rho=\frac12+\delta+i\gamma
\qquad 0<\delta<\frac12
\]

Its functional-equation and conjugation orbit gives two folded values

\[
a=\frac14-\delta^2+\gamma^2-2i\delta\gamma
\qquad \bar a
\tag{102.1}
\]

For \(t=a/x\) define

\[
v_N(t)=\sqrt t\,
\bigl((1+t)^{-1},\ldots,(1+t)^{-N}\bigr)^{\mathsf T}
\tag{102.2}
\]

Write \(v=u+iz\). The conjugate pair contributes

\[
v v^{\mathsf T}+\bar v\bar v^{\mathsf T}
=2uu^{\mathsf T}-2zz^{\mathsf T}
\tag{102.3}
\]

\begin{quote}
\textbf{Theorem 102.1 (quartet signature).} For \(N\ge2\), every
off-line quartet contributes a real rank-two form of signature
\((1,1)\): one positive and one negative direction.
\end{quote}

\textbf{Proof.} The vectors \(u\) and \(z\) are independent. Otherwise
\(v\) would be a complex scalar times a real vector, so every
consecutive ratio \(v_{j+1}/v_j=(1+t)^{-1}\) would be real. This is
impossible because \(t\) is nonreal. On the plane spanned by \(u,z\),
the determinant of \(2(uu^{\mathsf T}-zz^{\mathsf T})\) is

\[
-4\bigl(\|u\|^2\|z\|^2-\langle u,z\rangle^2\bigr)<0
\]

The two nonzero eigenvalues therefore have opposite signs. \(\square\)

Collect the on-line atoms and the positive halves of all off-line
quartets in \(P_N\succeq0\), and let \(W_N\) contain the negative
vectors. Whenever \(P_N\) is invertible, the matrix determinant lemma
gives

\[
\boxed{
\Delta_N^{\mathrm L}(x)
=x^{-N^2}\det P_N(x)
\det\!\left(I-W_N^{\mathsf T}P_N(x)^{-1}W_N\right)
}
\tag{102.4}
\]

For one quartet, with negative vector \(w\),

\[
\nu_N(x)=w_N^{\mathsf T}P_N(x)^{-1}w_N
\qquad
\Delta_N^{\mathrm L}(x)>0
\iff \nu_N(x)<1
\tag{102.5}
\]

The rank monotonicity is exact:

\[
\nu_N
=\max_{y\in\mathbb R^N}
\left(2y^{\mathsf T}w_N-y^{\mathsf T}P_Ny\right)
\tag{102.6}
\]

Embedding \(y\mapsto(y,0)\) preserves the objective, hence

\[
\boxed{\nu_{N+1}(x)\ge\nu_N(x)}
\tag{102.7}
\]

Thus the failure ranks form an up-set. This proves that the finite
ladder is ordered; it does not prove how high a fixed rank can see.

The PP133 scan reported that rank eight becomes blind to a maximally
displaced quartet above approximately \(\gamma=27\). That statement
remains \textbf{C}: PP134 was assigned to replace the 45-point \(x\)
scan by a rational cover using

\[
\nu_N\le \|w_N\|^2\operatorname{tr}(P_N^{-1})
\tag{102.8}
\]

but no PP134 certificate occurs in the recorded trace. Likewise, the
finite-ladder zero-localization conclusion remains a \textbf{proposed
no-go}, not a replayed theorem. Neither limitation affects the proved
global eighth gate (92.1).

\section{86. Rank seven: global certificate and
overlap}\label{rank-seven-global-certificate-and-overlap}

PP112 gives a three-lane global certificate with low range
\(0<x\le92116\), a \(d=1/20\) analytic middle lane, and a
safe-half-plane source tail. PP113 supplies an independent
\(d=1/28=1/T_7\) alternate with exact zero-box bridge
\(92116\le x\le157976\), an unsplit middle-core ratio
\(1.08765547236\ldots\), and a 128-box backup ratio
\(2.4012758008\ldots\).

The PP113 package hashes and terminal, bridge, splice, and cap receipts
have been reproduced twice byte-identically after direct extraction of
the canonical dependency. A separate full PP112 producing replay then
reproduced the middle, low, source-tail, coordinate, and splice receipts
byte-for-byte. Therefore

\[
\boxed{\Delta_7^{\mathrm L}(x)>0\ (x>0):\quad\textbf{A-CAT}}
\tag{86.1}
\]

The canonical receipt hashes are:

\noindent middle:
\path{2a90a75aca94fb4adee1177ef5e7124a4575f3a4c36a252335ca28059b09908c}

\noindent low:
\path{4c8f04809c8ea7a0fd358a35587f27d023c3f4b8ae1423112b33af272db6dadb}

\noindent tail:
\path{c5b6a3c5db3bfb5b2e72e0ec36569bdadf3ecd3c61ac2e42d8691cbd54a45ec8}

\noindent coordinate:
\path{b884ac47b2ce7d7ac0c271f88928e3739de98a0c53f69003fc888fb9d56974b3}

\noindent splice:
\path{317de1f7a05a16ce7fe8f42a9a83eedad2373b4e4b23cceb9116be6db9463e12}

The exact identity \(x(x-1)=2T_{x-1}\) and the repetend/triangular
observations may be retained as coordinate mnemonics. They do not select
the optimal band width and are not proof owners.

\section{\texorpdfstring{87. Rank eight: the \(q=57\) certificate
geometry}{87. Rank eight: the q=57 certificate geometry}}\label{rank-eight-the-q57-certificate-geometry}

PP114 constructs a four-lane rank-eight splice. Its \(q=56,L=6\)
adaptive cover has weakest ratio

\[
1.000006691152949\ldots
\tag{87.1}
\]

only \(6.7\) parts per million above one. The same exact receipt
contains a strictly stronger \(q=57,L=6\) cover:

\[
\begin{aligned}
&729\ \text{leaves},\qquad 11{,}649\ \text{corner evaluations}\\
&\text{weakest ratio }1.0001919333825133\ldots\\
&250000\le x\le750081\quad\text{for the finite bridge}\\
&75\ \text{bridge cells},\qquad861\ \text{authenticated boxes}\\
&x\le\frac{504058913542329082327040000}{457}
\quad\text{for the analytic middle lane}\\
&\text{middle/tail overlap }
\frac{180421979233767077267426}{457}>0
\end{aligned}
\tag{87.2}
\]

Every one of the 729 exact rational leaves and all 75 bridge cells has
passed an independent Fraction recomputation. The recorded Rail G and
Rail L receipts replay byte-identically; Rail T agrees in every
mathematical field and differs only in a recorded path string. The
completed global proof in §92 supplies the theorem authority for this
certificate geometry:

\[
\boxed{\Delta_8^{\mathrm L}(x)>0\ (x>0):\quad\textbf{A-CAT}\ \text{(§92)}}
\tag{87.3}
\]

For comparison, the \(q=80,L=7\) ratio \(1.100576995062874\ldots\) uses
one midpoint split in each of eight coordinates, hence \(2^8=256\)
rational subboxes. The unsplit ratios are \(0.8253553988179236\) for
\(L=7\) and \(0.9432633129347698\) for \(L=8\); the first passing
unsplit multiplicity is \(L=9\), ratio \(1.0611712270516161\). Moreover,
the proposed sparse bridge fails at \(x=250000\). The \(q=80\) lane is
therefore reconnaissance, not a theorem or preferred certificate.

Section 92 joins the low, bridge, middle, tail, and splice kernels into
the global eighth-gate proof. The arithmetic parameter \(q=57\) alone
does not establish any determinant sign.

\section{92. Global eighth local
determinant}\label{global-eighth-local-determinant}

\begin{quote}
\textbf{Theorem 92.1 (\texttt{D8-GLOBAL}, A-CAT).} For every real
\(x>0\), \[
\boxed{\Delta_8^{\mathrm L}(x)>0}
\tag{92.1}
\]
\end{quote}

Put

\[
X_{\rm cap}=\frac{504058913542329082327040000}{457}
\tag{92.2a}
\]

The authenticated PP131 splice is the gap-free union

\[
\begin{array}{rcl}
0<x\le250000 && \text{actual-node low theorem}\\
250000\le x\le337267 && \text{exact-box bridge}\\
337267\le x\le X_{\rm cap}
&& \text{$q=57,L=6$ ray and uniform core}\\
x\ge1102578756155569617614382 && \text{analytic source tail}
\end{array}
\tag{92.2}
\]

The provenance substitution is explicit in PP133 Section C.2: PP117's
independent clean-room replay occurs byte-for-byte inside PP131 as lane
\texttt{D8-Q57-L6-UNIFORM-CORE}. Thus PP131 consumes, rather than
replaces, the named PP117 rederivation.

The cap overlaps the tail by the exact positive amount

\[
\boxed{
\frac{180421979233767077267426}{457}>0
}
\tag{92.3}
\]

The independent replays reproduced all declared dependencies. The full
package and replay hashes are held with the certificate archive, not
here.

The finite verified-height input \(H=3000175332800\) is an unconditional
Platt--Trudgian theorem, not an RH assumption. The archived exact-box
chain and declared \texttt{zeros\_14.dat} MD5 are authenticated, but the
raw \texttt{.dat} bytes were not mounted for a fresh rehash or decode.
This provenance limitation is stated rather than hidden.

\section{93. Fixed-order Dobsch--Löwner
consequence}\label{fixed-order-dobschluxf6wner-consequence}

Let \(h(x)=xS_\xi(x)\) and \(W=h'\). With

\[
\mathcal H_8(x)=
\left[\frac{h^{(i+j-1)}(x)}{(i+j-1)!}\right]_{i,j=1}^{8}
\]

the fixed alternating-sign congruence gives

\[
\mathcal H_8(x)=D[M_{i+j}(x)]_{i,j=1}^{8}D
\qquad D=\operatorname{diag}(1,-1,\ldots,-1)
\tag{93.1}
\]

The global positivity of all leading determinants through rank eight and
Sylvester's criterion give \(\mathcal H_8(x)\succ0\). The audited
Dobsch--Donoghue theorem and Löwner characterization therefore yield:

\begin{quote}
\textbf{Theorem 93.1 (\texttt{DD-ORDER8} / \texttt{SRC-LOEWNER-N8},
A-CAT).} The folded owner \(h\) is matrix monotone of order eight on
\((0,\infty)\). Consequently, \[
\boxed{L_\Gamma-L_P\succeq0}
\tag{93.2}
\] for every positive node tuple of size at most eight.
\end{quote}

The conclusion is positive semidefinite, not universal strict positive
definiteness. The complete dependency digest is held with the
certificate archive. Order nine, arbitrary finite order, and the
operator intertwiner remain open.

\begin{figure}
\centering
\includegraphics{figures/v26c_rank_scope.png}
\caption{Global order and fixed-node rank are distinct directions. Order
eight is certified globally; rank twelve is certified at \(x=56\). The
fixed-node certificate does not fill the global order-nine region. See
§§92--93 and 138.}
\end{figure}

\section{138. Rank twelve at its exact finite
authority}\label{rank-twelve-at-its-exact-finite-authority}

PP190-T replaces the failed entrywise Gershgorin tail bound by a
correlation-preserving one-measure Schur argument. The controlling
manifest-bound production branch uses the exact rational unit-lower
preconditioner selected from its finite anchor. Its terminal polynomial
satisfies

\[
q_z(2+u)=\sum_{k=0}^{23}a_ku^k
\qquad a_k\in\mathbb Q_{\ge0},\qquad a_0>0
\tag{138.1}
\]

Hence \(q_z(t)>0\) for \(t\ge2\). The prime atoms themselves remain
indefinite; the positivity occurs only after their common-source
correlation is retained. At cutoff \(120{,}750{,}000\), the certified
Schur margin is

\[
\boxed{
\operatorname{Schur}
(\mathfrak B_{12}^{\mathrm{conf}}\mid
 \mathfrak B_{11}^{\mathrm{conf}})
\ge
1.2252441364765957718\times10^{-41}>0}
\tag{138.2}
\]

Therefore

\[
\boxed{\mathfrak B_{12}^{\mathrm{conf}}(56)\succ0}
\tag{138.3}
\]

Equation (138.3) is a finite rank-twelve interval certificate at one
basepoint, with the cutoff and margin in (138.2). Equation (138.3)
proves neither global matrix monotonicity of order nine (meaning global
in \(x\)), the all-node Löwner inequality, nor RH.

\Needspace{11\baselineskip}

\chapter{Technical chapter IX. Cyclic arithmetic, Riesz smoothing, and
zero
frontiers}\label{technical-chapter-ix.-cyclic-arithmetic-riesz-smoothing-and-zero-frontiers}

Cyclic denominator counts enter the zeta quotient through an Euler
product. Smoothing controls arithmetic jumps and exposes the
zero-frontier exponent.

\section{56. Cyclic arithmetic before
zeta}\label{cyclic-arithmetic-before-zeta}

To avoid collision with \(S_\xi\), define

\[
\Psi_{\rm cyc}(q)
=\sum_{a=1}^{q}\frac{q}{\gcd(a,q)}
\tag{56.1}
\]

Grouping residues by their order in the cyclic group gives

\[
\boxed{
\Psi_{\rm cyc}(q)
=\sum_{d\mid q}d\varphi(d)}
\tag{56.2}
\]

For a prime power,

\[
\Psi_{\rm cyc}(p^k)
=1+(p-1)\sum_{j=1}^{k}p^{2j-1}
=\boxed{\frac{p^{2k+1}+1}{p+1}}
\tag{56.3}
\]

Therefore the local Dirichlet factor is

\[
\sum_{k\ge0}\Psi_{\rm cyc}(p^k)p^{-ks}
=\frac{1-p^{1-s}}
{(1-p^{-s})(1-p^{2-s})}
\tag{56.4}
\]

and only then does the Euler product package the arithmetic:

\[
\boxed{
F_{\rm cyc}(s)
:=\sum_{q\ge1}\frac{\Psi_{\rm cyc}(q)}{q^s}
=\frac{\zeta(s)\zeta(s-2)}{\zeta(s-1)}}
\tag{56.5}
\]

\subsection{Singularity
classification}\label{singularity-classification}

The quotient has:

\begin{enumerate}
\def\labelenumi{\arabic{enumi}.}
\tightlist
\item
  a simple pole at \(s=3\) with residue \(\zeta(3)/\zeta(2)\);
\item
  a simple pole at \(s=1\) with residue \(1/6\);
\item
  for every nontrivial zero \(\rho\), a pole at \(s=1+\rho\) whose order
  is the multiplicity of \(\rho\);
\item
  no other poles in \(\Re s>1\).
\end{enumerate}

The inherited pole cannot cancel. At \(s=1+\rho\), the numerator factors
are \(\zeta(1+\rho)\) and \(\zeta(\rho-1)\). The first lies in the
zero-free half-plane \(\Re s>1\); the second is neither a nontrivial
zero nor a negative even integer. Thus both are nonzero.

\subsection{Sharp-cutoff no-go}\label{sharp-cutoff-no-go}

Let

\[
A(X)=\sum_{q\le X}\Psi_{\rm cyc}(q),\qquad
C=\frac{\zeta(3)}{3\zeta(2)}
\tag{56.6}
\]

and

\[
E(X)=A(X)-CX^3-\frac X6
\tag{56.7}
\]

For every prime \(p\), \(\Psi_{\rm cyc}(p)=p^2-p+1\), so

\[
\boxed{
E(p)-E(p-1)
=\left(1-\frac{\zeta(3)}{\zeta(2)}\right)(p^2-p)
+\frac56-C}
\tag{56.8}
\]

The quadratic coefficient is positive. For each prime, at least one of
\(|E(p)|\) and \(|E(p-1)|\) is at least half the jump. Since there are
infinitely many primes,

\[
\boxed{
E(X)=\Omega(X^2),\qquad
\limsup_{X\to\infty}\frac{|E(X)|}{X^2}
\ge\frac12\left(1-\frac{\zeta(3)}{\zeta(2)}\right)}
\tag{56.9}
\]

The unsmoothed \(O_\varepsilon(X^{3/2+\varepsilon})\) target is
therefore impossible, regardless of RH.

\begin{figure}
\centering
\includegraphics{figures/cyclic_sharp_vs_riesz.png}
\caption{Finite cyclic data show the quadratic sharp-cutoff cusps and
the continuous entry created by the first Riesz weight.}
\end{figure}

\section{57. First-Riesz cyclic
criterion}\label{first-riesz-cyclic-criterion}

Define the first Riesz mean

\[
\boxed{
\mathcal R_1(X)
=\sum_{q\le X}\Psi_{\rm cyc}(q)
\left(1-\frac qX\right)}
\tag{57.1}
\]

Mellin inversion gives, initially for \(c>3\),

\[
\boxed{
\mathcal R_1(X)
=\frac{1}{2\pi i}\int_{(c)}
F_{\rm cyc}(s)\frac{X^s}{s(s+1)}\,ds}
\tag{57.2}
\]

The residue at \(s=3\) is

\[
D X^3,\qquad
D=\frac{\zeta(3)}{12\zeta(2)}
\tag{57.3}
\]

The denominator has an exact adjacent-triangular reading:

\[
\boxed{
\left.\frac1{s(s+1)}\right|_{s=3}
=\frac1{3\cdot4}=\frac1{12}=\frac1{2T_3},\qquad
D=\frac1{2T_3}\frac{\zeta(3)}{\zeta(2)}}
\tag{57.3a}
\]

The continuation value \(\zeta(-1)=-1/12\) has the opposite sign but a
different analytic owner. The triangular reflection \(3\mapsto-4\) keeps
\(T_3=T_{-4}\) and therefore keeps the kernel factor positive; it does
not produce a separate identity \(-D\). The polar-shell picture of
{[}55{]} uses the same reciprocal-radius vocabulary but is not the proof
of this residue.

\begin{quote}
\textbf{Theorem 57.1 (first-Riesz RH equivalence).} \[
\boxed{
\mathrm{RH}
\iff
\mathcal R_1(X)
=DX^3+O_\varepsilon(X^{3/2+\varepsilon})
\quad\text{for every }\varepsilon>0}
\tag{57.4}
\]
\end{quote}

\textbf{RH implies the estimate.} Fix \(0<\delta<1/2\) and write
\(s=3/2+\delta+it\). Since \(\Re s>1\), \(\zeta(s)=O_\delta(1)\). The
functional equation gives

\[
\zeta\!\left(-\frac12+\delta+it\right)
\ll_\delta(1+|t|)^{1-\delta}
\tag{57.5}
\]

Under RH, Littlewood's reciprocal-zeta estimate in every fixed
half-plane to the right of the critical line gives, for every
\(\eta>0\),

\[
\frac1{\zeta(1/2+\delta+it)}
\ll_{\delta,\eta}(1+|t|)^\eta
\tag{57.6}
\]

Consequently

\[
F_{\rm cyc}(3/2+\delta+it)
\ll_{\delta,\eta}(1+|t|)^{1-\delta+\eta}
\tag{57.7}
\]

After the Riesz kernel, the vertical integrand is
\(O(|t|^{-1-\delta+\eta})\), which is absolutely integrable when
\(0<\eta<\delta\). The same fixed-strip bounds make the horizontal
integrals vanish. Shifting from \(c>3\) to \(\Re s=3/2+\delta\) crosses
only \(s=3\) and proves

\[
\mathcal R_1(X)
=DX^3+O_\delta(X^{3/2+\delta})
\tag{57.8}
\]

\textbf{The estimate implies RH.} For \(\Re s>3\), Tonelli's theorem and
a direct beta integral give

\[
\boxed{
\int_1^\infty\mathcal R_1(X)X^{-s-1}\,dX
=\frac{F_{\rm cyc}(s)}{s(s+1)}}
\tag{57.9}
\]

Therefore

\[
\int_1^\infty
\bigl(\mathcal R_1(X)-DX^3\bigr)X^{-s-1}\,dX
=\frac{F_{\rm cyc}(s)}{s(s+1)}-\frac{D}{s-3}
\tag{57.10}
\]

The assumed error makes the left side holomorphic in every half-plane
\(\Re s>3/2+\varepsilon\). An off-line zero with
\(\Re\rho>1/2+\varepsilon\) would create an uncancelled inherited pole
at \(s=1+\rho\) inside that half-plane. This is impossible. Letting
\(\varepsilon\downarrow0\) and applying functional-equation symmetry
gives RH. \(\square\)

The \(s=1\) residue would contribute \(X/12\) only after a contour shift
below \(\Re s=1\). It is below the stated error and is not crossed in
the RH proof.

\section{58. Fractional Riesz threshold and prime
cusp}\label{fractional-riesz-threshold-and-prime-cusp}

For \(\lambda>0\), define

\[
\mathcal R_\lambda(X)
=\frac1{\Gamma(\lambda+1)}
\sum_{q\le X}\Psi_{\rm cyc}(q)
\left(1-\frac qX\right)^\lambda
\tag{58.1}
\]

Its Mellin kernel is

\[
\boxed{
\frac{\Gamma(s)}{\Gamma(s+\lambda+1)}}
\tag{58.2}
\]

The \(s=3\) residue gives

\[
D_\lambda X^3,\qquad
D_\lambda=
\frac{2\zeta(3)}
{\zeta(2)\Gamma(\lambda+4)}
\tag{58.3}
\]

On \(\Re s=3/2+\delta\), Stirling's formula combines with (57.7) to give
an integrand of order

\[
|t|^{-\lambda-\delta+\eta}
\tag{58.4}
\]

Thus the direct absolute-contour proof closes for every \(\lambda\ge1\).
Conversely, for every \(\lambda>0\), an error
\(O_\varepsilon(X^{3/2+\varepsilon})\) would continue the Mellin
transform past all inherited poles and hence imply RH.

For \(0<\lambda<1\), a separate arithmetic obstruction appears. Fix
\(0<h<1/2\) and take the central second difference at a prime \(p\). The
new prime term occurs only at \(p+h\) and equals

\[
\frac{\Psi_{\rm cyc}(p)}{\Gamma(\lambda+1)}
\left(\frac{h}{p+h}\right)^\lambda
\asymp_\lambda h^\lambda p^{2-\lambda}
\tag{58.5}
\]

For terms already present at \(p-h\), Taylor's theorem and
\(\Psi_{\rm cyc}(n)\le n^2\) give

\[
\sum_{n\le p-1}\Psi_{\rm cyc}(n)
\left|\Delta_h^2
\left(1-\frac nX\right)^\lambda_{\!X=p}\right|
\le C_\lambda h^2p^{2-\lambda}
\tag{58.6}
\]

The summable factor is \(\sum_{r\ge1}r^{\lambda-2}<\infty\). Lower-order
differentiated factors and the second difference of the cubic main term
are \(O_\lambda(h^2p)\). Choose \(h=h(\lambda)\) so small that the
\(h^\lambda\) entering cusp dominates the \(h^2\) terms. Along every
sufficiently large prime,

\[
\boxed{
\left|
\Delta_h^2\bigl(\mathcal R_\lambda-D_\lambda X^3\bigr)(p)
\right|
\gg_\lambda p^{2-\lambda}}
\tag{58.7}
\]

If \(\lambda<1/2\), choose \(0<\varepsilon<1/2-\lambda\). Then
\(2-\lambda>3/2+\varepsilon\), contradicting the proposed RH-scale
error.

The rigorous phase diagram is

\[
\boxed{
\begin{array}{c|c}
0<\lambda<1/2&
\text{RH-scale estimate impossible}\\
1/2\le\lambda<1&
\text{estimate implies RH; converse open}\\
\lambda\ge1&
\text{RH-equivalent by the absolute-contour proof}
\end{array}}
\tag{58.8}
\]

The natural next boundary is \(\lambda=1/2\), not the unsmoothed cutoff.

\begin{figure}
\centering
\includegraphics{figures/fractional_riesz_phase_diagram.png}
\caption{The fractional Riesz phase diagram separates the prime-cusp
no-go region, the open middle interval, and the direct RH-equivalent
range.}
\end{figure}

\section{59. Denominator-moment
hierarchy}\label{denominator-moment-hierarchy}

For each integer \(j\ge1\), define

\[
\Psi_j(q)=\sum_{d\mid q}d^j\varphi(d)
\tag{59.1}
\]

The local factor and global quotient are

\[
\sum_{k\ge0}\Psi_j(p^k)p^{-ks}
=\frac{1-p^{j-s}}
{(1-p^{-s})(1-p^{j+1-s})}
\tag{59.2}
\]

and

\[
\boxed{
F_j(s)
=\sum_{q\ge1}\frac{\Psi_j(q)}{q^s}
=\frac{\zeta(s)\zeta(s-j-1)}{\zeta(s-j)}}
\tag{59.3}
\]

Every nontrivial zero \(\rho\) creates an uncancelled pole at
\(s=j+\rho\). The main pole at \(s=j+2\) has residue
\(\zeta(j+2)/\zeta(2)\).

Define

\[
\mathcal R_{j,1}(X)
=\sum_{q\le X}\Psi_j(q)\left(1-\frac qX\right)
\tag{59.4}
\]

The proof of Theorem 57.1 translates without loss: on
\(\Re s=j+1/2+\delta\), the two critical factors are again
\(\zeta(-1/2+\delta+it)\) and \(1/\zeta(1/2+\delta+it)\), while
\(\zeta(s)\) is bounded. Hence:

\begin{quote}
\textbf{Theorem 59.1 (cyclic moment hierarchy).} For every integer
\(j\ge1\), \[
\boxed{
\mathrm{RH}
\iff
\mathcal R_{j,1}(X)
=
\frac{\zeta(j+2)}
{\zeta(2)(j+2)(j+3)}X^{j+2}
+O_{\varepsilon,j}(X^{j+1/2+\varepsilon})}
\tag{59.5}
\]
\end{quote}

The half-power gap between the raw prime jump \(X^{j+1}\) and the RH
spectral scale \(X^{j+1/2}\) is universal across the hierarchy.

\begin{figure}
\centering
\includegraphics{figures/cyclic_moment_hierarchy.png}
\caption{The \(j=1\) cyclic theorem is the first member of a full
denominator-moment family with inherited poles \(j+\rho\).}
\end{figure}

\section{60. The zero-frontier
exponent}\label{the-zero-frontier-exponent}

Let

\[
\beta_*=\sup_\rho\Re\rho
\tag{60.1}
\]

and define

\[
\alpha_1
=\inf\left\{\alpha:
\mathcal R_1(X)-DX^3
=O_\varepsilon(X^{\alpha+\varepsilon})
\text{ for every }\varepsilon>0\right\}
\tag{60.2}
\]

The inherited poles give the unconditional lower bound

\[
\boxed{\alpha_1\ge1+\beta_*}
\tag{60.3}
\]

Any smaller exponent would analytically continue (57.10) across a pole
\(1+\rho\) with real part arbitrarily close to \(1+\beta_*\). The
reverse inequality follows from Littlewood's classical reciprocal-zeta
estimate in each fixed zero-free half-plane, combined with the same
first-Riesz contour argument. Thus, with that standard cited input,

\[
\boxed{\alpha_1=1+\beta_*}
\tag{60.4}
\]

This is a spectral-frontier identity, not a proof that \(\beta_*=1/2\).
It explains exactly why first Riesz smoothing measures the rightmost
zeta zeros.

\section{61. Cayley--Montgomery bridge: structural
only}\label{cayleymontgomery-bridge-structural-only}

Under RH, \(q_\gamma=1/4+\gamma^2\). At a large reference height \(T\),
set

\[
x_T=T^2+\frac14
\qquad
c_T(\gamma)=
\frac{q_\gamma-x_T}{q_\gamma+x_T}
=\frac{\gamma^2-T^2}{\gamma^2+T^2+1/2}
\tag{61.1}
\]

For \(\gamma=T+\delta\) with \(\delta=o(T)\),

\[
c_T(\gamma)
=\frac{\delta}{T}
+O\!\left(\frac{\delta^2}{T^2}
+\frac{|\delta|}{T^3}\right)
\tag{61.2}
\]

Define

\[
\eta_T(\gamma)
=\frac{T\log(T/2\pi)}{2\pi}c_T(\gamma)
\tag{61.3}
\]

For bounded offsets \(\gamma-T\) and \(\gamma'-T\) (in particular at the
local spacing scale \(O(1/\log T)\)), the nearby ordinates satisfy

\[
\eta_T(\gamma)-\eta_T(\gamma')
=\frac{\log(T/2\pi)}{2\pi}(\gamma-\gamma')+o(1)
\tag{61.4}
\]

which is the conventional local spacing normalization.

Both the sine kernel and the Postmaster Markov object admit
divided-difference forms, but no limit from the completed-source Pick
kernel to the sine kernel is proved here. Section 61A proves a limit for
the specified comparison projection of Section 12A. The present
statement remains a coordinate bridge, not a proof of pair correlation.

\begin{figure}
\centering
\includegraphics{figures/cayley_montgomery_unfolding.png}
\caption{At large \(T\), the local Cayley coordinate agrees to first
order with the conventional zero-spacing coordinate.}
\end{figure}

\section{61A. A comparison sine kernel at the ordinate
scale}\label{a.-a-comparison-sine-kernel-at-the-ordinate-scale}

The sine kernel of Section 12A can be matched to a specified high-height
coordinate. This is a theorem about that comparison projection, not the
missing completed-source Pick-kernel limit in Section 61.

With Fourier convention
\(\widehat f(t)=\int_{\mathbb R}f(u)e^{-2\pi itu}du\), interval
integration, convolution and Plancherel give respectively

\[
\operatorname{sinc}_\pi(u)=\int_{-1/2}^{1/2}e^{2\pi itu}dt,\qquad
\widehat{\operatorname{sinc}_\pi^2}(t)=(1-|t|)_+,\qquad
\int_{\mathbb R}\operatorname{sinc}_\pi(u)^2du=1.
\tag{61A.1}
\]

The continuous tent measures interval overlap. Its finite counterpart
(12A.4) counts pair differences exactly. Montgomery's original paper
{[}79{]} uses the corresponding distinct-point density
\(1-\operatorname{sinc}_\pi(u)^2\) in its pair-correlation conjecture.
Its theorem assumes RH and treats a restricted Fourier range. The
density is not the nearest-neighbor gap distribution; retaining
self-pairs adds a delta mass at zero. No conjectured pair-correlation
law is used as a proved arithmetic estimate here.

On the critical-line chart \(\rho=1/2+i\gamma\), \(\gamma>0\),

\[
 y_\gamma=\frac1{\gamma^2+1/4}
 =4\sin^2\frac{\theta(\gamma)}2,\qquad
 \theta(\gamma)=2\arctan\frac1{2\gamma},\qquad
 \theta'(\gamma)=-\frac1{\gamma^2+1/4}.
\tag{61A.2}
\]

These are algebraic identities. Using a real ordinate in this chart is
not the same as using the actual folded coordinate of an off-line zero.

\textbf{Theorem 61A.1 (joint bandwidth and height limit).} Put, for
\(T>2\pi\),

\[
 L_T=\log\frac{T}{2\pi},\qquad
 N_T=\left\lfloor\frac{(T^2+1/4)L_T}{2}\right\rfloor,\qquad
 \gamma_T(u)=T+\frac{2\pi u}{L_T}.
\tag{61A.3}
\]

For each fixed \(U\) and sufficiently large \(T\), uniformly for
\(|u|,|v|\le U\),

\[
\boxed{\frac\pi{N_T}\Pi_{N_T}
 \bigl(\theta(\gamma_T(u)),\theta(\gamma_T(v))\bigr)
 =\operatorname{sinc}_\pi(u-v)+O_U((TL_T)^{-1}).}
\tag{61A.4}
\]

\textbf{Proof.} Taylor's theorem applied to (61A.2), including the floor
error, gives
\(N_T[\theta(\gamma_T(u))-\theta(T)]=-\pi u+O_U((TL_T)^{-1})\). Here
\(\theta''(\gamma)=2\gamma/(\gamma^2+1/4)^2=O(T^{-3})\) on the required
interval. In the difference term of (12A.9), writing
\(z=N_T(\theta(\gamma_T(u))-\theta(\gamma_T(v)))\) gives
\(\sin z/[2N_T\sin(z/(2N_T))]\); the removable quotient converges
uniformly to \(\sin z/z\). In the reflected term the denominator is
asymptotic to \(\sin\theta(T)\sim1/T\), uniformly in \(u,v\). Its
normalized magnitude is \(O(T/N_T)=O((TL_T)^{-1})\). Together these
estimates prove (61A.4). \(\square\)

The relevant projection dimension is therefore
\(N\sim(T^2+1/4)\log(T/(2\pi))/2\). Fixed degree cannot resolve
arbitrarily high microscopic spacing in this coordinate. This \(N\) is
not automatically a Widder rung index or a count of zeta zeros.

\textbf{The source comparison retains a different functional.} In
\(\ell=r+s+1\), \(d=r-s\) coordinates, Section 155A gives

\[
 K_{r,s}=\lambda_\ell-\lambda_{|d|},\qquad
 y\mathcal P_r(y)\mathcal P_s(y)
 =2\cos(d\theta)-2\cos(\ell\theta).
\tag{61A.5}
\]

The same product-to-sum operation creates (12A.9), but the summation
data are different. The actual functional is
\(\mathcal L[f]=\sum_{[\rho]}m_\rho y_\rho f(y_\rho)\), and its identity
\(\lambda_{2m+1}=\mathcal L[\mathcal P_m^2]\) uses holomorphic squares,
not absolute squares. Neither \(\omega\) nor \(\nu_C\) is identified
with this functional. The all-rank condition does not occupy only
\(d=0\); it constrains the mixed entries. Section 155C supplies an
explicit positive reserve and retains its entire adjustment in the
arithmetic defect.

\textbf{Why a limiting spacing law is insufficient by itself.} Consider
a locally finite multiset of real ordinates and a bounded compactly
supported pair test of \((\gamma-\gamma')\log T\), divided by a
normalization of order \(T\log T\). Adding or removing finitely many
fixed ordinates changes its numerator by \(O(1)\): each fixed altered
point meets only finitely many unchanged points in the shrinking support
window. Its normalized limit is unchanged. Thus an abstract finite
off-line quartet can coexist with the same limiting ordinate statistic.
This is not a counterexample about zeta. It shows why the actual-source
inequality \(\mathcal E_{2m+1}\le\mathcal A_{2m+1}\) or the all-rank
defect bound must still be proved. The finite negative control of
Section 104A remains valid after this sine reformulation.

\Needspace{11\baselineskip}

\chapter{Technical chapter X. Prime wave packets and the cancellation
problem}\label{technical-chapter-x.-prime-wave-packets-and-the-cancellation-problem}

The moving shell leads to prime oscillation. Exact Bessel coefficients
improve the diagonal model; remaining correlation and uniformity
assumptions are stated at the end.

\section{45. Gate Z and the Lindelof
firewall}\label{gate-z-and-the-lindelof-firewall}

The finite octagon identity

\[
\omega=e^{2\pi i7/8}
\qquad
\sum_{j=0}^{7}\omega^j=0
\]

is exact. Its transfer to \(n^{-it}\) is not exact because

\[
\log(N+h)=\log N+\frac hN-\frac{h^2}{2N^2}+\cdots
\]

Any exponent-pair submission must retain uniformity in \(t,M,N\),
endpoints, integer-derivative danger blocks, and stationary phases. No
such new pair is proved here.

Lindelöf controls growth on the critical line. No argument here turns
that growth bound into exclusion of sparse off-line zeros. Gate Z is
useful analytic infrastructure, not the shortest zero-exclusion route.

\section{105. From the grid to the moving critical
shell}\label{from-the-grid-to-the-moving-critical-shell}

For each fixed \(x>0\) and fixed \(n\), PP157 obtains the RH-independent
asymptotic

\[
\boxed{
\begin{aligned}
G_{n,k}(x)
={}&
\frac{(1/2)_n}{8\sqrt{\pi x}}\,
k^{-n-1/2}
\left[
\log\frac{kx}{\pi^2}
+\gamma_E-(2H_{2n}-H_n)
\right] \\
&+O_{n,x}(k^{-n-1}\log k)
\end{aligned}
}
\tag{105.1}
\]

Thus every fixed \(n\) row is eventually positive. This is insufficient
for the full grid because the threshold is not uniform in \(n\).

The verified-height route gives a very large but finite A-CAT rectangle:

\[
G_{n,k}(56)>0
\qquad
(0\le n\le6{,}901{,}269{,}318,\ k\ge0)
\tag{105.2}
\]

Its size is not a logical substitute for all \(n\). The unresolved owner
moves with \(k\).

Put

\[
K=k+1,\qquad N=n+1,\qquad
\lambda=\frac{N}{\sqrt K}
\qquad
\Gamma^2+\frac14=\frac{xK}{N}
\tag{105.3}
\]

In the critical scaling \(\lambda\to\lambda_0\in(0,\infty)\), a zero
\(\rho=\tfrac12+\delta+i(\Gamma+c)\) contributes, after normalization,

\[
\exp\!\left[
\frac{2\lambda_0^2}{x}
\left(\rho-\frac12-i\Gamma\right)^2
\right]
\tag{105.4}
\]

Horizontal pairing gives

\[
2e^{-2\lambda_0^2c^2/x+2\lambda_0^2\delta^2/x}
\cos\frac{4\lambda_0^2\delta c}{x}
\tag{105.5}
\]

The flat-density Gaussian integral is independent of \(\delta\):

\[
\int_{\mathbb R}
e^{-2\lambda^2c^2/x+2\lambda^2\delta^2/x}
\cos\frac{4\lambda^2\delta c}{x}\,dc
=\sqrt{\frac{\pi x}{2\lambda^2}}
\tag{105.6}
\]

This is why isolated off-line displacement does not control the moving
shell; the prime oscillation does.

The even Gaussian test

\[
g_{\Gamma,\lambda}(u)
=e^{-2\lambda^2(u-\Gamma)^2/x}
+e^{-2\lambda^2(u+\Gamma)^2/x}
\tag{105.7}
\]

has Fourier transform

\[
\widehat g(y)
=\frac{\sqrt{2\pi x}}{\lambda}
e^{-\pi^2xy^2/(2\lambda^2)}
\cos(2\pi\Gamma y)
\tag{105.8}
\]

The exact explicit formula therefore isolates the prime cosine statistic

\[
P_x(\lambda,\Gamma)
=\sum_{m\ge2}
\frac{\Lambda(m)}{\sqrt m}
e^{-x(\log m)^2/(8\lambda^2)}
\cos(\Gamma\log m)
\tag{105.9}
\]

PP160 proves unconditionally that, for every fixed \(\varepsilon>0\),

\[
\lambda^2
\le2x(1-\varepsilon)\log\log\Gamma
\quad\Longrightarrow\quad
G_{n,k}(x)>0
\tag{105.10}
\]

for all sufficiently large admissible pairs. A convenient \(k\)-form is

\[
n<(\sqrt{2x}-o(1))\sqrt{k\log\log k}
\tag{105.11}
\]

The transition beyond this region is the first infinite part of the grid
not already covered by fixed-row asymptotics or source-side absolute
estimates.

\section{106. Fourier-saddle collision and the first-return
exponent}\label{fourier-saddle-collision-and-the-first-return-exponent}

Let

\[
A_x(\lambda)
=\sum_{m\ge2}
\frac{\Lambda(m)}{\sqrt m}
e^{-x(\log m)^2/(8\lambda^2)}
\tag{106.1}
\]

For \(\eta>0\), define the first positive near-return

\[
R_x^+(\lambda,\eta)
=\inf\left\{
\Gamma\ge1:
P_x(\lambda,\Gamma)\ge(1-\eta)A_x(\lambda)
\right\}
\tag{106.2}
\]

With \(\eta_\lambda=c_x/\lambda^2\), the return exponent is

\[
\boxed{
\Theta_x
=\liminf_{\lambda\to\infty}
\frac{x}{\lambda^2}
\log\log R_x^+(\lambda,\eta_\lambda)
}
\tag{106.3}
\]

For each fixed \(\lambda\), unique factorization makes the prime
logarithms rationally independent, so Kronecker approximation gives

\[
\limsup_{\Gamma\to\infty}
\frac{P_x(\lambda,\Gamma)}{A_x(\lambda)}=1
\tag{106.4}
\]

Fixed-dimensional recurrence is therefore unavoidable. It is harmless at
fixed \(\lambda\), where \(A_x(\lambda)\) is fixed and the Archimedean
main grows like \(\log\Gamma\). The moving dimension is the owner.

Truncating the active primes at

\[
\log p
\le2s+6\sqrt{s\log\lambda}
\qquad s=\lambda^2/x
\tag{106.5}
\]

leaves \(d_\lambda=\exp(2s+o(s))\) phases. Simultaneous Dirichlet
approximation gives

\[
\log\log R_x^+
\le2s+o(s)
\tag{106.6}
\]

On every bounded \(\Gamma\) interval the normalized cosine statistic
tends to zero, so \(R_x^+\to\infty\). Hence the sharp proved interval is
only

\[
\boxed{0\le\Theta_x\le2}
\tag{106.7}
\]

The Postmaster branch needs

\[
\boxed{\Theta_x>\frac12}
\tag{106.8}
\]

The analytic reason absolute values stop at this coefficient is visible
in the exact rational test

\[
H(t)
=x^{N-1}
\frac{(t^2+1/4)^K}
{(x+t^2+1/4)^{K+N}}
\tag{106.9}
\]

At its positive saddle,

\[
f''(\Gamma)
=-\frac{4\lambda^2}{x}
+O_x\!\left(\frac{\lambda^2}{\Gamma^2}\right)
\tag{106.10}
\]

For Fourier frequency \(L\), the complex shift is

\[
\sigma_F(L)=\frac{xL}{4\lambda^2}
\tag{106.11}
\]

The phase-blind prime-density exponent

\[
\Phi(L)=\frac L2-\frac{xL^2}{8\lambda^2}
\tag{106.12}
\]

is maximized at

\[
L_*=\frac{2\lambda^2}{x}
\qquad
\boxed{\sigma_F(L_*)=\frac12}
\tag{106.13}
\]

The prime coefficient \(m^{-1/2}\) and the saddle shift \(m^{-1/2}\)
collide exactly on \(\Re s=1\). Absolute convergence can reach the
boundary but cannot supply the missing sign. The surviving task is an
additive wave-packet anti-resonance theorem.

\section{107. The loose diagonal is exactly
self-limiting}\label{the-loose-diagonal-is-exactly-self-limiting}

A loose version of the moment method bounds the diagonal by a loose
expression of the shape \(r!S_2^r\). Write

\[
\log T=\exp(cs+o(s))
\qquad 0<c<1
\tag{107.1}
\]

The largest product length compatible with clean orthogonality is

\[
r_{\max}\sim\frac{e^{cs}}{2s}
\tag{107.2}
\]

while exclusion of one derivative-width near-\(L^1\) return requires

\[
r_{\rm req}\sim\frac{e^{cs}}{(1-c)s}
\tag{107.3}
\]

Therefore

\[
\boxed{
\frac{r_{\rm req}}{r_{\max}}
\longrightarrow\frac{2}{1-c}
}
\tag{107.4}
\]

Equivalently, with \(r=\beta\log T/(2s)\), the method needs

\[
\beta>\frac{2}{1-c}
\tag{107.5}
\]

The same central peak used in the upper bound forces

\[
\beta\le\frac{2}{1-c}
\tag{107.6}
\]

\begin{quote}
\textbf{Theorem 107.1 (loose-diagonal no-go).} The diagonal-quality
moment architecture cannot prove a strict return exponent larger than
\(c\) for any \(0<c<1\). Its reach and its failure are the same exponent
equality.
\end{quote}

The endpoint can have separate bookkeeping, but it cannot yield the
strict crossing (106.8). The exact-diagonal route uses the interior
saddle relation (108.9).

\section{118. Gaussian prime energy and the visibility
comparison}\label{gaussian-prime-energy-and-the-visibility-comparison}

The following source lemma separates an exact prime-energy asymptotic
from the still-unproved comparison between collar width and finite-rank
reach.

\subsection{Lemma 118.1 - Gaussian prime
energy}\label{lemma-118.1---gaussian-prime-energy}

For fixed \(s>0\), let

\[
a_p=\frac{\log p}{\sqrt p}
\exp\!\left(-\frac{(\log p)^2}{8s}\right)
\]

Then

\[
\boxed{
\mathcal S(s):=\sum_p a_p^2
=\sum_p\frac{(\log p)^2}{p}
\exp\!\left(-\frac{(\log p)^2}{4s}\right)<\infty
}
\tag{118.1}
\]

Moreover, as \(s\to\infty\), the prime number theorem and partial
summation give

\[
\boxed{\mathcal S(s)=2s+o(s)}
\tag{118.2}
\]

For convergence, the Gaussian in \(\log p\) eventually dominates every
fixed negative power of \(p\). For the asymptotic, \(d\pi(t)\) may be
replaced at main order by \(dt/\log t\); with \(u=\log t\) the resulting
main integral is

\[
\int_0^\infty u e^{-u^2/(4s)}\,du=2s
\]

The integral is exactly \(2s\); the discrete prime sum is only
asymptotic to it. Consequently the square-root and tilted peak locations
inferred from this density replacement are asymptotic diagnostics, not
exact finite-\(s\) identities.

The exact aperture formula from Section 48 survives PP175 unchanged. The
proposed comparison with PP133 does not receive theorem grade: PP133
reports an empirical rank threshold whose zero-count scale grows roughly
like \(\gamma\log\gamma\), whereas the Lill half-aperture shrinks like
\(1/\gamma\) for fixed horizontal displacement. Both say that high zeros
are harder to see, but no map from collar width to \(\nu_N\) has been
proved.

\section{108. The exact-diagonal theorem and its closure
geometry}\label{the-exact-diagonal-theorem-and-its-closure-geometry}

Set

\[
D_s(t)=\sum_pa_pp^{it}
\qquad
a_p=\frac{\log p}{\sqrt p}
e^{-(\log p)^2/(8s)}
\tag{108.1}
\]

and expand

\[
D_s(t)^r=\sum_nb_r(n)n^{it}
\tag{108.2}
\]

The exact diagonal is

\[
E_r=\sum_nb_r(n)^2
\tag{108.3}
\]

\begin{quote}
\textbf{Theorem 108.1 (Bessel-product diagonal).} \[
\boxed{
E_r=(r!)^2[z^r]\prod_pI_0(2a_p\sqrt z)
}
\tag{108.4}
\]
\end{quote}

\textbf{Proof.} For \(n=\prod_pp^{k_p}\) with \(\sum_pk_p=r\),

\[
b_r(n)=\frac{r!}{\prod_pk_p!}\prod_pa_p^{k_p}
\]

Since

\[
I_0(2a\sqrt z)
=\sum_{k\ge0}\frac{a^{2k}z^k}{(k!)^2}
\]

coefficient extraction in the Euler product gives exactly the sum of
\(b_r(n)^2/(r!)^2\). \(\square\)

\textbf{Relation to classical pseudomoments.} The quantity \(E_r\) is a
Gaussian-weighted all-prime pseudomoment, within the classical framework
of Conrey--Gamburd and subsequent work {[}30,32,38{]}. The local
identity is

\[
\sum_{k\ge0}\frac{a_p^{2k}z^k}{(k!)^2}
=I_0(2a_p\sqrt z)
\tag{108.4a}
\]

Wahl's prime product uses \(J_0\) factors {[}31, equation (1.4){]},
related by \(I_0(u)=J_0(iu)\). Theorem 108.1 is an exact coefficient
dictionary; the weight and growing-order remainder distinguish the later
analytic task. No priority claim is made for the Bessel local-factor
mechanism.

The first-order cancellation behind Wahl's renormalizer is explicit:
with \(a_p=p^{-1/2}\) and \(z=-u^2/4\),
\(I_0(2a_p\sqrt z)=J_0(u/\sqrt p)\), and the divergent \(p^{-1}\) term
in \(\log J_0\) is cancelled by \((1-1/p)^{-u^2/4}\).

The packet weight in (108.1) is an all-prime Gaussian weight. For each
fixed \(s\) its square is summable,

\[
\sum_p a_p^2
=\sum_p\frac{(\log p)^2}{p}
e^{-(\log p)^2/(4s)}<\infty
\tag{108.4b}
\]

so the bare local product converges, while its mass becomes concentrated
in a moving logarithmic packet as \(s\to\infty\). This damping
distinguishes the global object from Wahl's undamped product; the match
is at the Bessel local factor, not at the full weighted model.

Growing-order pseudomoment asymptotics occur in the literature; see
Bondarenko--Brevig--Saksman--Seip--Zhao {[}38{]}. The result below
concerns the stated Gaussian-weighted, coupled \((r,s)\) regime. Its
theorem is specific to that weight and uniform remainder; no general
novelty or priority claim for growing-order pseudomoment saddles is
made.

\begin{figure}
\centering
\includegraphics{figures/v26c_bessel_coefficients.png}
\caption{The exact Bessel coefficient dictionary keeps the factorial
square in each local prime multiplicity. Forming the product and
extracting \(z^r\) enforces total multiplicity \(r\); it does not
control the off-diagonal prime contribution.}
\end{figure}

The packet-specific analytic statement begins with the following saddle
estimate and its uniform remainder, not with (108.4).

If

\[
c_r=\frac{\log r}{s}
\]

stays in a compact subinterval of \((0,2)\), the symbolic saddle and
local central-limit argument give

\[
\boxed{
\log E_r
=r\left(c_r-\frac{c_r^2}{4}\right)s
+O(r\log s)
}
\tag{108.5}
\]

Also

\[
\log A_s^{2r}=rs+O(r\log s)
\qquad
A_s=\sum_pa_p
\tag{108.6}
\]

Thus

\[
\log\frac{A_s^{2r}}{E_r}
=r\left(1-\frac{c_r}{2}\right)^2s
+O(r\log s)
\tag{108.7}
\]

At

\[
r=\frac{\beta\log T}{2s}
\qquad
\frac{\log r}{s}\to c
\tag{108.8}
\]

an exact-diagonal translated bound of size \(T^{1+o(1)}E_r\) would rule
out a return whenever

\[
\boxed{
\frac{\beta}{2}\left(1-\frac c2\right)^2>1
\quad\hbox{equivalently}\quad
\beta>\frac{2}{(1-c/2)^2}
}
\tag{108.9}
\]

Solving the equality for \(c\) gives the exact-diagonal closure map

\[
\boxed{
c_{\rm ED}(\beta)
=2\left(1-\sqrt{\frac2\beta}\right)
}
\tag{108.10}
\]

At \(c=1/2\) the threshold is \(32/9\). Two margins are important:

\[
\frac{4}{2}\left(\frac34\right)^2-1=\frac18
\tag{108.11}
\]

and

\[
\frac{11/3}{2}\left(\frac34\right)^2-1=\frac1{32}
\tag{108.12}
\]

At the rational pilot

\[
\boxed{
\beta_0=\frac{11}{3}
\qquad
c_{\rm ED}(\beta_0)
=2\left(1-\sqrt{\frac6{11}}\right)
=0.5229021082\ldots>\frac12
}
\tag{108.13}
\]

This shows exponent room. It does not supply the translated theorem or
its uniformity in \(c\).

\begin{figure}
\centering
\includegraphics{figures/figure55_exact_diagonal_closure.png}
\caption{The exact diagonal permits coefficients above one half at the
rational pilot, but the shaded region is conditional exponent room.}
\end{figure}

The inherited rigorous exact-diagonal frontier remains
\(\beta_{\rm ED}\ge32/17\) in supremum sense. The pilot \(\beta_0=11/3\)
is a target, not an attained return theorem.

\section{109. Tilted diagonal, size owner, and the engineered
kernel}\label{tilted-diagonal-size-owner-and-the-engineered-kernel}

Define

\[
E_r(\theta)=\sum_nn^\theta b_r(n)^2
\qquad
\frac{\log r}{s}\to c\in(0,1)
\tag{109.1}
\]

PP168 proves

\[
\boxed{
\frac{\log E_r(\theta)}{rs}
=
\begin{cases}
(1+\theta)c-c^2/4,&0\le\theta\le c/2\\[1ex]
c+\theta^2,&\theta\ge c/2
\end{cases}
+O\!\left(\frac{\log s}{s}\right)
}
\tag{109.2}
\]

The value and first derivative agree at \(\theta=c/2\), while the second
derivatives are \(0\) and \(2\). The phase is \(C^1\) but not \(C^2\).
Its audited right-tail rate is

\[
\boxed{
I_c(\alpha)=\frac{\alpha^2-c^2}{4}
\qquad \alpha\ge c
}
\tag{109.3}
\]

\begin{figure}
\centering
\includegraphics{figures/figure54_tilted_diagonal_phase.png}
\caption{The two tilted branches meet with one derivative but not two.
The rate formula is asserted only on its audited right-tail domain.}
\end{figure}

At \(c=1/2\), with \(r=\beta\log T/(2s)\),

\[
\frac{E_r[n\asymp T^a]}{E_r}
=T^{-a^2/(2\beta)+\beta/32+o(1)}
\tag{109.4}
\]

is valid for

\[
a\ge\frac\beta4
\tag{109.5}
\]

After the weighted Montgomery--Vaughan penalty, the exponent is

\[
F_\beta(a)
=a-1-\frac{a^2}{2\beta}+\frac{\beta}{32}
=\frac{17\beta}{32}-1-\frac{(a-\beta)^2}{2\beta}
\tag{109.6}
\]

\begin{quote}
\textbf{Theorem 109.1 (Montgomery--Vaughan owner and the internal gap).}
At \(c=1/2\), assume the tilted asymptotic (109.2) in the regime
\(r=\beta\log T/(2s)\). Then \[
\int_0^T|D_s(t)^r|^2\,dt
=T E_r(0)+O(E_r(1))
\tag{109.6a}
\] and \[
\boxed{
\frac{E_r(1)}{T E_r(0)}
=T^{17\beta/32-1+o(1)}}
\tag{109.6b}
\] Consequently the classical mean-value theorem certifies a diagonal
asymptotic when \(\beta<32/17\) and ceases to be main-term quality when
\(\beta>32/17\). Since exact-diagonal return at \(c=1/2\) requires
\(\beta>32/9\), the classical mean-value range and the return range are
disjoint.
\end{quote}

\textbf{Proof.} Apply the Montgomery--Vaughan mean-value theorem first
to finite truncations of (108.2). Its error is bounded by
\(\sum_n n b_r(n)^2=E_r(1)\). The tilted formula gives

\[
\frac{\log E_r(0)}{rs}=\frac7{16}+O\!\left(\frac{\log s}s\right)
\qquad
\frac{\log E_r(1)}{rs}=\frac32+O\!\left(\frac{\log s}s\right)
\]

Thus

\[
\log\frac{E_r(1)}{E_r(0)}
=\frac{17}{16}rs+O(r\log s)
=\left(\frac{17\beta}{32}+o(1)\right)\log T
\]

which proves (109.6b). The Gaussian weight makes \(E_r(1)<\infty\) for
fixed \(s\), so the finite truncations converge in the required weighted
\(\ell^2\) norms and (109.6a) follows by passage to the limit. Finally
\(32/17<32/9\). \(\square\)

This is a no-overlap theorem for the \textbf{classical mean-value
estimate}, not a lower bound on the true off-diagonal error. When
\(\beta>32/17\), the theorem's allowance may dominate even if the actual
signed correlation is smaller. At the rational pilot,

\[
\frac{E_r(1)}{T E_r(0)}
=T^{91/96+o(1)}
\qquad \beta=\frac{11}{3}
\tag{109.6c}
\]

This quantifies why a correlation theorem, rather than another use of
the unmodified mean-value bound, is required.

It has one maximizer:

\[
\boxed{a=\beta,\qquad X=T^{\beta+o(1)}}
\tag{109.7}
\]

The active additive shifts satisfy

\[
\boxed{
L=\frac XT=X^{1-1/\beta}
}
\tag{109.8}
\]

The missing relative saving is independent of \(\beta\) in \(X\) units:

\[
\boxed{
T^{-\beta/4}=X^{-1/4}
}
\tag{109.9}
\]

\begin{figure}
\centering
\includegraphics{figures/figure56_size_owner.png}
\caption{The weighted exponent is a strict parabola whose unique owner
is the shell \(X=T^\beta\).}
\end{figure}

On this tilted owner a typical prime factor obeys

\[
\log p=2s+O_{\mathbb P}(\sqrt s)
\qquad
p=(\log X)^{4+o(1)}
\tag{109.10}
\]

The typical factor scale is not a hard whole-product cutoff. The extreme
factor for \(1-o(1)\) of the tilted mass reaches

\[
\log P^+(n)
=(2+\sqrt2+o_{\mathbb P}(1))s
\tag{109.11}
\]

or polylogarithmic exponent \(4+2\sqrt2\). Repeated factors are
negligible in the independent tilted product owner, but this does not
imply shifted-gcd independence.

For \(W\in C_c^\infty([-1,1])\) use

\[
\widehat W(\xi)=\int_{\mathbb R}W(t)e^{i\xi t}\,dt
\qquad
K_\pi(x)=\cos(\pi x)\widehat W(x)
\tag{109.12}
\]

The Fourier transform of \(K_\pi\) is supported away from the origin, so
all of its derivatives vanish there. Equivalently,

\[
\boxed{
\int_{\mathbb R}x^jK_\pi(x)\,dx=0
\qquad(j=0,1,2,\ldots)
}
\tag{109.13}
\]

Every finite smooth Taylor jet of a local shift density is invisible.
What remains must be nonsmooth arithmetic structure, lattice structure,
or a boundary term.

\section{110. The rational pilot and the exact local
factor}\label{the-rational-pilot-and-the-exact-local-factor}

At \(\beta_0=11/3\),

\[
\boxed{
X=T^{11/3},\qquad
H=T=X^{3/11},\qquad
L=X/T=X^{8/11}
}
\tag{110.1}
\]

The nonlinear-log and boundary collar is

\[
T^{-1}=X^{-3/11}
=X^{-1/4}X^{-1/44}
\tag{110.2}
\]

Thus it has a strict spare power \(X^{-1/44}\) relative to the target.

\begin{figure}
\centering
\includegraphics{figures/figure57_scale_ledger.png}
\caption{The pilot separates the window, shift, modulus, collar, and
target exponents on one common \(X\) scale.}
\end{figure}

The generating-function lift has an exact local \(p\)-adic owner. Put

\[
A_k=\frac{(za_p)^k}{k!}
\qquad
B_k=\frac{(wa_p)^k}{k!}
\qquad
\mu_k=(1-p^{-1})p^{-k}
\tag{110.3}
\]

For \(v=v_p(h)<\infty\),

\[
\boxed{
\begin{aligned}
L_{p,v}(z,w)
={}&
\sum_{k=0}^{v-1}\mu_kA_kB_k\\
&+p^{-v}\left[
(1-2/p)A_vB_v
+\sum_{j\ge1}\mu_j
\bigl(A_{v+j}B_v+A_vB_{v+j}\bigr)
\right]
\end{aligned}
}
\tag{110.4}
\]

\textbf{Proof.} If \(v_p(n)=k<v\), then \(v_p(n+h)=k\), giving the first
sum. On the remaining set write \(n=p^vm\) and \(h=p^vh'\) with
\(p\nmid h'\). The two classes \(m\equiv0\) and \(m\equiv-h'\pmod p\)
are distinct. With Haar probability \(1-2/p\) both are units; otherwise
exactly one has positive valuation \(j\), with measure \(\mu_j\).
Multiplication by \(p^{-v}\) gives (110.4). \(\square\)

In the symmetric squarefree truncation with activity \(\lambda\),

\[
L_{p,1}-L_{p,0}
=\frac{(1-\lambda)(1-(1-2/p)\lambda)}p
\qquad
L_{p,2}-L_{p,1}=\frac{\lambda^2}{p^2}
\tag{110.5}
\]

The sign change at \(\lambda=1\) does not improve an absolute
coefficient tail. The correct marginal Mellin activity is

\[
\boxed{\lambda_p=za_pp^{-\sigma}}
\tag{110.6}
\]

At the pilot saddle,

\[
\sigma_*=o(1)
\qquad
z_*=\frac{11}{3s}[1+o(1)]
\tag{110.7}
\]

and on the owner-prime shell,

\[
\lambda_p
=e^{-3s/2+o(s)}
=(\log X)^{-3+o(1)}
\longrightarrow0
\tag{110.8}
\]

The Mellin factor in (110.6) is essential to this small-activity regime.

For reciprocal smooth-modulus counting, the modeled tail exponent is
only \(1/4\) at the typical factor scale and \(1/(4+2\sqrt2)\) at the
whole-mass cutoff, while the pilot needs

\[
\delta>\frac{11}{32}
\tag{110.9}
\]

Because \(\lambda_p\to0\), the first local coefficient is
divisor-normalized like \(1/p\); a bounded sign or weight improvement is
only subpower over \(r\asymp\log X/\log\log X\). The smoothness-counting
subroute is therefore closed. This is not a no-go for every possible
signed extracted coefficient theorem.

The first missing global line inside this factorization route is

\[
R_{r,X}(h)
=M_{r,X}\mathfrak S_{r,s}(h)+\mathcal E_{r,X}(h)
\tag{110.10}
\]

with a kernel-tested error at relative scale \(X^{-1/4+o(1)}\) after
both growing-order coefficient extractions.

\section{111. The block owner and the total/off-diagonal Poisson
decomposition}\label{the-block-owner-and-the-totaloff-diagonal-poisson-decomposition}

Let

\[
f_X(n)=b_r(n)\phi(n/X)
\qquad
P_X(t)=\sum_nf_X(n)n^{it}
\tag{111.1}
\]

and define

\[
B_X=\sum_nf_X(n)
\qquad
E_X=\sum_nf_X(n)^2
\qquad
M_X=\frac{B_X^2}{X}
\tag{111.2}
\]

Expanding the smoothed moment gives

\[
\mathcal M_X(U,H)
=HK_\pi(0)E_X
+H\sum_{h\ne0}R_X(h)K_{\pi,n}(h)
\tag{111.3}
\]

where

\[
R_X(h)=\sum_nf_X(n)f_X(n+h)
\tag{111.4}
\]

and \(K_{\pi,n}(h)\) retains the exact logarithm \(H\log(1+h/n)\). The
diagonal is the block quantity \(E_X\), not the global \(E_r\).

The recorded exponent calculations give

\[
E_X=A_s^{2r}X^{-3/4+o(1)}
\qquad
A_s^{2r}=X^{1/2+o(1)}
\tag{111.5}
\]

No \(\ell^1\) localization at \(X\) has been proved. Positivity gives
the sufficient inequality

\[
B_X\le\sum_nb_r(n)=A_s^r
\tag{111.6}
\]

Therefore

\[
\boxed{
\frac{E_X}{M_X}
=\frac{E_XX}{B_X^2}
\ge
\frac{E_XX}{A_s^{2r}}
=X^{1/4+o(1)}
}
\tag{111.7}
\]

Equality is conditional on the separate localization
\(B_X=A_s^rX^{o(1)}\).

For the linearized divisor lattice put \(u=q/L\). Poisson summation
yields

\[
\boxed{
A^{\rm tot}(u)
=\frac{\pi}{u}
\sum_{k\in\mathbb Z}
\left[
W(2\pi k/u-\pi)+W(2\pi k/u+\pi)
\right]
}
\tag{111.8}
\]

The off-diagonal action is

\[
\boxed{
A^\times(u)=A^{\rm tot}(u)-K_\pi(0)
}
\tag{111.9}
\]

Let

\[
u_0=\frac{2\pi}{\pi+1}
\qquad
u_1=\frac{2\pi}{\pi-1}
\tag{111.10}
\]

For \(0<u<u_0\),

\[
\boxed{
A^{\rm tot}(u)=0
\qquad
A^\times(u)=-K_\pi(0)
}
\tag{111.11}
\]

Thus the engineered hole annihilates the total lattice action. It does
not annihilate the \(h\ne0\) correlation by itself.

\begin{figure}
\centering
\includegraphics{figures/figure58_poisson_total_offdiag.png}
\caption{For a representative smooth nonnegative bump, the exact
low-modulus identities and the distinction between total and
off-diagonal actions are visible directly.}
\end{figure}

For \(W\ge0\), \(A^{\rm tot}(u)\ge0\), while as \(u\to\infty\),

\[
A^{\rm tot}(u)=K_\pi(0)+O_N(u^{-N})
\qquad
A^\times(u)=O_N(u^{-N})
\tag{111.12}
\]

The high-modulus diagonal channel can therefore reconstruct the explicit
diagonal in total action while becoming invisible off diagonal.

\section{112. The matched high-modulus channel and the large-sieve
no-go}\label{the-matched-high-modulus-channel-and-the-large-sieve-no-go}

In the squarefree local specialization, the matched joint monomial has

\[
[\alpha\gamma](L_{p,1}-L_{p,0})
=\frac{p-2}{p^2}
\qquad
[\alpha\gamma](L_{p,\infty}-L_{p,1})
=\frac1{p^2}
\tag{112.1}
\]

Their sum is the positive diagonal coefficient \((p-1)/p^2\). A matched
prime is assigned modulus exponent one with relative probability
\(1-1/(p-1)\) and exponent two with relative probability \(1/(p-1)\).
Because

\[
p=(\log X)^{4+o(1)}
\qquad
r\sim\frac{\log X}{4\log\log X}
\qquad
\sum_{j\le r}\frac1{p_j-1}=o(1)
\tag{112.2}
\]

the positive matched squarefree local channel lies at

\[
\boxed{
q=\operatorname{rad}(n)=n=X^{1+o(1)}
\qquad
\frac qL=X^{3/11+o(1)}
}
\tag{112.3}
\]

This is A for the local matched channel and only candidate A for a
global extracted singular series. Other signed channels and their total
variation are not controlled.

For prime \(p\asymp L\), let

\[
F_p(a)=\sum_{n\equiv a\pmod p}f_X(n)
\qquad
B_{X,p}^*=\sum_{(n,p)=1}f_X(n)
\tag{112.4}
\]

Character orthogonality gives exactly

\[
\boxed{
\sum_{a\in(\mathbb Z/p\mathbb Z)^\times}
\left|
F_p(a)-\frac{B_{X,p}^*}{p-1}
\right|^2
=\frac1{p-1}
\sum_{\chi\ne\chi_0}
\left|\sum_nf_X(n)\chi(n)\right|^2
}
\tag{112.5}
\]

The classical multiplicative large sieve yields

\[
\sum_{p\asymp L}\operatorname{Var}_p(f_X)
\ll X^{o(1)}LE_X
\tag{112.6}
\]

The unsigned short-shift owner is

\[
\mathcal U_X=L\frac{B_X^2}{X}=LM_X
\tag{112.7}
\]

The desired absolute allowance is

\[
X^{-1/4}\mathcal U_X
\tag{112.8}
\]

Using inequality (111.7),

\[
\frac{LE_X}{X^{-1/4}\mathcal U_X}
=X^{1/4}\frac{E_X}{M_X}
\ge X^{1/2+o(1)}
\tag{112.9}
\]

\begin{quote}
\textbf{Theorem 112.1 (large-sieve normalization no-go).} The classical
blockwise character large sieve misses the kernel-tested short-shift
target by at least \(X^{1/2+o(1)}\). Equality in the exponent requires
the unproved \(\ell^1\) block localization; failure of that localization
only worsens the miss.
\end{quote}

The additive identity itself is exact. With

\[
F_X(\alpha)=\sum_nf_X(n)e(n\alpha)
\qquad e(t)=e^{2\pi it}
\tag{112.10}
\]

one has

\[
\boxed{
\sum_hK_\pi(h/L)R_X(h)
=\int_0^1|F_X(\alpha)|^2
\sum_hK_\pi(h/L)e(-h\alpha)\,d\alpha
}
\tag{112.11}
\]

For the Fourier convention

\[
\mathcal FK(\xi)=\int_{\mathbb R}K(x)e^{-2\pi ix\xi}\,dx
\tag{112.12}
\]

\[
\mathcal FK_\pi(\xi)
=\pi\left[
W(2\pi\xi-\pi)+W(2\pi\xi+\pi)
\right]
\tag{112.13}
\]

Poisson summation therefore selects, modulo one, two windows centered at

\[
\alpha=\pm\frac1{2L}
\tag{112.14}
\]

with half-width \(1/(2\pi L)\). Controlling the energy in these windows
after the correct diagonal/off-diagonal subtraction is the direct
surviving branch owner.

\section{129. Nonlinear logarithmic spectral-hole
persistence}\label{nonlinear-logarithmic-spectral-hole-persistence}

The exact short-shift kernel is evaluated at

\[
H\log(1+h/n)
\]

whereas the exact Poisson hole in Section 111 uses the linearized
coordinate \(Hh/n\). PP178 determines what finite-order nonlinearity
does inside the factorized total-lattice model.

Let

\[
\Phi_{H,v}(x)=H\log\left(1+\frac{x}{Hv}\right)
\qquad 1\le v\le2
\tag{129.1}
\]

where \(v\) is the smooth dyadic \(n/X\) variable. Let \(\rho\) be
smooth and supported in \([1,2]\), and let \(\chi\) equal one near zero
and be smoothly supported where \(1+x/(Hv)\) stays uniformly positive.
The factorized nonlinear total-lattice action is

\[
\mathcal A_H^{\rm tot}(u)
=\sum_{m\in\mathbb Z}\int_1^2
\rho(v)\chi(mu/H)K_\pi\!\left(\Phi_{H,v}(mu)\right)\,dv
\tag{129.1a}
\]

This is not the coefficient-weighted correlation \(R_X(h)\). For this
section use the angular-frequency transform
\(\widehat K_{\rm ang}(\omega)=\int_{\mathbb R}K(x)e^{-i\omega x}dx\),
so \(\widehat K_{\rm ang}(\omega)=\mathcal FK(\omega/(2\pi))\) under
(112.12). If \(K=K_\pi\) is defined by (109.12), then

\[
\operatorname{supp}\widehat K_{\rm ang}
\subset[-\pi-1,-\pi+1]\cup[\pi-1,\pi+1]
\tag{129.2}
\]

Writing \(y=x/v\), the cutoff composition through order three is

\[
\chi(x/H)K(\Phi_{H,v}(x))
=\sum_{j=0}^{3}H^{-j}\mathcal D_jK(y)
+\mathcal R_{4,H,v}(x)
\tag{129.3}
\]

with

\[
\begin{aligned}
\mathcal D_0K(y)&=K(y)\\
\mathcal D_1K(y)&=-\frac{y^2}{2}K'(y)\\
\mathcal D_2K(y)&=\frac{y^3}{3}K'(y)
                  +\frac{y^4}{8}K''(y)\\
\mathcal D_3K(y)&=-\frac{y^4}{4}K'(y)
                  -\frac{y^5}{6}K''(y)
                  -\frac{y^6}{48}K'''(y)
\end{aligned}
\tag{129.4}
\]

Every higher coefficient has the same form: a finite sum of polynomial
multiples of derivatives of \(K\). Polynomial multiplication and
differentiation preserve Fourier support. Scaling by \(v\) replaces the
bands by their contractions by \(1/v\); smooth integration over
\(1\le v\le2\) leaves the outer edge at \(\pi+1\) and does not fill the
central hole.

Here exact support belongs to the full jets \(\mathcal D_jK(x/v)\), not
to their products with \(\chi(x/H)\). Replacing a cut-off jet by its
full jet costs \(O_{A,N}(H^{-A}(1+|x|)^{-N})\) for every \(A,N\),
because the difference is supported where \(|x|\asymp H\) and every jet
is Schwartz. This cutoff discrepancy is absorbed into
\(\mathcal R_{J,H,v}\) before Poisson summation.

Consequently, with

\[
u_0=\frac{2\pi}{\pi+1}
\]

Poisson summation annihilates every finite smeared jet for \(0<u<u_0\):

\[
\boxed{
\sum_{m\in\mathbb Z}
\int_1^2\rho(v)\mathcal D_jK(mu/v)\,dv=0
\qquad(j\ge0)
}
\tag{129.5}
\]

A smooth dyadic cutoff therefore changes only the outer Schwartz tail.
For every fixed \(J,N\), after choosing the cutoff-tail order at least
\(J\), the combined composition-and-cutoff remainder satisfies

\[
|\mathcal R_{J,H,v}(x)|
\le C_{J,N}H^{-J}(1+|x|)^{-N}
\tag{129.6}
\]

uniformly for \(v\in[1,2]\). Since the lattice sum of the right side is
\(O_N(H^{-J}(1+u^{-1}))\), PP178 gives

\begin{quote}
\textbf{Theorem 129.1 (nonlinear persistence of the factorized total
hole).} For \(0<u<u_0\), \[
\boxed{
\mathcal A_H^{\rm tot}(u)
=O_J\!\left(H^{-J}(1+u^{-1})\right)
}
\tag{129.7}
\] and, after removing the zero lattice point, \[
\boxed{
\mathcal A_H^\times(u)
=-K_\pi(0)\int_1^2\rho(v)\,dv
+O_J\!\left(H^{-J}(1+u^{-1})\right)
}
\tag{129.8}
\]
\end{quote}

At the pilot, \(u\ge L^{-1}\), \(H=X^{3/11}\), and \(L=X^{8/11}\).
Keeping four jets gives

\[
H^{-4}(1+u^{-1})
\ll H^{-4}L
=X^{-4/11}
=X^{-1/4}X^{-5/44}
\tag{129.9}
\]

Thus truncation at \(J=4\) is the first one whose remainder clears the
required relative power, with spare \(X^{-5/44}\).

\begin{figure}
\centering
\includegraphics{figures/figure67_nonlinear_spectral_hole.png}
\caption{Every nonlinear logarithmic jet preserves the two Fourier
bands; at the rational pilot, the fourth remainder is the first to cross
below the target scale.}
\end{figure}

The scope is essential. Theorem 129.1 applies after factorization to the
unweighted total lattice, with smooth \(v\)-smearing and a harmless
dyadic cutoff. It does not replace \(R_X(qm)\) by a constant, construct
\(c_X(q)\), or control the growing-order factorization error. It proves
that finite-order kernel geometry is not the missing power. The
surviving obstruction is strictly arithmetic transport and extraction.

\section{113. The closure contract}\label{the-closure-contract}

The branch can advance only after the following conditions are met.

\begin{longtable}[]{@{}
  >{\raggedright\arraybackslash}p{(\columnwidth - 6\tabcolsep) * \real{0.2143}}
  >{\raggedleft\arraybackslash}p{(\columnwidth - 6\tabcolsep) * \real{0.2857}}
  >{\raggedleft\arraybackslash}p{(\columnwidth - 6\tabcolsep) * \real{0.2857}}
  >{\raggedright\arraybackslash}p{(\columnwidth - 6\tabcolsep) * \real{0.2143}}@{}}
\toprule\noalign{}
\begin{minipage}[b]{\linewidth}\raggedright
Arrow
\end{minipage} & \begin{minipage}[b]{\linewidth}\raggedleft
Direction
\end{minipage} & \begin{minipage}[b]{\linewidth}\raggedleft
Status
\end{minipage} & \begin{minipage}[b]{\linewidth}\raggedright
Missing uniformity or theorem
\end{minipage} \\
\midrule\noalign{}
\endhead
\bottomrule\noalign{}
\endlastfoot
all-rank folded positivity \(\leftrightarrow-\Theta_\xi'\) CM &
equivalence & exact & none \\
\(-\Theta_\xi'\) CM \(\leftrightarrow\) one-point Hausdorff grid &
equivalence & exact & none \\
full grid \(\Rightarrow\) every moving shell & necessary & exact &
converse not automatic \\
PP157--PP160 source region \(\Rightarrow G_{n,k}>0\) there & sufficient
& proved & transition shell remains \\
\(\Theta_x>1/2\Rightarrow\) transition-shell positivity & intended
sufficient & partial & one uniform theorem must be mounted \\
additive-window saving \(\Rightarrow\) return exclusion & conditional
sufficient & branch-open & extraction, nonlinear kernel, and coefficient
uniformity \\
return exclusion only at \(c=1/2\Rightarrow\Theta_x>1/2\) & false
without more & open & needs a \(c\)-neighborhood or stability theorem \\
transition-shell positivity \(\Rightarrow\) full grid & conditional
sufficient & open & all complementary regimes must be joined \\
\end{longtable}

Two owners must therefore remain named.

\textbf{GLOBAL-OWNER.} Prove all-order folded complete monotonicity,
equivalently the complete one-point Hausdorff grid, by a positive source
transport or another all-order source theorem.

\textbf{BRANCH-OWNER.} Construct the extracted finite-\(X\) object, or
bypass it, and prove the kernel-tested selected-window saving

\[
X^{-1/4+o(1)}
\tag{113.1}
\]

with the uniformity needed by the return problem.

The exact-diagonal map makes the needed neighborhood plausible but does
not provide it. If the translated theorem held uniformly for

\[
\frac12\le c\le\frac12+\varepsilon
<c_{\rm ED}(11/3)
\tag{113.2}
\]

then returns through that coefficient interval would be excluded with a
positive power margin, yielding a strict lower return exponent. A
theorem only at \(c=1/2\) gives a boundary statement unless a
quantitative stability lemma transports the \(1/32\) margin to an
interval.

Even after \(\Theta_x>1/2\), one final theorem must state that the
return bound closes every cell in the transition shell and that Sections
105--106 cover the complement. Only then does (104.12) return the
argument to the RH-facing criterion.

\begin{quote}
\textbf{Conditional closure proposition 113.1.} Assume:

\begin{enumerate}
\def\labelenumi{\arabic{enumi}.}
\tightlist
\item
  the selected additive-window estimate holds uniformly on a nontrivial
  \(c\)-interval above \(1/2\) at \(\beta_0=11/3\);
\item
  continuous return exclusion transfers uniformly to the integer
  Postmaster saddles;
\item
  the resulting strict return exponent makes every transition-shell cell
  positive; and
\item
  the proved fixed-row, finite-prefix, and source regions cover all
  remaining cells.
\end{enumerate}

Then the full grid is nonnegative, hence the complete-monotonicity
criterion holds. None of assumptions 1, 3, or the complete coverage in 4
is proved in the present corpus.
\end{quote}

This proposition records the proof contract. It is not an RH theorem.

\subsection{External-normalization
firewall}\label{external-normalization-firewall}

The symbol \(\beta\) in this branch is defined internally by

\[
r=\frac{\beta\log T}{2s}
\tag{113.3}
\]

PP174 supplies a useful \textbf{partial internal dictionary}. The unique
Montgomery--Vaughan error owner in (109.6)--(109.7) satisfies

\[
\boxed{
X=T^{\beta+o(1)}
\qquad
\beta=\frac{\log X}{\log T}+o(1)}
\tag{113.4}
\]

Thus \(\beta\) is the logarithmic length coordinate of the dominant
\(\theta=1\) coefficient shell. This identifies the internal physical
owner and explains the threshold \(32/17\) in Theorem 109.1.

It does \textbf{not} yet identify the full packet with every standard
external length parameter. The distinction is theorem-level:

\begin{enumerate}
\def\labelenumi{\arabic{enumi}.}
\tightlist
\item
  \(D_s(t)^r\) is an infinite Gaussian-weighted Dirichlet series,
  whereas a published large-value theorem is normally stated for a hard
  block \(N<n\le2N\);
\item
  passing from the global series to the block \(f_X\) requires
  quantitative \(\ell^2\) localization, and any \(\ell^1\) normalization
  remains separately open;
\item
  coefficient hypotheses such as \(|b_n|\le1\), value normalization,
  spacing, and the relation between \(N\) and \(T\) must be checked
  theorem by theorem;
\item
  Guth--Maynard's condition \(N\le T^{5/6-\varepsilon}\) describes one
  regime where their new estimate improves earlier bounds. It is not a
  universal definition of a literature parameter \(\alpha\) and does not
  prove that all large-value methods are irrelevant when \(X>T\).
\end{enumerate}

Accordingly, the firewall is narrow but remains necessary. No claim is
made that Guth--Maynard barriers, density-hypothesis exponents, or their
numerical improvements apply directly to (108.9). Conversely, the
internal exact-diagonal geometry is not advertised as an improvement to
those external records. The absence of \(32/9\) from a search is
expected for an internal threshold, but is not itself a novelty theorem.

Guth--Maynard arXiv:2405.20552 must be read at v2 (7 April 2026), and
Chen--Gupta--Li arXiv:2507.08296 at v2 (27 July 2026) proves the \(7/3\)
zero-density exponent for Dirichlet \(L\)-functions by an adaptation of
that method. These are current background inputs, not lemmas in the
Postmaster proof.

\Needspace{11\baselineskip}

\chapter{Technical chapter XI. Finite Weil models and controlled
limits}\label{technical-chapter-xi.-finite-weil-models-and-controlled-limits}

Finite dictionaries and positive cutoff tails have precise scope.
Ground-state and joint-limit hypotheses needed for spectral recovery
remain separate.

\section{19. The actual-Weil ground-state
criterion}\label{the-actual-weil-ground-state-criterion}

Let \(A_\lambda\) represent the full global Weil form, with bottom
eigenvalue \(\epsilon_0\), next eigenvalue \(\epsilon_1\), gap
\(\Delta_\lambda=\epsilon_1-\epsilon_0\), and normalized proxy
\(v_\lambda=k_\lambda/\|k_\lambda\|_2\). Define the closed-form Rayleigh
excess

\[
r_\lambda=q_{A,\lambda}[v_\lambda]-\epsilon_0
\]

If the bottom eigenvalue is simple, then

\[
1-|\langle v_\lambda,\phi_\lambda\rangle|^2
\le\frac{r_\lambda}{\Delta_\lambda}
\]

Because the proxy is even, \(r_\lambda/\Delta_\lambda<1\) forces the
simple ground state to be even. A sufficient scale is

\[
\boxed{
\frac{\|k_\lambda\|_2^2r_\lambda}{\Delta_\lambda}
\le C(1-\chi_2(\lambda))
}
\]

Together with bottom simplicity, proxy convergence, the published
real-zero theorem, and Hurwitz passage, this would imply RH. Bottom
simplicity and the gap-normalized estimate are open.

Every fixed finite Li/Hankel/localizing suite can be passed by a
synthetic off-line quartet. Therefore no fixed finite stopping point is
an RH certificate.

\section{44. Two-channel finite-Weil
sampling}\label{two-channel-finite-weil-sampling}

The coarse one-channel sampling cannot distinguish \(y=\log n\) from
\(L-y\). The two-channel construction samples

\[
T+kh
\qquad\text{and}\qquad
T+\left(k+\frac12\right)h
\]

or uses two polyphase classes on the fine lattice. The de-aliased prime
frequency is

\[
\frac{\pi\log n}{L}
\]

Necessary attainable-weight constraints include

\[
w_j\le\frac12
\qquad
\sum_{j\text{ even}}w_j
=\sum_{j\text{ odd}}w_j
=\frac12
\]

and, at the cubic endpoint,

\[
\sum_j(-1)^jj^rw_j=0
\qquad r=0,1,\ldots,6
\]

The construction has the following exact structure:

\begin{itemize}
\tightlist
\item
  seven-jet ownership \(J_0,\ldots,J_6\);
\item
  individual overlap order seven and operator endpoint order six;
\item
  Newton/central-factorial coordinates;
\item
  finite center/difference Weyl and Wigner identities;
\item
  scoped deep-packet safety with the stated frequency normalization.
\end{itemize}

The open seats are sharp attainable-range characterization,
multiplicative \(A_T/A_0\) comparison, full edge recombination, and
uniform bulk coercivity.

\begin{figure}
\centering
\includegraphics{figures/polyphase_aliasing.png}
\caption{Integer-grid aliasing and its exact half-grid/polyphase
separation.}
\end{figure}

\section{123. Exact finite Guinand--Weil
dictionary}\label{exact-finite-guinandweil-dictionary}

This section restates and checks the fixed-rank theorem of Groskin
arXiv:2607.02828v3. The attribution is load-bearing: no novelty claim is
made for the external theorem. To avoid collision with this volume's
Loewner order, write \(M\) for the external Galerkin band.

Fix \(c>1\) and put

\[
L=\log c
\qquad
\Delta=\frac{L}{2\pi}
\qquad
\rho=\frac{2\pi}{L}
\tag{123.1}
\]

For \(v=(v_0,\ldots,v_M)\in\mathbb R^{M+1}\), use the isometric
symmetric embedding

\[
u_0=v_0
\qquad
u_k=u_{-k}=\frac{v_k}{\sqrt2}\quad(1\le k\le M)
\tag{123.2}
\]

and define

\[
T_v(t)=\sum_{m=-M}^{M}u_me^{2\pi i mt}
\qquad
K_v(\omega)=2\int_0^\omega T_v(t)T_v(\omega-t)\,dt
\tag{123.3}
\]

The associated Fourier weight and test function are

\[
\widehat g_v(\xi)=
\begin{cases}
\pi K_v(1-|\xi|/\Delta),&|\xi|\le\Delta\\
0,&|\xi|>\Delta
\end{cases}
\qquad
g_v(z)=\int_{-\Delta}^{\Delta}\widehat g_v(\xi)e^{2\pi iz\xi}\,d\xi
\tag{123.4}
\]

The finite source calculus can be checked before invoking any explicit
formula. For

\[
\psi_{\alpha,\omega}(x)=\frac{\alpha}{\pi}\sin(2\pi\omega x)
\]

the divided-difference matrix satisfies

\[
\boxed{\langle v,Q_{\alpha,\omega}v\rangle
=\alpha K_v(\omega)}
\tag{123.5}
\]

Indeed, the off-diagonal coefficient is

\[
\alpha\frac{\sin(2\pi\omega m)-\sin(2\pi\omega n)}
{\pi(m-n)}
\]

the diagonal coefficient is \(2\alpha\omega\cos(2\pi\omega m)\), and
expansion of (123.3) reproduces both; the imaginary contributions cancel
under \(u_{-m}=u_m\). By linearity the source-to-form map factors
through the exact \(2M+1\) coordinates

\[
\int\omega\,d\mu
\qquad
\int\sin(2\pi k\omega)\,d\mu
\qquad
\int\omega\cos(2\pi k\omega)\,d\mu
\quad(1\le k\le M)
\tag{123.6}
\]

The entries \((0,0)\), \((k,0)\), and \((k,k)\) recover these
coordinates, proving that the quotient dimension is exact.

\begin{quote}
\textbf{Theorem 123.1 (Groskin's finite dictionary, audited).} For fixed
\(c>1\), fixed \(M\ge0\), and every real even Galerkin vector \(v\), \[
\boxed{
\langle v,Q_\infty^{(c,M)}v\rangle
=\sum_{z:\,\zeta(1/2+iz)=0}g_v(z)
}
\tag{123.7}
\] with multiplicity and absolute convergence.
\end{quote}

The proof is the Guinand--Weil explicit formula applied to (123.4). Its
Fourier support is \([-L/(2\pi),L/(2\pi)]\), so every prime power
\(q>c\) drops out exactly. The prime, pole, and archimedean terms are
then precisely those of the fixed-\((c,M)\) truncated matrix.

Three qualifications are part of the theorem statement in this volume.

\begin{enumerate}
\def\labelenumi{\arabic{enumi}.}
\tightlist
\item
  The map is coefficient vector to test function. It does not realize an
  arbitrary Guinand--Weil test function.
\item
  The theorem is finite in \(M\) and fixed in \(c\).
\item
  The zero sum ranges over the actual complex zero multiset. No reality
  of the parameters \(z\) and hence no RH assumption is used.
\end{enumerate}

\begin{figure}
\centering
\includegraphics{figures/v26c_finite_dictionary.png}
\caption{At finite band \(M\), one linear coordinate and two families of
\(M\) trigonometric coordinates give the exact \(2M+1\)-dimensional
Guinand--Weil dictionary. Finite coordinate recovery is distinct from
positivity of an infinite operator.}
\end{figure}

\section{124. Positive archimedean tail and the cutoff
certificate}\label{positive-archimedean-tail-and-the-cutoff-certificate}

For \(T>\rho M\), put \(a_T=T/\rho\) and, for
\(n\in I_M=\{-M,\ldots,M\}\),

\[
p_T(n)=\frac1{a_T-n}
\qquad
q_T(n)=\frac1{a_T+n}
\tag{124.1}
\]

The algebraic hinge is

\[
\frac{\dfrac{m}{a^2-m^2}-\dfrac{n}{a^2-n^2}}{m-n}
=\frac12\left(
\frac1{(a-m)(a-n)}+\frac1{(a+m)(a+n)}
\right)
\tag{124.2}
\]

with diagonal limit \((a^2+n^2)/(a^2-n^2)^2\). Direct integration of the
true archimedean source gives at each integer node

\[
\mathcal S(T,n,L)
=\frac{2\sin^2(LT/2)}{\rho}\frac{n}{a_T^2-n^2}
\tag{124.3}
\]

and differentiation before node evaluation gives exactly the same
diagonal limit. Thus the divided-difference matrix is formed from the
true source, not merely an integer-node interpolant.

Set

\[
h_+(T)=\Re\psi\!\left(\frac14+\frac{iT}{2}\right)-\log\pi
\tag{124.4}
\]

Its derivative has the positive partial-fraction expansion

\[
h_+'(T)=\frac{T}{2}\sum_{k=0}^{\infty}
\frac{k+1/4}{((k+1/4)^2+(T/2)^2)^2}>0
\tag{124.5}
\]

and \(h_+(7)>0\).

\begin{quote}
\textbf{Theorem 124.1 (Groskin's archimedean tail order, audited).} If
\(T_2>T_1>\max(\rho M,7)\), then \[
\boxed{
Q_{{\rm arch},T_2}-Q_{{\rm arch},T_1}
=\frac1{\pi^2}\int_{T_1}^{T_2}
h_+(T)\frac{\sin^2(LT/2)}{\rho}
\bigl(p_Tp_T^{\mathsf t}+q_Tq_T^{\mathsf t}\bigr)\,dT
}
\tag{124.6}
\] The increment is positive definite, and every ordered minor in the
natural node order is strictly positive.
\end{quote}

Positive definiteness follows without a compactness guess. If a vector
annihilates (124.6), the rational function \(\sum_{n\in I_M}x_n/(a-n)\)
vanishes on an interval. It therefore vanishes identically, and every
residue \(x_n\) is zero. For ordered minors, the pushed measure lies on
two intervals exterior to every node. Andreief's identity expresses a
minor as an integral of two Cauchy determinants. Their signs depend on
the same number of left-interval samples and cancel; positivity on open
sets makes the minor strict.

\begin{figure}
\centering
\includegraphics{figures/figure65_finite_weil_dictionary_tail.png}
\caption{At fixed prime cutoff and Galerkin band, the coefficient vector
has an exact zero-sum image and the omitted archimedean tail is an
ordered integral of rank-two Cauchy Gram atoms.}
\end{figure}

For the total fixed-\((c,M)\) matrix define

\[
B_T=\frac1{\pi^2}\int_T^\infty
h_+(r)\frac{\sin^2(Lr/2)}{\rho}
\bigl(\|p_r\|_2^2+\|q_r\|_2^2\bigr)\,dr
\tag{124.7}
\]

Then

\[
0\prec Q_\infty^{(c,M)}-Q_T^{(c,M)}\preceq B_TI
\tag{124.8}
\]

and hence

\[
\lambda_j(Q_T)<\lambda_j(Q_\infty)
\le\lambda_j(Q_T)+B_T
\tag{124.9}
\]

This produces the exact two-sided decision rule:

\begin{longtable}[]{@{}ll@{}}
\toprule\noalign{}
finite-\(T\) reading & cutoff-free conclusion at fixed \((c,M)\) \\
\midrule\noalign{}
\endhead
\bottomrule\noalign{}
\endlastfoot
\(\lambda_j(Q_T)\ge0\) & \(\lambda_j(Q_\infty)>0\) \\
\(\lambda_j(Q_T)<-B_T\) & \(\lambda_j(Q_\infty)<0\) \\
\(-B_T\le\lambda_j(Q_T)<0\) & inconclusive \\
\end{longtable}

For \(M\ge1\),

\[
B_T\le
\frac{2(2M+1)\rho}{\pi^2}
\left[
\frac{\log T}{T-\rho M}
+\frac1{\rho M}\log\frac{T}{T-\rho M}
\right]
\tag{124.10}
\]

and, at fixed \((c,M)\),

\[
B_T=
\frac{(2M+1)\rho}{\pi^2T}
\left(\log\frac{T}{2\pi}+1\right)(1+o(1))
\tag{124.11}
\]

At \((c,M,T)=(100,200,800)\), direct arithmetic gives

\[
\begin{aligned}
\rho M&=272.875270768\ldots\\
\text{bound (124.10)}&=1.5754527435\ldots\\
\text{leading (124.11)}&=0.4051372811\ldots
\end{aligned}
\tag{124.12}
\]

The sharper interval value reported by Groskin is not promoted because
that ancillary certificate is not mounted here.

\subsection{Three-limit firewall}\label{three-limit-firewall}

The tail theorem removes exactly one regulator:

\[
\boxed{
T\to\infty\text{ at fixed }(c,M)\ \text{is controlled}
\qquad
M\to\infty\ \text{and}\ c\to\infty\ \text{remain open}
}
\tag{124.13}
\]

It does not answer the general one-dimensional Putinar interpolation
question: it constructs one explicit positive Cauchy--Stieltjes tail
rather than a universal moment-extension theorem. Nor does it prove this
volume's all-order source inequality. The external \(M\) is a
Fourier/Galerkin band, not the matrix order in \(\Delta_N^{\mathrm L}\);
moreover the external prime and pole blocks are not made positive by the
archimedean increment.

\section{125. Connes--Suzuki theorem and limit
audit}\label{connessuzuki-theorem-and-limit-audit}

The direct technical sources clarify why spectacular finite spectra do
not yet form an RH proof.

\begin{longtable}[]{@{}
  >{\raggedright\arraybackslash}p{(\columnwidth - 4\tabcolsep) * \real{0.3333}}
  >{\raggedright\arraybackslash}p{(\columnwidth - 4\tabcolsep) * \real{0.3333}}
  >{\raggedright\arraybackslash}p{(\columnwidth - 4\tabcolsep) * \real{0.3333}}@{}}
\toprule\noalign{}
\begin{minipage}[b]{\linewidth}\raggedright
object
\end{minipage} & \begin{minipage}[b]{\linewidth}\raggedright
established result
\end{minipage} & \begin{minipage}[b]{\linewidth}\raggedright
missing identification
\end{minipage} \\
\midrule\noalign{}
\endhead
\bottomrule\noalign{}
\endlastfoot
Connes--van Suijlekom, Prop. 5.10 & finite real-root polynomial under
simple spectrum and a one-dimensional even kernel & applicability
requires the kernel hypotheses \\
Connes--van Suijlekom, Thm. 6.1 & \(\widehat\xi\) has only real zeros
when the lower-bounded operator has a simple isolated even ground state
& those spectral hypotheses must be proved for the target Weil form \\
CCM, \emph{Zeta Spectral Triples} & self-adjoint rank-one perturbations
and high-accuracy finite spectra & simple/even ground state and a
zero-carrying joint limit are explicitly missing \\
Suzuki, Thms. 1.3--1.5 & continuity of \(\lambda_a\); small-\(a\) simple
even ground state; real zeros of \(W(a,\theta;z)\) & \(a\to\infty\)
recovery of zeta ordinates is conjectural \\
Groskin numerical v4 & 16-cutoff computation and reported 307--329-digit
recovery at the largest cell & fixed-\(M\) bias and continuum limit are
not bounded \\
\end{longtable}

CCM Section 8 states the two essential missing steps: prove that the
localized Weil ground state is simple and even, then prove that the
proposed \(k_\lambda\) approximates that eigenvector strongly enough to
carry its zeros to the zeta-zero limit. This is the controlling
statement, not the expository letter's inaccessible continuation.

Suzuki's finite/local results do not require RH. Spectral recovery
remains a conjecture, and Section 7 is explicitly heuristic, begins by
assuming RH, and introduces \(A_\infty\) only as a formal symbol. Real
zeros of the finite characteristic function follow from
self-adjointness; their identification with zeta zeros does not.

The high-precision numerical paper considered here is
arXiv:2605.20224v4. It states that fixed-\(M\) Galerkin bias is not
independently measurable, labels its classifier verdict intermediate,
and makes no proof claim. At \(T=1200\), the \(c=100\) negative block
rearranges rather than disappears. These computations do not bound the
continuum limit.

\begin{figure}
\centering
\includegraphics{figures/figure66_connes_suzuki_limit_firewall.png}
\caption{The finite Weil-form results close the archimedean cutoff at
fixed parameters; simple-even control and the Galerkin/prime-cutoff
limits remain distinct open bridges before any zeta-zero
identification.}
\end{figure}

\Needspace{11\baselineskip}

\chapter{Technical chapter XII. Dilation, cyclic resolvents, and source
carriers}\label{technical-chapter-xii.-dilation-cyclic-resolvents-and-source-carriers}

Finite dilation complements fail a Stieltjes sign. The sharp dilation
frontier and cyclic-resolvent theorem identify the full conditions,
followed by source carrier formulas.

\section{139. Universal finite-dyadic Stieltjes
obstruction}\label{universal-finite-dyadic-stieltjes-obstruction}

Normalize

\[
\mathscr D^Nh(x)=\frac{x}{2^NR}\Phi(2^Ns)
\qquad
B_N=h-\mathscr D^Nh
\]

and

\[
G_N(z)=\frac{B_N(z)}z
=\frac1R\left[\Phi(s)-2^{-N}\Phi(2^Ns)\right]
\tag{139.1}
\]

On the upper lip of the folded cut, write

\[
z_v=-\frac{1+v^2}{4}+i0,\qquad R_+(z_v)=iv
\]

Functional reflection gives the exact boundary identity

\[
\boxed{
\Im G_{N,+}(z_v)
=\frac{2^{-N}}v
\Re\Phi\!\left(2^{N-1}(1+iv)\right)}
\tag{139.2}
\]

The positive \(\Xi\) kernel proves \(\Phi(\sigma)>0\) for every real
\(\sigma>1/2\). Hence the right side of (139.2) is positive for
sufficiently small \(v>0\). A Stieltjes function must map the upper
half-plane into the closed lower half-plane. Consequently

\[
\boxed{
\forall N\ge1:\quad G_N\text{ is not Stieltjes and }B_N
\text{ is not complete Bernstein}}
\tag{139.3}
\]

This does not contradict preservation of the complete-Bernstein cone by
\(\mathscr D\): cone preservation by \(\mathscr D\) does not make
\(I-\mathscr D^N\) positive.

The moving terminal has the exact triangular coordinate

\[
s_N=2^{N-1},\qquad n_N=2^{N-1}-1
\qquad x_N=2T_{2^{N-1}-1}
\tag{139.4}
\]

For six scales, \((s_6,n_6,x_6)=(32,31,992)\) and \(992=2T_{31}\). The
forbidden trace coefficient is

\[
\lim_{v\downarrow0}v\Im G_{N,+}(z_v)
=2^{-N}\Phi(2^{N-1})>0
\tag{139.5}
\]

It tends to zero as \(O(N2^{-N})\) but never vanishes at finite \(N\).
This is the precise finite/infinite boundary layer.

\section{140. Holomorphy, positivity, and the surviving
equivalence}\label{holomorphy-positivity-and-the-surviving-equivalence}

Let \(\Sigma=\mathbb C\setminus(-\infty,0]\) with the principal square
root. For every finite \(N\), the terminal term \(\Phi(2^Ns)\) is
Euler-safe on \(\Sigma\) and cannot cancel an off-cut pole of the first
term. Therefore

\[
\boxed{G_N\text{ holomorphic on }\Sigma\iff\mathrm{RH}}
\tag{140.1}
\]

Equation (140.1) is an exact equivalence, not a resolution. PP190-IH
proves that Stieltjes positivity is strictly stronger than this
holomorphy statement for every finite block and fails unconditionally.

The joint conclusions are:

\begin{enumerate}
\def\labelenumi{\arabic{enumi}.}
\tightlist
\item
  ranks through twelve can be certified at the fixed completed-source
  node;
\item
  no finite dyadic telescoping block can be the global CBF seed;
\item
  Gamma and prime contributions must couple before positivity is taken;
\item
  the required owner is infinite, nonlocal, or exactly
  completion-compensated.
\end{enumerate}

The dilation and cyclic-resolvent theorems below refine the current
proof map in Reading Volume §R24; their exact identities retain the
all-order source obligation.

\section{142. Dilation innovations, reflection, and no-gain
gates}\label{dilation-innovations-reflection-and-no-gain-gates}

The following exact statements organize the sharp dilation frontier. The
proof map in Reading Volume §R24 places their all-order quantifiers
beside the completed source sign. Numerical cutoffs and proof grades are
stated at their theorem owners.

\subsection{142.1 Divisor-square inversion and
innovation}\label{divisor-square-inversion-and-innovation}

The divisor-square carrier and its exact dyadic inversion are

\[
D(s)=\frac{\zeta(s)^4}{\zeta(2s)},\qquad
J(s)=4A(s)-2A(2s)
\]

\[
\boxed{
A(s)=\sum_{k\ge0}2^{-k-2}J(2^ks)}
\tag{142.1}
\]

The coefficient positivity and derivative-level convergence proved in
the source package justify (142.1) in its stated half-plane. The
normalized forward lift preserves the complete-Bernstein cone in its
proved domain, but its inverse difference does not. For the completed
innovation \(k=h-\mathscr D_2h\) the directed endpoint enclosure is

\[
\boxed{-0.011543<k'(0)<-0.011392<0}
\tag{142.2}
\]

The exact reconstruction

\[
\boxed{h=\sum_{j\ge0}\mathscr D_2^jk}
\tag{142.3}
\]

is convergent but signed. The finite Euler-prefix/counterterm and fixed
prime-translation tests likewise fail. Thus a positive forward lift
neither makes its difference positive nor licenses inversion of the
cone.

\subsection{142.2 PP190-T reflection, inertia, and no-gain
gates}\label{pp190-t-reflection-inertia-and-no-gain-gates}

For a nonreal folded shell \(q=|q|e^{i\theta}\), write

\[
\frac{e^{i\theta/2}}{x+q}=A_x+iB_x
\]

The conjugate pair contributes

\[
\boxed{
K_q(x,y)+K_{\bar q}(x,y)
=2|q|\bigl(A_xA_y-B_xB_y\bigr)}
\tag{142.4}
\]

This is the exact reflection/Krein difference of squares. Every nonreal
folded pair has shell-local real signature \((1,1)\); inertia is
invariant under invertible congruence, so no shell-local change of basis
can turn that pair into a positive atom. The negative-channel
contraction formulation is exact, but a Schur lift is a no-gain
tautology unless the larger completed block is independently proved
positive. The zero-tail member has both a certified negative direction
and an independent positive \(b_0\) direction and is therefore genuinely
indefinite.

\subsection{142.3 PP190-IH continuous
strengthening}\label{pp190-ih-continuous-strengthening}

For every finite real \(\lambda>1\), put

\[
G_\lambda^{(0)}(z)
=\frac1{R(z)}
\left[\Phi(s(z))-\lambda^{-1}\Phi(\lambda s(z))\right]
\]

The source package proves

\[
\boxed{
\lambda>1\Longrightarrow
G_\lambda^{(0)}\text{ is not Stieltjes}}
\tag{142.5}
\]

The printed dyadic theorem is the specialization \(\lambda=2^N\).
Independent positive-axis witnesses include
\(G_6^{(0)}(30)>G_6^{(0)}(12)\) and \((G_6^{(0)})'(12)>0\). There is no
contradiction with \(B_6'(0)>0\): the latter is a physical-axis
derivative, whereas (142.5) is forced by the opposite orientation on the
folded upper cut.

\section{143. Exact Stieltjes and complete-Bernstein dilation
frontier}\label{exact-stieltjes-and-complete-bernstein-dilation-frontier}

Let

\[
R(z)=\sqrt{1+4z},\qquad
s(z)=\frac{1+R(z)}2
\]

\[
S_\lambda(z)=\frac{\Phi(\lambda s(z))}{\lambda R(z)}
\qquad
T_\lambda(z)=zS_\lambda(z)
\qquad \lambda>0
\]

and define the horizontal zero frontier

\[
\beta_*=\sup_\rho\Re\rho
\]

\subsection{Theorem 143.1}\label{theorem-143.1}

For every \(\lambda>0\), the following are equivalent:

\begin{enumerate}
\def\labelenumi{\arabic{enumi}.}
\tightlist
\item
  \(S_\lambda\) is a Stieltjes function;
\item
  \(T_\lambda\) is a complete Bernstein function;
\item
  \(S_\lambda\) is holomorphic on the principal slit plane;
\item
  \(\lambda\ge2\beta_*\).
\end{enumerate}

Consequently,

\[
\boxed{
\{\lambda>0:T_\lambda\in\mathrm{CBF}\}
=[2\beta_*,\infty)}
\tag{143.1}
\]

\textbf{Proof.} Write \(\Xi(w)=\xi(1/2+w)\) and retain each reflected
zero pair. The genus-zero paired logarithmic derivative is

\[
\Phi(1/2+w)
=\sum_{\{\pm\alpha\}}
m_\alpha\left(\frac1{w-\alpha}+\frac1{w+\alpha}\right)
\tag{143.2}
\]

Here \(m_\alpha\) is the multiplicity of the paired zeros. The pairing
is essential: each paired summand is \(O(|\alpha|^{-2})\), while the two
separated reciprocal-zero series need not converge absolutely. If
\(\Re\sigma>\beta_*\), both reflected denominators in (143.2) have
positive real part, so \(\Re\Phi(\sigma+it)>0\). For
\(\lambda>2\beta_*\) the principal upper half-plane maps strictly into
that positive half-plane. On the upper folded lip,

\[
\Im S_{\lambda,+}(-r)
=-\frac{
\Re\Phi\!\left(\frac\lambda2
(1+i\sqrt{4r-1})\right)}
{\lambda\sqrt{4r-1}}<0
\tag{143.3}
\]

The keyhole maximum principle then gives the Stieltjes property. Locally
uniform closure gives the endpoint \(\lambda=2\beta_*\).

For necessity, if \(\rho=\beta+i\gamma\) and \(\lambda s(z_\rho)=\rho\),
then

\[
z_\rho=\frac{\rho(\rho-\lambda)}{\lambda^2}
\qquad
\Im z_\rho=\frac{\gamma(2\beta-\lambda)}{\lambda^2}
\tag{143.4}
\]

If \(\beta>\lambda/2\), the pole lies in the principal upper half-plane.
Whenever \(\lambda<2\beta_*\) such a zero exists, contradicting
slit-plane holomorphy. The standard Stieltjes/complete-Bernstein
quotient equivalence finishes the proof. \(\square\)

The source-normalized density is

\[
\boxed{
\omega_\lambda(r)=
\frac{\Re\Phi\!\left(
\frac{\lambda}{2}(1+i\sqrt{4r-1})\right)}
{\pi\lambda\sqrt{4r-1}}}
\tag{143.5}
\]

In the strict range,

\[
S_\lambda(x)
=\int_{1/4}^{\infty}\frac{\omega_\lambda(r)}{x+r}\,dr
\qquad
\omega_\lambda(r)>0
\tag{143.6}
\]

At the frontier, a zero of multiplicity \(m\) at
\(\rho=\lambda/2+i\gamma\) maps to

\[
r_\rho=\frac14+\frac{\gamma^2}{\lambda^2}
\]

and contributes the positive atom

\[
\boxed{\frac{m/\lambda^2}{x+r_\rho}}
\tag{143.7}
\]

The exact RH implication chain is therefore

\[
\boxed{
\begin{aligned}
\mathrm{RH}
&\iff \beta_*=\frac12\\
&\iff \inf\{\lambda:T_\lambda\in\mathrm{CBF}\}=1\\
&\iff T_\lambda\in\mathrm{CBF}\quad(\lambda>1)\\
&\iff T_{\lambda_j}\in\mathrm{CBF}
\quad\text{for some }\lambda_j\downarrow1
\end{aligned}}
\tag{143.8}
\]

Unconditionally, the frontier lies in \([1,2]\), \(T_2\) is complete
Bernstein, and every \(\lambda>2\) lies in the strict range. No sequence
approaching \(1\) has been forced independently of RH.

\begin{figure}
\centering
\includegraphics{figures/v26c_dilation_threshold.png}
\caption{The exact dilation threshold is \(\lambda=2\beta_*\), where
\(\beta_* = \sup_\rho\Re\rho\). A CBF sequence with
\(\lambda_j\downarrow1\) forces \(\beta_*=1/2\). The axis is schematic:
the threshold value itself remains unknown.}
\end{figure}

\section{144. Origin dilation and critical-line-centered
dilation}\label{origin-dilation-and-critical-line-centered-dilation}

Dilation about the origin and dilation about the critical-line center
have different pole geometry. Keep their basepoints explicit.

\subsection{144.1 Origin-dilation
complement}\label{origin-dilation-complement}

Define

\[
T_\lambda^{(0)}(z)
=\frac{z}{\lambda R(z)}\Phi(\lambda s(z))
\qquad
B_\lambda^{(0)}=h-T_\lambda^{(0)}
\tag{144.1}
\]

On \(\Omega=\mathbb C\setminus(-\infty,-1/4]\),

\[
\boxed{
B_\lambda^{(0)}\text{ is holomorphic on }\Omega
\iff
\begin{cases}
\text{false},&0<\lambda<1\\
\text{true},&\lambda=1\\
\mathrm{RH},&\lambda>1
\end{cases}}
\tag{144.2}
\]

For \(0<\lambda<1\), every zero with \(\Re\rho\ge1/2\) produces a
terminal preimage at \(s=\rho/\lambda\) in the principal domain.
Cancellation would force an outward zero orbit and eventually leave the
critical strip. At \(\lambda=1\), \(B_1^{(0)}=0\). For \(\lambda>1\), RH
moves both relevant pole families off the principal domain or onto the
removed cut; conversely, an off-line zero would require the impossible
outward orbit \(\rho,\lambda\rho,\lambda^2\rho,\ldots\).

For \(\lambda\ge2\), the terminal is Euler-safe throughout the principal
domain. For every \(\lambda>1\), the positive source value near the cut
endpoint gives the wrong Stieltjes orientation. For \(0<\lambda<1\), the
orientation reverses but the terminal poles enter the domain. This is
the origin-dilation dichotomy. It does not settle arbitrary positive
continuous mixtures, whose moving singularity curves may collide and
require a separate analysis.

\subsection{144.2 Critical-line-centered
complement}\label{critical-line-centered-complement}

Define instead

\[
\widetilde T_\lambda(z)
=\frac{z}{\lambda R(z)}
\Phi\!\left(\frac12+\lambda
\left(s(z)-\frac12\right)\right)
\qquad
\widetilde B_\lambda=h-\widetilde T_\lambda
\tag{144.3}
\]

The upper folded lip stays on the critical line for every \(\lambda>0\).
Hence the terminal cut contribution has zero imaginary part. The precise
holomorphy theorem is

\[
\boxed{
\widetilde B_1=0,\qquad
\widetilde B_\lambda\text{ holomorphic on }\Omega
\iff\mathrm{RH}\quad(\lambda\ne1)}
\tag{144.4}
\]

To see the converse, write each zero as \(\rho=1/2+\alpha\). Holomorphic
cancellation for \(\lambda\ne1\) would make the nonreal zero multiset
invariant under scaling \(\alpha\mapsto\lambda\alpha\). If
\(\lambda>1\), the orbit leaves the critical strip; if \(0<\lambda<1\),
the orbit accumulates at \(\alpha=0\), impossible for the discrete zero
set. Under RH, every centered preimage remains on the removed cut.

Equation (144.4) gives the centered-dilation geometry but supplies no
forcing sign. The same critical-line reality that removes the cut
obstruction also removes the boundary leverage.

\section{145. Cyclic resolvents and finite visibility of every off-line
shell}\label{cyclic-resolvents-and-finite-visibility-of-every-off-line-shell}

Let \(q_\rho=\rho(1-\rho)\), combine equal folded shells with
multiplicity \(M_q\), and write

\[
h(x)=\sum_qM_q\frac{x}{x+q}
\qquad
L_h(x,y)=\sum_qM_q
\frac{q}{(x+q)(y+q)}
\tag{145.1}
\]

On the weighted shell space

\[
\mathscr K=\ell^2(\mathcal Q,M_q|q|)
\qquad
r_x(q)=\frac1{x+q}
\tag{145.2}
\]

the positive-node resolvents are cyclic.

\subsection{Theorem 145.1}\label{theorem-145.1}

For every nonempty open interval \(I\subset(0,\infty)\),

\[
\boxed{
\overline{\operatorname{span}_{\mathbb C}
\{r_x:x\in I\}}=\mathscr K}
\tag{145.3}
\]

\textbf{Proof.} If \(a\in\mathscr K\) is orthogonal to every \(r_x\),
then the meromorphic Cauchy transform

\[
F_a(z)=\sum_q
\frac{M_q|q|\overline{a(q)}}{z+q}
\tag{145.4}
\]

vanishes on \(I\). Normal convergence gives a meromorphic function on
the complement of its poles. The identity theorem and the residues at
\(z=-q\) force every coordinate of \(a\) to vanish. Conjugation
symmetrization gives the corresponding real totality statement.
\(\square\)

If \(q_0\notin\mathbb R\), choose a conjugation-real vector supported on
\(q_0,\bar q_0\) whose phase makes

\[
\mathcal E(f)=\sum_qM_qqf(q)^2
\]

strictly negative. Approximation by (145.3) then yields

\[
\boxed{\begin{gathered}
\mathrm{RH\ false}\Longrightarrow\\
\text{a negative finite ordinary Löwner witness exists}\\
\text{inside every positive node interval}
\end{gathered}}
\tag{145.5}
\]

At any fixed \(x_0>0\), the powers \((x_0+q)^{-k-1}\) are likewise
total: orthogonality makes every derivative of (145.4) vanish at
\(x_0\). Therefore

\[
\boxed{
\mathrm{RH\ false}\Longrightarrow
\forall x_0>0\ \exists N<\infty:
\mathfrak D_N(h;x_0)\not\succeq0}
\tag{145.6}
\]

There is no rank bound in (145.5)--(145.6). Positive rank-nine,
rank-eleven, rank-twelve, and fifth-rung point certificates are fully
compatible with a later finite failure.

For a fixed finite node set \(X\) and a safe terminal,

\[
L_{h-T_\lambda}(X)=L_h(X)-L_{T_\lambda}(X)
\qquad
L_{T_\lambda}(X)\succ0
\]

while \(L_{T_\lambda}(X)\to0\) as \(\lambda\to\infty\). Hence

\[
\boxed{
\exists\lambda\ge2:
L_{h-T_\lambda}(X)\succeq0
\iff L_h(X)\succ0}
\tag{145.7}
\]

The identical statement holds for a fixed confluent block. This
rank-adaptive complement is a strictness reformulation, not an escape
from the global gate.

Combining the inherited one-point theorem with (143.8) and (145.3) gives
the exact implication chain

\[
\boxed{
\begin{aligned}
\mathrm{RH}
&\iff h\text{ is complete Bernstein}\\
&\iff L_h(X)\succeq0
\quad\text{for every finite positive }X\\
&\iff \mathfrak D_N(h;x_0)\succeq0
\quad\text{for all }N
\text{ at any one prescribed }x_0>0\\
&\iff T_{\lambda_j}\in\mathrm{CBF}
\quad\text{for some }\lambda_j\downarrow1
\end{aligned}}
\tag{145.8}
\]

None of the arrows in (145.8) supplies the missing premise.

\section{146. All-rank probability and common-carrier
identities}\label{all-rank-probability-and-common-carrier-identities}

At a fixed \(x_0>0\), let

\[
b_n(x_0)=x_0^n\frac{(-1)^n}{(n+1)!}h^{(n+1)}(x_0)
\]

Taylor's theorem gives the exact all-rank generating function

\[
\boxed{
\sum_{n\ge0}b_n(x_0)z^{n+1}
=\frac{h(x_0)-h(x_0(1-z))}{x_0}}
\tag{146.1}
\]

At \(x_0=56\),

\[
R(z)=\sqrt{225-224z}
\qquad
s(z)=\frac{1+R(z)}2
\]

\[
\boxed{
H_{56}(z)
=\frac{\Phi(8)}{15}
-(1-z)\frac{\Phi(s(z))}{R(z)}}
\tag{146.2}
\]

The terminal exponential has probability generating function

\[
\boxed{
\Pi_t(z)=\frac1{R(z)}
\exp\!\left[-\frac{R(z)-1}{2}t\right]
=\sum_{n\ge0}\pi_n(t)z^n
\qquad \Pi_t(1)=1}
\tag{146.3}
\]

Let \(\beta_k(t)\) be the PP189 source polynomials in their controlling
normalization. Coefficient extraction from \((1-z)\Pi_t(z)\) gives, for
every \(k\ge0\),

\[
\boxed{
e^{-8t}\beta_k(t)
=e^{-t}\bigl(\pi_{k+1}(t)-\pi_k(t)\bigr)}
\tag{146.4}
\]

The exact replay reconstructs (146.4) through \(\beta_{23}\), including
the degree-\(23\) controlling range. With

\[
\alpha=\frac{224}{225}
\qquad c=\frac{15t}{2}
\]

\[
\pi_n(t)
=\frac{e^{-7t}}{15}
\frac{\alpha^n}{2^nn!}\theta_n(c)
\qquad
\theta_n=(2n-1)\theta_{n-1}+c^2\theta_{n-2}
\tag{146.5}
\]

The Bessel recurrence is the analytic return of the same quadratic
Jacobian used in §R1 and technical §34. For \(x>0\) and \(t\ge0\), put
\(K_t(x)=e^{-t s(x)}/R(x)\), with \(s(x)=(1+R(x))/2\). Then

\[
\boxed{(-1)^rD_x^rK_t(x)
=\frac{2^r e^{-t s(x)}}{R(x)^{2r+1}}
\theta_r\!\left(\frac{tR(x)}2\right),\qquad r\ge0}
\tag{146.5a}
\]

For completeness its polynomial normalization is

\[
\theta_r(v)=\sum_{j=0}^r
\frac{(2r-j)!}{2^{r-j}(r-j)!j!}v^j
\qquad
\theta_{r+1}=(v+2r+1)\theta_r-v\theta_r'
\tag{146.5b}
\]

The last identity follows by coefficient comparison. Since \(D_xR=2/R\)
and \(D_x(tR/2)=t/R\), differentiating (146.5a) gives exactly (146.5b),
proving (146.5a) by induction from \(r=0\). The same coefficients give
\(\theta_{r+1}=(2r+1)\theta_r+v^2\theta_{r-1}\) for \(r\ge1\), which is
the shifted recurrence in (146.5). Finally
\(\Pi_t(z)=e^tK_t(x_0(1-z))\); Taylor extraction at \(x_0=56\) recovers
the coefficient formula in (146.5), including every scale factor.

All coefficients of \(\theta_r\) are positive, so \(K_t\) is completely
monotone. The multiplication by \(x\) in the completed source must still
be carried through:

\[
\frac{d}{dx}[xK_t(x)]
=\frac{e^{-t s(x)}}{R(x)^3}
\bigl(1+2x-txR(x)\bigr)
\tag{146.5c}
\]

For fixed \(x>0\) this changes sign as \(t\) increases. The exact
root-to-derivative connection therefore preserves the source problem:
positivity of the elementary kernel does not prove the completed matrix
sign.

The ratio recurrence proves that each \(\beta_{n-1}\) has exactly one
positive zero. Individual prime cells therefore change sign; completion,
not atomwise positivity, remains necessary.

The same law has the mixed-Poisson density

\[
\boxed{
f_t(\ell)=
\frac{e^{t/2}}{2\sqrt{\pi x_0\ell}}
\exp\!\left[-\frac{\ell}{4x_0}
-\frac{x_0t^2}{4\ell}\right]
\qquad t\ge0}
\tag{146.6}
\]

At \(x_0=56\),

\[
\mathbb E_tN=56(t+2)
\qquad
\operatorname{Var}_tN=6328t+25200
\tag{146.7}
\]

The endpoint estimates at \(\ell=0\) justify the Rodrigues formula at
\(t=0\) as well as \(t>0\):

\[
(-1)^r\Delta^{r+1}\pi_n(t)
=\frac{(-1)^r}{(n+r+1)!}
\int_0^\infty e^{-\ell}\ell^{n+r+1}
f_t^{(r+1)}(\ell)\,d\ell
\tag{146.8}
\]

For

\[
u_a(x)=\frac{x}{R(x)}e^{-s(x)a}
\qquad
K_{n,r}(a)=G_{n,r}[-u_a](x_0)
\]

local uniform derivative bounds justify every finite differentiation,
difference, Gamma integral, and prime sum. Thus for each fixed finite
\((n,r)\),

\[
\boxed{
\begin{aligned}
G_{n,r}[h](x_0)={}&-a_\Gamma K_{n,r}(0)\\
&+\int_0^\infty
\frac{K_{n,r}(a)-e^{-a/2}K_{n,r}(0)}
{1-e^{-2a}}\,da\\
&+\sum_{m\ge2}\Lambda(m)K_{n,r}(\log m)
\end{aligned}}
\tag{146.9}
\]

where \(a_\Gamma=[\psi(1/4)-\log\pi]/2\) is the same scalar anchor as in
§154. The local estimates depend on \((n,r)\); (146.9) is not uniform in
rank.

Finally, put

\[
H_R(t)=[\beta_{i+j}(t)]_{i,j=0}^{R-1}
\quad
A=\frac{56}{15}
\quad
D_R(t)=\operatorname{diag}
(1,At,\ldots,(At)^{R-1})
\]

The exact leading coefficient gives

\[
\boxed{
\frac{15}{At}D_R(t)^{-1}H_R(t)D_R(t)^{-1}
=M_R+O_R(t^{-1})
\qquad
M_R=\left[\frac1{(i+j+1)!}\right]}
\tag{146.10}
\]

Since \(M_R\) is nonsingular with alternating \(LDL^{\mathsf T}\)
pivots,

\[
\boxed{
\operatorname{inertia}H_R(t)
=\left(\left\lceil\frac R2\right\rceil,
\left\lfloor\frac R2\right\rfloor,0\right)}
\tag{146.11}
\]

for every fixed \(R\ge2\) and all sufficiently large \(t\). Rank one is
eventually positive. No effective or rank-uniform threshold is asserted.

\section{147. Exact carrier weights and finite-anchor
map}\label{exact-carrier-weights-and-finite-anchor-map}

The raw Groskin jump readout is

\[
\boxed{
-\frac{u}{2d^2}\langle J_N,dM_N(u)\rangle_F
=\sum_q\Lambda(q)q^{-1/2}
\delta_{\log q}(du)}
\tag{147.1}
\]

Because \(u_a=e^{-a/2}\kappa(\cdot,a)\), evaluation at \(a=\log q\)
absorbs the second square-root weight. At the one-point cell,

\[
\boxed{
\sum_q\Lambda(q)K_{n,r}(\log q)
=\sum_q\frac{\Lambda(q)}q
(-1)^r\Delta^{r+1}\pi_n(\log q)}
\tag{147.2}
\]

Equations (147.1)--(147.2) are exact arithmetic transport. The response
matrix is already indefinite at \(q=2\) on the tested node pair;
therefore the map is not a positive congruence and does not prove
\(L_\Gamma-L_P\succeq0\).

The finite anchors now have the following exact grades:

\begin{longtable}[]{@{}
  >{\raggedright\arraybackslash}p{(\columnwidth - 4\tabcolsep) * \real{0.3333}}
  >{\raggedright\arraybackslash}p{(\columnwidth - 4\tabcolsep) * \real{0.3333}}
  >{\raggedright\arraybackslash}p{(\columnwidth - 4\tabcolsep) * \real{0.3333}}@{}}
\toprule\noalign{}
\begin{minipage}[b]{\linewidth}\raggedright
anchor
\end{minipage} & \begin{minipage}[b]{\linewidth}\raggedright
certified statement
\end{minipage} & \begin{minipage}[b]{\linewidth}\raggedright
authority boundary
\end{minipage} \\
\midrule\noalign{}
\endhead
\bottomrule\noalign{}
\endlastfoot
local rank \(9\), \(x=56\) &
\(\mathfrak B_9^{\mathrm{conf}}[h;56]\succ0\); dyadic complement first
becomes positive at depth \(N=79\) with final pivot
\(1.5436100114\ldots\times10^{-27}\) & one basepoint; minimal only in
that dyadic family \\
local rank \(11\), \(x=56\) & directed confluent certificate & one
basepoint \\
local rank \(12\), \(x=56\) & controlling PP190-T certificate with
(138.2) & one basepoint and one manifest-bound cutoff/tail proof \\
fifth rung, \(x=90\) &
\(W_5(90)=3.2334949506154321\ldots\times10^{-8}>0\) & this anchor is one
point; global zero-assisted positivity is in Reading Volume §R13, while
the independent source proof remains open \\
Connes--CvS \(c=100\) cell & current package rebuilds the
\(501\times501\) cell, selected positive value, ten recovery checks, and
a directed normalized residual
\(\le9.1042746997704021\ldots\times10^{-501}\) & exact supplied finite
dyadic matrix only \\
\end{longtable}

The Connes--CvS bound is the normalized quantity

\[
\boxed{
\operatorname{dist}(\widehat\lambda,\sigma(A))
\le
\frac{\|Av-\widehat\lambda v\|_2}{\|v\|_2}
\le9.1042746997704021\ldots\times10^{-501}}
\tag{147.3}
\]

The implementation encloses \(\|v\|_2^2\) from below before division, so
the certificate controls the normalized residual. The eleven preceding
negative values are nondirected branch diagnostics; a directed inertia
or branch-minimality certificate is not included.

\Needspace{11\baselineskip}

\chapter{Technical chapter XIII. Completion forms, reflection contours,
and realization
limits}\label{technical-chapter-xiii.-completion-forms-reflection-contours-and-realization-limits}

Positive auxiliary forms are compared with the actual Gamma-minus-prime
form. Reflection, conjugation and the topology of a proposed realization
remain explicit.

\section{103. Why the source problem becomes
cancellation-first}\label{why-the-source-problem-becomes-cancellation-first}

The PP146--PP150 corridor closes a family of natural positive
transports. It does not close all conceivable transports.

\subsection{103.1 The Mills calibration and the Stieltjes-cone
obstruction}\label{the-mills-calibration-and-the-stieltjes-cone-obstruction}

The normalized Mills owner

\[
G_M(x)=\frac{m(\sqrt x)}{\sqrt x}
\]

is the Stieltjes transform of the \(\chi_1^2\) probability measure

\[
d\mu_M(t)=\frac{e^{-t/2}}{\sqrt{2\pi t}}\,dt
\tag{103.1}
\]

Its moments and Hankel determinants are

\[
\mu_n=(2n-1)!!,\qquad
\Delta_N=\prod_{j=0}^{N-1}(2j)!
\qquad
\Delta_N^{(1)}=\prod_{j=0}^{N-1}(2j+1)!
\tag{103.2}
\]

By contrast, one folded prime atom with log-frequency \(a>0\) is

\[
p_a(x)=\frac{e^{-a s(x)}}{R(x)}
\qquad
R(x)=\sqrt{1+4x}
\qquad
s(x)=\frac{1+R(x)}2
\tag{103.3}
\]

Its inverse Laplace density is

\[
\kappa_a(u)
=\frac1{2\sqrt{\pi u}}
\exp\!\left[-\frac{(u+a)^2}{4u}\right]
\tag{103.4}
\]

This density is positive, so \(p_a\) is completely monotone. But

\[
\frac{d}{du}\log\kappa_a(u)
=-\frac1{2u}-\frac14+\frac{a^2}{4u^2}
\tag{103.5}
\]

is positive near \(u=0\). Thus \(\kappa_a\) is not completely monotone
and \(p_a\) is not Stieltjes. Positive mixtures or subordinations that
remain inside the Stieltjes heat cone cannot map the Mills model to even
one prime atom.

\subsection{103.2 The triangular operator explains grading, not
positivity}\label{the-triangular-operator-explains-grading-not-positivity}

In Bargmann--Fock coordinates, with \(N=zD_z\),

\[
\mathcal T=\frac12N(N+1)
\qquad
2\mathcal T+\frac14I
=\left(N+\frac12\right)^2
\tag{103.6}
\]

The inverse on the zero-constant sector has the exact Beta
representation

\[
\mathcal T^{-1}f(z)
=2\int_0^1\frac{1-t}{t}f(tz)\,dt
\qquad
\frac1{T_{k+1}}=\int_0^1t^k\,2(1-t)\,dt
\tag{103.7}
\]

This is a positive Hausdorff multiplier. The forward multiplier is not a
moment or Löwner positivity preserver: applying \(T_{k+1}\) to the
Hausdorff sequence \(1,1,1,\ldots\) gives \(1,3,6,\ldots\) and

\[
\det\begin{pmatrix}1&3\\3&6\end{pmatrix}=-3
\tag{103.8}
\]

Because \(\mathcal T\) and \(\mathcal T^{-1}\) are functions of \(N\),
they commute with every mod-\(4\) projector. They explain the residue
grading but cannot supply the missing channel-changing positive
operator.

\subsection{103.3 One folded atom is not the matrix
owner}\label{one-folded-atom-is-not-the-matrix-owner}

For a prime log-atom \(y>0\), define

\[
h_y(x)=\frac{x}{R}e^{-sy}
\]

Exact differentiation gives

\[
h_y'(x)
=e^{-sy}\left[
\frac{R^2+1}{2R^3}
-\frac{xy}{R^2}
\right]
\tag{103.9}
\]

Hence \(h_y'(x)\ge0\) precisely when

\[
y\le\eta(R)
=\frac{2(R^2+1)}{R(R^2-1)}
\tag{103.10}
\]

At \(R=4\), \(\eta(4)=17/30<\log2\). Every prime atom is therefore
decreasing for \(x\ge15/4\). Prime-atom PSD cannot be the full source
mechanism.

Finally, the completed heat matrix has the cancellation-first identity

\[
D_NH_N(x)D_N
=\int_0^\infty e^{-xu}\Theta_\xi(u)
E_N(u)K_N(xu)E_N(u)\,du
\tag{103.11}
\]

Here \(H_N=[h^{(a+b+1)}/(a+b+1)!]_{a,b=0}^{N-1}\),
\(D_N=\operatorname{diag}((-1)^a)\), and
\(E_N(u)=\operatorname{diag}(u^a)_{a=0}^{N-1}\). The pencil is

\[
K_N(t)_{ab}
=\frac1{(a+b)!}-\frac{t}{(a+b+1)!}
\qquad 0\le a,b<N
\tag{103.12}
\]

The constant heat mode cancels only after the completed integral is
formed. The pencil itself is indefinite for every \(N\ge2\): for
\(N=2m\) its inertia is \((m,m,0)\) for \(t\ge0\); for \(N=2m+1\) there
is one real crossing and otherwise both signs remain. Therefore the
proof cannot be pointwise PSD. The owner is signed cancellation across
the completed Gamma and prime components.

\section{157. The endpoint-subtracted Beta--log Gram
bridge}\label{the-endpoint-subtracted-betalog-gram-bridge}

The Montgomery--Vaughan video factor is

\[
J_y(\delta)
=\int_0^y e^{2\pi i\delta x}\,dx
=\frac{e^{2\pi i\delta y}-1}{2\pi i\delta}
\qquad
J_y(0)=y
\tag{157.1}
\]

The terminal \(-1\) is the lower-endpoint subtraction. It removes the
apparent diagonal singularity and installs the derivative value \(y\).
Equivalently,

\[
\boxed{
J_y(\delta)=y\,e^{\pi i\delta y}\operatorname{sinc}(\delta y),\qquad
\operatorname{sinc}(u)=\frac{\sin\pi u}{\pi u}}
\tag{157.2}
\]

For \(L_n=\log n\) and \(F(x)=\sum_na_ne^{2\pi iL_nx}\),

\[
0\le\int_0^1\int_0^y|F(x)|^2\,dx\,dy
=\int_0^1(1-u)|F(u)|^2\,du
\tag{157.3}
\]

After normalization, the owner is the probability measure

\[
\boxed{d\nu_\triangle(u)=2(1-u)\,du,\qquad0\le u\le1}
\tag{157.4}
\]

It simultaneously owns the reciprocal-triangular moments

\[
\int_0^1u^k\,d\nu_\triangle(u)
=\frac{2}{(k+1)(k+2)}
=\frac1{T_{k+1}}
\tag{157.5}
\]

and the log-character Gram kernel

\[
\boxed{\begin{aligned}
\mathcal B(\delta)
&=\int_0^1e^{2\pi i\delta u}\,d\nu_\triangle(u)\\
&=\frac{2(1+ic-e^{ic})}{c^2}
\qquad c=2\pi\delta,\qquad \mathcal B(0)=1
\end{aligned}}
\tag{157.6}
\]

Its real part is

\[
\Re\mathcal B(\delta)
=\left(\frac{\sin\pi\delta}{\pi\delta}\right)^2
=\operatorname{sinc}^2(\delta)\ge0
\tag{157.7}
\]

With the inner product conjugate-linear in its first slot, for any
finite frequency set,

\[
B_{mn}=\mathcal B(L_n-L_m)
=\langle e^{2\pi iL_m u},e^{2\pi iL_n u}\rangle_{L^2(\nu_\triangle)}
\tag{157.8}
\]

is Hermitian positive semidefinite. Monomials and log characters may be
placed in the same Gram space, producing the positive block

\[
\begin{bmatrix}
[\langle u^i,u^j\rangle]&[\langle u^i,e^{2\pi iL_nu}\rangle]\\
[\langle e^{2\pi iL_mu},u^j\rangle]&[\mathcal B(L_n-L_m)]
\end{bmatrix}\succeq0
\tag{157.9}
\]

This is a genuine source-side sesquilinear form. It is not yet the
missing completed source inequality. Its exact firewall is the
accumulating gap

\[
L_{n+1}-L_n=\log\left(1+\frac1n\right)\sim\frac1n
\tag{157.10}
\]

No fixed compact positive window can orthogonalize infinitely many
frequencies whose gaps approach zero. The PP204 replay shows that
fixed-rank Jacobi projections resolve small frequencies but fail
uniformly as the log frequency grows, and a window capable of separating
adjacent labels near \(N\) must scale at least on the order of \(N\) in
the nearest-neighbor model.

\begin{figure}
\centering
\includegraphics{figures/figure77_beta_log_projection_and_scale_firewall.png}
\caption{Projection and scale diagnostics for the Beta-logarithm model.
Each curve is finite and replayed; the plotted trend is a firewall, not
an asymptotic proof beyond the stated theorem.}
\end{figure}

The surviving target is not a larger scalar window. It is a
source-defined, operator-weighted, nonstationary compression whose
\textbf{final} form has the translation covariance required by the Weil
functional while its internal lift still couples the two reflection
sheets.

\section{154. Completion-coupled translation defects and the PP97
carrier}\label{completion-coupled-translation-defects-and-the-pp97-carrier}

Reserve \(\Lambda_{\mathrm{vM}}(m)\) for the von Mangoldt function. For
a real finite resolvent combination \(F=\sum_i c_if_{x_i}\) define

\[
\begin{aligned}
\widehat F(t)&=\int_{\mathbb R}F(u)e^{-itu}\,du,\qquad
g_F(v)=\langle F,\tau_vF\rangle\\
D_F(v)&=g_F(0)-g_F(v)=\frac12\|(I-\tau_v)F\|_2^2
\end{aligned}
\tag{154.1}
\]

Here \((\tau_vF)(u)=F(u+v)\) and

\[
w_\infty(t)=\frac1{4\pi}
\left[\Re\psi\!\left(\frac14+\frac{it}{2}\right)-\log\pi\right]
=\frac1{4\pi}\left[\Re\psi\!\left(\frac{s}2\right)-\log\pi\right]_{s=1/2+it}
\tag{154.2}
\]

The PP200/PP201 Gamma representation, independently rebuilt in PP205 and
replayed in PP206, is

\[
\boxed{
\int_{\mathbb R}|\widehat F(t)|^2w_\infty(t)\,dt
=a_\Gamma\|F\|_2^2
+\int_0^\infty D_F(v)\,d\mu_\Gamma(v)}
\tag{154.3}
\]

where

\[
\boxed{\begin{aligned}
a_\Gamma&=\frac12\!\left(\psi\!\left(\frac14\right)-\log\pi\right)
=-2.686091709612832791\ldots\\
d\mu_\Gamma(v)&=\frac{e^{v/2}}{2\sinh v}\,dv>0
\end{aligned}}
\tag{154.4}
\]

The scalar anchor is negative; every translated defect is nonnegative.
The formula is therefore a signed anchor plus a positive continuum, not
a positive measure representation of the completed form by itself.

The density has the exact expansion

\[
\frac{e^{v/2}}{2\sinh v}=\sum_{k\ge0}e^{-(2k+1/2)v}
\tag{154.5}
\]

To read the poles in the expanded variable, analytically continue the
Fourier height \(t\). With \(a_k=2k+1/2\),

\begin{longtable}[]{@{}
  >{\raggedright\arraybackslash}p{(\columnwidth - 6\tabcolsep) * \real{0.2500}}
  >{\raggedright\arraybackslash}p{(\columnwidth - 6\tabcolsep) * \real{0.2500}}
  >{\raggedright\arraybackslash}p{(\columnwidth - 6\tabcolsep) * \real{0.2500}}
  >{\raggedright\arraybackslash}p{(\columnwidth - 6\tabcolsep) * \real{0.2500}}@{}}
\toprule\noalign{}
\begin{minipage}[b]{\linewidth}\raggedright
continued height
\end{minipage} & \begin{minipage}[b]{\linewidth}\raggedright
\(s=1/2+it\)
\end{minipage} & \begin{minipage}[b]{\linewidth}\raggedright
folded point \(q=s(1-s)\)
\end{minipage} & \begin{minipage}[b]{\linewidth}\raggedright
behavior of \(e^{-itv}\)
\end{minipage} \\
\midrule\noalign{}
\endhead
\bottomrule\noalign{}
\endlastfoot
\(t=+ia_k\) & \(s=-2k\) & \(-2k(2k+1)=-2T_{2k}\) & \(e^{+a_kv}\)
grows \\
\(t=-ia_k\) & \(s=2k+1\) & \(-2k(2k+1)=-2T_{2k}\) & \(e^{-a_kv}\)
decays \\
\end{longtable}

These are the two sheets of the folded Gamma pole. The \(k=0\) point
\(s=0\) is an endpoint of the completed factor, not a nontrivial zeta
zero.

The functional equation at its own fixed point supplies the exact
anchor:

\[
\boxed{
a_\Gamma=-\frac{\zeta'}{\zeta}\!\left(\frac12\right)}
\tag{154.6}
\]

Indeed, \(\xi'(1/2)=0\), the terms \(1/s\) and \(1/(s-1)\) cancel at
\(s=1/2+i0\), and the completed logarithmic derivative leaves (154.6).
This is analytic continuation. It is \textbf{not} permission to replace
the divergent series \(\sum_m\Lambda_{\mathrm{vM}}(m)/\sqrt m\) by
\(a_\Gamma\) inside a positive sum.

The prime channel is nevertheless exactly computable on the resolvent
span. Put

\[
a_i=\sqrt{x_i+\frac14},\qquad
s_i=\frac12+a_i=\sigma_i+i0>1
\tag{154.7}
\]

and

\[
\begin{aligned}
A_i&=\frac{x_i}{2a_i}\,2\sum_{j\ne i}\frac{c_ic_j}{x_i-x_j}
+c_i^2\frac{x_i+1/2}{4a_i^3}\\
B_i&=c_i^2\frac{x_i}{4a_i^2}
\end{aligned}
\tag{154.8}
\]

Then

\[
g_F(v)=\sum_i(A_i-B_iv)e^{-a_iv}
\tag{154.9}
\]

and, because every \(\sigma_i>1\),

\[
\boxed{
\sum_{m\ge2}\frac{\Lambda_{\mathrm{vM}}(m)}{\sqrt m}g_F(\log m)
=\sum_i\left[
A_i\left(-\frac{\zeta'}\zeta\right)(s_i)
-B_i\left(\frac{\zeta'}\zeta\right)'(s_i)
\right]}
\tag{154.10}
\]

No prime table, cutoff, zero table, or tail estimate is required in
(154.10). This is the practical source-side closure of the carrier.

For one node, it agrees exactly with PP97:

\begin{longtable}[]{@{}lrl@{}}
\toprule\noalign{}
quantity at \(x=0.05\) & replayed source value & role \\
\midrule\noalign{}
\endhead
\bottomrule\noalign{}
\endlastfoot
\((xG)'\) & \(-1.205641588399722174914\) & archimedean channel \\
\((xP)'\) & \(-1.228733588379453600229\) & prime channel \\
\(W_1=(xG)'-(xP)'\) & \(+0.02309199997973142531\) & completed
difference \\
\end{longtable}

The completion is visible in the cancellation, not in separate
positivity of the two channels.

For a cutoff \(M\) set

\[
\theta^\sharp(M)=\sum_{2\le m\le M}
\frac{\Lambda_{\mathrm{vM}}(m)}{\sqrt m}
\tag{154.11}
\]

Then

\[
-\sum_{m\le M}\frac{\Lambda_{\mathrm{vM}}(m)}{\sqrt m}g_F(\log m)
=\sum_{m\le M}\frac{\Lambda_{\mathrm{vM}}(m)}{\sqrt m}D_F(\log m)
-\theta^\sharp(M)\|F\|_2^2
\tag{154.12}
\]

Both terms on the right diverge. The tempting substitution
\(\theta^\sharp(M)\rightsquigarrow a_\Gamma\) is the named
\textbf{regularization trap}: it keeps the divergent positive-defect sum
and replaces only its compensator, manufacturing a false proof.

The geometric obstruction is sharper. The Gamma mass diverges at short
shift,

\[
d\mu_\Gamma(v)\sim\frac{dv}{2v}\quad(v\downarrow0)
\tag{154.13}
\]

whereas the smooth prime mass grows at long shift like \(e^{v/2}dv\).
The densities differ by the factor \(e^v\) asymptotically. A positive
scalar diagonal weight cannot move the required compensation from one
end to the other; an operator-weighted, genuinely two-point mechanism
remains the open class.

\section{158. Two measures that must not be
conflated}\label{two-measures-that-must-not-be-conflated}

Two exact measures have different triangular roles:

\begin{longtable}[]{@{}
  >{\raggedright\arraybackslash}p{(\columnwidth - 6\tabcolsep) * \real{0.2500}}
  >{\raggedright\arraybackslash}p{(\columnwidth - 6\tabcolsep) * \real{0.2500}}
  >{\raggedright\arraybackslash}p{(\columnwidth - 6\tabcolsep) * \real{0.2500}}
  >{\raggedright\arraybackslash}p{(\columnwidth - 6\tabcolsep) * \real{0.2500}}@{}}
\toprule\noalign{}
\begin{minipage}[b]{\linewidth}\raggedright
measure
\end{minipage} & \begin{minipage}[b]{\linewidth}\raggedright
support and mass
\end{minipage} & \begin{minipage}[b]{\linewidth}\raggedright
exact moments
\end{minipage} & \begin{minipage}[b]{\linewidth}\raggedright
role
\end{minipage} \\
\midrule\noalign{}
\endhead
\bottomrule\noalign{}
\endlastfoot
\(d\nu_\triangle(u)=2(1-u)du\) & \([0,1]\), probability &
\(\int u^k d\nu_\triangle=1/T_{k+1}\) for \(k\ge0\) & Beta--Hankel and
log-character Gram owner \\
\(d\mu_\Gamma(v)=e^{v/2}(2\sinh v)^{-1}dv\) & \((0,\infty)\), infinite
total mass & finite moments only for \(k\ge1\) & archimedean
translation-defect owner \\
\end{longtable}

For \(k\ge1\), expand (154.5) and integrate termwise:

\[
\boxed{
\begin{aligned}
\int_0^\infty v^k\,d\mu_\Gamma(v)
&=\sum_{m\ge0}\frac{k!}{(2m+1/2)^{k+1}}\\
&=\frac{k!}{2^{k+1}}\zeta\!\left(k+1,\frac14\right)\\
&=\frac{(-1)^{k+1}}{2^{k+1}}
\psi^{(k)}\!\left(\frac14\right)
\end{aligned}}
\tag{158.1}
\]

The first moment is

\[
\int_0^\infty v\,d\mu_\Gamma(v)
=\frac14\psi^{(1)}\!\left(\frac14\right)
=\frac{\pi^2+8G}{4}
\tag{158.2}
\]

where \(G\) is Catalan's constant. At \(k=0\) the integral diverges; the
finite quantity \(a_\Gamma\) in (154.4) is the renormalized scalar
anchor and is not the zeroth moment of a finite measure.

The measure comparison figure in Reading Volume §R18 shows these
distinct supports.

The phrase ``triangular rail'\,' therefore names two different exact
appearances: reciprocal-triangular \textbf{moments} on the compact side
and doubled-even-triangular \textbf{pole locations} on the Gamma side.
Their common arithmetic vocabulary is not a measure identity.

\section{137. Exact completed Weil/resolvent
identity}\label{exact-completed-weilresolvent-identity}

For positive nodes \(x_i\), real coefficients \(c_i\), and

\[
q(s)=s(1-s),\qquad
\Phi_c(s)=q(s)\left(\sum_i\frac{c_i}{x_i+q(s)}\right)^2
\]

PP190-T identifies the completed Weil form exactly. Let
\(\mathcal W_{\mathrm{pair}}\) denote the folded-orbit normalization
used in this volume and let \(\mathcal W_{GW}\) denote the classical
convention that counts every zero. The compulsory fold factor is

\[
\boxed{
\mathcal W_{\mathrm{pair}}(c)
=c^{\mathsf T}(L_\Gamma-L_P)c
=\frac12\mathcal W_{GW}[H_c]}
\tag{137.1}
\]

On \(s=\tfrac12+it\), put

\[
H_c(t)=\left(t^2+\frac14\right)
\left(\sum_i\frac{c_i}{t^2+x_i+\frac14}\right)^2\ge0
\]

and

\[
g_c(u)=\frac1{2\pi}\int_{\mathbb R}H_c(t)e^{-itu}\,dt
\]

Then

\[
\boxed{
\mathcal W_{\mathrm{pair}}(c)=
\frac1{4\pi}\int_{\mathbb R}H_c(t)
\left[\Re\psi\!\left(\frac14+\frac{it}{2}\right)-\log\pi\right]dt
-\sum_{m\ge2}\frac{\Lambda(m)}{\sqrt m}g_c(\log m)}
\tag{137.2}
\]

There is no endpoint remainder for this test family. If
\(a_x=\sqrt{x+\tfrac14}\) and

\[
f_x(u)=e^{-a_x|u|}
\left(\frac1{4a_x}+\frac12\operatorname{sgn}u\right)
\]

then \(H_c=|\widehat f_c|^2\) and
\(g_c(u)=\langle f_c,\tau_uf_c\rangle\). Thus (137.2) is a single
completion-coupled nonlocal quadratic form. It is not yet a Hilbert norm
square. Separating its two finite-cutoff norm channels produces
divergent common backgrounds at infinite cutoff; only their coupled
autocorrelation difference converges on the resolvent span.

At \(x=56\), \(a_x=15/2\), and the prime sampling factor is exactly

\[
m^{-1/2}e^{-(15/2)\log m}=m^{-8}
\tag{137.3}
\]

This is the exact adapter between the PP189 common-tail arithmetic and
the completed resolvent form.

\section{148. Resolvent autocorrelation and the zero-margin
firewall}\label{resolvent-autocorrelation-and-the-zero-margin-firewall}

For

\[
a_x=\sqrt{x+\frac14}
\qquad
f_x(u)=e^{-a_x|u|}
\left(\frac1{4a_x}+\frac12\operatorname{sgn}u\right)
\]

put

\[
\kappa(x,u)=\frac{x\,e^{-a_x|u|}}{2a_x}
\]

Parseval and one partial-fraction identity give the exact translation
divided difference

\[
\boxed{
\langle f_x,\tau_uf_y\rangle
=\frac{\kappa(x,u)-\kappa(y,u)}{x-y}}
\tag{148.1}
\]

At \(u=0\), \(\kappa(x,0)=x/\sqrt{1+4x}\) is operator monotone and
\(\kappa(x,0)/x\) is Stieltjes. The prime channel is

\[
\boxed{
xP(x)=\sum_{m\ge2}
\Lambda(m)m^{-1/2}\kappa(x,\log m)}
\tag{148.2}
\]

and the archimedean channel is the signed spectral integral

\[
\boxed{
xG(x)=\int_{\mathbb R}
\frac{x}{t^2+x+1/4}\,d\sigma_\infty(t)}
\tag{148.3}
\]

\[
d\sigma_\infty(t)
=\frac1{4\pi}
\left[
\Re\psi\!\left(\frac14+\frac{it}{2}\right)-\log\pi
\right]dt
\tag{148.4}
\]

Thus every prime power is a translate of one operator-monotone Gram
anchor. The difficulty is exact: translation destroys positive
definiteness, while completion restores only a coupled difference.

There is no pointwise symbol margin to exploit. Since
\(\Xi(t)=\xi(1/2+it)\) is real for real \(t\),

\[
\boxed{
\Re\frac{\xi'}{\xi}\!\left(\frac12+it\right)=0}
\tag{148.5}
\]

away from its poles, unconditionally. The real parts of the analytically
continued archimedean and prime symbols therefore agree on the critical
line. Any proposed proof based on a strict Gamma-over-prime boundary
margin is vacuous. The folded cut has no absolutely continuous jump from
\(h\); under RH the representing measure is purely atomic at the folded
zero locations.

The PNT-mode endpoint also survives:

\[
\int_0^\infty e^{u/2}\left(\frac14K(u)-K''(u)\right)\,du
=K'(0+)-\frac12K(0)
\tag{148.6}
\]

For the \(x=56\) resolvent, \(K'(0+)=0\) and \(K(0)=2/3375\), so

\[
\boxed{K'(0+)-K(0)/2=-1/3375\ne0}
\tag{148.7}
\]

A boundary counterterm or a test class with \(K(0)=0\) is compulsory.

\begin{figure}
\centering
\includegraphics{figures/v26c_reflection_pair.png}
\caption{For a real-coefficient test, reflection and conjugation use
different partners off the critical line. The full resolvent family
separates the distinction; a single multinode test may annihilate an
orbit. See §§159--160.}
\end{figure}

\section{149. Cauchy--Löwner contour
identities}\label{cauchyluxf6wner-contour-identities}

The exploratory contour calculation is relevant because it identifies
the precise complex-analytic owner of every finite Löwner and confluent
entry. Let \(\Gamma\) be a positively oriented contour inside a domain
of holomorphy of \(h\), enclosing the positive nodes \(x_i\) and no
singularities. Cauchy's formula gives

\[
\boxed{
L_h(x_i,x_j)
=\frac1{2\pi i}
\oint_\Gamma
\frac{h(z)}{(z-x_i)(z-x_j)}\,dz}
\tag{149.1}
\]

Consequently, for real \(c_i\),

\[
\boxed{
c^{\mathsf T}L_hc
=\frac1{2\pi i}
\oint_\Gamma h(z)
\left(\sum_i\frac{c_i}{z-x_i}\right)^2dz}
\tag{149.2}
\]

Coalescing all nodes at \(x\) yields

\[
\boxed{
\frac{h^{(i+j+1)}(x)}{(i+j+1)!}
=\frac1{2\pi i}
\oint_\Gamma
\frac{h(z)}{(z-x)^{i+j+2}}\,dz}
\tag{149.3}
\]

For the signed Dobsch normalization, replace \(z-x\) by \(x-z\); this
contributes \((-1)^{i+j}\). On the diagonal \(i=j=n\), (149.3) reaches
the native odd order

\[
i+j+1=2n+1
\]

which explains the repeated derivative order and factorial \((2n+1)!\)
in the confluent/Hausdorff hierarchy.

Equations (149.1)--(149.3) are exact but do not prove positivity. The
integrand contains a holomorphic square

\[
\left(\sum_i\frac{c_i}{z-x_i}\right)^2
\]

not an absolute square. A single contour therefore cannot create the
conjugation demanded by the nonholomorphic folded moment. The sharpened
next question is whether completion supplies a double-contour or
sesquilinear kernel of the form

\[
\iint
\mathcal K_{\mathrm{comp}}(z,w)
F(z)\overline{F(w)}\,dz\,d\bar w
\tag{149.4}
\]

whose compression is (149.2) and whose positivity is source-visible. No
such kernel is constructed here.

\section{159. Stationary-symbol blindness and the missing
conjugation}\label{stationary-symbol-blindness-and-the-missing-conjugation}

Let

\[
\Phi(s)
:=\frac{\xi'(s)}{\xi(s)},\qquad
s=\sigma+it
\tag{159.1}
\]

Away from zero ordinates, \(\Xi(t)=\xi(1/2+it)\) is real. Hence
\(i\Phi(1/2+it)=d(\log\Xi(t))/dt\) is real and

\[
\boxed{\Re\Phi(1/2+it)=0}
\tag{159.2}
\]

pointwise wherever the quotient is finite. There is no positive
archimedean-over-prime margin on the critical boundary; their real parts
cancel exactly after completion.

The distributional boundary value is also lossy. An off-line functional-
equation pair at \(\beta=1/2\pm\delta\) and height \(\gamma\)
contributes, with \(y=t-\gamma\),

\[
-\frac{\delta}{\delta^2+y^2}
+\frac{\delta}{\delta^2+y^2}=0
\tag{159.3}
\]

Thus the consistent distributional symbol is

\[
\boxed{
\Omega(t)=\pi\sum_{\rho:\,\beta_\rho=1/2}
m_\rho\,\delta_{\gamma_\rho}(t)\ge0}
\tag{159.4}
\]

unconditionally, while the pointwise regularization is simply zero. Both
versions erase the off-line content. A stationary difference kernel
\(K(u,v)=k(u-v)\) that compresses only to this scalar multiplier is
therefore \textbf{blind}, not RH-equivalent.

For the meromorphic continuation

\[
\mathcal F_-(s)=(1-s)\Psi_X(q(s))
\tag{159.5a}
\]

the exact completed zero form uses reflection and one representative per
functional-equation orbit:

\[
Q_X(c)=\sum_{[\rho]}m_\rho\,\mathcal F_-(\rho)\mathcal F_-(1-\rho)
\tag{159.5}
\]

whereas the corresponding Hermitian comparison is normalized as

\[
\frac12\sum_{\rho\ \mathrm{all}}m_\rho\,
|\mathcal F_-(\rho)|^2
\tag{159.6}
\]

For one fixed real-coefficient test, pointwise equality of the
reflection product with the Hermitian square holds exactly when

\[
\boxed{\Re\rho=\frac12\quad\text{or}\quad\Psi_X(q_\rho)=0}
\tag{159.7}
\]

A multinode test can annihilate an off-line orbit, as the exact
synthetic example in reader (20.6b) demonstrates. A separating family is
needed for the RH equivalence. For the nonzero one-node carrier, no such
annihilator exists, and equality is strict off the line. This names the
missing operation precisely:

\[
\boxed{\begin{gathered}
\text{Create the identification}\\
s\mapsto1-s\quad\text{with}\quad s\mapsto\bar s\\
\text{without assuming RH}
\end{gathered}}
\tag{159.8}
\]

The construction must retain reflection information lost by the scalar
boundary compression (159.4). Internal nonstationarity is one design
strategy, not a proved universal classification. Final translation
covariance is compatible with the exact Weil form. Section 161 gives a
separate precise obstruction for regular measurable multiplier norms; it
leaves singular spectral measures and form limits open.

A finite conjugation-symmetric feature family can turn an aggregate
holomorphic-square sum into a tensor norm; Lamzouri's Proposition 2.1
{[}66{]} provides a concrete example. This does not make individual
off-line summands nonnegative, and it does not identify the resulting
finite form with the completed source form. Its zeta application
controls asymptotic proportions of zeros. The present fixed-scale
criterion instead requires an independent bound through unbounded
degree, with multiplicities and the complete source retained.

\subsection{159.1 What the sesquilinear language
contributes}\label{what-the-sesquilinear-language-contributes}

Use the same convention as Reading Volume §R18A and technical §161:
\(\mathfrak h(u,v)\) is conjugate-linear in its first argument and
linear in its second. It is Hermitian when

\[
\mathfrak h(v,u)=\overline{\mathfrak h(u,v)}
\tag{159.9}
\]

and positive when \(\mathfrak h(u,u)\ge0\). In the complex setting the
diagonal quadratic form determines a Hermitian form through
polarization:

\[
\boxed{
\begin{aligned}
4\mathfrak h(u,v)
={}&Q(u+v)-Q(u-v)\\
&-iQ(u+iv)+iQ(u-iv)
\qquad Q(w)=\mathfrak h(w,w)
\end{aligned}}
\tag{159.10}
\]

This elementary fact is useful because it makes the type mismatch
visible. The reflected zero pairing

\[
\mathfrak b_{\mathrm{ref}}(F,G)
=\frac12\sum_{\rho\ \mathrm{all}}m_\rho\,
\mathcal F(\rho)\mathcal G(1-\rho)
\tag{159.11}
\]

is complex \textbf{bilinear} as printed; no conjugate appears. A
Hermitian extension in the declared convention conjugates its first test
factor and keeps its second factor linear. On RH, real-coefficient test
functions obey

\[
1-\rho=\bar\rho,\qquad
\mathcal G(\bar\rho)=\overline{\mathcal G(\rho)}
\tag{159.12}
\]

so (159.11), restricted to real-coefficient tests, agrees with the
Hermitian pairing. A complex Hermitian extension conjugates the
coefficients in the first slot under our convention; a complex bilinear
form does not become sesquilinear by a change of notation. Off RH the
reflection/conjugation agreement is not available for a separating
family. The lecture material on sesquilinear forms therefore helps by
identifying the exact algebraic type a successful lift must have; it
does not itself construct the missing lift. Any proposed polarization
argument must first exhibit a genuinely positive diagonal \(Q\) on an
enlarged two-sheet space. Applying (159.10) directly to (159.11) would
merely conceal the unproved step.

\section{160. The reflection-contour normal
form}\label{the-reflection-contour-normal-form}

For a finite positive node set \(X=(x_1,\ldots,x_N)\) and real
coefficients \(c_i\), define

\[
\Psi_X(q)=\sum_{i=1}^N\frac{c_i}{x_i+q},\qquad
q(s)=s(1-s)
\tag{160.1}
\]

and reserve

\[
\boxed{
\mathscr L_F(s)
:=q(s)\Psi_X(q(s))^2}
\tag{160.2}
\]

This is PP205's test symbol \(\Lambda(s)\) under a collision-free name.
With \(s=\sigma+it\), its argument is explicitly

\[
q(s)
=\sigma(1-\sigma)+t^2+i\,t(1-2\sigma)
\tag{160.3}
\]

Functional-equation symmetry gives

\[
\mathscr L_F(1-s)=\mathscr L_F(s)
\qquad
\Phi(1-s)=-\Phi(s)
\tag{160.4}
\]

so \(\mathscr L_F\Phi\) is antisymmetric under reflection. Put

\[
\sigma_*=\min_i\frac{1+\sqrt{1+4x_i}}2>1
\tag{160.5}
\]

For every \(1<\sigma_0<\sigma_*\), the exact single-contour formula is

\[
\boxed{
\begin{aligned}
Q_X(c)
&=\frac1{2\pi}\int_{\mathbb R}
\mathscr L_F(\sigma_0+it)\Phi(\sigma_0+it)\,dt\\
&=\frac1{2\pi i}\int_{(\sigma_0)}
\mathscr L_F(s)\Phi(s)\,ds
\end{aligned}}
\tag{160.6}
\]

On this line \(\sigma_0>1\), the source expansion of \(\Phi\) is
absolutely convergent. The poles of \(\Psi_X(q(s))\) lie outside the
contour window, the horizontal edges vanish like \(O(\log|t|/|t|^2)\) or
faster under moment cancellation, and the enclosed poles are the
nontrivial zeros. Reflection turns the two vertical integrals into twice
(160.6), while the residue sum is twice the folded form.

At the terminal line \(s=1/2+it\),

\[
\boxed{
\mathscr L_F(1/2+it)=|\widehat F(t)|^2\ge0,\qquad
\Phi(1/2+it)\in i\mathbb R}
\tag{160.7}
\]

The positive object and the pairing are exactly orthogonal there. With
the poles of \(\Phi\) treated by symmetric indentation and symmetric
truncation, the critical-line statement is
\(\operatorname{PV}I_F(1/2)=0\); the integrand is not zero. The RH
content lives in the contour shift and its residues, not in a pointwise
boundary sign.

Define

\[
I_F(\sigma)=\frac1{2\pi i}\int_{(\sigma)}
\mathscr L_F(s)\Phi(s)\,ds
\tag{160.8}
\]

On any open interval containing no zero abscissae or test poles, \(I_F\)
is constant. On its admissible contour lines it is odd about \(1/2\).
For endpoint lines avoiding poles and shifts crossing no test poles, the
residue argument gives the exact jump law

\[
\boxed{
I_F(\sigma_2)-I_F(\sigma_1)
=\sum_{\sigma_1<\beta_\rho<\sigma_2}
m_\rho\mathscr L_F(\rho)}
\tag{160.9}
\]

This is a classical contour/residue mechanism in project-native
packaging. It is not yet a positivity theorem: \(\mathscr L_F(\rho)\) is
a holomorphic square at an off-line zero and has no manifest sign.

The contour figure in Reading Volume §R20 illustrates this jump law.

Three fast screening tests now govern every candidate:

\begin{enumerate}
\def\labelenumi{\arabic{enumi}.}
\tightlist
\item
  \textbf{Reflection test.} If it produces \(|\mathcal F_X(\rho)|^2\) in
  place of \(\mathcal F_X(\rho)\mathcal F_X(1-\rho)\), it must prove the
  required equality for a separating family; a fixed test can have
  accidental zeros.
\item
  \textbf{Stationary-loss test.} If the internal construction reduces to
  the boundary multiplier (159.4), it has erased every off-line pair.
\item
  \textbf{Terminal test.} If it claims positivity from the integrand at
  \(\sigma=1/2\), it claims a sign for a contour integral forced to
  vanish.
\end{enumerate}

The correct next object is a positive form before compression whose
compression reproduces (160.6) on \(\sigma_0>1\). Section 161 adds a
regularity test to this completion-coupled, two-sheet target: an
ordinary strong tensor-norm limit cannot supply the required spectral
owner.

\section{161. Regular translation energy: an
obstruction}\label{regular-translation-energy-an-obstruction}

The varying-endpoint Gram in Reading Volume §R18B is a positive
auxiliary object. Its natural tensor lift, however, has a further
obstruction. The distinction is between an ordinary convergent vector
series and a singular limit of forms.

Use the inner product conjugate-linear in the first slot. Put
\(L_n=\log n\), let \(\phi_n\) be (EG.1), and take a finite index set
\(I\subset\{2,3,\ldots\}\). Define

\[
\mathcal R_{b,I}F=\sum_{n\in I}b_n\phi_n\otimes(I-\tau_{L_n})F
\qquad g_F(v)=\langle F,\tau_vF\rangle
\tag{161.1}
\]

Translation unitarity gives the complete finite expansion

\[
\begin{aligned}
\|\mathcal R_{b,I}F\|^2
={}&\sum_{m,n\in I}\overline b_m b_nK_{mn}
\big[g_F(0)-g_F(L_n)-g_F(-L_m)\\
&\hspace{43mm}+g_F(L_n-L_m)\big]
\end{aligned}
\tag{161.2}
\]

The final shifts are \(\log(n/m)\). Pairs with the same reduced ratio
must be grouped before a coefficient can be declared noncancellable. For
a reduced ratio \(a/b\), those pairs are \((m,n)=(\ell b,\ell a)\). In
an infinite construction the four displayed channels must not be
separated without an appropriate convergence argument.

\begin{quote}
\textbf{Theorem 161.1 --- regular translation-energy obstruction.} Fix
the coefficients and auxiliary vectors independently of the input test
and its node. Suppose (161.1) is a finite sum, or its partial sums
converge strongly in the fixed ordinary Hilbert space
\(\mathcal H\otimes L^2(\mathbb R)\) on the resolvent tests. Here
\(\mathcal H=L^2(0,\infty)\) for (EG.1). Then its norm cannot satisfy
\(\|\mathcal R_bf_x\|^2=W_1(x)\) for every \(x>0\). In particular, it
cannot reproduce the completed source form on every rational-resolvent
test.
\end{quote}

\textbf{Proof.} Under the forward transform \(e^{-itu}\), translation by
\(L_n\) acts as multiplication by \(e^{itL_n}\). For finite \(I\) set

\[
A_I(t)=\sum_{n\in I}b_n\phi_n(1-e^{itL_n})
\qquad M_I(t)=\|A_I(t)\|_{\mathcal H}^2\ge0
\]

Plancherel gives

\[
\|\mathcal R_{b,I}F\|^2
=\frac1{2\pi}\int_{\mathbb R}|\widehat F(t)|^2M_I(t)dt
\tag{161.3}
\]

The same regular multiplier representation survives strong convergence.
Fix \(x_0>0\). The transform

\[
\widehat f_{x_0}(t)=\frac{1/2-it}{t^2+x_0+1/4}
\]

does not vanish on the real line. Strong convergence says that
\(A_I\widehat f_{x_0}\) converges in \(L^2(\mathbb R;\mathcal H)\).
Divide its measurable limit by \(\widehat f_{x_0}\) to define \(A(t)\).
For every other \(x>0\), the ratio \(\widehat f_x/\widehat f_{x_0}\) is
bounded, so the same multiplier works on all these tests. Thus (161.3)
holds with \(M(t)=\|A(t)\|^2\ge0\).

Push forward \(M(t)dt/(2\pi)\) under \(q=t^2+1/4\) and call the
resulting positive measure \(\nu\). It is absolutely continuous on
\((1/4,\infty)\) and has no atom at \(1/4\). Equality of the one-node
diagonal would give

\[
W_1(x)=\int_{[1/4,\infty)}\frac{q}{(x+q)^2}\,d\nu(q)
\tag{161.4}
\]

Integrate from \(0\) to \(x\). Tonelli applies to the nonnegative
integrand; \(W_1=(xS_\xi)'\) and \(xS_\xi(x)\to0\) at the origin.
Consequently

\[
\boxed{S_\xi(x)=\int_{[1/4,\infty)}\frac{d\nu(q)}{x+q}}
\tag{161.5}
\]

This integral is finite: its kernel is comparable to the one in (161.4),
since \(q\ge1/4\).

The completed fold has a different analytic requirement. Reflection
removes the square-root branch in

\[
\Phi_0(x)=\frac{\xi(s_x)}{\xi(1)}
\qquad S_\xi(x)=\frac{\Phi_0'(x)}{\Phi_0(x)}
\tag{161.6}
\]

Thus \(\Phi_0\) is entire with real coefficients, and \(S_\xi\) is
meromorphic in the full plane, with discrete poles and real values on
real intervals away from them. A positive Stieltjes measure with this
transform must be purely atomic. Indeed, on any compact interval whose
reflected interval contains no pole, the two boundary values of
\(S_\xi\) agree and are real. Stieltjes inversion gives zero measure
there. Equivalently, integrate its Poisson-kernel form

\[
-\frac1\pi\Im S_\xi(-y+i\varepsilon)
=\int\frac{\varepsilon\,d\nu(q)}
{\pi[(q-y)^2+\varepsilon^2]}
\tag{161.7}
\]

over such an interval and let \(\varepsilon\downarrow0\). Local
finiteness and the weighted tail integrability supplied by (161.5)
justify this inversion. Hence \(\nu\) is supported on the discrete
reflected pole set. Any off-cut pole would already contradict analytic
continuation of (161.5).

The measure is thus both absolutely continuous and discretely supported,
hence zero. That would imply \(S_\xi=0\), contradicting
\(S_\xi(0)=\lambda_1>0\). No RH assumption was used. \(\square\)

\begin{figure}
\centering
\includegraphics{figures/v26c_measure_obstruction.png}
\caption{The regular-energy obstruction compares two measure types. An
ordinary measurable multiplier produces an absolutely continuous
measure; a positive Stieltjes representation of the completed
meromorphic resolvent would require an atomic measure. Their only common
measure is zero. Singular form limits lie outside the stated class.}
\end{figure}

\textbf{Comparison with the completed source.} Section 154's full
Gamma-minus-prime form has exactly the one-node diagonal \(W_1\) used
above, so the candidate's regular Fourier density cannot reproduce it,
even if mixed rational-log translations are reorganized. Its
integer-label Gram may be nonstationary, yet its action on \(F\) reduces
to a scalar Fourier density.

\textbf{Scope and surviving route.} This theorem excludes finite sums,
strongly convergent series in the fixed ordinary tensor norm, and more
generally regular measurable vector-multiplier realizations with finite
test energies. It does not exclude singular spectral measures, a
different source-defined norm, a discrete realization, or weak limits of
quadratic forms that fail strong vector convergence. Absolutely
continuous measures can converge weakly to atoms. Such a construction
must specify its domain and topology, retain the Gamma/prime
compensation, and prove exact equality on the full resolvent span.
Sections 163--165 construct a source-defined discrete sequence and prove
a sufficient variation criterion; that bound remains unproved.

The new requirement is therefore an atomic positive spectral owner, or a
rigorously justified singular form limit, derived from the completed
source. The trace target (23.3) and the all-node source sign remain
open.

\Needspace{11\baselineskip}

\chapter{Technical chapter XIV. Completed Bernstein rows and the
fixed-scale
energy}\label{technical-chapter-xiv.-completed-bernstein-rows-and-the-fixed-scale-energy}

Source derivatives determine the rows and their exact Pascal/Laguerre
norm. The growth theorem leads to a fixed-scale prime estimate; the
elementary/Gamma part is controlled.

\section{163. Row mass, refinement, and the generating
function}\label{row-mass-refinement-and-the-generating-function}

Fix \(a>0\) and abbreviate \(S=S_\xi\). All zero sums below are folded
sums, with multiplicities and conjugate folded values retained. Since
\(\Re q_\rho>0\) and the zero count is \(O(T\log T)\),
\(\sum_{[\rho]}m_\rho/|a+q_\rho|<\infty\). The existing Widder and
shifted-zeta identities give

\[
\begin{aligned}
F_{n,k}(x)&=(-1)^nD^{n+k}[x^kS(x)]
=(n+k)!\sum_{[\rho]}\frac{m_\rho q_\rho^k}{(x+q_\rho)^{n+k+1}}\\
Z_r(x)&=\frac{(-1)^{r-1}}{(r-1)!}S^{(r-1)}(x)
\end{aligned}
\tag{163.1}
\]

Substitution into (BR.1) proves its Bernstein expansion. For each
complex \(\theta\), the binomial theorem and degree-elevation identity
imply

\[
\begin{aligned}
\sum_{j=0}^NB_{N,j}&=S(a)\\
B_{N,j}&=\frac{N+1-j}{N+1}B_{N+1,j}
+\frac{j+1}{N+1}B_{N+1,j+1}
\end{aligned}
\tag{163.2}
\]

Taking absolute values, then summing \(j\), proves
\(S(a)\le\mathcal V_N(a)\le\mathcal V_{N+1}(a)\): each coefficient of
the finer row is counted with total weight one. A negative entry at one
degree forces a negative entry at every subsequent degree. This
statement concerns a finite negative witness, not a method that
guarantees one.

For the row polynomial \(P_N(z)=\sum_jB_{N,j}z^j\), finite Leibniz
expansion gives an independent formula entirely in source derivatives:

\[
\begin{aligned}
P_N(z)&=\left.\frac1{N!}D_x^N[(x-az)^NS(x)]\right|_{x=a}\\
&=\sum_{r=0}^N\binom Nr\frac{S^{(r)}(a)}{r!}
a^r(1-z)^r
\end{aligned}
\tag{163.3}
\]

Summing for sufficiently small \(u\), and then continuing
meromorphically, produces

\[
\begin{aligned}
\sum_{N\ge0}P_N(z)u^N
&=\frac1{1-u}S\!\left(a\frac{1-uz}{1-u}\right)\\
&=\sum_{[\rho]}\frac{m_\rho}{a+q_\rho-u(q_\rho+az)}
\end{aligned}
\tag{163.4}
\]

The first equality follows also from
\(\sum_{N\ge r}\binom Nr u^N=u^r/(1-u)^{r+1}\). The completed source,
including its prime-dependent part, occurs in (163.3)--(163.4). Its
asymptotic coefficients alone do not determine these rows.

For later use, the falling moments satisfy

\[
\sum_{j=0}^N(j)_rB_{N,j}(a)
=(N)_r a^rZ_{r+1}(a)\quad(0\le r\le N)
\tag{163.5}
\]

This is the binomial falling-moment identity, valid algebraically for
complex \(\theta\). At \(N\ge1\), Stirling's finite change of basis
between \(j^r\) and \((j)_\ell\) then shows

\[
\sum_{j=0}^N(j/N)^rB_{N,j}(a)\longrightarrow a^rZ_{r+1}(a)
\quad\text{for each fixed }r\ge0
\tag{163.6}
\]

This moment convergence is unconditional. Weak convergence of the signed
measures requires a further bound on total variation.

\section{163A. The exact Bernstein--Laguerre
energy}\label{a.-the-exact-bernsteinlaguerre-energy}

The row coefficients admit a norm that keeps the completed source inside
one polynomial. Let \(L_j^{(0)}(t)=\sum_{r=0}^j\binom jr(-t)^r/r!\) be
the ordinary Laguerre polynomials. Their coefficients are precisely the
signed grid entries \(c_{j,r}\) divided by \(r!\), including \(L_0=1\).
Define

\[
\mathcal Q_N(t;a)=\sum_{j=0}^N B_{N,j}(a)L_j^{(0)}(t)
=\sum_{r=0}^N\binom Nr\frac{a^rS^{(r)}(a)}{(r!)^2}t^r
\tag{163A.1}
\]

To prove the second equality, expand each \(L_j\) and use (163.5) in the
form \(\sum_j\binom jrB_{N,j}=\binom Nr a^rZ_{r+1}\). The sign in
\(Z_{r+1}=(-1)^rS^{(r)}/r!\) cancels the Laguerre sign. Both factorials
remain. Independently, finite coefficient comparison gives the classical
polynomial identity

\[
L_N^{(0)}(\theta t)=\sum_{j=0}^N\binom Nj
\theta^j(1-\theta)^{N-j}L_j^{(0)}(t)
\tag{163A.2}
\]

It holds for complex \(\theta\). Thus (163A.1) also equals
\(\sum_q m_q L_N^{(0)}(at/(a+q))/(a+q)\), with equal folded values
combined into \(m_q\). Every such sum is normally convergent for fixed
degree, by the summable resolvent weights.

Rodrigues' formula \(L_j=e^tD^j(e^{-t}t^j)/j!\) and integration by parts
give \(\int_0^\infty e^{-t}L_jL_k\,dt=\delta_{jk}\). The boundary terms
vanish; when \(j>k\), the \(j\)-th derivative of \(L_k\) vanishes, and
when \(j=k\), its leading coefficient gives norm one. Together with
circle Parseval this proves

\[
\boxed{\begin{aligned}
\mathscr E_N(a)&:=\int_0^\infty e^{-t}|\mathcal Q_N(t;a)|^2\,dt
=\sum_{j=0}^N B_{N,j}(a)^2\\
&=\frac1{2\pi}\int_{-\pi}^{\pi}|P_N(e^{i\vartheta};a)|^2\,d\vartheta
\end{aligned}}
\tag{163A.3}
\]

The finite linear map \(z^j\mapsto L_j^{(0)}(t)\) is therefore an
isometry between these two polynomial spaces. It sends \((1-z)^r\) to
\(t^r/r!\); this is the exact connection between the source polynomial
in (163.3) and the second factorial in (163A.1). No equality with the
earlier Laguerre flux hierarchy is asserted.

The elementary coefficient-norm inequalities give

\[
\frac{S(a)^2}{N+1}\le\mathscr E_N(a),\qquad
\sqrt{\mathscr E_N(a)}\le\mathcal V_N(a)
\le\sqrt{N+1}\sqrt{\mathscr E_N(a)}
\tag{163A.4}
\]

Consequently (BR.4) is equivalent to the following estimate at the same
fixed \(a\): for every \(\varepsilon>0\) there is a finite
\(C_{a,\varepsilon}\) such that
\(\mathscr E_N(a)\le C_{a,\varepsilon}(1+\varepsilon)^{2N}\) for every
\(N\). The polynomial factor \(N+1\) is absorbed by first using a
smaller positive epsilon. Positivity of the norm is automatic; its upper
growth bound is the missing assertion.

A finite algebraic form, useful for checking normalization, is

\[
\mathscr E_N(a)=\sum_{r,s=0}^N\binom Nr\binom Ns
\frac{(r+s)!}{(r!)^2(s!)^2}
a^{r+s}S^{(r)}(a)S^{(s)}(a)
\tag{163A.5}
\]

\subsection{The triangular row inside the energy
matrix}\label{the-triangular-row-inside-the-energy-matrix}

The factorial coefficients in (163A.5) are the entries of a symmetric
Pascal matrix. For a fixed \(a>0\), collect the real completed-source
derivatives into a column \(d^{(N)}\) and define

\[
\begin{aligned}
d_r^{(N)}(a)&=\binom Nr\frac{a^rS^{(r)}(a)}{r!}
\qquad(0\le r\le N)\\
(M_N^{\rm Pas})_{r,s}&=\binom{r+s}{r},\qquad
(C_N)_{r,j}=\begin{cases}\binom rj&j\le r\\0&j>r\end{cases}
\end{aligned}
\tag{163A.6}
\]

In the signed grid notation, \((C_N)_{r,j}=(-1)^j c_{r,j}\) and
\((M_N^{\rm Pas})_{r,s}=(-1)^r c_{r+s,r}\). The figure shows cells with
row index at most 17; these matrix identities use the same array for
arbitrary finite indices.

The derivatives are real because the completed folded sum retains
conjugate pairs. Direct regrouping of (163A.5), followed by
Vandermonde's identity, gives

\[
\boxed{\begin{aligned}
\mathscr E_N(a)&=(d^{(N)})^{\mathsf T}M_N^{\rm Pas}d^{(N)}\\
M_N^{\rm Pas}&=C_NC_N^{\mathsf T},\qquad
\det M_N^{\rm Pas}=1
\end{aligned}}
\tag{163A.7}
\]

Indeed, \(\sum_j\binom rj\binom sj=\binom{r+s}{r}\) and \(C_N\) is
triangular with unit diagonal. This classical Pascal factorization is
described in {[}65{]}. Its role here is to identify the exact matrix
already present in the completed-source energy. The first four rows and
columns make the earlier triangular numbers visible:

\[
M_3^{\rm Pas}=\begin{pmatrix}
1&1&1&1\\
1&2&3&4\\
1&3&6&10\\
1&4&10&20
\end{pmatrix},\qquad
\boxed{(M_N^{\rm Pas})_{2,s}
=\binom{s+2}{2}=T_{s+1}}
\tag{163A.8}
\]

The last identity uses zero-based indices, \(N\ge2\), and
\(0\le s\le N\). Thus the third row and column are triangular; the
neighboring rows contain the constant, integer, and tetrahedral
sequences. These are the Pascal coefficients of §171, now acting as a
Gram matrix on source derivatives.

There is also an exact reflection of polynomial coefficients. Set
\((U_N)_{j,r}=(-1)^j\binom rj\) for \(j\le r\) and zero otherwise. The
identity \(P_N(z;a)=\sum_r d_r^{(N)}(a)(1-z)^r\) in (163.3) gives

\[
\boxed{B_N=U_Nd^{(N)},\qquad U_N^2=I,\qquad
U_N^{\mathsf T}U_N=M_N^{\rm Pas}}
\tag{163A.9}
\]

Here \(B_N=(B_{N,0},\ldots,B_{N,N})^{\mathsf T}\) and
\((U_N)_{j,r}=c_{r,j}\) is the transposed signed grid. The involution is
the substitution \(z\mapsto1-z\) in the polynomial variable: applying it
twice returns the original polynomial. This variable is independent of
the source coordinate \(s\) and of the Cayley coordinate in Reading
Volume §R7. For \(N\ge1\) the change of basis is involutive but not
orthogonal in the ordinary coefficient norm; its Gram matrix is
precisely \(M_N^{\rm Pas}\).

The smallest nontrivial case reconnects this matrix to the reciprocal
sequence of §171:

\[
M_1^{\rm Pas}=\begin{pmatrix}1&1\\1&2\end{pmatrix},\qquad
\lambda_\pm=\frac{3\pm\sqrt5}{2},\qquad
\operatorname{tr}((M_1^{\rm Pas})^m)=\mathcal P_m(3)
\tag{163A.10}
\]

Indeed \(\lambda_+\lambda_-=1\) and \(\lambda_++\lambda_-=3\). The trace
sequence is \(2,3,7,18,47,123,\ldots\), with the same recurrence as
(RP.2). The calligraphic \(\mathcal P_m\) is the existing reciprocal
polynomial; \(P_N\) continues to denote the source polynomial.

For general \(N\), (163A.9) gives

\[
(M_N^{\rm Pas})^{-1}=U_N M_N^{\rm Pas}U_N,\qquad
\lambda_j\lambda_{N-j}=1\quad(0\le j\le N),
\tag{163A.11}
\]

where \(\lambda_0\le\cdots\le\lambda_N\) are the positive eigenvalues.
The first identity follows by multiplying \(U_N^{\mathsf T}U_N\) and
using \(U_N^2=I\); similarity then proves reciprocal pairing. The final
diagonal entry and the Rayleigh bound imply

\[
\operatorname{cond}_2(M_N^{\rm Pas})
=\lambda_N^2\ge\binom{2N}{N}^{\!2}
\tag{163A.12}
\]

Thus involution does not make this change of coefficients orthogonal or
well-conditioned at large degree. This is the matrix counterpart of
keeping the reflected coefficients and their attached basis functions
distinct in (SW.4c).

The Pascal matrix is positive definite for every \(N\), regardless of
the zero locations. Its positivity explains the norm identity, while the
degree-dependent vector \(d^{(N)}\) carries the source information
needed for an upper growth estimate. The diagonal entry
\(\binom{2N}{N}\) itself grows exponentially, so the factorization alone
supplies no such bound.

In particular \(\mathscr E_1=S(a)^2+2aS(a)S'(a)+2a^2S'(a)^2\). The
symbol \(\mathscr E_N\) denotes a polynomial energy. The off-line heat
defect \(\mathcal E_H\) in (HG.2) is a different quantity.

\begin{figure}
\centering
\includegraphics{figures/v27b_energy_isometry.png}
\caption{The completed row has two exact polynomial realizations with
the same squared norm. The open assertion concerns the growth of that
norm at one fixed source scale.}
\end{figure}

\section{164. Exact growth and the heat
comparison}\label{exact-growth-and-the-heat-comparison}

Fix \(a>0\). For a folded point \(q\) put

\[
R_q(a)=\frac{a+|q|}{|a+q|},\qquad
\mathcal R(a)=\sup_{[\rho]}R_{q_\rho}(a)\ge1.
\tag{164.1}
\]

The completed Bernstein row of §163 has variation \(\mathcal V_N(a)\)
and Laguerre energy \(\mathscr E_N(a)\). A verified zero height \(H\)
gives the uniform finite-height estimate

\[
\mathcal V_N(a)\le S(a)\sqrt{1+H^{-2}}
\left(1+\frac{a}{(a+H^2)^2}\right)^{N/2}.
\tag{164.4b}
\]

This is extremely sharp at fixed \(a\) but still has one fixed
exponential factor greater than one, so it is not the
arbitrary-\(\varepsilon\) source bound.

\subsection{Theorem 164.2. Energy asymptotics and ordinary
limits}\label{theorem-164.2.-energy-asymptotics-and-ordinary-limits}

There is \(0<\mathcal K(a)<\infty\) such that

\[
\boxed{\mathscr E_N(a)\sim
\mathcal K(a)\frac{\mathcal R(a)^{2N}}{\sqrt N}.}
\tag{164.5}
\]

Consequently

\[
\boxed{\begin{aligned}
\lim_{N\to\infty}\left(\frac{\mathcal V_N(a)}{S(a)}\right)^{1/N}
&=\lim_{N\to\infty}\left(\frac{\mathscr E_N(a)}{S(a)^2}\right)^{1/(2N)}
=\mathcal R(a),\\
\lim_{N\to\infty}\frac{\mathscr E_{N+1}(a)}{\mathscr E_N(a)}
&=\mathcal R(a)^2.
\end{aligned}}
\tag{164.6}
\]

The proof reduces to collisions of Bernstein rows. For \(0<p<1\),

\[
I_N(p)=\sum_j\binom Nj^2p^{2j}(1-p)^{2N-2j}
=\frac1{2\pi}\int_{-\pi}^{\pi}
[1-4p(1-p)\sin^2(\vartheta/2)]^N d\vartheta,
\tag{164.7}
\]

so Laplace scaling at \(\vartheta=0\) gives

\[
I_N(p)\sim\frac1{2\sqrt{\pi Np(1-p)}},
\qquad I_N(p)\le\frac{C}{\sqrt{Np(1-p)}}.
\tag{164.8}
\]

Distinct centers have a strictly smaller exponential collision rate. The
finitely many maximizing folded orbits therefore determine the leading
term (164.5), with equal folded values aggregated before squaring their
weights.

\subsection{Corollary 164.3. Bounded energy at one fixed
scale}\label{corollary-164.3.-bounded-energy-at-one-fixed-scale}

Since \(\mathcal R(a)=1\) exactly when every folded zero is positive
real,

\[
\boxed{\mathrm{RH}\iff
\mathscr E_N[S_\xi](a)\to0
\iff\sup_N\mathscr E_N[S_\xi](a)<\infty.}
\tag{164.13}
\]

A bounded subsequence at unbounded degrees also suffices because the
ordinary root limit exists. Moving the basepoint with \(N\) is a
different limit and can hide the obstruction. The longer
collision-kernel and multiplicity proof is preserved in earlier
editions.

\section{165. Bounded variation and the atomic source
representation}\label{bounded-variation-and-the-atomic-source-representation}

\begin{quote}
\textbf{Theorem 165.1.} If \(\sup_N\mathcal V_N(a)<\infty\) at one
prescribed \(a>0\), then the measures
\(\eta_N=\sum_jB_{N,j}(a)\delta_{j/N}\), \(N\ge1\), converge weakly on
\([0,1]\) to the unique positive measure with moments \(a^rZ_{r+1}(a)\).
Its Stieltjes transform is the completed \(S_\xi\). The hypothesis is
unproved; the result is an implication from a source bound.
\end{quote}

\textbf{Proof.} A uniform total-variation bound makes these signed
measures weak-* relatively compact as functionals on \(C[0,1]\). Any
subsequential limit has the moments (163.6). Polynomial approximation
implies uniqueness, so the whole sequence converges. Initially for
\(|x-a|<a\), the geometric series in the following integral is uniformly
convergent:

\[
S(x)=\int_{[0,1]}\frac{d\eta(t)}{1+(x/a-1)t}
\tag{165.1}
\]

Its Taylor coefficients at \(a\) are exactly those of \(S\). The right
side is analytic on \(\mathbb C\setminus(-\infty,0]\), even for signed
finite \(\eta\). Analytic continuation of (165.1) excludes every off-cut
pole of \(S\). Since its poles are \(-q_\rho\), RH follows. Under RH the
explicit positive measure

\[
\eta=\sum_{[\rho]}\frac{m_\rho}{a+q_\rho}
\delta_{a/(a+q_\rho)}
\tag{165.2}
\]

has precisely the same moments and therefore is that unique limit. No
positivity of the original signed rows was assumed. \(\square\)

The correct matrix normalization follows from \(h=xS\). Put
\(A_x(t)=1+(x/a-1)t\). Its divided differences, including the diagonal
by continuity, are

\[
L_h(x,y)=\int_{[0,1]}\frac{1-t}{A_x(t)A_y(t)}\,d\eta(t)
\tag{165.3}
\]

For arbitrary complex coefficients this gives the Hermitian Gram norm
\(\int(1-t)|\sum_i c_i/A_{x_i}(t)|^2d\eta(t)\ge0\). Thus the
representation recovers all nodes at once. For the finite-vector
normalization of §R10A, weight \(\eta\) by \(1/(1-t)\) and push forward
under \(d=t/[a(1-t)]\). The resulting measure \(\nu\) satisfies

\[
S(x)=\int\frac{d\nu(d)}{1+xd},\qquad
\nu([0,\infty))=S(0)=\lambda_1
\tag{165.4}
\]

There is no extra numerator \(d\) in this finite-vector formula. The
discrete measure (165.2) has no atom at \(t=1\); its weighted mass in
(165.4) is finite by \(S(0)<\infty\).

\textbf{Completion before variation.} The elementary source term \(1/x\)
alone produces \(B_{N,j}=a^{-1}\delta_{j,N}\), a persistent endpoint
atom at \(t=1\). The completed function removes the corresponding pole
at \(x=0\). The exact row mass, endpoint cancellation, and limiting
positive measure require the completed combination. Section 166A proves
when separate channel bounds can nevertheless establish the qualitative
subexponential energy criterion. The source basepoint \(a=2=2T_1\) gives
\(s_a=2\), an available Euler-series location; it does not change this
requirement or prove the variation bound.

\section{166. What topology remains after the regular-energy
obstruction}\label{what-topology-remains-after-the-regular-energy-obstruction}

Section 161 excludes a regular measurable multiplier in the fixed
ordinary \(L^2\) norm. The discrete rows above are instead measures on a
compact interval. A weak limit of measures is not a strong vector limit
in that old norm. This difference is mathematical, not terminological.

For example, consider a sampling quadratic form
\(Q(f)=\sum_jc_j|f(t_j)|^2\) with \(c_j>0\), discrete support, and a
domain containing compactly supported smooth functions. Choose an
isolated atom \(t_0\) and shrinking smooth bumps with \(f_n(t_0)=1\) and
support avoiding the other atoms. Then

\[
\|f_n\|_{L^2(dt)}\to0,\qquad Q(f_n-f_m)=0
\qquad Q(f_n)=c_0>0
\tag{166.1}
\]

This violates closability in ordinary \(L^2(dt)\). It does not exclude
the same sampling form on a different domain with a norm controlling
point values. A surviving source realization must state that domain and
norm, and prove its source identity. Referring generally to a singular
limit does not supply those data.

The existing dilation criterion and higher-moment support theorem give a
second exact crosswalk. Write \(\beta_*:=\sup_\rho\Re\rho\) and
\(r_\xi:=\sup_\rho|\Re\rho-1/2|\). Functional-equation symmetry implies
\(2\beta_*=1+2r_\xi\). The criterion of §143 and (GM.6) therefore give
(HG.5). The normalized \(L^{2m}\) norms increase to this support radius,
so finite moments supply lower bounds. A source proof reducing the
dilation threshold to \(1\) remains missing.

The regular-energy obstruction, the discrete-measure proposal, and the
dilation limit can therefore be read together: they specify which object
must converge, in which topology, and which sign estimate has not yet
been established. They do not turn the calibrated auxiliary model into
the completed source.

\section{166A. The remaining norm estimate and the fixed source
scale}\label{a.-the-remaining-norm-estimate-and-the-fixed-source-scale}

Extend (163A.1) linearly to \(\mathcal Q_N[f]\), with \(S\) replaced by
\(f\), and let
\(\mathscr E_N[f]=\|\mathcal Q_N[f]\|_{L^2(e^{-t}dt)}^2\). Use the same
replacement in (163.3) to define \(P_N[f]\). The exact source in §54
gives

\[
\mathcal Q_N[S]=a^{-1}L_N^{(0)}+\mathcal Q_N[G]-\mathcal Q_N[P]
\tag{166A.1}
\]

Indeed \(\mathcal Q_N[1/x]=a^{-1}L_N^{(0)}\) by direct differentiation.
The exact energy includes all cross terms. This identity concerns a
degree-dependent source polynomial; it does not assert the existence of
the fixed regular multiplier excluded in §161.

\subsection{A cut-plane estimate closes the Gamma
part}\label{a-cut-plane-estimate-closes-the-gamma-part}

\textbf{Lemma 166A.1.} If \(f\) is holomorphic on
\(D=\mathbb C\setminus(-\infty,0]\), then, at each fixed \(a>0\) and
each \(0<r<1\), a finite \(C_{f,a,r}\) satisfies

\[
\mathscr E_N[f](a)\le C_{f,a,r}^2 r^{-2N}\quad(N\ge0)
\tag{166A.2}
\]

\textbf{Proof.} The source-polynomial generating function is
\((1-u)^{-1}f(a(1-uz)/(1-u))\). For \(|u|<1\), \(|z|\le1\), its argument
lies in \(D\). In fact, an argument \(-v\) with \(v\ge0\) would require
\(u=(a+v)/(az+v)\), whose modulus is at least one; a zero denominator
cannot give a finite solution. For \(|u|=r\) the image is a compact
subset of \(D\). Its maximum, multiplied by \((1-r)^{-1}\), defines
\(C_{f,a,r}\). Uniform Cauchy coefficient estimates give
\(|P_N[f](z)|\le C_{f,a,r}r^{-N}\) on \(|z|=1\). Circle Parseval proves
(166A.2). \(\square\)

This lemma applies unconditionally to \(A(x)=1/x+G(x)\). On the
principal branch \(R(x)=\sqrt{1+4x}\) has positive real part in \(D\),
so \((1+R)/4\) has real part greater than \(1/4\). The digamma function
in \(G=[\psi((1+R)/4)-\log\pi]/(2R)\) has no pole there. Taking
\(r=(1+\varepsilon)^{-1}\) proves the required subexponential estimate
for the whole elementary/Gamma channel.

Consequently the triangle inequality, now using this proved channel
bound, gives an exact qualitative reduction:

\[
\boxed{\begin{aligned}
&\forall\varepsilon>0\ \exists C_{a,\varepsilon}<\infty:
\quad\mathscr E_N[S](a)\le C_{a,\varepsilon}(1+\varepsilon)^{2N}
\quad\forall N\\
&\hspace{12mm}\iff
\forall\varepsilon>0\ \exists C'_{a,\varepsilon}<\infty:
\quad\mathscr E_N[P](a)\le C'_{a,\varepsilon}(1+\varepsilon)^{2N}
\quad\forall N
\end{aligned}}
\tag{166A.3}
\]

Use \(S=A-P\) in one direction and \(P=A-S\) in the other; constants
absorb the squared triangle bound. No positivity of \(A\) or \(P\) is
required. Channel separation is therefore legitimate for this growth
criterion after the Gamma estimate is proved. Exact row masses, endpoint
cancellation, the limiting measure, and sharper numerical bounds still
use the completed combination (166A.1).

The first unproved assertion is now the right-hand estimate in (166A.3)
for the full prime term, at a fixed basepoint such as \(a=2\). The
Dirichlet series for \(P\) and all its finite derivatives converge
there. That local convergence gives no uniform estimate at arbitrary
degree. For example, if \(|u|<r<1/3\), then \(|x-a|\le2ar/(1-r)<a\) and
the transformed source remains in \(\Re s>1\), where the Euler series is
available. Extending this argument to every \(r<1\) requires control
outside that Euler-safe region. One cannot assume \(P\) holomorphic
throughout \(D\) in order to apply the lemma: excluding its off-cut
poles is the unresolved zero-location statement.

Each finite prime sum separately satisfies the lemma. Its constants need
not be uniform in the prime cutoff, and the raw Euler partial sums do
not converge throughout the larger domain. Thus neither finite prime
truncation nor an exchange of the degree and cutoff limits proves
(166A.3). This locates the remaining analytic task without replacing it
by positivity of an auxiliary norm.

\subsection{A moving basepoint erases the
diagnostic}\label{a-moving-basepoint-erases-the-diagnostic}

For every nontrivial zero, \(q=A+ib\) satisfies \(b^2\le A\). Writing
\(v=|q|\) gives \(v-A=b^2/(v+A)\le1/2\); hence
\(R_q(a)^2-1\le a/(a+A)^2\le1/a\). Also

\[
\frac{|a+q|^{-1}}{\Re(a+q)^{-1}}
=\sqrt{1+\frac{b^2}{(a+A)^2}}
\le\sqrt{1+\frac1{4a}}
\tag{166A.4}
\]

The last bound uses \(A/(a+A)^2\le1/(4a)\). Sum the positive real parts
and apply (164.1) to obtain, without assuming RH,

\[
\boxed{1\le\frac{\mathcal V_N(a)}{S(a)}
\le\sqrt{1+\frac1{4a}}\left(1+\frac1a\right)^{N/2}}
\tag{166A.5}
\]

Therefore every sequence \(a_N\to\infty\) satisfies

\[
\lim_{N\to\infty}
\left(\frac{\mathcal V_N(a_N)}{S(a_N)}\right)^{1/N}=1
\qquad
\frac{\mathcal V_N(N)}{S(N)}\le e^{1/2}\sqrt{1+\frac1{4N}}
\quad(N\ge1)
\tag{166A.6}
\]

These conclusions are unconditional and cannot prove RH. As in the
fixed-rung large-\(x\) theorem of §55, changing the source scale can
hide the unbounded-order question. A sparse sequence of degrees at one
fixed \(a\) is allowed by (164.6); sending \(a\) to infinity with degree
is a different limit. The required fixed-scale source norm estimate
remains open.

\section{166B. Gamma Stieltjes repair and the bounded full-prime
criterion}\label{b.-gamma-stieltjes-repair-and-the-bounded-full-prime-criterion}

The Gamma-pole analysis gives a shorter derivation of the Gamma measure
used in this criterion. Put \(R_x=\sqrt{1+4x}\), \(s_x=(1+R_x)/2\),
\(b(x)=x/R_x\), and \(a_\Gamma=[\psi(1/4)-\log\pi]/2<0\). The
quarter-centered digamma split is

\[
\boxed{h_\Gamma(x)=a_\Gamma b(x)+\sum_{m\ge0}g_m(x),
\qquad
g_m(x)=\frac{2x}{(4m+1)(4m+1+R_x)}.}
\tag{166B.1}
\]

Each term has a positive Stieltjes representation. With
\(y=\sqrt{4u-1}\),

\[
\frac1{R_x}=\int_{1/4}^\infty\frac{du}{\pi y(x+u)},\qquad
\frac{g_m(x)}x=\int_{1/4}^\infty
\frac{2y\,du}{\pi(4m+1)((4m+1)^2+y^2)(x+u)}.
\tag{166B.2}
\]

Thus

\[
\boxed{G(x)=\int_{1/4}^\infty\frac{d\sigma_\Gamma(u)}{x+u}}
\tag{166B.3}
\]

with signed density

\[
\frac{d\sigma_\Gamma}{du}
=\frac{a_\Gamma}{\pi y}
+\frac2\pi\sum_{m\ge0}
\frac{y}{(4m+1)((4m+1)^2+y^2)}.
\tag{166B.4}
\]

The total baseline shift between the completion-centered pole terms and
these positive quarter-centered terms is exact:

\[
\sum_{m\ge0}\frac1{(4m+1)(2m+1)}
=\frac\pi4+\frac{\log2}{2}.
\tag{166B.4a}
\]

The repaired positive part pushes forward exactly to the
Catalan-resolvent reserve of (155C.5), plus the \(m=0\) Catalan measure.
This is a reconciliation, not a new reserve.

At fixed \(a>0\), the weighted total variation is finite; its Bernstein
endpoint is geometrically small and the Gamma collision energy tends to
zero. In particular,

\[
\mathcal V_N[G](a)\le C_\Gamma(a),\qquad
|B_{N,N}[G](a)|\le C_\Gamma(a)\beta_a^N,
\quad \beta_a=\frac{a}{a+1/4}<1,
\tag{166B.7}
\]

and

\[
\mathscr E_N[1/x+G](a)
=a^{-2}+\mathscr E_N[G](a)+2a^{-1}B_{N,N}[G](a)
\longrightarrow a^{-2}.
\tag{166B.8}
\]

\subsection{Theorem 166B.1. Bounded full-prime energy at one fixed
scale}\label{theorem-166b.1.-bounded-full-prime-energy-at-one-fixed-scale}

For \(P(x)=-(\zeta'/\zeta)(s_x)/R_x\),

\[
\boxed{\begin{aligned}
\mathrm{RH}&\iff\sup_{N\ge0}\mathscr E_N[P](a)<\infty\\
&\iff\mathscr E_N[P](a)\longrightarrow a^{-2}\\
&\iff\mathscr E_N[P-1/x](a)\longrightarrow0.
\end{aligned}}
\tag{166B.9}
\]

The proof is the identity \(S_\xi=1/x+G-P\) plus the decay above and
Corollary 164.3. No separate sign is assigned to the Gamma and prime
pieces.

At \(a=2\) this becomes the explicit all-degree criterion

\[
\boxed{\mathrm{RH}\iff
\forall N\ge0:\quad
\mathscr E_N[P](2)\le\left(1+\log2+\frac\pi6\right)^2.}
\tag{166B.11}
\]

The constant has the triangular reading

\[
1+\log2+\frac\pi6
=2-\sum_{n\ge1}\frac1{2T_{2n}}
+\frac12\sqrt{1+\frac1{12}\sum_{n\ge1}\frac1{T_n^2}},
\tag{166B.12}
\]

and the elementary bounds from §R4 imply the convenient rational
equivalent

\[
\boxed{\mathrm{RH}\iff
\forall N:\ \mathscr E_N[P](2)<5
\iff\forall N:\ \mathscr E_N[P](2)\le5.}
\tag{166B.12a}
\]

\subsection{Theorem 166B.2. Prime
plateau}\label{theorem-166b.2.-prime-plateau}

With the same \(\mathcal K(a),\mathcal R(a)\) as Theorem 164.2,

\[
\boxed{\mathscr E_N[P](a)-a^{-2}
\sim\mathscr E_N[P-1/x](a)
\sim\mathcal K(a)\frac{\mathcal R(a)^{2N}}{\sqrt N}.}
\tag{166B.13}
\]

The exact endpoint decomposition is

\[
\mathscr E_N[P](a)=a^{-2}+\mathscr E_N[P-1/x](a)
+2a^{-1}B_{N,N}[P-1/x](a),
\tag{166B.16}
\]

and the completed resolvent endpoint satisfies

\[
|B_{N,N}[S_\xi](a)|\le
\left(\sum_{[\rho]}\frac{m_\rho}{|a+q_\rho|}\right)
\left(\frac a{d_a}\right)^N,
\quad d_a=\min_{[\rho]}|a+q_\rho|>a.
\tag{166B.17}
\]

Hence the cross term is geometrically smaller than the leading spectral
energy. The theorem describes the two alternatives; it does not select
the bounded one.

The adversarial Gamma-pole analysis also closes a tempting shortcut:
splitting the integer pole lattice into prime-power and complementary
indices does not yield a positive Löwner reserve. The exact rank-two
witness is recorded in §48C; adding the full repaired complement merely
returns the RH-equivalent completed comparison above.

\section{166C. Prime cutoffs, the persistent endpoint, and exponential
cutoff
depth}\label{c.-prime-cutoffs-the-persistent-endpoint-and-exponential-cutoff-depth}

Fix \(a>0\) and truncate the actual von-Mangoldt source by integer
value, including prime powers:

\[
P_X(x)=\frac1{R_x}\sum_{2\le n\le X}\Lambda(n)n^{-s_x},
\qquad R_x=\sqrt{1+4x},\quad s_x=\frac{1+R_x}{2}.
\tag{166C.1}
\]

Every fixed cutoff converges at each fixed degree but loses the
persistent endpoint at large degree. The recovery transition occurs only
when the cutoff depth is exponential in \(N\): the central-limit
calculation retained in earlier editions gives the threshold
\(\log X_N\sim (N+1)/a\). Consequently no single cutoff-uniform
subexponential cap can prove the full-prime criterion. The point is
uniformity in degree, not failure of any fixed finite approximation.

The support-count shortcut fails even earlier. Only \(40\) prime powers
lie at or below \(113\), while the true von-Mangoldt mass satisfies
\(\Psi(113)>117.3867>114\). Sparse support is therefore not small
weighted mass. The Gamma measure is continuous whereas \(d\Psi(e^u)\) is
atomic, so no finite constant gives literal measure domination of the
prime source by the Gamma lattice.

A surviving cutoff argument must match the cutoff to the degree and
control both pieces. One sufficient formulation is to construct \(X_N\)
and nonnegative \(A_N,D_N\) with

\[
\boxed{\ \|\mathcal Q_N[P_{X_N}]\|_2\le A_N,\qquad
\|\mathcal Q_N[P-P_{X_N}]\|_2\le D_N,\qquad
A_N+D_N\le\sqrt5.\ }
\tag{166C.25}
\]

Section 166D proves the remainder half for one explicit exponential
schedule; the matched main-term bound remains open. The exact endpoint
coefficient, central-limit window, and raw-cutoff negative control are
archived in earlier editions (Zenodo version chain).

\section{166D. An explicit full-row remainder for a growing prime
cutoff}\label{d.-an-explicit-full-row-remainder-for-a-growing-prime-cutoff}

Section 166C proves that no bound is available uniformly over
independent pairs \((N,X)\), and (166C.25) names the two estimates a
surviving cutoff argument needs. One of the two is now proved,
unconditionally, for an explicit schedule. Throughout,
\(\alpha=(\sqrt5-1)/2\) and the norm is the weighted one of (163A.3),
\(\lVert\mathcal Q_N[f]\rVert_2^2=\int_0^\infty e^{-t}|\mathcal Q_N[f](t;2)|^2dt
=\sum_{j=0}^N|B_{N,j}[f](2)|^2\) --- not the unweighted \(L^2(dt)\) norm
and not the Euclidean norm of the monomial coefficients.

\begin{quote}
\textbf{Theorem 166D.1.} For every integer \(N\ge0\) and every integer
\(X\ge2\), with no hypothesis, \[
\boxed{\ \bigl\lVert\mathcal Q_N[P-P_X](\cdot;2)\bigr\rVert_2
\le\frac{5^N}{\sqrt5}\,X^{-\alpha}
\Bigl(\frac{\log X}{\alpha}+\frac1{\alpha^2}\Bigr)\ }
\tag{166D.1}
\] In particular the schedule \(X_N=20^{N+1}\) gives \[
\begin{gathered}
\bigl\lVert\mathcal Q_N[P-P_{X_N}]\bigr\rVert_2\le D_N
:=\frac{20^{-\alpha}}{\sqrt5}
\Bigl(\frac{(N+1)\log20}{\alpha}+\frac1{\alpha^2}\Bigr)\kappa^N
\longrightarrow0,\\[2pt]
\kappa=\frac5{20^\alpha}=0.785035410093857\ldots<1
\end{gathered}
\]
\end{quote}

\emph{Proof.} On the closed complex disk \(|x-2|\le1\) put
\(\omega=1+4x\). Since \(\Re\omega\ge5\), the principal root satisfies
\(\Re\sqrt\omega=\sqrt{(|\omega|+\Re\omega)/2}\ge\sqrt5\) and
\(|\sqrt\omega|\ge\sqrt5\), so \(\Re s_x\ge1+\alpha>1\) --- and this is
tight, the minimum \(1+\alpha\) being attained on the boundary. The
Euler tail is therefore absolutely and uniformly convergent on a
neighbourhood of the disk. Since \(0\le\Lambda(n)\le\log n\) and prime
powers enter by their integer value as in (166C.1),

\[
M_X:=\sup_{|x-2|\le1}\bigl|P(x)-P_X(x)\bigr|
\le\frac1{\sqrt5}\sum_{n>X}\frac{\log n}{n^{1+\alpha}}
\le\frac{X^{-\alpha}}{\sqrt5}\Bigl(\frac{\log X}\alpha+\frac1{\alpha^2}\Bigr)
\]

the comparison integrating \((\log t)t^{-1-\alpha}\), which decreases
for \(t\ge2\) because \(1-(1+\alpha)\log t<0\) there. Cauchy's estimate
gives \(|(P-P_X)^{(r)}(2)|/r!\le M_X\). With the exact Pascal conversion
(BE.1a), \(d_r=\binom Nr2^rf^{(r)}(2)/r!\) and
\(B_{N,j}=\sum_{r\ge j}(-1)^j\binom rjd_r\), so

\[
\bigl\lVert\mathcal Q_N[f]\bigr\rVert_2\le\sum_{j=0}^N|B_{N,j}|
\le\sum_{r=0}^N2^r|d_r|
\le M_X\sum_{r=0}^N\binom Nr4^r=5^NM_X
\]

which controls all \(N+1\) coefficients simultaneously, both endpoints
and degree zero included. For \(\kappa<1\) exactly: \(\alpha>3/5\)
follows from \((11/5)^2<5\), and \(20^3=8000>3125=5^5\) gives
\(20^\alpha>5\). \(\square\)

\textbf{The surviving obligation is now one prescribed diagonal
sequence.} Put
\(A_N=\lVert\mathcal Q_N[P_{20^{N+1}}](\cdot;2)\rVert_2\). The triangle
and reverse triangle inequalities give
\(\bigl|A_N-\lVert\mathcal Q_N[P]\rVert_2\bigr|\le D_N\), so by Theorem
166B.1

\[
\boxed{\ \mathrm{RH}\iff\sup_{N\ge0}A_N<\infty\ }
\tag{166D.2}
\]

This does not contradict Theorem 166C.4, which rules out bounds uniform
over \emph{independent} cutoffs and degrees; here the cutoff is
prescribed as a function of the degree. The exponentially small error
also transfers the prime-energy asymptotics of Theorems 164.2 and 166B.2
to this schedule: under RH,
\(A_N^2=\tfrac14+\mathcal K(2)N^{-1/2}+o(N^{-1/2})\), using
\(|A_N^2-e_N^2|\le D_N(2e_N+D_N)=o(N^{-1/2})\) with
\(e_N=\lVert\mathcal Q_N[P]\rVert_2\); and if RH fails,
\(A_N^2\sim\mathcal K(2)\mathcal R(2)^{2N}N^{-1/2}\) with
\(\mathcal R(2)>1\), since \(D_N/e_N\to0\). It is the \textbf{norm}
error that is exponentially small; the energy statement needs that
displayed extra factor. Neither alternative has been selected by any
source estimate.

\textbf{Precise advance and remaining gap.} The remainder half of
(166C.25) is now explicit for one generous exponential schedule;
(166D.1) does \textbf{not} establish \(A_N+D_N\le\sqrt5\), and no bound
\(\sup_NA_N<\infty\) is proved. A finite set of computed degrees, a
positive endpoint, a bound depending uncontrollably on \(N\), or one
fixed exponential allowance above one does not discharge (166D.2). The
schedule \(X_N=20^{N+1}\) is a mathematical construction and is
\textbf{not} a proposal to enumerate \(20^{N+1}\) integers.

\emph{Scope note on generalization.} If the fixed basepoint \(a=2\) and
radius one are replaced by a general Cauchy radius \(h\in(0,a)\) with
\(\sigma_h=(1+\sqrt{1+4(a-h)})/2\) and matrix factor \((1+2a/h)^N\), the
monotonicity step above needs \(X\ge e^{1/\sigma_h}\); the uniform
integer assumption \(X\ge3\) suffices. Near \(\sigma_h=1\) the integrand
need not decrease from \(X=2\). This qualifies that generalization only,
not the concrete \(a=2\) theorem.

\subsection{A smaller sufficient whole-row
cutoff}\label{a-smaller-sufficient-whole-row-cutoff}

\textbf{Theorem 166D.2 (radius four-fifths).} Put
\(\alpha_4=(\sqrt{29/5}-1)/2\). For all integers \(N\ge0\), \(X\ge2\),

\[
\boxed{\|\mathcal Q_N[P-P_X](\cdot;2)\|_2
\le\frac{6^NX^{-\alpha_4}}{\sqrt{29/5}}
 \left(\frac{\log X}{\alpha_4}+\frac1{\alpha_4^2}\right).}
\tag{166D.3}
\]

\textbf{Proof.} On \(|x-2|\le4/5\) the principal root \(\sqrt{1+4x}\)
has real part and modulus at least \(\sqrt{29/5}\). The preceding tail
argument gives that disk supremum with \(6^N\) omitted. Cauchy's
estimate for the \(r\)-th Taylor coefficient gains \((5/4)^r\), and the
exact Pascal conversion now has sum \(\sum_r\binom Nr5^r=6^N\). The tail
comparison is decreasing from \(X=2\). \(\square\)

For the prescribed schedule \(X_N=13^{N+1}\), the error tends to zero
with an explicit linear prefactor and exponential base

\[
\kappa_{13}=\frac6{13^{\alpha_4}}=0.985728931213\ldots<1,\qquad
\boxed{\mathrm{RH}\iff
 \sup_N\|\mathcal Q_N[P_{13^{N+1}}](\cdot;2)\|_2<\infty.}
\tag{166D.4}
\]

Exactly, \((12/5)^2<29/5\) implies \(\alpha_4>7/10\), and
\(13^7=62{,}748{,}517>60{,}466{,}176=6^{10}\) implies
\(13^{\alpha_4}>6\). The equivalence follows from the triangle and
reverse triangle inequalities and Theorem 166B.1. No uniform norm bound
is proved. The earlier base-twenty schedule remains valid and has faster
decay in its displayed majorant. Neither base nor disk is asserted
optimal. The base-five result below concerns recovered outputs and mixed
derivative blocks, not the whole row.

\section{166E. Common-source recovery and the positive-scale error
transfer}\label{e.-common-source-recovery-and-the-positive-scale-error-transfer}

Use \(\mathcal P_m\) from Section 155A,
\(c_{m,j}=[y^j]\mathcal P_m(y)^2\), \(J=2m\), and the recovery weights
of (104A.6) for rates \(b_0,\ldots,b_J\). If all observations share one
moment perturbation \(\delta\mu\), finite summation proves

\[
\delta\Lambda_m(k)=\sum_{j=0}^Jc_{m,j}b_j^k\delta\mu_j
\quad\Longrightarrow\quad
\boxed{\sum_{k=1}^{J+1}w_k\delta\Lambda_m(k)
 =\sum_{j=0}^Jc_{m,j}\delta\mu_j.}
\tag{166E.1}
\]

Indeed \(\sum_kw_kb_j^k=1\). A common error is combined before absolute
bounds are taken. Independent rounding residuals \(e_k\) still
contribute \(\sum_k|w_k||e_k|\): if \(|\delta\mu_j|\le\epsilon_j\), the
complete bound is \(\sum_j|c_{m,j}|\epsilon_j+\sum_k|w_k||e_k|\). The
correlation must be established and preserved, not assigned after
rounding independent outputs.

\textbf{Theorem 166E.1 (a row-norm transfer at one positive scale).} Fix
\(a>0\), \(N\ge J\), and a real analytic source \(f\) near \(a\). Define
\(\mu_j^{(a)}[f]=(-1)^jf^{(j)}(a)/j!\) and use exactly the row operator
of (163.3), with coefficients \(B_{N,r}[f](a)\). Then

\[
\begin{aligned}
\mu_j^{(a)}[f]&=\frac1{\binom Nj a^j}
 \sum_{r=j}^N\binom rj B_{N,r}[f](a),\\
t_{m,N,a,k}(r)&=\sum_{j=0}^{\min(r,J)}
 c_{m,j}b_j^k\frac{\binom rj}{\binom Nj a^j},\\
\Lambda_{m,a}^f(k)&:=\sum_{j=0}^Jc_{m,j}b_j^k\mu_j^{(a)}[f]
 =t_{m,N,a,k}^{\mathsf T}B_N[f],\\
\left|\sum_{k=1}^{J+1}w_k\Lambda_{m,a}^f(k)\right|
 &\le\|t_{m,N,a,0}\|_2
   \|\mathcal Q_N[f](\cdot;a)\|_{L^2(e^{-t}dt)}.
\end{aligned}
\tag{166E.2}
\]

\textbf{Proof.} Differentiate the polynomial in (163.3) at \(z=1\);
equivalently use its falling-moment identity (163.5). This gives the
first line. Substitution gives the next two, and (166E.1) gives
\(\sum_kw_kt_{m,N,a,k}=t_{m,N,a,0}\). Cauchy--Schwarz with the actual
isometry (163A.3) proves the last line. No sign assumption is used.
\(\square\)

For \(a=2\), \(f=P-P_X\), and \(\alpha=(\sqrt5-1)/2\), Theorem 166D.1
gives

\[
\left|\sum_{k=1}^{2m+1}w_k\Lambda_{m,2}^{P-P_X}(k)\right|
\le\|t_{m,N,2,0}\|_2\frac{5^N}{\sqrt5}X^{-\alpha}
 \left(\frac{\log X}{\alpha}+\frac1{\alpha^2}\right).
\tag{166E.3}
\]

This keeps the integer-value prime-power cutoff, not a substituted Euler
product. The completed positive-scale approximation \(S_X=1/x+G-P_X\)
satisfies \(S_\xi-S_X=-(P-P_X)\), so the same absolute bound applies.
Furthermore

\[
\|t_{m,N,a,0}\|_2\le\sqrt{N+1}\,\mathcal P_m(-1/a)^2,
\qquad\mathcal P_m(-1/2)=2^{m+1}-2^{-m}.
\tag{166E.4}
\]

The binomial ratio is at most one and \(c_{m,j}\) has sign \((-1)^j\);
these give the first bound. The recurrence (155A.1) gives the second.
Thus the prescribed schedule \(X_N=20^{N+1}\) makes the transferred
error tend to zero for each fixed \(m\) as \(N\to\infty\). Equations
(166E.5)--(166E.9) below prove uniform growing-degree versions.

The shifted output at \(a>0\) is not an actual Li coefficient. The full
\(P\) has a pole \(1/x\) at the origin, while finite \(P_X\) is regular
there; \(S_X\) retains the uncancelled elementary pole. Setting \(a=0\)
in (166E.2) is invalid. Section 155B supplies the correctly regularized
boundary functional, but no directed all-order boundary error bound or
source sign is inferred from (166E.3). The uniform cutoff-row norm in
(166D.2) remains unproved. The finite-polynomial reconstruction in
Section 166G is a different operation, proved using a completed-source
disk and the full mixed-block amplification. It is not inferred from
(166E.3); independent rounding retains a separate budget.

\subsection{Growing degree and direct
recovery}\label{growing-degree-and-direct-recovery}

Write \(\mathcal T_{m,N}[f]=\sum_{j=0}^{2m}c_{m,j}(-1)^jf^{(j)}(2)/j!\),
realized through row degree \(N\ge2m\). The earlier row-norm bound
already gives a useful growing window. Put
\(C_\alpha=20^{-\alpha}(\log20/\alpha+\alpha^{-2})/\sqrt5\). For
\(0\le m\le\lfloor N/8\rfloor\), (166E.4) and (166D.1) imply

\[
|\mathcal T_{m,N}[P-P_{20^{N+1}}]|
\le4C_\alpha(N+1)^{3/2}(\kappa4^{1/8})^N\longrightarrow0.
\tag{166E.5}
\]

Here \(\alpha>3/5\) and \(5^8<2^{19}\) prove \(\kappa4^{1/8}<1\)
exactly. More generally this row-norm argument works for a fixed
\(m\le\theta N\) when \(\kappa4^\theta<1\) and \(2m\le N\). It remains
available for analytic errors known through their row norms. The
prime-tail functional itself admits a sharper direct estimate.

\textbf{Theorem 166E.2 (the full admissible recovery range).} Put
\(\alpha_3=(\sqrt6-1)/2\). On the disk \(|x-2|\le3/4\),

\[
M_X:=\sup_{|x-2|\le3/4}|P(x)-P_X(x)|
\le\frac{X^{-\alpha_3}}{\sqrt6}
 \left(\frac{\log X}{\alpha_3}+\frac1{\alpha_3^2}\right).
\tag{166E.6}
\]

For all integers \(N\ge0\), \(X\ge2\),

\[
\boxed{\max_{0\le m\le\lfloor N/2\rfloor}
 |\mathcal T_{m,N}[P-P_X]|
\le\frac{9\,3^NX^{-\alpha_3}}{4\sqrt6}
 \left(\frac{\log X}{\alpha_3}+\frac1{\alpha_3^2}\right).}
\tag{166E.7}
\]

\textbf{Proof.} On this disk the principal root has real part and
modulus at least \(\sqrt6\), proving (166E.6) by the same decreasing
tail comparison. Cauchy's bound is \(|f^{(j)}(2)|/j!\le M_X(4/3)^j\).
The alternating coefficients of \(\mathcal P_m^2\) give
\(|\mathcal T_{m,N}[f]|\le M_X\mathcal P_m(-4/3)^2\). The polynomial
recurrence has characteristic roots \(3\) and \(1/3\) at this argument;
its initial values give

\[
\mathcal P_m(-4/3)=\frac{3^{m+1}-3^{-m}}2,
\qquad \mathcal P_m(-4/3)^2\le\frac94\,9^m\le\frac94\,3^N.
\tag{166E.8}
\]

This proves the theorem, including \(m=N=0\). There is no whole-row
amplification because \(N\) enters the recovered functional only through
its admissibility condition \(2m\le N\). \(\square\)

For any integer \(b\ge5\), define

\[
\begin{gathered}
q_b=\frac3{b^{\alpha_3}},\qquad
C_b=\frac{9b^{-\alpha_3}}{4\sqrt6}
 \left(\frac{\log b}{\alpha_3}+\frac1{\alpha_3^2}\right),\\
\max_{2m\le N}|\mathcal T_{m,N}[P-P_{b^{N+1}}]|
\le C_b(N+1)q_b^N\longrightarrow0.
\end{gathered}
\tag{166E.9}
\]

Indeed \((12/5)^2<6\) and \(5^7>3^{10}\) give \(\alpha_3>7/10\) and
\(q_b<1\). For orientation, \(q_5=0.934429044142\ldots\) and
\(q_{20}=0.342142037896\ldots\); the exact inequalities carry
convergence. The result covers the full admissible range, not just the
coarser window (166E.5). It controls one common analytic source error.
Independent observation residuals still cost \(\sum_k|w_k||e_k|\). No
actual Li coefficient at the boundary or sign of a recovered output is
certified by this positive-scale estimate.

\section{166F. Matched cutoff matrices at one fixed positive
scale}\label{f.-matched-cutoff-matrices-at-one-fixed-positive-scale}

The derivative cutoff can be attached to an equivalent target at its own
basepoint. This defines a separate positive-scale moment form; it does
not identify that form with the boundary Li matrix or take a limit
\(a\downarrow0\).

For real analytic \(f\) near \(a\) and a polynomial
\(p(y)=\sum_jp_jy^j\), put

\[
\mathcal K_a[f;p]=\sum_j[y^j]p(y)^2\,
 \frac{(-1)^jf^{(j)}(a)}{j!},\qquad
(K_{a,d}[f])_{rs}=\mathcal K_a[f;\mathcal P_r,\mathcal P_s].
\tag{166F.1}
\]

The bilinear notation in the matrix replaces \(p^2\) by
\(\mathcal P_r\mathcal P_s\). The reserve \(\mathsf R_d\) remains
exactly (155C.8), in this same polynomial basis.

\textbf{Theorem 166F.1 (exact common-error block constant).} Let \(f\)
be holomorphic on a neighborhood of \(|x-2|\le3/4\) and real on the real
axis, with disk supremum at most \(M\). For every \(d\ge1\),

\[
\boxed{\left\|\mathsf R_d^{-1/2}K_{2,d}[f]\mathsf R_d^{-1/2}\right\|_{\rm op}
\le125\Sigma_d M,\qquad
\Sigma_d=\frac9{128}(9^d-9^{-d})+\frac{3d}{8}.}
\tag{166F.2}
\]

\textbf{Proof.} The third-kind polynomials (155D.4) have alternating
coefficients, since their \(j\) roots lie in \((0,4)\). At \(-4/3\)
their recurrence and initial values give
\(V_j(-4/3)=(3^{j+1}+3^{-j})/4\). For \(p=\sum_{j<d}a_jV_j\), the
coefficient absolute sum evaluated at \(4/3\) is at most
\(\sum_{j<d}|a_j|V_j(-4/3)\). Cauchy's derivative bound and
multiplication of these coefficient majorants give

\[
|\mathcal K_2[f;p]|
\le M\left(\sum_{j<d}|a_j|V_j(-4/3)\right)^2
\le M\Sigma_d\|p\|_{L^2(\nu_C)}^2
\le125M\Sigma_d\mathcal R[p].
\]

Summing the geometric terms in \(V_j(-4/3)^2\) gives the stated exact
\(\Sigma_d\). The supremum over all nonzero real polynomials of degree
below \(d\) proves the normalized operator-norm bound. \(\square\)

The coarser \((9^d-1)/8\) also bounds \(\Sigma_d\), but is unnecessary
here. This is the full mixed block: its highest derivative order is
\(2d-2\). For \(f=P-P_X\) use \(M_X\) in (166E.6). With

\[
X_d=5^{2d-1},\qquad
\varepsilon_d=125\Sigma_d
\frac{X_d^{-\alpha_3}}{\sqrt6}
 \left(\frac{\log X_d}{\alpha_3}+\frac1{\alpha_3^2}\right),
\tag{166F.3}
\]

one has \(\varepsilon_d=O(dq_5^{2d-2})\to0\). The parameter match is
\(N=2d-2\), so \(X_d=5^{N+1}\). This controls every admissible mixed
entry together in reserve norm, not the entire Bernstein--Laguerre row
at base five.

\subsection{The equivalent target at scale
two}\label{the-equivalent-target-at-scale-two}

Fix \(a=2\) once and for all, and write

\[
\begin{aligned}
y_{a,\rho}&=(a+q_\rho)^{-1},\\
\mathcal K_a[p]&=\mathcal K_a[S_\xi;p]
 =\sum_{[\rho]}m_\rho y_{a,\rho}p(y_{a,\rho})^2,\\
\mathcal D_a[p]&=\mathcal R[p]-\mathcal K_a[p],\\
\beta_{a,d}^+&=\sup_{\substack{p\ne0\\\deg p<d}}
 \mathcal D_a[p]/\mathcal R[p].
\end{aligned}
\tag{166F.4}
\]

The actual shifted weights have finite absolute mass. Their bounded
support accumulates only at zero, and

\[
\Im y_{a,\rho}=-\frac{\Im q_\rho}{|a+q_\rho|^2}.
\tag{166F.5}
\]

Thus the shift preserves nonreality. Under RH the support is positive
real and contained strictly inside \([0,4]\).

\textbf{Theorem 166F.2 (positive-scale criterion).} A finite uniform
upper bound on \(\beta_{2,d}^+\), a fixed polynomial upper bound, or an
upper bound at every exponential rate in \(d\) is equivalent to RH.

\textbf{Proof.} Under RH \(\mathcal K_2\) is positive and
\(\beta_{2,d}^+\le1\). For the finite-bound converse repeat Theorem
155D.1 with shifted moments: their generating function is \(S_\xi(a-z)\)
and has genuine poles at \(z=a+q_\rho\). Compact support of the
hypothetical positive representing measure excludes a nonreal such pole.
For the subexponential converse define the exterior root by
\(w+w^{-1}=2-y_{2,\rho}\) and let \(\mathfrak L_2=\sup r(y_{2,\rho})\).
If a nonreal point exists, \(\mathfrak L_2>1\) is attained finitely
often, with a strict gap to the remaining radii. The finite annihilator
and real affine phase adjustment of Theorem 155D.3 apply to these
shifted points and their summable weights. They yield
\(\beta_{2,d}^+\ge c\mathfrak L_2^{2d}\) at every sufficiently large
rank. This contradicts any all-rate upper estimate. The universal
polynomial upper estimate on test values proves the matching rate,
exactly as in (155D.10). The rate is \(\mathfrak L_2\), not an
identification with the boundary Li modulus \(\mathfrak L\). \(\square\)

\subsection{The cutoff approximates this very
target}\label{the-cutoff-approximates-this-very-target}

Retain the full completion and define

\[
S_X(x)=\frac1x+G(x)-P_X(x),\qquad
\widehat{\mathsf D}_{2,d}=\mathsf R_d-K_{2,d}[S_{X_d}],\qquad
\widehat\beta_{2,d}^+=
\lambda_{\max}(\mathsf R_d^{-1/2}\widehat{\mathsf D}_{2,d}\mathsf R_d^{-1/2}).
\tag{166F.6}
\]

The elementary and Gamma channels are present, not truncated out. Since
\(S_X-S_\xi=P-P_X\), the sign of the exact error is

\[
\mathsf D_{2,d}-\widehat{\mathsf D}_{2,d}=K_{2,d}[P-P_{X_d}],\qquad
|\widehat\beta_{2,d}^+-\beta_{2,d}^+|\le\varepsilon_d\longrightarrow0.
\tag{166F.7}
\]

The second assertion is the variational eigenvalue bound applied to
(166F.2). Hence

\[
\boxed{\begin{aligned}
\mathrm{RH}&\iff\sup_d\widehat\beta_{2,d}^+<\infty\\
&\iff\exists C,A\ge0\ \forall d:
 \widehat\beta_{2,d}^+\le C(1+d)^A\\
&\iff\forall\epsilon>0\ \exists C_\epsilon<\infty\ \forall d:
 \widehat\beta_{2,d}^+\le C_\epsilon e^{\epsilon d}.
\end{aligned}}
\tag{166F.8}
\]

The vanishing error cannot hide the eventual exponential escape. A fixed
unbounded rank subsequence suffices for the all-rate condition by that
same eventual escape, with rates measured in the actual rank. The
varying-cutoff matrices need not be nested, so their extreme eigenvalues
need not be monotone.

No upper bound in (166F.8) is proved from the prime source. This matched
positive-scale approximation bypasses a boundary transfer for this
route; it is not itself that transfer. Section 166G separately
reconstructs the actual boundary matrix with Taylor degree \(15d-1\) and
cutoff \(8^{15d}\). The raw cutoff retains its origin pole. Independent
entry-rounding errors need their own matrix estimate in addition to the
common analytic bound.

\section{166G. Controlled reconstruction of the actual boundary
matrix}\label{g.-controlled-reconstruction-of-the-actual-boundary-matrix}

The raw cutoff \(S_X=1/x+G-P_X\) has an uncancelled origin pole. The
construction below instead evaluates a finite Taylor polynomial at zero.
It uses the same boundary moments and compensated reserve as Section
155C, not the shifted target \(K_{2,d}\) of Section 166F:

\[
\mu_j=\frac{(-1)^jS_\xi^{(j)}(0)}{j!},\qquad
K_d=\mathsf C_d[\mu_{r+s}]\mathsf C_d^{\mathsf T},\qquad
\mathsf D_d=\mathsf R_d-K_d.
\tag{166G.1}
\]

Here \(\mathsf C_d\) is the coefficient matrix of \(\mathcal P_0,\ldots,
\mathcal P_{d-1}\). The cutoff includes full elementary/Gamma completion
and prime powers by their integer value; it is not a truncated Euler
product. At the safe point two, \(S_X-S_\xi=P-P_X\).

\textbf{Lemma 166G.1 (completed disk).} The carried low-height exclusion
\(|\Im\rho|>1\) gives holomorphy on a neighborhood of \(|z-2|\le23/8\)
and

\[
\sup_{|z-2|\le23/8}|S_\xi(z)|\le16\lambda_1<1.
\tag{166G.2}
\]

\textbf{Proof.} For \(q=A+iB\) one has \(A=\gamma^2+\beta(1-\beta)>1\)
and \(|B|<|\gamma|<A\). On this disk \(\Re z\ge-7/8\), so
\(|z+q|^{-1}\le(A-7/8)^{-1}\le8/A\le16\Re(1/q)\). The folded sum
converges normally by this summable majorant, whose sum is
\(16\lambda_1\). Section 155C.1 gives \(\lambda_1<1/16\). Every pole
lies strictly outside a neighborhood of the disk. This uses the
low-height exclusion, not RH or an unknown positive folded measure.
\(\square\)

For rank \(d\ge1\) set \(L_d=15d-1\), \(X_d=8^{15d}\) and

\[
\boxed{\widehat S_d(z)=\sum_{\ell=0}^{L_d}
\frac{S_{X_d}^{(\ell)}(2)}{\ell!}(z-2)^\ell.}
\tag{166G.3}
\]

Write \(\alpha=(\sqrt6-1)/2\), \(r=18/23\), \(q_T=3/8^\alpha\) and

\[
\begin{aligned}
E_d^\partial&=\frac{23}{5}r^{15d}
 +\frac1{2\sqrt6}\left(\frac{15d\log8}{\alpha}+\frac1{\alpha^2}\right)q_T^{15d},\\
\Theta&=36(18/23)^{15}<1.
\end{aligned}
\tag{166G.4}
\]

\textbf{Theorem 166G.2 (boundary disk and mixed block).} Put
\(\widehat\mu_j^{(d)}=(-1)^j\widehat S_d^{(j)}(0)/j!\) and
\(\widehat K_d^\partial=\mathsf C_d[\widehat\mu_{r+s}^{(d)}]
\mathsf C_d^{\mathsf T}\). Then

\[
\begin{aligned}
\sup_{|z|\le1/4}|S_\xi(z)-\widehat S_d(z)|&\le E_d^\partial,\\
\|\mathsf R_d^{-1/2}(\widehat K_d^\partial-K_d)
\mathsf R_d^{-1/2}\|_{\rm op}
&\le\varepsilon_d^\partial:=125\frac{36^d-1}{35}E_d^\partial
 =O(d\Theta^d)\longrightarrow0.
\end{aligned}
\tag{166G.5}
\]

\textbf{Proof.} Let \(T_L\) denote Taylor truncation at two. The
completed remainder on \(|z|\le1/4\) is at most
\(r^{L+1}/(1-r)=(23/5)r^{L+1}\), since \(|z-2|\le9/4\). For
\(f_X=P-P_X\), the disk-\(3/4\) estimate (166E.6) is
\(M_X=X^{-\alpha}(\log X/\alpha+\alpha^{-2})/\sqrt6\). Consequently
\(|T_Lf_X(z)|\le M_X\sum_{\ell=0}^L3^\ell
\le M_X3^{L+1}/2\). The exact polynomial identity

\[
S_\xi-T_LS_X=(S_\xi-T_LS_\xi)-T_L(P-P_X)
\tag{166G.6}
\]

proves the first bound after \(L+1=15d\) and \(X=8^{15d}\). Neither term
continues the raw cutoff through its pole.

For a holomorphic boundary error \(f\) of supremum \(E\), Cauchy bounds
\(|f^{(j)}(0)|/j!\) by \(E4^j\). Write \(p=\sum_{j<d}a_jV_j\) in the
Catalan-orthonormal third-kind basis of Section 155D. Its alternating
coefficients and the exact value at \(-4\) give

\[
 V_j(-4)=\frac{\ell^{j+1}+\ell^{-j}}{\ell+1}\le6^j,
 \quad\ell=3+2\sqrt2,\qquad
 |\mathcal K_0[f;p]|\le
 E\frac{36^d-1}{35}\|p\|_{L^2(\nu_C)}^2,
\tag{166G.7}
\]

where \(\mathcal K_0[f;p]=\sum_j[y^j]p^2(-1)^jf^{(j)}(0)/j!\).
Cauchy--Schwarz proves the last inequality from
\(E(\sum|a_j|V_j(-4))^2\). Apply
\(\|p\|_{L^2(\nu_C)}^2\le125\mathcal R[p]\) and take the absolute
Rayleigh supremum. This controls all mixed entries simultaneously.
Finally \(\alpha>2/3\) gives \(q_T<3/4<18/23\), and the exact integer
comparison \(36\,18^{15}<23^{15}\) proves the claimed decay. \(\square\)

The required moment translation is finite:

\[
 c_\ell^X=\frac{(-1)^\ell S_X^{(\ell)}(2)}{\ell!},\qquad
 \widehat\mu_j^{(d)}=\sum_{\ell=j}^{15d-1}
 \binom\ell j2^{\ell-j}c_\ell^{X_d}.
\tag{166G.8}
\]

The highest boundary moment is \(\mu_{2d-2}\), while the source degree
is \(15d-1\). For \(m<d\) the odd-Li diagonal error is at most
\(\varepsilon_d^\partial(\mathsf R_d)_{mm}\). This is a signed
polynomial identity, not a replacement positive measure.

\subsection{The matched boundary criterion and its error
budget}\label{the-matched-boundary-criterion-and-its-error-budget}

Define
\(\widehat{\mathsf D}_d^\partial=\mathsf R_d-\widehat K_d^\partial\) and
its largest normalized eigenvalue \(\widehat\beta_d^{\partial,+}\). The
error sign and variational estimate are

\[
 \mathsf D_d-\widehat{\mathsf D}_d^\partial=\widehat K_d^\partial-K_d,
 \qquad |\widehat\beta_d^{\partial,+}-\beta_d^+|
 \le\varepsilon_d^\partial\longrightarrow0.
\tag{166G.9}
\]

Section 155D therefore gives

\[
\boxed{\begin{aligned}
\mathrm{RH}&\iff\sup_d\widehat\beta_d^{\partial,+}<\infty\\
&\iff\exists C,A\ge0\ \forall d:\quad
 \widehat\beta_d^{\partial,+}\le C(1+d)^A\\
&\iff\forall\epsilon>0\ \exists C_\epsilon<\infty\ \forall d:\quad
 \widehat\beta_d^{\partial,+}\le C_\epsilon e^{\epsilon d}.
\end{aligned}}
\tag{166G.10}
\]

No bound in (166G.10) is proved from the arithmetic source. The order of
operations is complete finite source at two, finite Taylor polynomial,
boundary moments, then reflected mixed matrix. Holding \(X\) fixed while
sending \(L\) to infinity would revive the original pole obstruction.
Section 166F's smaller positive-scale cutoff is a separate result.

If the coefficients \(S_X^{(\ell)}(2)/\ell!\) have independent errors
\(e_\ell\), add
\(125(36^d-1)\sum_{\ell=0}^{15d-1}|e_\ell|(9/4)^\ell/35\) to the matrix
error in (166G.5). The common analytic error is not a rounding budget.
The bundled exact rational regression uses the retained synthetic model
\(q_0=17/4\) (multiplicity two), \(q_\pm=147/16\pm3i/2\). Its positive
shifted moments through degree \(284\) reconstruct a negative
\(\lambda_{37}^*\in(-0.598,-0.597)\), with rational error below
\(10^{-30}\). This tests the adapter, not the actual zeta source or an
all-order upper bound.

\chapter{Technical chapter XV. Remaining proof obligations and
countermodels}\label{technical-chapter-xv.-remaining-proof-obligations-and-countermodels}

The full criteria and proved finite results have been separated above.
These final references identify the remaining source implication and its
countermodels.

\section{66. Finite results within the exact
criteria}\label{finite-results-within-the-exact-criteria}

Reader §R12 gives the current scalar proof channels; §13 gives the
zero-assisted finite rectangle. Technical §§92--93 and 138 distinguish
global order eight from rank twelve at one node. These finite theorems
are summarized once in Reading Volume §R22.

Their exact converse has no finite stopping rank: §49 gives an open
interval with infinitely many negative rungs if RH fails, while §§51 and
145 give finite witnesses at every prescribed scale. A uniform source
theorem is still needed to establish the positive side at every order.

\section{67. The remaining source
implication}\label{the-remaining-source-implication}

The all-node inequality (54.8), restated at Reader (23.1), remains
unproved. Sections 155D--155E reduce equivalent source tasks to the
compensated form's one-sided growth or to ordinary-Laguerre curvature. A
polynomial upper bound on either parity of \(I_n\) suffices under the
explicit progression theorem; no such bound is proved. The compensation
has a proved polynomial bound, but the signed prime integral has not.

At fixed positive scale two, Section 166F aligns a complete cutoff
matrix with its own RH-equivalent target, with vanishing error at
\(5^{2d-1}\). Section 166G separately reconstructs the actual boundary
matrix by a controlled finite Taylor polynomial. Both one-sided source
upper bounds remain open. Whole-row convergence at \(13^{N+1}\) or
\(20^{N+1}\) does not supply them. Sections 155F--155G show why an
envelope, positive prime-power weights and a long fixed exact prime
prefix are insufficient. Sections 155H--155I restore the exact divisor
law and identify its Mellin error with the denominator residual's Hardy
norm. The remaining sufficient target is (155H.6), below every positive
power of the actual prefix length; the two tested families do not meet
it.

The independent seventh-rung source certificate is the next finite-order
task. Global \(W_7>0\) is not established by the independent source
channel; \(W_7\rightsquigarrow F^{(14)}\) is the first open source rung.
No historical certificate has been promoted by these all-order
reformulations.

\section{69. Smallest countermodels that remain
decisive}\label{smallest-countermodels-that-remain-decisive}

\begin{itemize}
\tightlist
\item
  A coarse one-channel sample cannot distinguish a frequency from its
  Nyquist mirror.
\item
  A nonnegative function can have a Fourier transform of changing sign.
\item
  A fixed finite moment/Weil compression cannot determine a meromorphic
  target Cauchy functional.
\item
  A finite set of Li/Hankel inequalities can be passed by an off-axis
  synthetic quartet.
\item
  Any fixed positive Widder prefix can coexist with a later negative
  rung in a finite synthetic Stieltjes test model.
\item
  The recurrence \(W_{k+1}=-(2k+1)W_k'-xW_k''\) does not preserve
  positivity without derivative control.
\item
  Matching minimum eigenvalues does not determine ground-state parity.
\item
  Exact roots-of-unity cancellation does not survive nonlinear
  logarithmic phase without curvature control.
\end{itemize}

The fixed counting and prime-supported countermodels in Sections
155F--155G additionally retain the specified positivity and PNT
hypotheses without the actual arithmetic cancellation. The divergent
balanced primitive product in Section 155H rejects that one
construction, not all inverse-factor methods.

These countermodels should remain beside the positive theory. They
prevent the next branch from silently rebuilding an already closed
route.

\newpage

\chapter{Summary of the technical
boundary}\label{summary-of-the-technical-boundary}

The Dossier starts with exact triangular arithmetic and reaches the
completed source without identifying its comparison measures with the
actual folded zero data. Its common index set is the reflection-pair
set, with algebraic multiplicity retained. The canonical product and
resolvent therefore count an off-line quartet as two conjugate folded
nodes. This indexing is sufficient to recover the first-Li normalization
and the Widder-zero-heat identities without an RH assumption.

The inherited finite source certificates and exact arithmetic identities
remain the results recorded in the Reading Volume's status key. Their
receipt provenance is v5.0, not a newly claimed replay. The
independently unproved next source rung is \(W_7\); the all-order,
all-node and all-degree requirements in the corresponding criteria
remain unproved. Countermodels exclude the listed shortcuts under their
stated hypotheses, not every possible route to the required source
inequalities.

Thus the proofs establish their own domains, the certificates establish
their own finite ranges, and the empirical appendix establishes neither.
Version 5.1 clarifies those interfaces but advances no RH status. No
entropy construction or semantic interpretation is a premise of this
Dossier.

\chapter{Part E. Empirical context and historical
transcriptions}\label{part-e.-empirical-context-and-historical-transcriptions}

\EVID{} Under the single master status key in the Reading Volume,
everything in this Part is numerical or transcribed evidence. It
supports no theorem in Chapters I--XV, and no result anywhere in the
program depends on it. It is retained so that the empirical claims that
circulated alongside this work stay attached to the exact statements
that replaced them.

Section 168 gives the exact counting and triangular-coordinate
identities, including the transport rule (168.6): an unchanged event
requires transporting the complete selection rule -- windows, merging
and support conditions -- and applying \(T_x\) to the cumulative smooth
zero count is a different operation. The historical comparisons below
provide context and possible future tests; they are not an RH proof.

The broader report {[}62{]} scans \(2.30\)--\(3.20\) with lower controls
and describes an enrichment shelf near \(2.70\)--\(2.81\). The rational
marker \(11/4=2.75\) is inside that reported interval. It is not
established as an endpoint of that shelf, and it is not the mean
normalized spacing. The statistic concerns random-surrogate
preceding-marker hits near small-gap valleys. Different supported-block
sets must not be treated as a common-support ranking. The five-marker
figure retained below is a local comparison and cannot show the full
shelf. The raw gap arrays and complete randomization replay are not
included here, so no significance level and no shelf boundary is
asserted.

The final three pages reproduce original report pages 11--13 as
historical pages: not a fresh replay, and not a confirmation of every
interpretation in that report. This Part uses its own \texttt{S} label
namespace, and §168 retains tags (168.2)--(168.6).

\subsection{Historical comparison, as
transcribed}\label{historical-comparison-as-transcribed}

\subsection{Historical section 168. What the reported spacing
experiments
contribute}\label{historical-section-168.-what-the-reported-spacing-experiments-contribute}

The two historical reports {[}61--62{]} are retained as empirical
instruments. They test different coordinates and have different
controls. The present revision transcribes their reported counts and
arithmetic comparisons; their raw gap arrays and full randomization
replay are not included here. No fresh search-adjusted significance
level is claimed.

\subsubsection{Row-sector audit near the proposed modulus
window}\label{row-sector-audit-near-the-proposed-modulus-window}

The orthogonal follow-up {[}61{]} scans \(401\) moduli from \(800\) to
\(1200\) using six row sectors and \(500\) permutations. Its label
``base-6'' means six bins, not positional base-six notation. The windows
give:

\begin{longtable}[]{@{}
  >{\raggedright\arraybackslash}p{(\columnwidth - 8\tabcolsep) * \real{0.2000}}
  >{\raggedright\arraybackslash}p{(\columnwidth - 8\tabcolsep) * \real{0.2000}}
  >{\raggedright\arraybackslash}p{(\columnwidth - 8\tabcolsep) * \real{0.2000}}
  >{\raggedright\arraybackslash}p{(\columnwidth - 8\tabcolsep) * \real{0.2000}}
  >{\raggedright\arraybackslash}p{(\columnwidth - 8\tabcolsep) * \real{0.2000}}@{}}
\toprule\noalign{}
\begin{minipage}[b]{\linewidth}\raggedright
window
\end{minipage} & \begin{minipage}[b]{\linewidth}\raggedright
best reported modulus
\end{minipage} & \begin{minipage}[b]{\linewidth}\raggedright
rank of \(1009\)
\end{minipage} & \begin{minipage}[b]{\linewidth}\raggedright
adjusted \(p\) at best modulus
\end{minipage} & \begin{minipage}[b]{\linewidth}\raggedright
adjusted \(p\) at \(1009\)
\end{minipage} \\
\midrule\noalign{}
\endhead
\bottomrule\noalign{}
\endlastfoot
\(900\) & \(1133\) & \(48\) & \(0.079840\) & \(1\) \\
\(1000\) & \(1087\) & \(151\) & \(0.235529\) & \(1\) \\
\(1100\) & \(1133\) & \(143\) & \(1\) & \(1\) \\
\(1300\) & \(1022\) & \(58\) & \(0.938124\) & \(1\) \\
\end{longtable}

This audit does not validate a special \(990/1009\) row-sector
mechanism. It also does not test every possible gcd or TFG partition. A
reduced-ray experiment should declare \(q=n/\gcd(n,k)\) and its
reciprocal weights in advance, then retain the scan correction.
Arithmetic equality of totients or a useful decimal period remains true
independently of this outcome.

\subsubsection{\texorpdfstring{A localized-enrichment comparison near
\(2.75\)}{A localized-enrichment comparison near 2.75}}\label{a-localized-enrichment-comparison-near-2.75}

The second report {[}62{]} studies unfolded nearest-neighbor gaps
\(s_n=N_0(\gamma_{n+1})-N_0(\gamma_n)\), with
\(N_0(T)=\frac{T}{2\pi}\log\frac{T}{2\pi}-\frac{T}{2\pi}+7/8\). Its
active marker is

\[
a_+=\frac\pi2+\frac2\pi+\cos1
\approx2.747718405031
\qquad a_-=\frac\pi2+\frac2\pi-\cos1
\tag{S.1}
\]

For \(r=\pi-3\), the identity \(\cos1=\sin[(1+r)/2]\) is exact. The
approximation \(\cos1\approx1/2+2r^2\) has error about
\(0.000205346767\) at this fixed \(r\); it is not the Taylor expansion
of that sine function about \(r=0\). Both the exact marker and the
useful approximation can be retained with this distinction stated.

The original reported 22-block comparison totals \(736\) observed
preceding hits versus \(218\) random-surrogate hits, ratio
\(3.376146789\); each of its printed block ratios exceeds one. The
broader follow-up shows a band, with several neighboring markers:

\begin{longtable}[]{@{}lllll@{}}
\toprule\noalign{}
marker & supported blocks & observed hits & surrogate hits & ratio \\
\midrule\noalign{}
\endhead
\bottomrule\noalign{}
\endlastfoot
\(a_+\) & \(21\) & \(702\) & \(218\) & \(3.22018\) \\
\(11/4\) & \(22\) & \(743\) & \(224\) & \(3.31696\) \\
\(7\pi/8\) & \(21\) & \(692\) & \(190\) & \(3.64211\) \\
\(2.740\) & \(22\) & \(784\) & \(218\) & \(3.59633\) \\
\(2.780\) & \(21\) & \(570\) & \(158\) & \(3.60759\) \\
\end{longtable}

The primary summary reports \(736\) observed hits and \(218\) surrogate
hits over \(22\) supported blocks. The follow-up row for the same active
marker reports \(702\) and \(218\) over \(21\) supported blocks. These
are different reported aggregates. The printed tables do not establish
that the difference is exactly one removed block, or that the repeated
surrogate count means the same random sample was reused. Those questions
require the block ledger and the randomization record, neither of which
is included here.

\begin{figure}
\centering
\includegraphics{figures/v27_empirical_marker_comparison.png}
\caption{Reported follow-up counts near 2.75, transcribed from the
empirical report {[}62{]}. Each label shows its own supported-block
count. Observed/surrogate ratios use different block sets and cannot by
themselves rank the markers on common support.}
\end{figure}

The denominator in this table counts random-surrogate hits, not hits at
the sign-reversed marker \(a_-\). The displayed support sets differ, so
these ratios do not select a unique winner. The report's statistic
searches a preceding window for a marker near a small-gap valley; it is
a conditional gap-pattern statistic. An unconditional pair-correlation
value at separation \(2.75\) does not test the same event. Its
random-center controls likewise do not preserve all local gap
correlations.

The useful next comparison is specific: use a common set of blocks for
\(a_+\), \(7\pi/8\), \(11/4\), and nearby prespecified controls; retain
the same windows, merging rules, exclusions, and support thresholds;
then run the complete statistic on held-out zeros and an appropriate
correlated null, such as a matched sine-kernel/GUE simulation. Include
the full selection procedure in any significance calculation. A decimal
or triangular coordinate can organize this experiment without
predetermining its result.

The exact arithmetic of §§98 and 167, the measured enrichment band, and
the source positivity frontier therefore remain available side by side.
Their proposed connection is a research question requiring a transport
identity or a successful independent test. RH remains open.

\subsection{Sources for this Part}\label{sources-for-this-part}

{[}61{]} \emph{A Search-Adjusted Audit of Row-Sector Modulus Signals in
Terminal Zeta-Zero Data}, unpublished historical report. The table above
preserves the transcribed audit; it has not been rerun.

{[}62{]} \emph{A Finite-Count Search for Localized Enrichment Near 2.75
in Unfolded Zeta-Zero Gaps}, unpublished historical report, 24 pages.
The final three pages of this Part reproduce pages 11--13 of that
report; those pages match the original copy in both extracted text and
rendered pixels.

\subsection{Original wider-scan pages
follow}\label{original-wider-scan-pages-follow}

The next three pages reproduce pages 11, 12 and 13 of the original
24-page report, in order. They include the active-marker follow-up,
alternative-marker support, and the broader threshold scan. Their
original titles and page furniture are preserved so the historical
source can be identified.

Their original page numbers, figure numbers, table numbers and internal
references --- including ``Appendix D'', the ``R141H'' and ``R142A''
ledgers and the named companion rails --- belong to that report, not to
this volume and not to this supplement. Consult the complete report for
those references. The facsimiles preserve historical evidence; they are
not a fresh replay and not a confirmation of every interpretation in
that report.

\newpage

\chapter{References}\label{references}

\begingroup
\small

\begin{enumerate}
\def\labelenumi{\arabic{enumi}.}
\item
  E. Bombieri, \emph{The Riemann Hypothesis - official problem
  description}, Clay Mathematics Institute, Millennium Prize Problems,
  \href{https://www.claymath.org/millennium/riemann-hypothesis/}{Clay
  problem description}.
\item
  NIST Digital Library of Mathematical Functions, Sections 5.5,
  5.7--5.9, 5.11, 24.4, 25.4, 25.10, 25.11, and 25.16; in particular
  \href{https://dlmf.nist.gov/25.11}{Hurwitz special values and
  derivatives}, \href{https://dlmf.nist.gov/25.10}{zero counting}, and
  \href{https://dlmf.nist.gov/5.11}{Gamma and digamma asymptotics}.
\item
  E. C. Titchmarsh, revised by D. R. Heath-Brown, \emph{The Theory of
  the Riemann Zeta-Function}, 2nd ed., Oxford University Press, 1986.
\item
  X.-J. Li, ``The positivity of a sequence of numbers and the Riemann
  hypothesis,'' \emph{Journal of Number Theory} 65 (1997), no. 2,
  325--333, \href{https://doi.org/10.1006/jnth.1997.2137}{DOI
  10.1006/jnth.1997.2137}.
\item
  E. Bombieri and J. C. Lagarias, ``Complements to Li's criterion for
  the Riemann hypothesis,'' \emph{Journal of Number Theory} 77 (1999),
  no. 2, 274--287, \href{https://doi.org/10.1006/jnth.1999.2392}{DOI
  10.1006/jnth.1999.2392}.
\item
  L. de Branges, \emph{Hilbert Spaces of Entire Functions},
  Prentice-Hall, 1968.
\item
  G. Csordas, T. S. Norfolk, and R. S. Varga, ``The Riemann hypothesis
  and the Turan inequalities,'' \emph{Transactions of the American
  Mathematical Society} 296 (1986), 521-541.
\item
  J. Huckstead, \emph{Divisor-Band Bijections and Totient Shells in the
  Triangular-Fractional Grid}, version 5, Zenodo (2026), DOI
  10.5281/zenodo.21440877.
\item
  J. Huckstead, \emph{Triangular Addressing, Null Coordinates, and
  Reciprocal-Metallic Locks in a Scale-Covariant Plane Geometry},
  Working Draft III, Zenodo (2026), DOI 10.5281/zenodo.21444929.
\item
  J. Huckstead, \emph{From Measure to Flux: Energy-Conjugate Bookkeeping
  and Exact Semidiscrete Scattering in Finite Radial Wave Models},
  publication layer v0.6.13, Zenodo (2026), DOI 10.5281/zenodo.21444339.
\item
  J. Huckstead, \emph{Two-Term Asymptotics for the Total Variation of
  Exponentially Weighted Laguerre Polynomials}, v0.1, Zenodo (2026), DOI
  10.5281/zenodo.21861619.
\item
  J. Sondow and C. Dumitrescu, ``A monotonicity property of Riemann's xi
  function and a reformulation of the Riemann hypothesis,''
  \emph{Periodica Mathematica Hungarica} 60 (2010), 37-40, DOI
  10.1007/s10998-010-1037-3; arXiv:1005.1104.
\item
  Y. Matiyasevich, F. Saidak, and P. Zvengrowski, ``Horizontal
  monotonicity of the modulus of the Riemann zeta function and related
  functions,'' arXiv:1205.2773 (2012).
\item
  J. C. Lagarias, ``On a positivity property of the Riemann
  xi-function,'' \emph{Acta Arithmetica} 89 (1999), 217-234.
\item
  A. D. Sokal, ``Real-variables characterization of generalized
  Stieltjes functions,'' \emph{Expositiones Mathematicae} 28 (2010),
  179-185; arXiv:0902.0065.
\item
  D. V. Widder, ``The Stieltjes transform,'' \emph{Transactions of the
  American Mathematical Society} 43 (1938), 7-60.
\item
  F. Johansson, ``Rigorous high-precision computation of the Hurwitz
  zeta function and its derivatives,'' arXiv:1309.2877 (2013).
\item
  GNU MPFR Project, \emph{GNU MPFR 4.2.2 Manual},
  https://www.mpfr.org/mpfr-current/mpfr.html.
\item
  F. Hausdorff, ``Summationsmethoden und Momentfolgen. I,''
  \emph{Mathematische Zeitschrift} 9 (1921), 74-109, DOI
  10.1007/BF01378337.
\item
  F. Hausdorff, ``Summationsmethoden und Momentfolgen. II,''
  \emph{Mathematische Zeitschrift} 9 (1921), 280-299, DOI
  10.1007/BF01279032.
\item
  A. Pringsheim, ``Ueber Functionen, welche in gewissen Punkten endliche
  Differentialquotienten jeder endlichen Ordnung, aber keine Taylor'sche
  Reihenentwickelung besitzen,'' \emph{Mathematische Annalen} 44 (1894),
  41-56, EuDML 157700.
\item
  K. Löwner, ``Über monotone Matrixfunktionen,'' \emph{Mathematische
  Zeitschrift} 38 (1934), 177-216, DOI 10.1007/BF01170633.
\item
  R. L. Schilling, R. Song, and Z. Vondraček, \emph{Bernstein Functions:
  Theory and Applications}, 2nd ed., De Gruyter, 2012, DOI
  10.1515/9783110269338.
\item
  M. Riesz, ``Sur l'hypothèse de Riemann,'' \emph{Acta Mathematica} 40
  (1916), 185-190, DOI 10.1007/BF02418544.
\item
  L. Báez-Duarte, ``A sequential Riesz-like criterion for the Riemann
  hypothesis,'' \emph{International Journal of Mathematics and
  Mathematical Sciences} 2005, 3527-3537, DOI 10.1155/IJMMS.2005.3527.
\item
  A. Agarwal, M. Garg, and B. Maji, ``Riesz-type criteria for the
  Riemann hypothesis,'' \emph{Proceedings of the American Mathematical
  Society} 150 (2022), DOI 10.1090/proc/16064.
\item
  H. L. Montgomery and R. C. Vaughan, ``Hilbert's inequality,''
  \emph{Journal of the London Mathematical Society} (2) 8 (1974),
  73--82.
\item
  H. L. Montgomery, \emph{Ten Lectures on the Interface Between Analytic
  Number Theory and Harmonic Analysis}, CBMS Regional Conference Series
  in Mathematics 84, American Mathematical Society, 1994.
\item
  H. Iwaniec and E. Kowalski, \emph{Analytic Number Theory}, American
  Mathematical Society Colloquium Publications 53, 2004.
\item
  J. B. Conrey and A. Gamburd, ``Pseudomoments of the Riemann
  zeta-function and pseudomagic squares,'' arXiv:math/0307213v2;
  \emph{Journal of Number Theory} 117 (2006).
\item
  M. Wahl, ``On the mod-Gaussian convergence of a sum over primes,''
  arXiv:1201.5295v3 (2013), especially equation (1.4).
\item
  M. Gerspach and Y. Lamzouri, ``Low pseudomoments of Euler products,''
  arXiv:2103.03870 (2021).
\item
  A. Connes, ``The Riemann Hypothesis: Past, Present and a Letter
  Through Time,'' arXiv:2602.04022v1 (2026).
\item
  D. S. P. Salazar, ``Order-Moment Transport and Hankel Determinants in
  Special-Function Inequalities,'' arXiv:2606.31647v2 (2026).
\item
  M. Bhowmik, A. Chatterjee, and M. Putinar, ``Holomorphic Interpolation
  of Multivariate Completely Monotone Functions,'' arXiv:2606.12102
  (2026).
\item
  L. Guth and J. Maynard, ``New large value estimates for Dirichlet
  polynomials,'' \emph{Annals of Mathematics} 203 (2026), no. 2,
  623--675, DOI 10.4007/annals.2026.203.2.6; arXiv:2405.20552v2.
\item
  B. Chen, V. Gupta, and Y. C. Li, ``Large Value Estimates for Dirichlet
  Polynomials with Characters and Zero Density of Dirichlet
  \(L\)-Functions,'' arXiv:2507.08296v2 (2026).
\item
  A. Bondarenko, O. F. Brevig, E. Saksman, K. Seip, and J. Zhao,
  ``Pseudomoments of the Riemann zeta function,'' \emph{Bulletin of the
  London Mathematical Society} 50 (2018), no. 4, 709--724, DOI
  10.1112/blms.12183; arXiv:1701.06842v3.
\item
  E. Lill, ``Résolution graphique des équations numériques de tous les
  degrés à une seule inconnue, et description d'un instrument inventé
  dans ce but,'' \emph{Nouvelles Annales de Mathématiques}, 2e série, 6
  (1867), 359--362.
\item
  E. Lill, ``Résolution graphique des équations algébriques qui ont des
  racines imaginaires,'' \emph{Nouvelles Annales de Mathématiques}, 2e
  série, 7 (1868), 363--367.
\item
  B. Bellotti and P.-J. Wong, ``Improved estimates for the argument and
  zero-counting function of the Riemann zeta-function,''
  arXiv:2412.15470v2 (2025), accepted by \emph{Mathematics of
  Computation}.
\item
  D. J. Platt and T. S. Trudgian, ``The Riemann hypothesis is true up to
  \(3\cdot10^{12}\),'' \emph{Bulletin of the London Mathematical
  Society} 53 (2021), 792--797; arXiv:2004.09765.
\item
  O. Dobsch, ``Matrixfunktionen beschränkter Schwankung,''
  \emph{Mathematische Zeitschrift} 43 (1938), 353--388.
\item
  W. F. Donoghue, Jr., \emph{Monotone Matrix Functions and Analytic
  Continuation}, Grundlehren der mathematischen Wissenschaften 207,
  Springer, 1974.
\item
  Euclid, \emph{The Thirteen Books of The Elements}, translated with
  introduction and commentary by T. L. Heath, 2nd ed., Dover, 1956, Book
  VII, Definition 22, and Book IX, Proposition 36.
\item
  L. E. Dickson, \emph{History of the Theory of Numbers}, Volume I:
  Divisibility and Primality, Carnegie Institution of Washington, 1919,
  Chapter I.
\item
  A. Groskin, ``A finite Guinand--Weil dictionary and archimedean tail
  order for the truncated Weil quadratic form,'' arXiv:2607.02828v3
  (2026), archival DOI 10.5281/zenodo.21124802.
\item
  A. Connes, C. Consani, and H. Moscovici, ``Zeta Spectral Triples,''
  arXiv:2511.22755v1 (2025), forthcoming in the EMS Lecture Notes in
  Mathematics volume \emph{Applications of Noncommutative Geometry to
  Gauge Theories, Field Theories, and Quantum Space-Time}.
\item
  A. Connes and W. D. van Suijlekom, ``Quadratic Forms, Real Zeros and
  Echoes of the Spectral Action,'' arXiv:2511.23257v1 (2025).
\item
  A. Connes, C. Consani, and H. Moscovici, ``Zeta zeros and prolate wave
  operators,'' arXiv:2310.18423v2 (2024).
\item
  M. Suzuki, ``Weil's quadratic form via the screw function,''
  arXiv:2606.09096v2 (2026).
\item
  A. Groskin, ``High-Precision Approximation of Riemann Zeros via the
  Truncated Weil Form,'' arXiv:2605.20224v4 (2026), archival DOI
  10.5281/zenodo.19546514.
\item
  E. W. Weisstein, ``Gamma Function,'' \emph{MathWorld--A Wolfram Web
  Resource},
  \href{https://mathworld.wolfram.com/GammaFunction.html}{MathWorld
  entry}, page captured 27 August 2026.
\item
  \emph{Hermitian Spaces}, Chapter 12 lecture slides, CIS 5150,
  University of Pennsylvania, Sections 12.1--12.2, especially the
  definitions of sesquilinear/Hermitian forms and the complex
  polarization identities,
  https://www.cis.upenn.edu/\textasciitilde cis5150/cis515-20-sl9.pdf.
\item
  J. Huckstead, \emph{The Polar Shell Renderer: A Lens, Not a Law},
  revised 9 July 2026, Zenodo, DOI 10.5281/zenodo.21280464.
\item
  A. Voros, ``Zeta functions for the Riemann zeros,'' \emph{Annales de
  l'Institut Fourier} 53 (2003), no. 3, 665--699; first preprint 2001,
  \href{https://arxiv.org/abs/math/0104051v4}{arXiv:math/0104051v4}.
  Equations (40)--(41), (59), (74), (81), and (91)--(95), and Table 1
  supply the classical spectral-zeta formulas used here.
\item
  A. Voros, ``Erratum - Zeta functions for the Riemann zeros,''
  \emph{Annales de l'Institut Fourier} 54 (2004), no. 4, 1139,
  \href{https://aif.centre-mersenne.org/articles/10.5802/aif.2046/}{DOI
  10.5802/aif.2046}.
\item
  J. B. Keiper, ``Power series expansions of Riemann's \(\xi\)
  function,'' \emph{Mathematics of Computation} 58 (1992), no. 198,
  765--773, \href{https://doi.org/10.2307/2153215}{DOI 10.2307/2153215}.
\item
  NIST, \emph{Digital Library of Mathematical Functions}, §§5.4--5.5,
  Gamma special values, recurrence, reflection, and duplication.
  \href{https://dlmf.nist.gov/5.5}{Official formulas}.
\item
  S. M. Eisenberg, ``Moment Sequences and the Bernstein Polynomials,''
  \emph{Canadian Mathematical Bulletin} 12 (1969), no. 4, 401--411.
  \href{https://doi.org/10.4153/CMB-1969-050-8}{DOI
  10.4153/CMB-1969-050-8}. Classical background; the source-specific
  growth argument is printed independently in §164, without a claim of
  literature novelty.
\item
  \emph{A Search-Adjusted Audit of Row-Sector Modulus Signals in
  Terminal Zeta-Zero Spacing Residuals}, unpublished historical report
  (2026), 10 pages. Only its reported tables are used; its raw arrays
  and permutation replay are not part of this work.
\item
  \emph{A Finite-Count Search for Localized Enrichment Near 2.75 in
  Unfolded Riemann Zeta-Zero Gap Coordinates}, unpublished historical
  report, May 2026, 24 pages. Its reported counts and stated design are
  compared in §168.
\item
  NIST, \emph{Digital Library of Mathematical Functions}, sine
  \href{https://dlmf.nist.gov/4.19.E1}{Maclaurin series} and
  \href{https://dlmf.nist.gov/4.22.E1}{Euler infinite product}, and the
  \href{https://dlmf.nist.gov/18.5.E4}{fourth-kind Chebyshev
  representation}. The triangular factorization in (PR.6) is a direct
  regrouping of the factorial coefficients; Gamma reflection is covered
  by {[}59{]}.
\item
  L. Euler, \emph{De formulis exponentialibus replicatis}, E489 (1778),
  \href{https://scholarlycommons.pacific.edu/euler-works/489/}{original
  paper in the Euler Archive}. The fixed-point differentiation and
  stability calculation used here are displayed in §R5E.
\item
  A. Edelman and G. Strang, \emph{Pascal Matrices},
  \href{https://web.mit.edu/18.06/www/Essays/pascal-work.pdf}{MIT course
  essay}. Classical symmetric Pascal matrix, Vandermonde factorization,
  and determinant one. Section 163A identifies this matrix inside the
  completed-source energy.
\item
  Y. Lamzouri, \emph{A new proof that more than two thirds of the zeros
  of the Riemann zeta function are simple or on the critical line},
  \href{https://arxiv.org/abs/2609.02882}{arXiv:2609.02882v2}, 17 pages.
  Proposition 2.1, Lemma 3.1 and Theorem 1.1 are the statements compared
  in §159. Its notation is local to that paper; in particular its
  critical-line counting symbol is not the smooth approximation \(N_0\)
  of §168.
\end{enumerate}

Shared references {[}67{]}--{[}78{]} are listed in the Reading Volume;
the numbering below retains that common bibliography.

\begin{enumerate}
\def\labelenumi{\arabic{enumi}.}
\setcounter{enumi}{78}
\item
  H. L. Montgomery, ``The pair correlation of zeros of the zeta
  function,'' in \emph{Analytic Number Theory}, Proceedings of Symposia
  in Pure Mathematics 24, American Mathematical Society, 1973, 181--193.
  The historical theorem and conjectured density are distinguished in
  Section 61A; the comparison does not import the conjecture as an
  actual-source estimate.
\item
  L. Elaissaoui and Z. E.-A. Guennoun, ``Fourier expansion of the
  Riemann zeta function and applications,'' \emph{Journal of Number
  Theory} 211 (2020), 113--138,
  \href{https://doi.org/10.1016/j.jnt.2019.09.025}{DOI
  10.1016/j.jnt.2019.09.025}; preprint arXiv:1809.02829. Theorem 1.1 and
  Corollary 2.4 of the preprint (v2): the zeta Fourier-coefficient and
  Parseval identities. Section 155G re-derives the Mellin normalization
  independently.
\end{enumerate}

Entries 83--120 were added in the previous edition. They are the
classical primary sources, and the modern accounts, for results that
this volume uses or re-derives. Each is cited for the statement named
where it is used.

\begin{enumerate}
\def\labelenumi{\arabic{enumi}.}
\setcounter{enumi}{82}
\tightlist
\item
  Nicomachus of Gerasa, \emph{Introduction to Arithmetic} (c.~100 CE),
  end of Chapter 20; used at (2.3).
\item
  J. Faulhaber, \emph{Academia Algebrae}, Augsburg, 1631; used at (2.1).
\item
  C. G. J. Jacobi, ``De usu legitimo formulae summatoriae
  Maclaurinianae,'' \emph{Journal für die reine und angewandte
  Mathematik} 12 (1834), 263--272,
  \href{https://doi.org/10.1515/crll.1834.12.263}{DOI
  10.1515/crll.1834.12.263}. The first general proof of Faulhaber's
  assertion.
\item
  D. E. Knuth, ``Johann Faulhaber and sums of powers,''
  \emph{Mathematics of Computation} 61 (1993), no. 203, 277--294,
  \href{https://doi.org/10.2307/2152953}{DOI 10.2307/2152953}. The
  inverse form and derivative rule of §R2.
\item
  R. Wituła, K. Kaczmarek, P. Lorenc, E. Hetmaniok, and M. Pleszczyński,
  ``Jordan numbers, Stirling numbers and sums of powers,''
  \emph{Discussiones Mathematicae -- General Algebra and Applications}
  34 (2014), 155--166, \href{https://doi.org/10.7151/dmgaa.1225}{DOI
  10.7151/dmgaa.1225}. The case \(r=1\) of their product decomposition
  (13) is \(3S_2=j_nT_n\), the summed form of (2.2a).
\item
  N. Derby, ``A search for sums of powers,'' \emph{The Mathematical
  Gazette} 99 (2015), no. 546, 416--421,
  \href{https://doi.org/10.1017/mag.2015.77}{DOI 10.1017/mag.2015.77}.
  Used in §§R2 and R5C; the printed quadratic should read
  \((p-T_{2n})(p+n)=0\).
\item
  D. Singmaster, ``How often does an integer occur as a binomial
  coefficient?'' (Research Problems), \emph{American Mathematical
  Monthly} 78 (1971), no. 4, 385--386,
  \href{https://doi.org/10.2307/2316907}{DOI 10.2307/2316907}.
\item
  B. Pascal, \emph{Traité du triangle arithmétique}, Paris, 1665
  (written 1654).
\item
  P. Mengoli, \emph{Novae quadraturae arithmeticae, seu de additione
  fractionum}, Bologna, 1650; the telescope (4.1).
\item
  L. Euler, ``De summis serierum reciprocarum,'' \emph{Commentarii
  academiae scientiarum Petropolitanae} 7 (1740), 123--134; Euler
  Archive E41,
  \href{https://scholarlycommons.pacific.edu/euler-works/41/}{scholarlycommons.pacific.edu/euler-works/41}.
\item
  L. Euler, ``De progressionibus harmonicis observationes,''
  \emph{Commentarii academiae scientiarum Petropolitanae} 7 (1740),
  150--161; Euler Archive E43,
  \href{https://scholarlycommons.pacific.edu/euler-works/43/}{scholarlycommons.pacific.edu/euler-works/43}.
  The paper that introduces Euler's constant.
\item
  J. Stirling, \emph{Methodus Differentialis: sive Tractatus de
  Summatione et Interpolatione Serierum Infinitarum}, London, 1730.
\item
  J. L. Raabe, ``Angenäherte Bestimmung der Factorenfolge
  \(1\cdot2\cdot3\cdot4\cdot5\cdots n=\Gamma(1+n)=\int x^ne^{-x}\,dx\),
  wenn \(n\) eine sehr grosse Zahl ist,'' \emph{Journal für die reine
  und angewandte Mathematik} 25 (1843), 146--159,
  \href{https://doi.org/10.1515/crll.1843.25.146}{DOI
  10.1515/crll.1843.25.146}. The source of Raabe's integral
  \(\int_0^1\log\Gamma=\tfrac12\log2\pi\).
\item
  E. Cesàro, ``Sur la série harmonique,'' \emph{Nouvelles Annales de
  Mathématiques} (3) 4 (1885), 295--296.
\item
  M. B. Villarino, ``Ramanujan's harmonic number expansion into negative
  powers of a triangular number,'' preprint (2007),
  \href{https://arxiv.org/abs/0707.3950}{arXiv:0707.3950}; see also
  ``Ramanujan's harmonic number expansion,'' preprint (2005),
  \href{https://arxiv.org/abs/math/0511335}{arXiv:math/0511335}. Derives
  (GA.2), with an error estimate, from {[}120{]}.
\item
  C. Mortici and M. B. Villarino, ``On the Ramanujan--Lodge harmonic
  number expansion,'' \emph{Applied Mathematics and Computation} 251
  (2015), 423--430, \href{https://doi.org/10.1016/j.amc.2014.11.088}{DOI
  10.1016/j.amc.2014.11.088}.
\item
  C.-P. Chen, ``On the coefficients of asymptotic expansion for the
  harmonic number by Ramanujan,'' \emph{The Ramanujan Journal} 40
  (2016), 279--290, \href{https://doi.org/10.1007/s11139-015-9670-3}{DOI
  10.1007/s11139-015-9670-3}; and ``Ramanujan's formula for the harmonic
  number,'' \emph{Applied Mathematics and Computation} 317 (2018),
  121--128, \href{https://doi.org/10.1016/j.amc.2017.08.053}{DOI
  10.1016/j.amc.2017.08.053}.
\item
  D. W. DeTemple, ``A quicker convergence to Euler's constant,''
  \emph{American Mathematical Monthly} 100 (1993), no. 5, 468--470,
  \href{https://doi.org/10.2307/2324300}{DOI 10.2307/2324300}.
\item
  L. J. Boya, ``Another relation between \(\pi\), \(e\), \(\gamma\) and
  \(\zeta(n)\),'' \emph{RACSAM} 102 (2008), no. 2, 199--202,
  \href{https://doi.org/10.1007/BF03191819}{DOI 10.1007/BF03191819}. The
  reciprocal-triangular series (GA.3), obtained from Stirling's formula.
\item
  E. T. Whittaker and G. N. Watson, \emph{A Course of Modern Analysis},
  5th ed., Cambridge University Press, 2021, Chapter 12 (the Gamma
  function), \href{https://doi.org/10.1017/9781009004091}{DOI
  10.1017/9781009004091}.
\item
  J. C. Lagarias, ``Euler's constant: Euler's work and modern
  developments,'' \emph{Bulletin of the American Mathematical Society}
  50 (2013), no. 4, 527--628,
  \href{https://doi.org/10.1090/S0273-0979-2013-01423-X}{DOI
  10.1090/S0273-0979-2013-01423-X}.
\item
  O. Espinosa and V. H. Moll, ``On some integrals involving the Hurwitz
  zeta function: Part 1,'' \emph{The Ramanujan Journal} 6 (2002),
  159--188, \href{https://doi.org/10.1023/A:1015706300169}{DOI
  10.1023/A:1015706300169}.
\item
  H. Kinkelin, ``Ueber eine mit der Gammafunction verwandte
  Transcendente und deren Anwendung auf die Integralrechnung,''
  \emph{Journal für die reine und angewandte Mathematik} 57 (1860),
  122--138, \href{https://doi.org/10.1515/crll.1860.57.122}{DOI
  10.1515/crll.1860.57.122}; J. W. L. Glaisher, ``On the product
  \(1^1\cdot2^2\cdot3^3\cdots n^n\),'' \emph{Messenger of Mathematics} 7
  (1878), 43--47.
\item
  B. Nyman, \emph{On Some Groups and Semigroups of Translations},
  thesis, Uppsala, 1950.
\item
  A. Beurling, ``A closure problem related to the Riemann
  zeta-function,'' \emph{Proceedings of the National Academy of
  Sciences} 41 (1955), no. 5, 312--314,
  \href{https://doi.org/10.1073/pnas.41.5.312}{DOI
  10.1073/pnas.41.5.312}.
\item
  L. Báez-Duarte, M. Balazard, B. Landreau, and E. Saias, ``Notes sur la
  fonction \(\zeta\) de Riemann, 3,'' \emph{Advances in Mathematics} 149
  (2000), 130--144, \href{https://doi.org/10.1006/aima.1999.1861}{DOI
  10.1006/aima.1999.1861}.
\item
  L. Báez-Duarte, M. Balazard, B. Landreau, and E. Saias, ``Étude de
  l'autocorrélation multiplicative de la fonction « partie fractionnaire
  »,'' \emph{The Ramanujan Journal} 9 (2005), 215--240,
  \href{https://doi.org/10.1007/s11139-005-0834-4}{DOI
  10.1007/s11139-005-0834-4}.
\item
  J. Hadamard, ``Étude sur les propriétés des fonctions entières et en
  particulier d'une fonction considérée par Riemann,'' \emph{Journal de
  Mathématiques Pures et Appliquées} (4) 9 (1893), 171--215,
  \href{http://www.numdam.org/item/JMPA_1893_4_9__171_0/}{Numdam}.
\item
  H. von Mangoldt, ``Zur Verteilung der Nullstellen der Riemannschen
  Funktion \(\xi(t)\),'' \emph{Mathematische Annalen} 60 (1905), 1--19,
  \href{https://doi.org/10.1007/BF01447494}{DOI 10.1007/BF01447494}.
\item
  H. Davenport, \emph{Multiplicative Number Theory}, 3rd ed., revised by
  H. L. Montgomery, Graduate Texts in Mathematics 74, Springer, 2000,
  Chapter 12, \href{https://doi.org/10.1007/978-1-4757-5927-3}{DOI
  10.1007/978-1-4757-5927-3}. The Hadamard constant
  \(B=-\tfrac12\gamma_{\mathrm E}-1+\tfrac12\log4\pi\) and
  \(\sum_\rho\Re\rho^{-1}=-B\).
\item
  T.-J. Stieltjes, ``Recherches sur les fractions continues,''
  \emph{Annales de la Faculté des Sciences de Toulouse} (1) 8 (1894),
  no. 4, J1--J122,
  \href{http://www.numdam.org/item/AFST_1894_1_8_4_J1_0/}{Numdam}.
\item
  S. Bernstein, ``Sur les fonctions absolument monotones,'' \emph{Acta
  Mathematica} 52 (1929), 1--66,
  \href{https://doi.org/10.1007/BF02592679}{DOI 10.1007/BF02592679}.
\item
  D. V. Widder, \emph{The Laplace Transform}, Princeton University
  Press, 1941.
\item
  A. Weil, ``Sur les « formules explicites » de la théorie des nombres
  premiers,'' \emph{Communications du Séminaire Mathématique de
  l'Université de Lund}, tome supplémentaire (1952), 252--265.
\item
  E. C. Titchmarsh, \emph{Introduction to the Theory of Fourier
  Integrals}, 2nd ed., Oxford University Press, 1948. Mellin inversion
  and the Mellin--Plancherel theorem used at (GA.6).
\item
  H. M. Edwards, \emph{Riemann's Zeta Function}, Academic Press, 1974.
\item
  B. C. Berndt, \emph{Ramanujan's Notebooks, Part V}, Springer-Verlag,
  New York, 1998, Chapter 38, Entry 9, p.~521; the expansion (GA.2).
\item
  D. W. DeTemple and S.-H. Wang, ``Half-integer approximations for the
  partial sums of the harmonic series,'' \emph{Journal of Mathematical
  Analysis and Applications} 160 (1991), no. 1, 149--156,
  \href{https://doi.org/10.1016/0022-247X(91)90296-C}{DOI
  10.1016/0022-247X(91)90296-C}.
\end{enumerate}

\endgroup

\backmatter
\end{document}
